% Generated mechanically from frozen input p01/ko/categories.tex.
% Do not edit this generated file; edit the bound locale source instead.


% OK, start here.
%

\setcounter{chapter}{3}
\chapter{범주}



\phantomsection
\label{categories-section-phantom}


\section{서론}
\label{categories-section-introduction}

\noindent
범주는 \cite{GenEqui}에서 처음 도입되었다. 범주들의 범주(이는 고유 모임이다)는
$2$-범주이다. 마찬가지로 스택들의 범주도 $2$-범주를 이룬다. 범주는 이미
알지만 $2$-범주는 모른다면, 뒤에 나오는 형식적 정의의 입문으로 절
\ref{categories-section-formal-cat-cat}을 읽기 바란다.

\section{정의}
\label{categories-section-definition-categories}

\noindent
표기를 정할 겸 정의들을 상기한다.

\begin{definition}
\label{categories-definition-category}
{\it 범주} $\mathcal{C}$는 다음 자료로 이루어진다.
\begin{enumerate}
\item 대상들의 집합 $\Ob(\mathcal{C})$.
\item 각 쌍 $x, y \in \Ob(\mathcal{C})$에 대한 사상들의 집합
$\Mor_\mathcal{C}(x, y)$.
\item 각 삼중항 $x, y, z\in \Ob(\mathcal{C})$에 대한 합성 사상
$ \Mor_\mathcal{C}(y, z) \times \Mor_\mathcal{C}(x, y)
\to \Mor_\mathcal{C}(x, z) $. 이는 $(\phi, \psi) \mapsto
\phi \circ \psi$로 표기한다.
\end{enumerate}
이 자료는 다음 규칙들을 만족해야 한다.
\begin{enumerate}
\item 모든 원소 $x\in \Ob(\mathcal{C})$에 대하여 사상
$\text{id}_x\in \Mor_\mathcal{C}(x, x)$가 존재하여, 합성이 의미를 갖을
때마다 $\text{id}_x \circ \phi = \phi$이고
$\psi \circ \text{id}_x = \psi $이다.
\item 합성은 결합적이다. 즉 합성이 의미를 갖을 때마다
$(\phi \circ \psi) \circ \chi =
\phi \circ ( \psi \circ \chi)$이다.
\end{enumerate}
\end{definition}

\noindent
모든 사상 집합 $\Mor_\mathcal{C}(x, y)$이 서로소라고 요구하는 것이 관례이다.
이렇게 하면 사상 $\phi : x \to y$는 유일한 {\it 정의역} $x$와 유일한
{\it 공역} $y$를 갖는다. 이는 엄밀히 필요한 조건은 아니지만, 그렇지 않은
경우 위 조건 (2)를 서술할 때 주의해야 한다. 이 조건은 편리하므로 흔히
가정할 것이다. 이 경우 $\phi$의 정의역과 $\psi$의 공역이 같으면 $\phi$와
$\psi$가 {\it 합성 가능}하다고 하며, 이때 $\phi \circ \psi$가 정의된다.
동치인 정의로, 범주를 대상들의 집합, 사상(화살표)들의 집합, 정의역, 공역,
합성으로 이루어진 오중항 $(\text{Ob}, \text{Arrows}, s, t, \circ)$으로 잡고
긴 공리 목록을 부과할 수도 있다. 때때로 이 관점을 사용할 것이다.

\begin{remark}
\label{categories-remark-big-categories}
큰 범주. 어떤 문헌에서는 범주가 대상들의 고유 모임을 가져도 된다고 한다.
이 글에서도 이를 허용하되, 다음 목록의 경우에만 허용한다(필요에 따라 목록을
갱신할 것이다). 특히 ``$\mathcal{C}$를 범주라 하자''라고 말하면
$\Ob(\mathcal{C})$가 집합이라는 뜻으로 이해한다.
\begin{enumerate}
\item 집합들의 범주 $\textit{Sets}$.
\item 아벨 군들의 범주 $\textit{Ab}$.
\item 군들의 범주 $\textit{Groups}$.
\item 군 $G$가 주어졌을 때 왼쪽 $G$-작용을 갖는 집합들의 범주
$G\textit{-Sets}$.
\item 환 $R$이 주어졌을 때 $R$-가군들의 범주 $\text{Mod}_R$.
\item 체 $k$가 주어졌을 때 $k$ 위 벡터공간들의 범주.
\item 환들의 범주.
\item 분할멱 환들의 범주. 분할멱 대수학, 절
\ref{dpa-section-divided-power-rings}을 참조하라.
\item 스킴들의 범주.
\item 위상공간들의 범주 $\textit{Top}$.
\item 위상공간 $X$가 주어졌을 때 $X$ 위 집합의 준층들의 범주
$\textit{PSh}(X)$.
\item 위상공간 $X$가 주어졌을 때 $X$ 위 집합의 층들의 범주 $\Sh(X)$.
\item 위상공간 $X$가 주어졌을 때 $X$ 위 아벨 군의 준층들의 범주
$\textit{PAb}(X)$.
\item 위상공간 $X$가 주어졌을 때 $X$ 위 아벨 군의 층들의 범주
$\textit{Ab}(X)$.
\item 작은 범주 $\mathcal{C}$가 주어졌을 때 $\mathcal{C}$에서
$\textit{Sets}$로 가는 함자들의 범주.
\item 범주 $\mathcal{C}$가 주어졌을 때 $\mathcal{C}$ 위 집합의 준층들의 범주.
\item 사이트 $\mathcal{C}$가 주어졌을 때 $\mathcal{C}$ 위 집합의 층들의 범주.
\end{enumerate}
이 목록을 적은 이유 가운데 하나는 ``큰'' 범주들의 ``모임'' 같은 것을 다루는
일을 피하려는 것이다. 그것은 모든 클래스들의 모임을 다루는 것과 같으며,
이는 분명 메타수학적 대상이라고 생각한다.
\end{remark}






\begin{remark}
\label{categories-remark-unique-identity}
정의에서 바로, $\mathcal{A}$의 대상 $x$에 대한 임의의 두 항등사상은 서로
같다는 것이 따른다. 따라서 $x$의 {\it 그} 항등사상
$\text{id}_x$라고 말할 수 있고 그렇게 할 것이다.
\end{remark}

\begin{definition}
\label{categories-definition-isomorphism}
범주 $\mathcal{C}$의 사상 $\phi : x \to y$에 대하여, 사상
$\psi : y \to x$가 존재하여 $\phi \circ \psi = \text{id}_y$이고
$\psi \circ \phi = \text{id}_x$이면 그 사상을 {\it 동형}이라고 한다.
\end{definition}

\noindent
동형 $\phi$를 때때로 {\it 가역} 사상이라고도 하며, 정의에 나오는 사상
$\psi$를 {\it 역}이라 하고 $\phi^{-1}$로 표기한다. 역이 존재하면 유일하다.
범주 $\mathcal{A}$의 대상 $x$가 주어졌을 때
$\Mor_\mathcal{A}(x, x)$의 가역 원소들의 집합
$\text{Aut}_\mathcal{A}(x)$는 합성에 관하여 군을 이룬다. 이 군을
$\mathcal{A}$에서 $x$의 {\it 자기동형군}이라고 한다.

\begin{definition}
\label{categories-definition-groupoid}
모든 사상이 동형인 범주를 {\it 준군}이라고 한다.
\end{definition}

\begin{example}
\label{categories-example-group-groupoid}
군 $G$로부터 대상이 하나인 준군을 얻는다. 그 대상을 $x$라 하면 사상들은
$\Mor(x, x) = G$이고, 합성 규칙은 $G$의 군 연산으로 주어진다. 대상이 하나인
모든 준군은 이 꼴이다.
\end{example}

\begin{example}
\label{categories-example-set-groupoid}
집합 $C$로부터 다음과 같이 정의되는 준군 $\mathcal{C}$를 얻는다. 대상으로는
$\Ob(\mathcal{C}) := C$를 택하고, 사상 집합 $\Mor(x, y)$는 $x\neq y$이면
공집합, $x = y$이면 $\{\text{id}_x\}$로 택한다.
\end{example}

\begin{definition}
\label{categories-definition-functor}
두 범주 $\mathcal{A}, \mathcal{B}$ 사이의 {\it 함자}
$F : \mathcal{A} \to \mathcal{B}$는 다음 자료로 주어진다.
\begin{enumerate}
\item 사상 $F : \Ob(\mathcal{A}) \to \Ob(\mathcal{B})$.
\item 모든 $x, y \in \Ob(\mathcal{A})$에 대한 사상
$F : \Mor_\mathcal{A}(x, y) \to \Mor_\mathcal{B}(F(x), F(y))$.
이는 $\phi \mapsto F(\phi)$로 표기한다.
\end{enumerate}
이 자료는 다음과 같이 합성과 항등사상에 양립해야 한다. $\mathcal{A}$의
합성 가능한 사상쌍 $(\phi, \psi)$에 대하여
$F(\phi \circ \psi) = F(\phi) \circ F(\psi)$이고,
$F(\text{id}_x) = \text{id}_{F(x)}$이다.
\end{definition}

\noindent
모든 범주 $\mathcal{A}$에는 {\it 항등} 함자 $\text{id}_\mathcal{A}$가 있음을
주목하자. 또한 함자 $G : \mathcal{B} \to \mathcal{C}$와 함자
$F : \mathcal{A} \to \mathcal{B}$가 주어지면, 자명한 방식으로 정의되는
{\it 합성} 함자 $G \circ F : \mathcal{A} \to \mathcal{C}$가 있다.

\begin{definition}
\label{categories-definition-faithful}
$F : \mathcal{A} \to \mathcal{B}$를 함자라 하자.
\begin{enumerate}
\item 임의의 대상 $x, y \in \Ob(\mathcal{A})$에 대하여 사상
$$
F : \Mor_\mathcal{A}(x, y) \to \Mor_\mathcal{B}(F(x), F(y))
$$
이 단사이면 $F$를 {\it 충실}하다고 한다.
\item 이 사상들이 모두 전단사이면 $F$를 {\it 충만충실}하다고 한다.
\item
임의의 대상 $y \in \Ob(\mathcal{B})$에 대하여 $F(x)$가 $\mathcal{B}$에서
$y$와 동형이 되는 대상 $x \in \Ob(\mathcal{A})$가 존재하면 함자 $F$를
{\it 본질적으로 전사}라고 한다.
\end{enumerate}
\end{definition}

\begin{definition}
\label{categories-definition-subcategory}
범주 $\mathcal{B}$의 {\it 부분범주}란, 그 대상과 화살표가 각각
$\mathcal{B}$의 대상과 화살표의 부분집합을 이루는 범주 $\mathcal{A}$로서,
$\mathcal{A}$의 출발점, 도착점, 합성이 $\mathcal{B}$의 것과 일치하고
$\mathcal{A}$의 각 대상의 항등사상도 $\mathcal{B}$에서의 항등사상과
일치하는 것을 말한다. 모든 $x, y \in \Ob(\mathcal{A})$에 대하여
$\Mor_\mathcal{A}(x, y) = \Mor_\mathcal{B}(x, y)$이면 $\mathcal{A}$를
$\mathcal{B}$의 {\it 충만한 부분범주}라고 한다. $\mathcal{A}$가 충만한
부분범주이고, $x \in \Ob(\mathcal{A})$가 주어졌을 때 $x$와 동형인
$\mathcal{B}$의 모든 대상도 $\mathcal{A}$에 속하면 이를
$\mathcal{B}$의 {\it 엄밀히 충만한} 부분범주라고 한다.
\end{definition}

\noindent
$\mathcal{A} \subset \mathcal{B}$가 부분범주이면 항등 사상은
$\mathcal{A}$에서 $\mathcal{B}$로 가는 함자이다. 더 나아가 부분범주
$\mathcal{A} \subset \mathcal{B}$가 충만할 필요충분조건은 포함 함자가
충만충실한 것이다. 범주 $\mathcal{B}$가 주어졌을 때 $\mathcal{B}$의
충만한 부분범주들의 집합은 $\Ob(\mathcal{B})$의 부분집합들의 집합과
같다는 점에 유의하자.

\begin{remark}
\label{categories-remark-functor-into-sets}
$\mathcal{A}$가 범주라고 하자. $\mathcal{A}$에서 $\textit{Sets}$로 가는
함자 $F$는 집합의 범주가 ``크더라도'' 수학적 대상이다(즉, 클래스나
집합론의 공식이 아니라 집합이다. Sets, Section
\ref{sets-section-sets-everything}을 보라). 실제로 대상들에 대한 $F$의
상은 집합 $F(\Ob(\mathcal{A}))$이며, 따라서 $F$를 $\mathcal{A}$와,
대상들이 $F(\Ob(\mathcal{A}))$의 원소들인 집합 범주의 충만한
부분범주 사이의 함자로 생각할 수 있다.
\end{remark}

\begin{example}
\label{categories-example-group-homomorphism-functor}
군의 준동형 $p : G\to H$는 Example \ref{categories-example-group-groupoid}의 대응하는
준군들 사이의 함자를 준다. 이 함자가 충실할(각각 충만충실할)
필요충분조건은 $p$가 단사일(각각 동형일) 것이다.
\end{example}

\begin{example}
\label{categories-example-category-over-X}
범주 $\mathcal{C}$와 대상 $X\in \Ob(\mathcal{C})$가 주어졌다고 하자.
{\it $X$ 위의 대상들의 범주}를 다음과 같이 정의하고 $\mathcal{C}/X$로
표기한다. $\mathcal{C}/X$의 대상은 어떤 $Y\in \Ob(\mathcal{C})$에 대한
사상 $Y\to X$이다. 대상 $Y\to X$와 $Y'\to X$ 사이의 사상은 자명한
도표를 가환하게 하는 $\mathcal{C}$의 사상 $Y\to Y'$이다. 사상을 단순히
잊는 함자 $p_X : \mathcal{C}/X\to \mathcal{C}$가 있음에 유의하자.
또한 $\mathcal{C}$의 사상 $f : X'\to X$가 주어지면, $f$와 합성하여
얻는 유도 함자 $F : \mathcal{C}/X' \to \mathcal{C}/X$가 있으며
$p_X\circ F = p_{X'}$이다.
\end{example}

\begin{example}
\label{categories-example-category-under-X}
범주 $\mathcal{C}$와 대상 $X\in \Ob(\mathcal{C})$가 주어졌다고 하자.
{\it $X$ 아래의 대상들의 범주}를 다음과 같이 정의하고 $X/\mathcal{C}$로
표기한다. $X/\mathcal{C}$의 대상은 어떤 $Y\in \Ob(\mathcal{C})$에 대한
사상 $X\to Y$이다. 대상 $X\to Y$와 $X\to Y'$ 사이의 사상은 자명한
도표를 가환하게 하는 $\mathcal{C}$의 사상 $Y\to Y'$이다. 사상을 단순히
잊는 함자 $p_X : X/\mathcal{C}\to \mathcal{C}$가 있음에 유의하자.
또한 $\mathcal{C}$의 사상 $f : X'\to X$가 주어지면, $f$와 합성하여
얻는 유도 함자 $F : X/\mathcal{C} \to X'/\mathcal{C}$가 있으며
$p_{X'}\circ F = p_X$이다.
\end{example}




\begin{definition}
\label{categories-definition-transformation-functors}
$F, G : \mathcal{A} \to \mathcal{B}$를 함자라 하자. {\it 자연변환}, 또는
{\it 함자의 사상} $t : F \to G$란 다음 조건을 만족하는 모음
$\{t_x\}_{x\in \Ob(\mathcal{A})}$이다.
\begin{enumerate}
\item $t_x : F(x) \to G(x)$는 범주 $\mathcal{B}$의 사상이다.
\item $\mathcal{A}$의 모든 사상 $\phi : x \to y$에 대하여 다음 도표가
가환한다.
$$
\xymatrix{
F(x) \ar[r]^{t_x} \ar[d]_{F(\phi)} & G(x) \ar[d]^{G(\phi)} \\
F(y) \ar[r]^{t_y} & G(y) }
$$
\end{enumerate}
\end{definition}

\noindent
때때로 $t$가 $F$에서 $G$로 가는 사상임을 나타내기 위해 도표
$$
\xymatrix{
\mathcal{A}
\rtwocell^F_G{t}
&
\mathcal{B}
}
$$
를 사용한다.

\medskip\noindent
모든 함자 $F$에는 {\it 항등} 변환 $\text{id}_F : F \to F$가 딸려 있다.
또한 함자의 사상 $t : F \to G$와 함자의 사상 $s : E \to F$가 주어지면,
그 {\it 합성} $t \circ s$를 규칙
$$
(t \circ s)_x = t_x \circ s_x : E(x) \to G(x)
$$
에 따라 모든 $x \in \Ob(\mathcal{A})$에 대해 정의한다. 이것이 실제로
$E$에서 $G$로 가는 함자의 사상임은 쉽게 확인된다. 이와 같이 범주
$\mathcal{A}$와 $\mathcal{B}$가 주어지면 새로운 범주, 즉 $\mathcal{A}$와
$\mathcal{B}$ 사이의 함자들의 범주를 얻는다.

\begin{remark}
\label{categories-remark-functors-sets-sets}
$\mathcal{A}$가 ``큰'' 범주일 때에는 같은 일이 성립하지 않는 한 가지
경우가 바로 이것이다. 예를 들어 함자 $\textit{Sets} \to \textit{Sets}$를
생각하자. 현재의 정의에 따르면 그러한 함자는 집합이 아니라 클래스이다.
다시 말해, 그것은 (일부 변수들이 지정된 집합들과 같은) 집합론의 공식으로
주어진다! 가능한 모든 집합론 공식을 수학적 대상의 정의 일부로 간주하려는
것은 좋은 생각이 아니다. 예를 들어 스킴들의 범주 위의 층들을 생각할 때에도
같은 문제가 생긴다. 이 점은 나중에 다시 다룬다.
\end{remark}

\begin{definition}
\label{categories-definition-equivalence-categories}
{\it 범주의 동치} $F : \mathcal{A} \to \mathcal{B}$란, 함자
$G : \mathcal{B} \to \mathcal{A}$가 존재하여 합성 $F \circ G$와
$G \circ F$가 각각 항등 함자 $\text{id}_\mathcal{B}$와
$\text{id}_\mathcal{A}$에 동형인 함자를 말한다. 이 경우 $G$를 $F$의
{\it 준역}이라고 한다.
\end{definition}

\begin{lemma}
\label{categories-lemma-construct-quasi-inverse}
$F : \mathcal{A} \to \mathcal{B}$를 충만충실한 함자라 하자. 모든
$X \in \Ob(\mathcal{B})$에 대하여 $\mathcal{A}$의 대상 $j(X)$와 동형
$i_X : X \to F(j(X))$가 주어졌다고 하자. 그러면 대상들에 대한 이 규칙을
확장하는 유일한 함자 $j : \mathcal{B} \to \mathcal{A}$가 존재하여 $j$가
동형들 $i_X$는 함자의 동형 $\text{id}_\mathcal{B} \to F \circ j$를
정의한다. 더욱이 $j$와 $F$는 서로 준역인 범주의 동치이다.
\end{lemma}

\begin{proof}
$j : \mathcal{B} \to \mathcal{A}$를 구성하는 데에는 두 단계가 있다.
먼저 $X \in \mathcal{B}$에 $j(X)$를 대응시키는 사상
$j : \Ob(\mathcal{B}) \to \Ob(\mathcal{A})$를 정의한다. 둘째,
$X,Y \in \Ob(\mathcal{B})$이고 $\phi : X \to Y$이면
$\phi' := i_Y \circ\phi\circ i_X^{-1}$를 생각한다. $F$가 충만충실하므로
$F(\varphi) = \phi'$를 만족하는 유일한 $\varphi$가 있다. $j(\phi) = \varphi$로
정의한다. $j$가 함자라는 확인은 생략한다. 구성에 의해 도표
$$
\xymatrix{
X \ar[r]_-{i_X} \ar[d]_{\phi}
&
F(j(X)) \ar[d]^{F\circ j(\phi)}
\\
Y \ar[r]^-{i_Y}
&
F(j(Y))
}
$$
는 가환한다. 따라서 각 $i_X$가 동형이므로 $\{i_X\}_X$는 함자의 동형
$\text{id}_\mathcal{B} \to F\circ j$이다. 결론을 내리려면 또한
$j \circ F$가 $\text{id}_\mathcal{A}$와 동형임을 보여야 한다. 그러나
$F$가 충만충실하므로, 이를 보이려면 $F$를 뒤에서 합성한 뒤에 보이는
것으로 충분하다. 즉, $F \circ j \circ F$가
$F \circ \text{id}_\mathcal{A}$와 동형임을 보이면 충분하다(작은 세부
사항은 생략한다). 이는 $F \circ j \cong \text{id}_\mathcal{B}$이므로
명백하다.
\end{proof}

\begin{lemma}
\label{categories-lemma-equivalence-categories}
함자가 범주의 동치일 필요충분조건은 그 함자가 충만충실하고
본질적으로 전사인 것이다.
\end{lemma}

\begin{proof}
$F : \mathcal{A} \to \mathcal{B}$가 본질적으로 전사이고 충만충실하다고
하자. 관례상 모든 범주는 작고 $F$는 본질적으로 전사이므로, 선택 공리를
사용하여 모든 $X \in \Ob(\mathcal{B})$에 대해 $\mathcal{A}$의 대상
$j(X)$와 동형 $i_X : X \to F(j(X))$를 택할 수 있다. 이제 $F$가
충만충실하다는 사실을 사용하여 Lemma \ref{categories-lemma-construct-quasi-inverse}를
적용한다.
\end{proof}

\begin{definition}
\label{categories-definition-product-category}
$\mathcal{A}$, $\mathcal{B}$를 범주라 하자. {\it 곱범주}
$\mathcal{A} \times \mathcal{B}$를 대상이
$\Ob(\mathcal{A} \times \mathcal{B}) =
\Ob(\mathcal{A}) \times \Ob(\mathcal{B})$이고
$$
\Mor_{\mathcal{A} \times \mathcal{B}}((x, y), (x', y'))
:=
\Mor_\mathcal{A}(x, x')\times
\Mor_\mathcal{B}(y, y').
$$
인 범주로 정의한다. 합성은 성분별로 정의한다.
\end{definition}


\section{반대범주와 요네다 보조정리}
\label{categories-section-opposite}

\begin{definition}
\label{categories-definition-opposite}
범주 $\mathcal{C}$가 주어졌을 때 {\it 반대범주}
$\mathcal{C}^{opp}$는 $\mathcal{C}$와 같은 대상들을 가지되 모든 사상의
방향을 뒤집은 범주이다.
\end{definition}

\noindent
다시 말해
$\Mor_{\mathcal{C}^{opp}}(x, y) = \Mor_\mathcal{C}(y, x)$이다.
$\mathcal{C}^{opp}$에서의 합성은 방향만 반대일 뿐 $\mathcal{C}$에서와
같다. 즉, $\phi : y \to z$와 $\psi : x \to y$가
$\mathcal{C}^{opp}$의 사상, 다시 말해 $\mathcal{C}$에서의 화살표
$z \to y$와 $y \to x$라면, $\phi \circ^{opp} \psi$는
$\mathcal{C}$에서의 합성 $z \to y \to x$에 대응하는
$\mathcal{C}^{opp}$의 사상 $x \to z$이다.

\begin{definition}
\label{categories-definition-contravariant}
$\mathcal{C}$, $\mathcal{S}$를 범주라 하자. $\mathcal{C}$에서
$\mathcal{S}$로 가는 {\it 반변} 함자 $F$란 함자
$\mathcal{C}^{opp}\to \mathcal{S}$를 말한다.
\end{definition}

\noindent
구체적으로, 반변 함자 $F$는 사상
$F : \Ob(\mathcal{C}) \to \Ob(\mathcal{S})$와
$\mathcal{C}$의 각 사상 $\psi : x \to y$에 대한 사상
$F(\psi) : F(y) \to F(x)$로 주어진다. 이들은 다음 성질을 만족해야 한다.
또 다른 사상 $\phi : y \to z$가 주어지면,
$F(\phi \circ \psi) = F(\psi) \circ F(\phi)$가
$F(z) \to F(x)$의 사상으로서 성립한다. (순서가 뒤집힘에 유의하라.)

\begin{definition}
\label{categories-definition-presheaf}
$\mathcal{C}$를 범주라 하자.
\begin{enumerate}
\item $\mathcal{C}$ 위의 {\it 집합 준층}, 또는 간단히 {\it 준층}은
$\mathcal{C}$에서 $\textit{Sets}$로 가는 반변 함자 $F$이다.
\item 준층들의 범주는 $\textit{PSh}(\mathcal{C})$로 표기한다.
\end{enumerate}
\end{definition}

\noindent
물론 준층들의 범주는 고유 모임이다.

\begin{example}
\label{categories-example-hom-functor}
점 함자.
임의의 $U\in \Ob(\mathcal{C})$에 대하여 반변 함자
$$
\begin{matrix}
h_U & : & \mathcal{C}
&
\longrightarrow
&
\textit{Sets} \\
& &
X
&
\longmapsto
&
\Mor_\mathcal{C}(X, U)
\end{matrix}
$$
가 있으며, 이 함자는 대상 $X$를 집합 $\Mor_\mathcal{C}(X, U)$로
보낸다. 다시 말해 $h_U$는 준층이다. 사상 $f : X\to Y$가 주어지면
대응하는 사상
$h_U(f) : \Mor_\mathcal{C}(Y, U)\to \Mor_\mathcal{C}(X, U)$는
$\phi$를 $\phi\circ f$로 보낸다. 이 준층은 항상
$h_U : \mathcal{C}^{opp} \to \textit{Sets}$로 표기한다. 이를 $U$에
대응하는 {\it 표현 가능한 준층}이라고 한다. $\mathcal{C}$가 스킴들의
범주이면 이 함자를 때때로 $U$의 \emph{점 함자}라고 한다.
\end{example}

\noindent
$\mathcal{C}$의 사상 $\phi : U \to V$가 주어지면, 사상 $U \to V$와
합성하여 정의되는 함자의 자연변환 $h(\phi) : h_U \to h_V$를 얻는다.
이는 $\mathcal{C}$에서의 사상 합성을 함자의 변환 합성으로 보낸다.
다시 말해 함자
$$
h :
\mathcal{C}
\longrightarrow
\textit{PSh}(\mathcal{C})
$$
를 얻는다. 그 공역은 “큰” 범주임에 유의하라. Remark
\ref{categories-remark-big-categories}를 보라. 한편 $h$는 실제 수학적 대상
(즉, 집합)이다. Remark \ref{categories-remark-functor-into-sets}와 비교하라.

\begin{lemma}[요네다 보조정리]
\label{categories-lemma-yoneda}
\begin{reference}
\cite{Yoneda-homology}에 어떤 형태로 등장했으며, 그로텐디크가
\cite{Gr-II}에서 일반화된 형태로 사용하였다.
\end{reference}
$U, V \in \Ob(\mathcal{C})$라 하자. 함자의 사상 $s : h_U \to h_V$가
주어지면 $h(\phi) = s$를 만족하는 유일한 사상
$\phi : U \to V$가 존재한다. 다시 말해 함자 $h$는 충만충실하다.
더 일반적으로, 임의의 반변 함자 $F$와 $\mathcal{C}$의 임의의 대상
$U$에 대하여 자연스러운 전단사
$$
\Mor_{\textit{PSh}(\mathcal{C})}(h_U, F) \longrightarrow F(U),
\quad
s \longmapsto s_U(\text{id}_U).
$$
가 있다.
\end{lemma}

\begin{proof}
첫 번째 명제에서는 단지
$\phi = s_U(\text{id}_U) \in \Mor_\mathcal{C}(U, V)$로 두면 된다.
두 번째 명제에서는 $\xi \in F(U)$가 주어졌을 때, $h_U(V) =
\Mor_\mathcal{C}(V, U)$의 원소 $f : V \to U$를 $F(f)(\xi)$로 보내도록
$s_V : h_U(V) \to F(V)$를 정하여 $s$를 정의한다.
\end{proof}

\begin{definition}
\label{categories-definition-representable-functor}
반변 함자 $F : \mathcal{C}\to \textit{Sets}$가 $\mathcal{C}$의 어떤 대상
$U$에 대한 점 함자 $h_U$와 동형이면 이를 {\it 표현 가능}하다고 한다.
\end{definition}

\noindent
$\mathcal{C}$를 범주라 하고
$F : \mathcal{C}^{opp} \to \textit{Sets}$를 표현 가능한 함자라 하자.
$\mathcal{C}$의 대상 $U$와 동형 $s : h_U \to F$를 택하자. 요네다
보조정리에 따르면 쌍 $(U, s)$는 유일한 동형을 제외하고 유일하다.
대상 $U$를 $F$를 {\it 표현하는} 대상이라고 한다. 요네다 보조정리에 의해
변환 $s$는 유일한 원소 $\xi \in F(U)$에 대응한다. 이 원소를
{\it 보편 원소}라고 한다. 이 원소는 $V \in \Ob(\mathcal{C})$에 대하여
사상
$$
\Mor_\mathcal{C}(V, U) \longrightarrow F(V),\quad
(f : V \to U) \longmapsto F(f)(\xi)
$$
이 전단사라는 성질을 갖는다. 따라서 $\xi$는 $F(V)$의 모든 원소가
유일한 $\mathcal{C}$의 사상 $f : V \to U$에 대한 $F(f)$에 의한 $\xi$의 상과 같다는
뜻에서 보편적이다.


\section{두 대상의 곱}
\label{categories-section-products-pairs}

\begin{definition}
\label{categories-definition-products}
$x, y\in \Ob(\mathcal{C})$라 하자. $x$와 $y$의 {\it 곱}은 대상
$x \times y \in \Ob(\mathcal{C})$와 사상
$p\in \Mor_{\mathcal C}(x \times y, x)$ 및
$q\in\Mor_{\mathcal C}(x \times y, y)$로 이루어지고 다음 보편 성질을
만족하는 것이다. 임의의 $w\in \Ob(\mathcal{C})$와 사상
$\alpha \in \Mor_{\mathcal C}(w, x)$ 및
$\beta \in \Mor_\mathcal{C}(w, y)$에 대하여 도표
$$
\xymatrix{
w \ar[rrrd]^\beta \ar@{-->}[rrd]_\gamma \ar[rrdd]_\alpha & & \\
& & x \times y \ar[d]_p \ar[r]_q & y \\
& & x &
}
$$
를 가환하게 하는 유일한
$\gamma\in \Mor_{\mathcal C}(w, x \times y)$가 존재한다.
\end{definition}

\noindent
곱이 존재하면 유일한 동형을 제외하고 유일하다. 이는 정의가
$x \times y$를 $w$에 대하여 함자적으로
$$
h_{x \times y}(w) = h_x(w) \times h_y(w)
$$
를 만족하는 $\mathcal{C}$의 대상으로 요구하므로 요네다 보조정리에서
따른다. 다시 말해 곱 $x \times y$는 함자
$w \mapsto h_x(w) \times h_y(w)$를 표현하는 대상이다.

\begin{definition}
\label{categories-definition-has-products-of-pairs}
임의의 $x, y \in \Ob(\mathcal{C})$에 대하여 곱 $x \times y$가 존재하면
범주 $\mathcal{C}$가 {\it 대상쌍의 곱을 갖는다}고 한다.
\end{definition}

\noindent
이 용어는 보통 다른 뜻을 갖는 “곱을 갖는다” 또는 “유한 곱을 갖는다”는
개념과 구별하기 위해 사용한다(특히 후자는 언제나 끝 대상의 존재를
함의한다).




\section{두 대상의 쌍대곱}
\label{categories-section-coproducts-pairs}

\begin{definition}
\label{categories-definition-coproducts}
$x, y \in \Ob(\mathcal{C})$라 하자. $x$와 $y$의 {\it 쌍대곱}, 또는
{\it 합병합}은 대상 $x \amalg y \in \Ob(\mathcal{C})$와 사상
$i \in \Mor_{\mathcal C}(x, x \amalg y)$ 및
$j \in \Mor_{\mathcal C}(y, x \amalg y)$로 이루어지고 다음 보편 성질을
만족하는 것이다. 임의의 $w \in \Ob(\mathcal{C})$와 사상
$\alpha \in \Mor_{\mathcal C}(x, w)$ 및
$\beta \in \Mor_\mathcal{C}(y, w)$에 대하여 도표
$$
\xymatrix{
& y \ar[d]^j \ar[rrdd]^\beta \\
x \ar[r]^i \ar[rrrd]_\alpha & x \amalg y \ar@{-->}[rrd]^\gamma \\
& & & w
}
$$
를 가환하게 하는 유일한
$\gamma \in \Mor_{\mathcal C}(x \amalg y, w)$가 존재한다.
\end{definition}

\noindent
쌍대곱이 존재하면 유일한 동형을 제외하고 유일하다. 정의가
$x \amalg y$를 $w$에 대하여 함자적으로
$$
\Mor_\mathcal{C}(x \amalg y, w) =
\Mor_\mathcal{C}(x, w) \times \Mor_\mathcal{C}(y, w)
$$
를 만족하는 $\mathcal{C}$의 대상으로 요구하므로, 이는 반대범주에
적용한 요네다 보조정리에서 따른다.

\begin{definition}
\label{categories-definition-has-coproducts-of-pairs}
임의의 $x, y \in \Ob(\mathcal{C})$에 대하여 쌍대곱
$x \amalg y$가 존재하면 범주 $\mathcal{C}$가
{\it 대상쌍의 쌍대곱을 갖는다}고 한다.
\end{definition}

\noindent
이 용어는 보통 다른 뜻을 갖는 “쌍대곱을 갖는다” 또는 “유한 쌍대곱을
갖는다”는 개념과 구별하기 위해 사용한다(특히 후자는 언제나
$\mathcal{C}$에 초기 대상이 존재함을 함의한다).





\section{섬유곱}
\label{categories-section-fibre-products}

\begin{definition}
\label{categories-definition-fibre-products}
$x, y, z\in \Ob(\mathcal{C})$라 하고,
$f\in \Mor_\mathcal{C}(x, y)$ 및
$g\in \Mor_{\mathcal C}(z, y)$라 하자. $f$와 $g$의 {\it 섬유곱}은
대상 $x \times_y z\in \Ob(\mathcal{C})$와 사상
$p \in \Mor_{\mathcal C}(x \times_y z, x)$ 및
$q \in \Mor_{\mathcal C}(x \times_y z, z)$로 이루어지고 도표
$$
\xymatrix{
x \times_y z \ar[r]_q \ar[d]_p & z \ar[d]^g \\
x \ar[r]^f & y
}
$$
를 가환하게 하며 다음 보편 성질을 만족하는 것이다. 임의의
$w\in \Ob(\mathcal{C})$와, $f \circ \alpha = g \circ \beta$를 만족하는
사상 $\alpha \in \Mor_{\mathcal C}(w, x)$ 및
$\beta \in \Mor_\mathcal{C}(w, z)$에 대하여 도표
$$
\xymatrix{
w \ar[rrrd]^\beta \ar@{-->}[rrd]_\gamma \ar[rrdd]_\alpha & & \\
& & x \times_y z \ar[d]^p \ar[r]_q & z \ar[d]^g \\
& & x \ar[r]^f & y
}
$$
를 가환하게 하는 유일한
$\gamma \in \Mor_{\mathcal C}(w, x \times_y z)$가 존재한다.
\end{definition}

\noindent
섬유곱이 존재하면 유일한 동형을 제외하고 유일하다. 이는 정의가
$x \times_y z$를 $w$에 대하여 함자적으로
$$
h_{x \times_y z}(w) = h_x(w) \times_{h_y(w)} h_z(w)
$$
를 만족하는 $\mathcal{C}$의 대상으로 요구하므로 요네다 보조정리에서
따른다. 다시 말해 섬유곱 $x \times_y z$는 함자
$w \mapsto h_x(w) \times_{h_y(w)} h_z(w)$를 표현하는 대상이다.

\begin{definition}
\label{categories-definition-cartesian}
범주의 가환 도표
$$
\xymatrix{
w \ar[r] \ar[d] &
z \ar[d] \\
x \ar[r] &
y
}
$$
에서 $w$와 사상 $w \to x$, $w \to z$가 사상 $x \to y$와
$z \to y$의 섬유곱을 이루면 이 도표를 {\it 카르테시안}이라고 한다.
\end{definition}

\begin{definition}
\label{categories-definition-has-fibre-products}
임의의 $f\in \Mor_{\mathcal C}(x, y)$와
$g\in \Mor_{\mathcal C}(z, y)$에 대하여 섬유곱이 존재하면 범주
$\mathcal{C}$가 {\it 섬유곱을 갖는다}고 한다.
\end{definition}

\begin{definition}
\label{categories-definition-representable-morphism}
범주 $\mathcal{C}$의 사상 $f : x \to y$에 대하여, $\mathcal{C}$의
모든 사상 $z \to y$에 대한 섬유곱 $x \times_y z$가 존재하면
그 사상을 {\it 표현 가능}하다고 한다.
\end{definition}

\begin{lemma}
\label{categories-lemma-composition-representable}
$\mathcal{C}$를 범주라 하자. $f : x \to y$와 $g : y \to z$가 표현
가능하다고 하자. 그러면 $g \circ f : x \to z$도 표현 가능하다.
\end{lemma}

\begin{proof}
$t \in \Ob(\mathcal C)$이고 $ \varphi \in \Mor_{\mathcal C}(t,z) $라
하자. $g$와 $f$가 표현 가능하므로 보편 성질이
Definition \ref{categories-definition-fibre-products}의 것인 가환 도표
$$
\xymatrix{
y \times_z t \ar[r]_q \ar[d]_p &
t \ar[d]^{\varphi} \\
y \ar[r]^{g} & z
}
\quad\quad
\xymatrix{
x \times_y (y\! \times_z\! t) \ar[r]_{q'} \ar[d]_{p'} &
y \times_z t \ar[d]^p \\
x \ar[r]^f & y
}
$$
를 얻는다. 사상 $q \circ q' : x \times_z t \to t$와
$p' : x \times_z t \to x$를 갖춘
$x \times_z t = x \times_y (y \times_z t)$가 섬유곱이라고 주장한다.
먼저 위 도표들의 가환성에서
$\varphi \circ q \circ q' = g \circ f \circ p'$가 따른다. 보편 성질을
확인하기 위해 $w \in \Ob(\mathcal C)$라 하고, 사상
$\alpha : w \to x$와 $\beta : w \to t$가
$\varphi \circ \beta = g \circ f \circ \alpha$를 만족한다고 하자.
섬유곱의 정의에 따라 도표
$$
\xymatrix{
w \ar[rrrd]^\beta \ar@{-->}[rrd]_\delta \ar[rrdd]_{f\circ\alpha} & & \\
& & y \times_z t \ar[d]_p \ar[r]_q & t \ar[d]^{\varphi} \\
& & y \ar[r]^{g} & z
}
$$
와
$$
\xymatrix{
w \ar[rrrd]^\delta \ar@{-->}[rrd]_\gamma \ar[rrdd]_{\alpha} & & \\
& & x \times_y (y\!\times_z\! t) \ar[d]_{p'} \ar[r]_{q'} &
y \times_z t \ar[d]^{p} \\
& & x \ar[r]^{f} & y
}
$$
를 가환하게 하는 유일한 사상 $\delta$와 $\gamma$가 있다. 그러면
$\gamma$는 도표
$$
\xymatrix{
w \ar[rrrd]^\beta \ar@{-->}[rrd]_\gamma \ar[rrdd]_{\alpha} & & \\
& & x \times_z t \ar[d]_{p'} \ar[r]_{q\circ q'} & t \ar[d]^{\varphi} \\
& & x \ar[r]^{g\circ f} & z
}
$$
를 가환하게 한다. 그 유일성을 보이기 위해
$ q\circ q'\circ\gamma' = \beta $ 및
$ p'\circ \gamma' = \alpha $를 만족하는 $\gamma'$를 택하자.
$\gamma$가 유일하므로 결론을 내리려면
$ q'\circ\gamma' = \delta $ 및 $ p'\circ\gamma' = \alpha $를 보이면
충분하다. 두 번째 등식은 가정하였다. 첫 번째 등식에는 그 사상의
유일성도 사용해야 한다. $\delta$는
$ q\circ\delta = \beta $ 및 $ p\circ\delta = f\circ\alpha $를 만족하는
유일한 사상임에 유의하자. 이미
$ q\circ (q'\circ\gamma') = \beta $를 가정하였다. 또한 섬유곱의
정의에서 $ f\circ p' = p\circ q' $임을 안다. 따라서
$$
p\circ (q'\circ\gamma') =
(p\circ q')\circ\gamma' =
(f\circ p')\circ\gamma' =
f\circ (p'\circ\gamma') = f\circ\alpha.
$$
그러므로 $ q'\circ\gamma' = \delta $이고, 이로써 증명이 끝난다.
\end{proof}

\begin{lemma}
\label{categories-lemma-base-change-representable}
$\mathcal{C}$를 범주라 하자. $f : x \to y$를 표현 가능한 사상이라 하고
$y' \to y$를 $\mathcal{C}$의 사상이라 하자. 그러면 사상
$x' := x \times_y y' \to y'$도 표현 가능하다.
\end{lemma}

\begin{proof}
$z \to y'$를 사상이라 하자. 섬유곱 $x' \times_{y'} z$는 함자
\begin{eqnarray*}
w & \mapsto & h_{x'}(w)\times_{h_{y'}(w)} h_z(w) \\
& = & (h_x(w) \times_{h_y(w)} h_{y'}(w)) \times_{h_{y'}(w)} h_z(w) \\
& = & h_x(w) \times_{h_y(w)} h_z(w)
\end{eqnarray*}
를 표현해야 하는데, 이 함자는 가정에 의해 표현 가능하다.
\end{proof}

\section{섬유곱의 예}
\label{categories-section-example-fibre-products}

\noindent
이 절에서는 섬유곱의 예들을 열거하고 설명한다.

\medskip\noindent
첫 번째로 정말 자명한 예를 들면, 집합들의 범주는 섬유곱을 가지므로 모든
사상이 표현 가능하다. 실제로 $f : X \to Y$와 $g : Z \to Y$가 집합의
사상이면, $X \times_Y Z$를 $f(x) = g(z)$를 만족하는 쌍 $(x, z)$들로
이루어진 $X \times Z$의 부분집합으로 정의한다. 사상
$p : X \times_Y Z \to X$와 $q : X \times_Y Z \to Z$는 각각 사영 사상
$(x, z) \mapsto x$와 $(x, z) \mapsto z$이다. 마지막으로 사상
$\alpha : W \to X$와 $\beta : W \to Z$가
$f \circ \alpha = g \circ \beta$를 만족하면, 사상
$W \to X \times Z$, $w\mapsto (\alpha(w), \beta(w))$의 상은 자명하게
원하는 대로 $X \times_Y Z$에 들어간다.

\medskip\noindent
대상이 특정 종류의 대수 구조를 갖춘 집합인 많은 범주에서는, 기저 집합들의
섬유곱이 그 범주에서의 섬유곱도 준다. 예를 들어 위의 $X$, $Y$, $Z$가
군이고 $f$, $g$가 군 준동형이라고 하자. 그러면 집합론적 섬유곱
$X \times_Y Z$는 두 쌍의 곱을
$(x, z) \cdot (x', z') = (xx', zz')$로 정의함으로써 군 구조를 물려받는다.
비슷한 논법이 성립하는 범주들을 열거하면 다음과 같다.
\begin{enumerate}
\item 군들의 범주 $\textit{Groups}$.
\item 고정된 군 $G$에 대하여 왼쪽 $G$-작용을 갖춘 집합들의 범주
$G\textit{-Sets}$.
\item 환들의 범주.
\item 환 $R$이 주어졌을 때 $R$-가군들의 범주.
\end{enumerate}




\section{섬유곱과 표현 가능성}
\label{categories-section-representable-map-presheaves}

\noindent
이 절에서는 한 범주에서 집합들의 범주로 가는 반변 함자들의 범주에서
섬유곱을 구한다. 이 논의는 나중에 사이트, 준층, 층을 다룰 때 더 일반적인
논의로 대체된다. 이러한 함자 사이의 “표현 가능한 사상” 개념도 흥미롭다.

\begin{lemma}
\label{categories-lemma-fibre-product-presheaves}
$\mathcal{C}$를 범주라 하자. $F, G, H : \mathcal{C}^{opp} \to
\textit{Sets}$를 함자라 하고, $a : F \to G$와 $b : H \to G$를 함자의
변환이라 하자. 그러면 범주 $\textit{PSh}(\mathcal{C})$에서의 섬유곱
$F \times_{a, G, b} H$가 존재하고, $\mathcal{C}$의 임의의 대상 $X$에
대하여 공식
$$
(F \times_{a, G, b} H)(X) =
F(X) \times_{a_X, G(X), b_X} H(X)
$$
으로 주어진다.
\end{lemma}

\begin{proof}
생략한다.
\end{proof}












\noindent
특수한 경우로 사상 $a : F \to G$, 대상
$U \in \Ob(\mathcal{C})$, 원소 $\xi \in G(U)$가 주어졌다고 하자.
Yoneda Lemma \ref{categories-lemma-yoneda}에 따르면 이는 변환
$\xi : h_U \to G$를 준다. 이 경우 섬유곱은 규칙
$$
(h_U \times_{\xi, G, a} F)(X) =
\{ (f, \xi') \mid f : X \to U, \ \xi' \in F(X), \ G(f)(\xi) = a_X(\xi')\}
$$
으로 기술된다. $F$, $G$도 표현 가능하다면, 이는 섬유곱이 존재할 경우
그 섬유곱을 표현하는 함자이다. Section \ref{categories-section-fibre-products}를
보라. Definition \ref{categories-definition-representable-morphism}과의 유사성에서
표현 가능한 변환의 개념을 다음과 같이 정의하게 된다.

\begin{definition}
\label{categories-definition-representable-map-presheaves}
$\mathcal{C}$를 범주라 하고 $F, G : \mathcal{C}^{opp} \to
\textit{Sets}$를 함자라 하자. 모든 $U \in \Ob(\mathcal{C})$와 모든
$\xi \in G(U)$에 대하여 함자 $h_U \times_{\xi, G, a} F$가 표현
가능하면 사상 $a : F \to G$를 {\it 표현 가능}하다고 하며, 또는
{\it $F$가 $G$ 위에서 상대적으로 표현 가능하다}고 한다.
\end{definition}

\begin{lemma}
\label{categories-lemma-representable-over-representable}
$\mathcal{C}$를 범주라 하고, $a : F \to G$를 $\mathcal{C}$에서
$\textit{Sets}$로 가는 반변 함자들의 사상이라 하자. $a$가 표현 가능하고
$G$가 표현 가능한 함자이면 $F$도 표현 가능하다.
\end{lemma}

\begin{proof}
생략한다.
\end{proof}

\begin{lemma}
\label{categories-lemma-representable-diagonal}
$\mathcal{C}$를 범주라 하고
$F : \mathcal{C}^{opp} \to \textit{Sets}$를 함자라 하자.
$\mathcal{C}$가 대상쌍의 곱과 섬유곱을 갖는다고 가정하자.
다음 조건들은 서로 동치이다.
\begin{enumerate}
\item 대각 사상 $\Delta : F \to F \times F$가 표현 가능하다.
\item $\mathcal{C}$의 모든 $U$와 임의의 $\xi \in F(U)$에 대하여 사상
$\xi : h_U \to F$가 표현 가능하다.
\item $\mathcal{C}$의 모든 쌍 $U, V$와 임의의 $\xi \in F(U)$,
$\xi' \in F(V)$에 대하여 섬유곱
$h_U \times_{\xi, F, \xi'} h_V$가 표현 가능하다.
\end{enumerate}
\end{lemma}

\begin{proof}
계속해서 요네다 보조정리를 사용하여 $F(U)$를 함자의 변환
$h_U \to F$와 동일시한다.

\medskip\noindent
(2)와 (3)의 동치. $U, \xi, V, \xi'$가 (3)에서와 같다고 하자.
(2)와 (3)은 모두 정확히 $h_U \times_{\xi, F, \xi'} h_V$가 표현
가능하다고 말한다. 유일한 차이는 (3)이 $U$와 $V$에 대칭이지만 (2)는
그렇지 않다는 것이다.

\medskip\noindent
조건 (1)을 가정하자. $U, \xi, V, \xi'$가 (3)에서와 같다고 하자.
$h_U \times h_V = h_{U \times V}$는 표현 가능함에 유의하라.
곱 $\xi \times \xi' : h_U \times h_V \to F \times F$에 대응하는 사상을
$\eta : h_{U \times V} \to F \times F$로 표기하자. 가정에 따라 섬유곱
$F \times_{\Delta, F \times F, \eta} h_{U \times V}$는 표현 가능하다.
이는 $W \in \Ob(\mathcal{C})$와 사상 $W \to U$, $W \to V$,
$h_W \to F$가 존재하여 도표
$$
\xymatrix{
h_W \ar[d] \ar[r] & h_U \times h_V \ar[d]^{\xi \times \xi'} \\
F \ar[r] & F \times F
}
$$
가 카르테시안이라는 뜻이다. Lemma
\ref{categories-lemma-fibre-product-presheaves}의 섬유곱에 대한 명시적 기술을
사용하면, 이것이 원하는 대로
$h_W = h_U \times_{\xi, F, \xi'} h_V$를 함의함을 알 수 있다.

\medskip\noindent
서로 동치인 조건 (2)와 (3)을 가정하자. $U$를 $\mathcal{C}$의 대상으로
하고 $(\xi, \xi') \in (F \times F)(U)$라 하자. (3)에 따라 섬유곱
$h_U \times_{\xi, F, \xi'} h_U$는 표현 가능하다. 대상 $W$와 동형
$h_W \to h_U \times_{\xi, F, \xi'} h_U$를 택하자. 두 사영
$\text{pr}_i : h_U \times_{\xi, F, \xi'} h_U \to h_U$는 요네다에 의해
사상 $p_i : W \to U$에 대응한다. 이제
$W' = W \times_{(p_1, p_2), U \times U} U$라 하자. 등식
$$
h_{W'} = h_W \times_{h_U \times h_U} h_U
= (h_U \times_{\xi, F, \xi'} h_U) \times_{h_U \times h_U} h_U
= F \times_{F \times F} h_U.
$$
때문에 $W'$가 $F \times_{\Delta, F \times F} h_U$를 표현한다는 것은
형식적으로 확인된다. 따라서 $\Delta$는 표현 가능하고 증명이 끝난다.
\end{proof}









\section{밂}
\label{categories-section-pushouts}

\noindent
섬유곱의 쌍대 개념은 밂이다.

\begin{definition}
\label{categories-definition-pushouts}
$x, y, z\in \Ob(\mathcal{C})$라 하고,
$f\in \Mor_\mathcal{C}(y, x)$ 및
$g\in \Mor_{\mathcal C}(y, z)$라 하자. $f$와 $g$의 {\it 밂}은 대상
$x\amalg_y z\in \Ob(\mathcal{C})$와 사상
$p\in \Mor_{\mathcal C}(x, x\amalg_y z)$ 및
$q\in\Mor_{\mathcal C}(z, x\amalg_y z)$로 이루어지고 도표
$$
\xymatrix{
y \ar[r]_g \ar[d]_f & z \ar[d]^q \\
x \ar[r]^p & x\amalg_y z
}
$$
를 가환하게 하며 다음 보편 성질을 만족하는 것이다. 임의의
$w\in \Ob(\mathcal{C})$와, $\alpha \circ f = \beta \circ g$를 만족하는
사상 $\alpha \in \Mor_{\mathcal C}(x, w)$ 및
$\beta \in \Mor_\mathcal{C}(z, w)$에 대하여 도표
$$
\xymatrix{
y \ar[r]_g \ar[d]_f & z \ar[d]^q \ar[rrdd]^\beta & & \\
x \ar[r]^p \ar[rrrd]^\alpha & x \amalg_y z \ar@{-->}[rrd]^\gamma & & \\
& & & w
}
$$
를 가환하게 하는 유일한
$\gamma\in \Mor_{\mathcal C}(x\amalg_y z, w)$가 존재한다.
\end{definition}

\noindent
삼중항 $(x\amalg_y z, p, q)$가 존재한다면 유일한 동형을 제외하고
유일함을 직접적인 논증으로 증명하는 것은 가능하고 간단하다. 다른 방법은
밂을 반대범주에서의 섬유곱으로 생각하여 Section
\ref{categories-section-fibre-products}의 논의에서 이 유일성을 바로 얻는 것이다.

\begin{definition}
\label{categories-definition-cocartesian}
범주의 가환 도표
$$
\xymatrix{
y \ar[r] \ar[d] & z \ar[d] \\
x \ar[r] & w
}
$$
에서 $w$와 사상 $x \to w$, $z \to w$가 사상 $y \to x$와
$y \to z$의 밂을 이루면 이 도표를 {\it 쌍대카르테시안}이라고 한다.
\end{definition}


\section{등화자}
\label{categories-section-equalizers}

\begin{definition}
\label{categories-definition-equalizers}
$X$, $Y$가 범주 $\mathcal{C}$의 대상이고 $a, b : X \to Y$가
사상이라고 하자. $a \circ e = b \circ e$이고 쌍 $(Z, e)$가 다음 보편
성질을 만족하면 사상 $e : Z \to X$를 쌍 $(a, b)$의 {\it 등화자}라고
한다. $\mathcal{C}$의 모든 사상 $t : W \to X$가
$a \circ t = b \circ t$를 만족할 때, $t = e \circ s$를 만족하는
유일한 사상 $s : W \to Z$가 존재한다.
\end{definition}

\noindent
위의 섬유곱과 마찬가지로 등화자가 존재하면 유일한 동형을 제외하고
유일하다. 이 정의는 사상이 $2$개보다 많은 경우로 곧바로 일반화된다.

\section{쌍대등화자}
\label{categories-section-coequalizers}

\begin{definition}
\label{categories-definition-coequalizers}
$X$, $Y$가 범주 $\mathcal{C}$의 대상이고 $a, b : X \to Y$가
사상이라고 하자. $c \circ a = c \circ b$이고 쌍 $(Z, c)$가 다음 보편
성질을 만족하면 사상 $c : Y \to Z$를 쌍 $(a, b)$의 {\it 쌍대등화자}라고
한다. $\mathcal{C}$의 모든 사상 $t : Y \to W$가
$t \circ a = t \circ b$를 만족할 때, $t = s \circ c$를 만족하는
유일한 사상 $s : Z \to W$가 존재한다.
\end{definition}

\noindent
위의 밂과 마찬가지로 쌍대등화자가 존재하면 유일한 동형을 제외하고
유일하며, 이는 반대범주를 생각하면 등화자의 유일성에서 따른다.
이 정의는 사상이 $2$개보다 많은 경우로 곧바로 일반화된다.

\section{초기 대상과 끝 대상}
\label{categories-section-initial-final}

\begin{definition}
\label{categories-definition-initial-final}
$\mathcal{C}$를 범주라 하자.
\begin{enumerate}
\item 범주 $\mathcal{C}$의 대상 $x$에서 $\mathcal{C}$의 모든 대상
$y$로 가는 사상 $x \to y$가 정확히 하나씩 존재하면 그 대상을
{\it 초기} 대상이라고 한다.
\item 범주 $\mathcal{C}$의 대상 $x$에 대하여 $\mathcal{C}$의 모든 대상
$y$에서 그 대상으로 가는 사상 $y \to x$가 정확히 하나씩 존재하면 그 대상을
{\it 종말} 대상이라고 한다.
\end{enumerate}
\end{definition}

\noindent
집합들의 범주에서 공집합 $\emptyset$은 초기 대상이며, 사실 유일한 초기
대상이다. 또한 임의의 {\it 한원소 집합}, 즉 원소가 하나인 집합은 종말
대상이다(따라서 끝 대상 자체는 유일하지 않다).




\section{단사사상과 전사사상}
\label{categories-section-mono-epi}

\begin{definition}
\label{categories-definition-mono-epi}
$\mathcal{C}$를 범주라 하고 $f : X \to Y$를 $\mathcal{C}$의 사상이라
하자.
\begin{enumerate}
\item 모든 대상 $W$와, $f \circ a = f \circ b$를 만족하는 모든 사상쌍
$a, b : W \to X$에 대하여 $a = b$이면 $f$를 {\it 단사사상}이라고 한다.
\item 모든 대상 $W$와, $a \circ f = b \circ f$를 만족하는 모든 사상쌍
$a, b : Y \to W$에 대하여 $a = b$이면 $f$를 {\it 전사사상}이라고 한다.
\end{enumerate}
\end{definition}

\begin{example}
\label{categories-example-mono-epi-sets}
집합들의 범주에서 단사사상은 단사 함수에, 전사사상은 전사 함수에 대응한다.
\end{example}

\begin{lemma}
\label{categories-lemma-characterize-mono-epi}
$\mathcal{C}$를 범주라 하고 $f : X \to Y$를 $\mathcal{C}$의 사상이라
하자. 그러면 다음이 성립한다.
\begin{enumerate}
\item $f$가 단사사상일 필요충분조건은 $X$가 섬유곱
$X \times_Y X$인 것이다.
\item $f$가 전사사상일 필요충분조건은 $Y$가 밂
$Y \amalg_X Y$인 것이다.
\end{enumerate}
\end{lemma}

\begin{proof}
$f$가 단사사상이라고 하자. $W$를 $\mathcal C$의 대상이라 하고
$\alpha, \beta \in \Mor_\mathcal C(W,X)$가
$f\circ \alpha = f\circ \beta$를 만족한다고 하자. $f$가 단사사상이므로
$\alpha = \beta$이다. 또한 가환 도표
$$
\xymatrix{
X \ar[r]^{\text{id}_X} \ar[d]_{\text{id}_X} & X \ar[d]^{f} \\
X \ar[r]^{f} & Y
}
$$
가 있으며, 이는 $\gamma := \alpha = \beta$에 대하여 보편 성질을
확인한다. 따라서 $X$는 실제로 섬유곱 $X\times_Y X$이다.

\medskip\noindent
$X \times_Y X \cong X $라고 하자. 도표
$$
\xymatrix{
X \ar[r]^{\text{id}_X} \ar[d]_{\text{id}_X} & X \ar[d]^{f} \\
X \ar[r]^{f} & Y }
$$
는 가환한다. 또한 $W \in \Ob(\mathcal C)$이고
$\alpha, \beta : W \to X$가 $f \circ \alpha = f \circ \beta$를
만족하면,
$$
\gamma = \text{id}_X\circ\gamma = \alpha = \beta
$$
를 만족하는 유일한 $\gamma$가 있다. 따라서 $\alpha = \beta$이다.

\medskip\noindent
두 번째 항목의 증명도 완전히 같으며, 다만 밂
$Y\amalg_X Y = Y$를 사용한다.
\end{proof}





\section{극한과 여극한}
\label{categories-section-limits}

\noindent
$\mathcal{C}$를 범주라 하자. $\mathcal{C}$에서의 {\it 도표}는 단순히
함자 $M : \mathcal{I} \to \mathcal{C}$이다. $\mathcal{I}$를
{\it 지표 범주}라고 하며, 또는 $M$을 $\mathcal{I}$-도표라고 한다.
$\mathcal{I}$의 대상 $i$의 상은 $M_i$로 표기한다. 따라서
$\mathcal{I}$의 사상 $\phi : i \to i'$에 대하여
$M(\phi) : M_i \to M_{i'}$이다.

\begin{definition}
\label{categories-definition-limit}
범주 $\mathcal{C}$에서 $\mathcal{I}$-도표 $M$의 {\it 극한}은
$\mathcal{C}$의 대상 $\lim_\mathcal{I} M$과 사상
$p_i : \lim_\mathcal{I} M \to M_i$로 주어지며 다음을 만족한다.
\begin{enumerate}
\item $\mathcal{I}$의 사상 $\phi : i \to i'$에 대하여
$p_{i'} = M(\phi) \circ p_i$이다.
\item $\mathcal{C}$의 임의의 대상 $W$와
$i \in \Ob(\mathcal{I})$로 지표화된 임의의 사상족
$q_i : W \to M_i$가 주어지고, $\mathcal{I}$의 모든
$\phi : i \to i'$에 대하여 $q_{i'} = M(\phi) \circ q_i$이면,
$\mathcal{I}$의 모든 대상 $i$에 대하여 $q_i = p_i \circ q$를 만족하는
유일한 사상 $q : W \to \lim_\mathcal{I} M$이 존재한다.
\end{enumerate}
\end{definition}

\noindent
극한 $(\lim_\mathcal{I} M, (p_i)_{i\in \Ob(\mathcal{I})})$이 존재하면
정의의 유일성 조건에 의해 유일한 동형을 제외하고 유일하다. 대상쌍의 곱,
섬유곱, 등화자는 극한의 예이다. 공도표에 대한 극한은 $\mathcal{C}$의
끝 대상이다. 집합들의 범주에서는 모든 극한이 존재한다. 쌍대 개념은
여극한이다.

\begin{definition}
\label{categories-definition-colimit}
범주 $\mathcal{C}$에서 $\mathcal{I}$-도표 $M$의 {\it 여극한}은
$\mathcal{C}$의 대상 $\colim_\mathcal{I} M$과 사상
$s_i : M_i \to \colim_\mathcal{I} M$로 주어지며 다음을 만족한다.
\begin{enumerate}
\item $\mathcal{I}$의 사상 $\phi : i \to i'$에 대하여
$s_i = s_{i'} \circ M(\phi)$이다.
\item $\mathcal{C}$의 임의의 대상 $W$와
$i \in \Ob(\mathcal{I})$로 지표화된 임의의 사상족
$t_i : M_i \to W$가 주어지고, $\mathcal{I}$의 모든
$\phi : i \to i'$에 대하여 $t_i = t_{i'} \circ M(\phi)$이면,
$\mathcal{I}$의 모든 대상 $i$에 대하여 $t_i = t \circ s_i$를 만족하는
유일한 사상 $t : \colim_\mathcal{I} M \to W$가 존재한다.
\end{enumerate}
\end{definition}

\noindent
여극한 $(\colim_\mathcal{I} M, (s_i)_{i\in \Ob(\mathcal{I})})$이
존재하면 정의의 유일성 조건에 의해 유일한 동형을 제외하고 유일하다.
대상쌍의 쌍대곱, 밂, 쌍대등화자는 여극한의 예이다. 공도표에 대한 여극한은
$\mathcal{C}$의 초기 대상이다. 집합들의 범주에서는 모든 여극한이 존재한다.

\begin{remark}
\label{categories-remark-diagram-small}
(여)극한의 지표 범주는 대상들의 고유 모임을 갖는 것이 허용되지 않는다.
이 프로젝트에서는 Remark \ref{categories-remark-big-categories}에 열거된 범주들
중 하나일 수 없다는 뜻이다.
\end{remark}

\begin{remark}
\label{categories-remark-limit-colim}
흔히 $\mathcal{I}$로 지표화한 표기 대신
$\lim_i M_i$, $\colim_i M_i$,
$\lim_{i\in \mathcal{I}} M_i$, 또는
$\colim_{i\in \mathcal{I}} M_i$를 쓴다. 이 표기와 아래 Section
\ref{categories-section-limit-sets}의 집합의 극한 및 여극한에 대한 기술을 사용하면
다음과 같이 말할 수 있다. $M : \mathcal{I} \to \mathcal{C}$를
도표라 하자.
\begin{enumerate}
\item 대상 $\lim_i M_i$가 존재하면 다음 성질을 만족한다.
$$
\Mor_\mathcal{C}(W, \lim_i M_i)
=
\lim_i \Mor_\mathcal{C}(W, M_i)
$$
여기서 우변의 극한은 집합들의 범주에서 취한다.
\item 대상 $\colim_i M_i$가 존재하면 다음 성질을 만족한다.
$$
\Mor_\mathcal{C}(\colim_i M_i, W)
=
\lim_{i \in \mathcal{I}^{opp}} \Mor_\mathcal{C}(M_i, W)
$$
여기서 우변은 반대범주에 대한, 집합들의 범주에 값을 갖는 극한이다.
\end{enumerate}
요네다 보조정리와 그 쌍대에 의해 이 공식은 각각 극한과 여극한을 완전히
결정한다.
\end{remark}

\begin{remark}
\label{categories-remark-cones-and-cocones}
$M : \mathcal{I} \to \mathcal{C}$를 도표라 하자. 이 상황에서 $M$의
{\it 원뿔}은 대상 $W$와 사상족
$q_i : W \to M_i$, $i \in \Ob(\mathcal{I})$로 주어지며,
$\mathcal{I}$의 모든 사상 $\phi : i \to i'$에 대하여 도표
$$
\xymatrix{
& W \ar[dl]_{q_i} \ar[dr]^{q_{i'}} \\
M_i \ar[rr]^{M(\phi)} & & M_{i'}
}
$$
가 가환한다. 원뿔들의 모음은 자명한 사상 개념과 함께 범주를 이룬다.
명백히 $M$의 극한이 존재하면 그것은 원뿔들의 범주의 끝 대상이다.
쌍대적으로 $M$의 {\it 쌍대원뿔}은 대상 $W$와 사상족
$t_i : M_i \to W$로 주어지며, $\mathcal{I}$의 모든 사상
$\phi : i \to i'$에 대하여 도표
$$
\xymatrix{
M_i \ar[rr]^{M(\phi)} \ar[dr]_{t_i} & & M_{i'} \ar[dl]^{t_{i'}} \\
& W
}
$$
가 가환한다. 쌍대원뿔들의 모음은 자명한 사상 개념과 함께 범주를 이룬다.
위와 마찬가지로 $M$의 여극한이 존재할 필요충분조건은 쌍대원뿔들의 범주가
초기 대상을 갖는 것이다.
\end{remark}

\noindent
극한과 여극한의 개념을 응용하여 곱과 쌍대곱을 정의한다.

\begin{definition}
\label{categories-definition-product}
$I$를 집합이라 하고, 각 $i \in I$에 대하여 범주 $\mathcal{C}$의 대상
$M_i$가 주어졌다고 하자. {\it 곱} $\prod_{i\in I} M_i$는 정의상
$\lim_\mathcal{I} M$이다(존재하는 경우). 여기서 $\mathcal{I}$는
$I$의 원소들을 대상으로 하고 사상으로는 항등사상들만 갖는 범주이다.
\end{definition}

\noindent
중요한 특수한 경우는 $I = \emptyset$인 경우이며, 이때 곱은 범주의 종말
대상이다. 사상 $p_i : \prod M_i \to M_i$를 {\it 사영 사상}이라고 한다.

\begin{definition}
\label{categories-definition-coproduct}
$I$를 집합이라 하고, 각 $i \in I$에 대하여 범주 $\mathcal{C}$의 대상
$M_i$가 주어졌다고 하자. {\it 쌍대곱} $\coprod_{i\in I} M_i$는 정의상
$\colim_\mathcal{I} M$이다(존재하는 경우). 여기서 $\mathcal{I}$는
$I$의 원소들을 대상으로 하고 사상으로는 항등사상들만 갖는 범주이다.
\end{definition}

\noindent
중요한 특수한 경우는 $I = \emptyset$인 경우이며, 이때 쌍대곱은 범주의
초기 대상이다. 쌍대곱에는 사상 $M_i \to \coprod M_i$가 딸려 있음에
유의하라. 이들을 때때로 {\it 쌍대사영}이라고 한다.

\begin{lemma}
\label{categories-lemma-functorial-colimit}
$M : \mathcal{I} \to \mathcal{C}$와
$N : \mathcal{J} \to \mathcal{C}$를 여극한이 존재하는 도표라 하자.
$H : \mathcal{I} \to \mathcal{J}$를 함자라 하고
$t : M \to N \circ H$를 함자의 변환이라 하자. 그러면 유일한 사상
$$
\theta :
\colim_\mathcal{I} M
\longrightarrow
\colim_\mathcal{J} N
$$
가 존재하여 모든 도표
$$
\xymatrix{
M_i \ar[d]_{t_i} \ar[r]
&
\colim_\mathcal{I} M \ar[d]^{\theta}
\\
N_{H(i)} \ar[r]
&
\colim_\mathcal{J} N
}
$$
가 가환한다.
\end{lemma}

\begin{proof}
생략한다.
\end{proof}

\begin{lemma}
\label{categories-lemma-functorial-limit}
$M : \mathcal{I} \to \mathcal{C}$와
$N : \mathcal{J} \to \mathcal{C}$를 극한이 존재하는 도표라 하자.
$H : \mathcal{I} \to \mathcal{J}$를 함자라 하고
$t : N \circ H \to M$을 함자의 변환이라 하자. 그러면 유일한 사상
$$
\theta :
\lim_\mathcal{J} N
\longrightarrow
\lim_\mathcal{I} M
$$
가 존재하여 모든 도표
$$
\xymatrix{
\lim_\mathcal{J} N \ar[d]^{\theta} \ar[r]
&
N_{H(i)} \ar[d]_{t_i}
\\
\lim_\mathcal{I} M \ar[r]
&
M_i
}
$$
가 가환한다.
\end{lemma}

\begin{proof}
생략한다.
\end{proof}


\begin{lemma}
\label{categories-lemma-colimits-commute}
$\mathcal{I}$, $\mathcal{J}$를 지표 범주라 하고
$M : \mathcal{I} \times \mathcal{J} \to \mathcal{C}$를 함자라 하자.
모든 $i$에 대하여 $M_{i, \infty} = \colim_j M_{i,j}$가 존재한다고
가정하자. 그러면 이렇게 얻은 함자
$M_{-, \infty} : \mathcal{I} \to \mathcal{C}$가 여극한을 가질
필요충분조건은 $M$이 여극한을 갖는 것이며, 이때 두 여극한은 일치한다.
특히 표시된 모든 여극한이 존재하면
$$
\colim_i \colim_j M_{i, j}
=
\colim_{i, j} M_{i, j}
=
\colim_j \colim_i M_{i, j}
$$
이다. 극한에 대해서도 마찬가지이다.
\end{lemma}

\begin{proof}
생략한다.
\end{proof}

\begin{lemma}
\label{categories-lemma-limits-products-equalizers}
$M : \mathcal{I} \to \mathcal{C}$를 도표라 하자.
$I = \Ob(\mathcal{I})$, $A = \text{Arrows}(\mathcal{I})$로 쓰고,
출발점과 도착점 사상을 $s, t : A \to I$로 표기하자.
$\prod_{i \in I} M_i$와 $\prod_{a \in A} M_{t(a)}$가 존재한다고 하자.
또한
$$
\xymatrix{
\prod_{i \in I} M_i
\ar@<1ex>[r]^\phi \ar@<-1ex>[r]_\psi
&
\prod_{a \in A} M_{t(a)}
}
$$
의 등화자가 존재한다고 하자. 여기서 사상들은 성분별로
$p_a \circ \psi = M(a) \circ p_{s(a)}$ 및
$p_a \circ \phi = p_{t(a)}$에 의해 결정된다. 그러면 이 등화자가
도표의 극한이다.
\end{lemma}

\begin{proof}
생략한다.
\end{proof}


\begin{lemma}
\label{categories-lemma-colimits-coproducts-coequalizers}
\begin{slogan}
모든 쌍대곱과 쌍대등화자가 존재하면 모든 여극한이 존재한다.
\end{slogan}
$M : \mathcal{I} \to \mathcal{C}$를 도표라 하자.
$I = \Ob(\mathcal{I})$, $A = \text{Arrows}(\mathcal{I})$로 쓰고,
출발점과 도착점 사상을 $s, t : A \to I$로 표기하자.
$\coprod_{i \in I} M_i$와 $\coprod_{a \in A} M_{s(a)}$가 존재한다고 하자.
또한
$$
\xymatrix{
\coprod_{a \in A} M_{s(a)}
\ar@<1ex>[r]^\phi \ar@<-1ex>[r]_\psi
&
\coprod_{i \in I} M_i
}
$$
의 쌍대등화자가 존재한다고 하자. 여기서 사상들은 다음과 같이 성분별로
결정된다. 성분 $M_{s(a)}$는 $\psi$에 의해 사상 $M(a)$를 거쳐 성분
$M_{t(a)}$로 간다. 성분 $M_{s(a)}$는 $\phi$에 의해 항등사상을 거쳐
성분 $M_{s(a)}$로 간다. 그러면 이 쌍대등화자가 도표의 여극한이다.
\end{lemma}

\begin{proof}
생략한다.
\end{proof}







\section{집합들의 범주에서의 극한과 여극한}
\label{categories-section-limit-sets}

\noindent
$\textit{Sets}$에서는 극한과 여극한이 모두 존재할 뿐 아니라 기술하기도
쉽다. 실제로 $M : \mathcal{I} \to \textit{Sets}$,
$i \mapsto M_i$를 집합들의 도표라 하고 $I = \Ob(\mathcal{I})$로
표기하자. 극한은 다음과 같이 기술된다.
$$
\lim_\mathcal{I} M
=
\{
(m_i)_{i\in I} \in \prod\nolimits_{i\in I} M_i
\mid
\forall \phi : i \to i' \text{ in }\mathcal{I},
M(\phi)(m_i) = m_{i'}
\}.
$$
따라서 극한의 원소를 모든 집합 $M_i$의 원소들로 이루어진 양립 가능한
계로 생각한다.

\medskip\noindent
한편 여극한은
$$
\colim_\mathcal{I} M
=
(\coprod\nolimits_{i\in I} M_i)/\sim
$$
이다. 여기서 동치관계 $\sim$는 $m_i \in M_i$,
$m_{i'} \in M_{i'}$이고 어떤 $\phi : i \to i'$에 대하여
$M(\phi)(m_i) = m_{i'}$일 때 $m_i \sim m_{i'}$로 놓아 생성되는
동치관계이다. 다시 말해 $m_i \in M_i$와 $m_{i'} \in M_{i'}$가
동치일 필요충분조건은 $\mathcal{I}$ 안에 사상의 사슬
$$
\xymatrix{
&
i_1 \ar[ld] \ar[rd] & &
i_3 \ar[ld] & &
i_{2n-1} \ar[rd] & \\
i = i_0 & &
i_2 & &
\ldots & &
i_{2n} = i'
}
$$
이 존재하고, 원소 $m_{i_j} \in M_{i_j}$들이 위 $\mathcal{I}$의
사상들에서 유도되는 사상 $M_{i_{2k-1}} \to M_{i_{2k-2}}$ 및
$M_{i_{2k-1}} \to M_{i_{2k}}$ 아래에서 서로에게 대응하는 것이다.

\medskip\noindent
이것은 다루기에 그다지 편한 종류의 대상이 아니다. 그러나 도표가 필터이면
훨씬 쉽게 기술할 수 있다. Section \ref{categories-section-directed-colimits}에서
이를 설명한다.



\section{연결 극한}
\label{categories-section-connected-limits}

\noindent
지표 범주가 연결이면 (여)극한을 연결이라고 한다.

\begin{definition}
\label{categories-definition-category-connected}
$x \sim y \Leftrightarrow \Mor_\mathcal{I}(x, y) \not = \emptyset$로부터
생성되는 동치관계가 동치류를 정확히 하나 가지면 범주 $\mathcal{I}$를
{\it 연결}이라고 한다.
\end{definition}

\noindent
여기서는 연결 공간이 공집합이 아니라는 Topology, Definition
\ref{topology-definition-connected-components}의 관례를 따른다.
다음 결과는 어떤 느슨한 의미에서 연결 극한을 특징짓는다.

\begin{lemma}
\label{categories-lemma-connected-limit-over-X}
$\mathcal{C}$를 범주라 하고 $X$를 $\mathcal{C}$의 대상이라 하자.
$M : \mathcal{I} \to \mathcal{C}/X$를 $X$ 위의 대상들의 범주에서의
도표라 하자. 지표 범주 $\mathcal{I}$가 연결이고 $M$의 극한이
$\mathcal{C}/X$에 존재하면, 합성
$\mathcal{I} \to \mathcal{C}/X \to \mathcal{C}$의 극한도 존재하고
같은 대상이다.
\end{lemma}

\begin{proof}
$L \to X$를 $\mathcal{C}/X$에서의 극한을 표현하는 대상이라 하자.
함자
$$
W \longmapsto \lim_i \Mor_\mathcal{C}(W, M_i).
$$
를 생각하자. $(\varphi_i)$를 우변 집합의 원소라 하자. 각 $M_i$에는
사상 $s_i : M_i \to X$가 딸려 있으므로 사상
$f_i = s_i \circ \varphi_i : W \to X$를 얻는다. 그런데
$\mathcal{I}$가 연결이므로 모든 $f_i$는 같다. $\mathcal{I}$는
공집합이 아니므로 적어도 하나의 $f_i$가 있다. 따라서 이 공통값
$W \to X$는 $W$에 $\mathcal{C}/X$의 대상 구조를 주며, $(\varphi_i)$는
$\lim_i \Mor_{\mathcal{C}/X}(W, M_i)$의 원소를 정의한다. 그러므로 원하는
대로 $\varphi_i$가 $\phi$와 $L \to M_i$의 합성이 되게 하는 유일한
사상 $\phi : W \to L$을 얻는다.
\end{proof}

\begin{lemma}
\label{categories-lemma-connected-colimit-under-X}
$\mathcal{C}$를 범주라 하고 $X$를 $\mathcal{C}$의 대상이라 하자.
$M : \mathcal{I} \to X/\mathcal{C}$를 $X$ 아래의 대상들의 범주에서의
도표라 하자. 지표 범주 $\mathcal{I}$가 연결이고 $M$의 여극한이
$X/\mathcal{C}$에 존재하면, 합성
$\mathcal{I} \to X/\mathcal{C} \to \mathcal{C}$의 여극한도 존재하고
같은 대상이다.
\end{lemma}

\begin{proof}
생략한다. 힌트: 이 보조정리는 Lemma
\ref{categories-lemma-connected-limit-over-X}의 쌍대이다.
\end{proof}




\section{공종 범주와 초기 범주}
\label{categories-section-cofinal}

\noindent
문헌에서는 다음 정의의 공종이라는 말 대신 ``종말''이라는 말을
사용하기도 한다.

\begin{definition}
\label{categories-definition-cofinal}
$H : \mathcal{I} \to \mathcal{J}$를 범주 사이의 함자라 하자.
$\mathcal{I}$가 $\mathcal{J}$에서 {\it 공종}이라고 하거나
$H$가 {\it 공종}이라고 함은 다음 두 조건이 성립한다는 뜻이다.
\begin{enumerate}
\item 모든 $y \in \Ob(\mathcal{J})$에 대하여 어떤
$x \in \Ob(\mathcal{I})$와 사상 $y \to H(x)$가 존재하고,
\item $y \in \Ob(\mathcal{J})$, $x, x' \in \Ob(\mathcal{I})$ 및
사상 $y \to H(x)$와 $y \to H(x')$가 주어지면, $\mathcal{I}$ 안의
사상들의 열
$$
x = x_0 \leftarrow x_1 \rightarrow x_2 \leftarrow x_3 \rightarrow \ldots
\rightarrow x_{2n} = x'
$$
과 $\mathcal{J}$ 안의 사상들 $y \to H(x_i)$가 존재하여,
$y \to H(x_0)$와 $y \to H(x_{2n})$는 주어진 사상들이고 도표
$$
\xymatrix{
& y \ar[ld] \ar[d] \ar[rd] \\
H(x_{2k}) & H(x_{2k + 1}) \ar[l] \ar[r] & H(x_{2k + 2})
}
$$
가 $k = 0, \ldots, n - 1$에 대하여 가환한다.
\end{enumerate}
달리 말해, $\mathcal{J}$의 대상 $y$를 고정하고 $x$가
$\mathcal{I}$의 대상이며 $b : y \to H(x)$가 사상인 순서쌍
$(x, b)$들의 집합 $S$를 생각하자. $b' = H(a) \circ b$를 만족하는
사상 $a : x \to x'$가 존재할 때 $(x, b) \sim (x', b')$라 하여
생성되는 $S$ 위의 동치관계를 생각하자. 그러면 $S$는 정확히 하나의
동치류로 이루어져야 한다.
\end{definition}

\begin{lemma}
\label{categories-lemma-cofinal}
$H : \mathcal{I} \to \mathcal{J}$를 범주의 함자라 하자.
$\mathcal{I}$가 $\mathcal{J}$에서 공종이라고 가정하자. 그러면 모든 도표
$M : \mathcal{J} \to \mathcal{C}$에 대하여 어느 한쪽이 존재하면
표준 동형
$$
\colim_\mathcal{I} M \circ H
=
\colim_\mathcal{J} M
$$
이 있다.
\end{lemma}

\begin{proof}
생략한다.
\end{proof}

\begin{definition}
\label{categories-definition-initial}
$H : \mathcal{I} \to \mathcal{J}$를 범주 사이의 함자라 하자.
$\mathcal{I}$가 $\mathcal{J}$에서 {\it 초기}라고 하거나
$H$가 {\it 초기}라고 함은 다음 두 조건이 성립한다는 뜻이다.
\begin{enumerate}
\item 모든 $y \in \Ob(\mathcal{J})$에 대하여 어떤
$x \in \Ob(\mathcal{I})$와 사상 $H(x) \to y$가 존재하고,
\item 임의의 $y \in \Ob(\mathcal{J})$, $x , x' \in \Ob(\mathcal{I})$ 및
$\mathcal{J}$ 안의 사상 $H(x) \to y$, $H(x') \to y$에 대하여,
$\mathcal{I}$ 안의 사상들의 열
$$
x = x_0 \leftarrow x_1 \rightarrow x_2 \leftarrow x_3 \rightarrow \ldots
\rightarrow x_{2n} = x'
$$
과 $\mathcal{J}$ 안의 사상들 $H(x_i) \to y$가 존재하여 도표
$$
\xymatrix{
H(x_{2k}) \ar[rd] &
H(x_{2k + 1}) \ar[l] \ar[r] \ar[d] &
H(x_{2k + 2}) \ar[ld] \\
& y
}
$$
가 $k = 0, \ldots, n - 1$에 대하여 가환한다.
\end{enumerate}
\end{definition}

\noindent
이는 ``공종'' 함자의 쌍대 개념일 뿐이다.

\begin{lemma}
\label{categories-lemma-initial}
$H : \mathcal{I} \to \mathcal{J}$를 범주의 함자라 하자.
$\mathcal{I}$가 $\mathcal{J}$에서 초기라고 가정하자.
그러면 모든 도표 $M : \mathcal{J} \to \mathcal{C}$에 대하여
어느 한쪽이 존재하면 표준 동형
$$
\lim_\mathcal{I} M \circ H = \lim_\mathcal{J} M
$$
이 있다.
\end{lemma}

\begin{proof}
생략한다.
\end{proof}

\begin{lemma}
\label{categories-lemma-colimit-constant-connected-fibers}
$F : \mathcal{I} \to \mathcal{I}'$를 함자라 하자. 다음을 가정하자.
\begin{enumerate}
\item $\mathcal{I}'$ 위의 $\mathcal{I}$의 섬유 범주들(Definition
\ref{categories-definition-fibre-category} 참조)은 모두 연결이고,
\item $\mathcal{I}'$의 모든 사상 $\alpha' : x' \to y'$에 대하여
$F(\alpha) = \alpha'$를 만족하는 $\mathcal{I}$의 사상
$\alpha : x \to y$가 존재한다.
\end{enumerate}
그러면 모든 도표 $M : \mathcal{I}' \to \mathcal{C}$에 대하여
여극한 $\colim_\mathcal{I} M \circ F$가 존재할 필요충분조건은
$\colim_{\mathcal{I}'} M$이 존재하는 것이며, 이 경우 두 여극한은 같다.
\end{lemma}

\begin{proof}
$\mathcal{I}$가 $\mathcal{I}'$에서 공종임을 보이고 Lemma
\ref{categories-lemma-cofinal}를 적용하여 증명할 수 있다. 다음과 같이 직접
증명할 수도 있다. $\mathcal{C}$의 임의의 대상 $T$에 대하여
$$
\lim_{\mathcal{I}^{opp}} \Mor_\mathcal{C}(M_{F(i)}, T)
=
\lim_{(\mathcal{I}')^{opp}} \Mor_\mathcal{C}(M_{i'}, T)
$$
임을 보이면 충분하다. $(g_{i'})_{i' \in \Ob(\mathcal{I}')}$가 우변의
원소이면 $f_i = g_{F(i)}$로 놓아 좌변의 원소
$(f_i)_{i \in \Ob(\mathcal{I})}$를 얻는다. 거꾸로,
$(f_i)_{i \in \Ob(\mathcal{I})}$를 좌변의 원소라 하자. 각 (연결인)
섬유 범주 $\mathcal{I}_{i'}$ 위에서 함자 $M \circ F$는 값이
$M_{i'}$인 상수 함자임에 유의하자. 따라서 $F(i) = i'$인
$i \in \Ob(\mathcal{I})$에 대한 사상 $f_i$들은 모두 같으며,
잘 정의된 사상 $g_{i'} : M_{i'} \to T$를 결정한다. 가정 (2)에 의해
모음 $(g_{i'})_{i' \in \Ob(\mathcal{I}')}$는 우변의 원소를 정의한다.
\end{proof}

\begin{lemma}
\label{categories-lemma-product-with-connected}
$\mathcal{I}$와 $\mathcal{J}$를 범주라 하고
$p : \mathcal{I} \times \mathcal{J} \to \mathcal{J}$를 사영이라 하자.
$\mathcal{I}$가 연결이면 도표 $M : \mathcal{J} \to \mathcal{C}$에 대하여
여극한 $\colim_\mathcal{J} M$이 존재할 필요충분조건은
$\colim_{\mathcal{I} \times \mathcal{J}} M \circ p$가 존재하는 것이며,
이 경우 두 여극한은 같다.
\end{lemma}

\begin{proof}
이는 Lemma \ref{categories-lemma-colimit-constant-connected-fibers}의 특수한 경우이다.
\end{proof}






\section{유한 극한과 여극한}
\label{categories-section-finite-limits}

\noindent
지표 범주가 유한한, 즉 대상과 사상의 수가 모두 유한한 (여)극한을
{\it 유한} (여)극한이라고 한다. 지표 범주가 공집합이 아니면
(여)극한을 {\it 비공집합}이라고 한다. 지표 범주가 연결이면
(여)극한을 {\it 연결}이라고 한다. Definition
\ref{categories-definition-category-connected}을 참조하라. 유한 지표 범주는
어떤 의미에서 ``충분히 많이'' 있음이 밝혀진다.

\begin{lemma}
\label{categories-lemma-finite-diagram-category}
$\mathcal{I}$를 다음 조건을 만족하는 범주라 하자.
\begin{enumerate}
\item $\Ob(\mathcal{I})$는 유한하고,
\item 유한 개의 사상 $f_1, \ldots, f_m \in \text{Arrows}(\mathcal{I})$가
존재하여 $\mathcal{I}$의 모든 사상은
$f_{j_1} \circ f_{j_2} \circ \ldots \circ f_{j_k}$ 꼴의 합성이다.
\end{enumerate}
그러면 다음을 만족하는 함자 $F : \mathcal{J} \to \mathcal{I}$가
존재한다.
\begin{enumerate}
\item[(a)] $\mathcal{J}$는 유한 범주이고,
\item[(b)] 임의의 도표 $M : \mathcal{I} \to \mathcal{C}$에 대하여
$\mathcal{I}$ 위에서 $M$의 (여)극한이 존재할 필요충분조건은
$\mathcal{J}$ 위에서 $M \circ F$의 (여)극한이 존재하는 것이며,
이 경우 두 (여)극한은 표준적으로 동형이다.
\end{enumerate}
더욱이 $\mathcal{J}$가 연결인 것(각각 공집합이 아닌 것)과
$\mathcal{I}$가 그러한 것은 서로 동치이다.
\end{lemma}

\begin{proof}
$\Ob(\mathcal{I}) = \{x_1, \ldots, x_n\}$이라 하자.
$f_j : x_{s(j)} \to x_{t(j)}$를 만족하는 함수들을
$s, t : \{1, \ldots, m\} \to \{1, \ldots, n\}$로 나타내자.
$\Ob(\mathcal{J}) = \{y_1, \ldots, y_n, z_1, \ldots, z_n\}$으로 둔다.
항등사상 외에 사상들
$g_j : y_{s(j)} \to z_{t(j)}$, $j = 1, \ldots, m$과 사상들
$h_i : y_i \to z_i$, $i = 1, \ldots, n$을 도입한다.
$\mathcal{J}$의 항등사상이 아닌 모든 사상은 어떤 $y$에서 어떤
$z$로 가므로 정의할 합성도, 확인할 결합법칙도 없다.
$F(y_i) = F(z_i) = x_i$로 둔다. $F(g_j) = f_j$ 및
$F(h_i) = \text{id}_{x_i}$로 둔다. $F$가 함자임은 명백하다.
$\mathcal{J}$가 유한함도 명백하다. $\mathcal{J}$가 연결인 것,
각각 공집합이 아닌 것과 $\mathcal{I}$가 그러한 것은 서로 동치임도
명백하다.

\medskip\noindent
$M : \mathcal{I} \to \mathcal{C}$를 도표라 하자. $\mathcal{C}$의 대상
$W$와 Definition \ref{categories-definition-limit}에서와 같은 사상들
$q_i : W \to M(x_i)$를 생각하자. 그러면
$q_i : W \to M(F(y_i)) = M(F(z_i)) = M(x_i)$로 놓아 도표
$M \circ F$에 대한 Definition \ref{categories-definition-limit}에서와 같은
사상들의 족을 얻는다. 거꾸로 도표 $M \circ F$에 대한
Definition \ref{categories-definition-limit}에서와 같은 사상들
$qy_i : W \to M(F(y_i))$와 $qz_i : W \to M(F(z_i))$가 주어졌다고 하자.
다음 등식으로부터
$$
M(F(h_i)) = \text{id} : M(F(y_i)) = M(x_i) \longrightarrow M(x_i) = M(F(z_i))
$$
모든 $i$에 대하여 $qy_i = qz_i$임을 안다. 이 공통값을 $q_i$로
놓자. $q_{s(j)} = qy_{s(j)}$와 $q_{t(j)} = qz_{t(j)}$가 사상
$M(f_j)$와 양립한다는 사실은 족 $q_i$가 $\mathcal{I}$의 모든
사상과 양립함을 보장한다. 실제로 가정에 의해 그러한 모든 사상은
사상들 $f_j$의 합성이기 때문이다. 따라서 표준 전단사
$$
\lim_{B \in \Ob(\mathcal{J})} \Mor_\mathcal{C}(W, M(F(B)))
=
\lim_{A \in \Ob(\mathcal{I})} \Mor_\mathcal{C}(W, M(A))
$$
를 얻었으며, 이는 보조정리의 극한에 관한 명제를 함의한다.
여극한에 관한 명제도 같은 방식으로 증명된다(증명 생략).
\end{proof}

\begin{lemma}
\label{categories-lemma-fibre-products-equalizers-exist}
$\mathcal{C}$를 범주라 하자. 다음은 서로 동치이다.
\begin{enumerate}
\item $\mathcal{C}$에 연결 유한 극한이 존재한다.
\item $\mathcal{C}$에 등화자와 섬유곱이 존재한다.
\end{enumerate}
\end{lemma}

\begin{proof}
등화자와 섬유곱은 유한 연결 극한이므로 (1)은 (2)를 함의한다.
역으로 $\mathcal{I}$를 유한 연결 지표 범주라 하자. Lemma
\ref{categories-lemma-finite-diagram-category}의 증명에서 구성한 지표 범주의 함자를
$F : \mathcal{J} \to \mathcal{I}$라 하자. 그러면 $\mathcal{I}$를
$\mathcal{J}$로 바꾸어도 됨을 알 수 있다. 따라서
$\Ob(\mathcal{I}) = \{x_1, \ldots, x_n\} \amalg \{y_1, \ldots, y_m\}$이고
$n, m \geq 1$이며, $\mathcal{I}$의 항등사상이 아닌 모든 사상은
어떤 $i$, $j$에 대한 사상 $f : x_i \to y_j$라고 가정할 수 있다.

\medskip\noindent
$n > 1$이라고 하자. $\mathcal{I}$가 연결이므로 지표들
$i_1, i_2$, $j_0$와 사상들 $a : x_{i_1} \to y_{j_0}$,
$b : x_{i_2} \to y_{j_0}$가 존재한다. 다음 범주를 생각하자.
$$
\mathcal{I}' =
\{x\} \amalg \{x_1, \ldots, \hat x_{i_1}, \ldots, \hat x_{i_2}, \ldots x_n\}
\amalg \{y_1, \ldots, y_m\}
$$
여기서
$$
\Mor_{\mathcal{I}'}(x, y_j) = \Mor_\mathcal{I}(x_{i_1}, y_j)
\amalg \Mor_\mathcal{I}(x_{i_2}, y_j)
$$
이고, 다른 모든 사상 집합은 $\mathcal{I}$에서와 같다. 임의의 함자
$M : \mathcal{I} \to \mathcal{C}$에 대하여
$$
M'(x) = M(x_{i_1}) \times_{M(a), M(y_{j_0}), M(b)} M(x_{i_2})
$$
로 놓아 함자 $M' : \mathcal{I}' \to \mathcal{C}$를 구성할 수 있다.
예를 들어 $f : x_{i_1} \to y_j$에 대응하는 사상
$f' : x \to y_j$에 대해서는 $M'(f) = M(f) \circ \text{pr}_1$로 둔다.
함자 $M$이 극한을 가지는 것과 함자 $M'$가 극한을 가지는 것은
서로 동치이다(증명 생략). 따라서 귀납법으로 $n = 1$인 경우로
환원된다.

\medskip\noindent
$n = 1$이면 임의의 $M : \mathcal{I} \to \mathcal{C}$의 극한은
사상쌍들 $x_1 \to y_j$의 연속적인 등화자이므로 가정에 의해 존재한다.
\end{proof}

\begin{lemma}
\label{categories-lemma-almost-finite-limits-exist}
$\mathcal{C}$를 범주라 하자. 다음은 서로 동치이다.
\begin{enumerate}
\item $\mathcal{C}$에 비공집합 유한 극한이 존재한다.
\item $\mathcal{C}$에 두 대상의 곱과 등화자가 존재한다.
\item $\mathcal{C}$에 두 대상의 곱과 섬유곱이 존재한다.
\end{enumerate}
\end{lemma}

\begin{proof}
두 대상의 곱, 섬유곱, 등화자는 지표 범주가 공집합이 아닌 극한이므로
(1)은 (2)와 (3)을 모두 함의한다. (2)를 가정하자. 그러면 유한 비공집합
곱과 등화자가 존재한다. 따라서 Lemma
\ref{categories-lemma-limits-products-equalizers}에 의해 유한 비공집합 극한이
존재한다. 즉 (1)이 성립한다. (3)을 가정하자.
$a, b : A \to B$가 $\mathcal{C}$의 사상이면 $a, b$의 등화자는
$$
(A \times_{a, B, b} A)\times_{(\text{pr}_1, \text{pr}_2), A \times A, \Delta} A.
$$
이다. 따라서 (3)은 (2)를 함의하며 보조정리가 증명되었다.
\end{proof}

\begin{lemma}
\label{categories-lemma-finite-limits-exist}
$\mathcal{C}$를 범주라 하자. 다음은 서로 동치이다.
\begin{enumerate}
\item $\mathcal{C}$에 유한 극한이 존재한다.
\item 유한 곱과 등화자가 존재한다.
\item 범주에 끝 대상이 있고 섬유곱이 존재한다.
\end{enumerate}
\end{lemma}

\begin{proof}
유한 곱, 섬유곱, 등화자, 끝 대상은 유한 지표 범주 위의 극한이므로
(1)은 (2)와 (3)을 모두 함의한다. 위의 Lemma
\ref{categories-lemma-limits-products-equalizers}에 의해 (2)는 (1)을 함의한다.
(3)을 가정하자. 곱 $A \times B$는 끝 대상 위의 섬유곱임에 유의하자.
$a, b : A \to B$가 $\mathcal{C}$의 사상이면 $a, b$의 등화자는
$$
(A \times_{a, B, b} A)\times_{(\text{pr}_1, \text{pr}_2), A \times A, \Delta} A.
$$
이다. 따라서 (3)은 (2)를 함의하며 보조정리가 증명되었다.
\end{proof}

\begin{lemma}
\label{categories-lemma-push-outs-coequalizers-exist}
$\mathcal{C}$를 범주라 하자. 다음은 서로 동치이다.
\begin{enumerate}
\item $\mathcal{C}$에 연결 유한 여극한이 존재한다.
\item $\mathcal{C}$에 쌍대등화자와 밂이 존재한다.
\end{enumerate}
\end{lemma}

\begin{proof}
생략한다. 힌트: 이는 Lemma
\ref{categories-lemma-fibre-products-equalizers-exist}의 쌍대이다.
\end{proof}

\begin{lemma}
\label{categories-lemma-almost-finite-colimits-exist}
$\mathcal{C}$를 범주라 하자. 다음은 서로 동치이다.
\begin{enumerate}
\item $\mathcal{C}$에 비공집합 유한 여극한이 존재한다.
\item $\mathcal{C}$에 두 대상의 쌍대곱과 쌍대등화자가 존재한다.
\item $\mathcal{C}$에 두 대상의 쌍대곱과 밂이 존재한다.
\end{enumerate}
\end{lemma}

\begin{proof}
생략한다. 힌트: 이는 Lemma
\ref{categories-lemma-almost-finite-limits-exist}의 쌍대이다.
\end{proof}

\begin{lemma}
\label{categories-lemma-colimits-exist}
$\mathcal{C}$를 범주라 하자. 다음은 서로 동치이다.
\begin{enumerate}
\item $\mathcal{C}$에 유한 여극한이 존재한다.
\item $\mathcal{C}$에 유한 쌍대곱과 쌍대등화자가 존재한다.
\item 범주에 초기 대상이 있고 밂이 존재한다.
\end{enumerate}
\end{lemma}

\begin{proof}
생략한다. 힌트: 이는 Lemma \ref{categories-lemma-finite-limits-exist}의 쌍대이다.
\end{proof}





\section{여과 여극한}
\label{categories-section-directed-colimits}

\noindent
여과 도표 위의 여극한은 계산하거나 기술하기가 더 쉽다. 정의는 다음과 같다.

\begin{definition}
\label{categories-definition-directed}
도표 $M : \mathcal{I} \to \mathcal{C}$가 다음 조건을 만족하면
{\it 유향} 또는 {\it 여과}라고 한다.
\begin{enumerate}
\item 범주 $\mathcal{I}$에는 적어도 하나의 대상이 있고,
\item $\mathcal{I}$의 모든 대상쌍 $x, y$에 대하여 어떤 대상 $z$와
사상들 $x \to z$, $y \to z$가 존재하며,
\item $\mathcal{I}$의 모든 대상쌍 $x, y$와 $\mathcal{I}$의 모든 사상쌍
$a, b : x \to y$에 대하여 $\mathcal{C}$의 사상으로서
$M(c \circ a) = M(c \circ b)$를 만족하는 $\mathcal{I}$의 사상
$c : y \to z$가 존재한다.
\end{enumerate}
지표 범주 $\mathcal{I}$가 {\it 유향} 또는 {\it 여과}라고 함은
$\text{id} : \mathcal{I} \to \mathcal{I}$가 여과라는 뜻이다
(달리 말해 위의 (3)에서 $M$을 지운다).
\end{definition}

\noindent
여과 지표 범주를 가진 모든 도표는 여과이며, 여과 여극한은 보통 이렇게
나타난다. 실제로 $M : \mathcal{I} \to \mathcal{C}$가 여과 도표이면
$M$을 $\mathcal{I} \to \mathcal{I}' \to \mathcal{C}$로 분해할 수 있다.
여기서 $\mathcal{I}'$는 여과 지표 범주이고\footnote{구체적으로,
$\mathcal{I}'$의 대상들은 $\mathcal{I}$와 같게 두되,
$\Mor_{\mathcal{I}'}(x, y)$는 $M(a) = M(b)$일 때
$a, b : x \to y$를 동일시하는 동치관계에 의한
$\Mor_\mathcal{I}(x, y)$의 몫으로 둔다.},
$\colim_\mathcal{I} M$이 존재할 필요충분조건은
$\colim_{\mathcal{I}'} M'$이 존재하는 것이며, 이 경우 두 여극한은
표준적으로 동형이다.

\medskip\noindent
$M : \mathcal{I} \to \textit{Sets}$가 여과 도표라고 하자. 이 경우 공식
$$
\colim_\mathcal{I} M
=
(\coprod\nolimits_{i\in I} M_i)/\sim
$$
의 동치관계는 간단히 다음과 같이 기술된다.
$$
m_i \sim m_{i'}
\Leftrightarrow
\exists i'', \phi : i \to i'', \phi': i' \to i'',
M(\phi)(m_i) = M(\phi')(m_{i'}).
$$
달리 말해 두 원소가 여극한에서 같은 것과 이들이 ``결국 같아지는'' 것은
서로 동치이다.

\begin{lemma}
\label{categories-lemma-directed-commutes}
$\mathcal{I}$와 $\mathcal{J}$를 지표 범주라 하자. $\mathcal{I}$는
여과이고 $\mathcal{J}$는 유한하다고 가정하자. 집합 도표들의 도표를
$M : \mathcal{I} \times \mathcal{J} \to \textit{Sets}$,
$(i, j) \mapsto M_{i, j}$라 하자. 이 경우
$$
\colim_i \lim_j M_{i, j}
=
\lim_j \colim_i M_{i, j}.
$$
특히 $\mathcal{I}$ 위의 여극한은 집합의 유한 곱, 섬유곱, 등화자와
가환한다.
\end{lemma}

\begin{proof}
생략한다. 실제로 범주가 여과일 필요충분조건은 그 범주 위의 여극한이
유한 극한(집합의 범주로 가는 극한)과 가환하는 것임을 증명하는 것은
재미있는 연습문제이다.
\end{proof}

\noindent
$\mathcal{J}$가 무한인 경우에는 보조정리의 반례가 있다. 구체적으로
$\mathcal{I}$는 $\mathbf{N} = \{1, 2, 3, \ldots\}$로 이루어지고
$i \leq i'$일 때 유일한 사상 $i \to i'$를 가지게 하자.
$\mathcal{J}$는 이산 범주 $\mathbf{N} = \{1, 2, 3, \ldots\}$라 하자
(사상은 항등사상뿐이다). $M_{i, j} = \{1, 2, \ldots, i\}$로 놓고
$i \leq i'$일 때 자명한 포함사상 $M_{i, j} \to M_{i', j}$를 취하자.
이 경우 $\colim_i M_{i, j} = \mathbf{N}$이므로
$$
\lim_j \colim_i M_{i, j}
=
\prod\nolimits_j \mathbf{N}
=
\mathbf{N}^\mathbf{N}
$$
이다. 한편 $\lim_j M_{i, j} = \prod\nolimits_j M_{i, j}$이므로
$$
\colim_i \lim_j M_{i, j}
=
\bigcup\nolimits_i \{1, 2, \ldots, i\}^{\mathbf{N}}
$$
이고, 이는 다른 극한보다 작다.

\begin{lemma}
\label{categories-lemma-cofinal-in-filtered}
\begin{reference}
\cite[Proposition 3.2.4]{KS}
\end{reference}
$\mathcal{I}$를 범주라 하고 $\mathcal{J}$를 충만한 부분범주라 하자.
$\mathcal{I}$가 여과라고 가정하자. 또한 $\mathcal{I}$의 모든 대상
$i$에 대하여 $\mathcal{J}$의 어떤 대상 $j$로 가는 사상
$i \to j$가 존재한다고 가정하자. 그러면 $\mathcal{J}$는 여과이며
$\mathcal{I}$에서 공종이다.
\end{lemma}

\begin{proof}
생략한다. 관련 개념을 익히기 좋은 연습문제이다.
\end{proof}

\noindent
때로는 지표 범주에 부과할 수 있는 조건들을 더 세밀하게 제어해야 한다.
따라서 요구할 수 있는 조건들에 관한 몇 가지 보조정리를 덧붙인다.

\begin{lemma}
\label{categories-lemma-preserve-products}
$\mathcal{I}$를 지표 범주, 즉 범주라 하자. $\mathcal{I}$의 모든 대상쌍
$x, y$에 대하여 어떤 대상 $z$와 사상들 $x \to z$, $y \to z$가
존재한다고 가정하자. 그러면 다음이 성립한다.
\begin{enumerate}
\item $M$과 $N$이 $\mathcal{I}$ 위의 집합 도표이면 사상
$\colim (M_i \times N_i) \to \colim M_i \times \colim N_i$는 전사이고,
\item 일반적으로 $\mathcal{I}$ 위의 집합 도표의 여극한은 유한 비공집합
곱과 가환하지 않는다.
\end{enumerate}
\end{lemma}

\begin{proof}
(1)의 증명. $(\overline{m}, \overline{n})$을
$\colim M_i \times \colim N_i$의 원소라 하자. 그러면 어떤
$x, y \in \Ob(\mathcal{I})$에 대하여 $m \in M_x$와 $n \in N_y$를
찾을 수 있으며, $m$은 $\overline{m}$으로, $n$은 $\overline{n}$으로
사상된다. Section \ref{categories-section-limit-sets}을 참조하라.
$\mathcal{I}$에서 $a : x \to z$와 $b : y \to z$를 택하자. 그러면
$(M(a)(m), N(b)(n))$은 $(M \times N)_z$의 원소이고,
$\colim (M_i \times N_i)$에서의 그 상은 원하는 대로
$(\overline{m}, \overline{n})$으로 사상된다.

\medskip\noindent
(2)의 증명. $G$를 자명하지 않은 군이라 하고 $\mathcal{I}$를 자기준동형
모노이드가 $G$인 한 대상 범주라 하자. 그러면 $\mathcal{I}$는 보조정리의
조건을 자명하게 만족한다. 이제 $G$가 평행이동으로 자기 자신에 작용하게
하고 $G$-집합 $G$를 집합값 $\mathcal{I}$-도표로 보자. 그러면
$$
\colim_\mathcal{I} G \times \colim_\mathcal{I} G \cong G/G \times G/G
$$
는 다음과 동형이 아니다.
$$
\colim_\mathcal{I} (G \times G) \cong (G \times G)/G
$$
이 예는 유한 곱만을 다루더라도 Lemma
\ref{categories-lemma-directed-commutes}의 추가 조건을 단순히 버릴 수 없음을
보여 준다.
\end{proof}

\begin{lemma}
\label{categories-lemma-colimits-abelian-as-sets}
$\mathcal{I}$를 지표 범주, 즉 범주라 하자. $\mathcal{I}$의 모든 대상쌍
$x, y$에 대하여 어떤 대상 $z$와 사상들 $x \to z$, $y \to z$가
존재한다고 가정하자. $M : \mathcal{I} \to \textit{Ab}$를
$\mathcal{I}$ 위의 아벨군 도표라 하자. 그러면 집합의 범주에서의
$M$의 여극한에서 아벨군의 범주에서의 $M$의 여극한으로 가는
전사사상이 존재한다.
\end{lemma}

\begin{proof}
집합의 범주에서의 여극한은 관계에 의한 서로소 합집합
$\coprod M_i$의 몫임을 상기하자. Section
\ref{categories-section-limit-sets}을 참조하라. 마찬가지로 아벨군의 범주에서의
여극한은 직합 $\bigoplus M_i$의 몫이다. 보조정리의 가정은
$i, j \in \Ob(\mathcal{I})$, $m \in M_i$, $n \in M_j$가 주어지면
어떤 대상 $k$와 사상들 $a : i \to k$, $b : j \to k$를 찾을 수 있다는
뜻이다. 따라서 $m + n$은 여극한에서 $M_k$의 원소
$M(a)(m) + M(b)(n)$으로 표현된다. 그러므로 $\coprod M_i$에서
여극한으로 가는 사상은 전사이다.
\end{proof}

\begin{lemma}
\label{categories-lemma-split-into-connected}
$\mathcal{I}$를 지표 범주, 즉 범주라 하자. 모든 실선 도표
$$
\xymatrix{
x \ar[d] \ar[r] & y \ar@{..>}[d] \\
z \ar@{..>}[r] & w
}
$$
에 대하여 도표를 가환하게 하는 대상 $w$와 점선 화살표들이
$\mathcal{I}$ 안에 존재한다고 가정하자. 그러면 $\mathcal{I}$는
공집합이거나, 같은 성질을 가지는 공집합이 아닌 연결 범주들의 서로소
합집합이다.
\end{lemma}

\begin{proof}
$\mathcal{I}$가 공집합 범주이면 보조정리는 참이다. 그렇지 않으면
어떤 $z$와 사상들 $x \to z$, $y \to z$가 존재할 때
$x \sim y$라고 하여 $\mathcal{I}$의 대상들 위에 관계를 정의한다.
보조정리의 가정에 의해 이는 동치관계이다. 따라서
$\Ob(\mathcal{I})$는 동치류들의 서로소 합집합이다. 이 동치류들에
대응하는 충만한 부분범주들을 $\mathcal{I}_j$라 하자. 그러면 원하는
대로 $\mathcal{I}_j$가 공집합이 아닌 상태에서
$\mathcal{I} = \coprod \mathcal{I}_j$이다.
\end{proof}

\begin{lemma}
\label{categories-lemma-preserve-injective-maps}
$\mathcal{I}$를 지표 범주, 즉 범주라 하자. 모든 실선 도표
$$
\xymatrix{
x \ar[d] \ar[r] & y \ar@{..>}[d] \\
z \ar@{..>}[r] & w
}
$$
에 대하여 도표를 가환하게 하는 대상 $w$와 점선 화살표들이
$\mathcal{I}$ 안에 존재한다고 가정하자. 그러면 다음이 성립한다.
\begin{enumerate}
\item $\mathcal{I}$ 위의 집합 도표들의 단사 사상 $M \to N$은 집합의
단사 사상 $\colim M_i \to \colim N_i$를 유도하고,
\item 일반적으로 아벨군의 도표들과 그 여극한에 대해서는 같은 명제가
성립하지 않는다.
\end{enumerate}
\end{lemma}

\begin{proof}
$\mathcal{I}$가 공집합 범주이면 보조정리는 참이다. 따라서
$\mathcal{I}$가 공집합이 아니라고 가정해도 된다. 이 경우 Lemma
\ref{categories-lemma-split-into-connected}에 의해 각 $\mathcal{I}_j$가 공집합이
아니고 같은 성질을 만족하도록
$\mathcal{I} = \coprod \mathcal{I}_j$로 쓸 수 있다. 다음 등식 때문에
$\colim_\mathcal{I} M = \coprod_j \colim_{\mathcal{I}_j} M|_{\mathcal{I}_j}$
(1)의 증명은 연결인 경우로 환원된다.

\medskip\noindent
$\mathcal{I}$가 연결이고 $M \to N$이 단사, 즉 모든 사상
$M_i \to N_i$가 단사라고 가정하자. $M_i$를 $M_i \to N_i$의 상과
동일시하여 $M_i$를 $N_i$의 부분집합으로 생각한다. 앞으로 따로
언급하지 않고 Section \ref{categories-section-limit-sets}에 주어진 여극한의
기술을 사용한다. $s, s' \in \colim M_i$가 $\colim N_i$의 같은
원소로 사상된다고 하자. $s$는 $M_i$의 원소 $m$에서 오고 $s'$는
$M_{i'}$의 원소 $m'$에서 온다고 하자. 그러면 $\mathcal{I}$의 대상들의 열
$i = i_0, i_1, \ldots, i_n = i'$와 사상들
$$
\xymatrix{
&
i_1 \ar[ld] \ar[rd] & &
i_3 \ar[ld] & &
i_{2n-1} \ar[rd] & \\
i = i_0 & &
i_2 & &
\ldots & &
i_{2n} = i'
}
$$
및 원소들 $n_{i_j} \in N_{i_j}$을 찾을 수 있다. 이 원소들은 위의
$\mathcal{I}$의 사상들에서 유도되는 사상
$N_{i_{2k-1}} \to N_{i_{2k-2}}$와
$N_{i_{2k-1}} \to N_{i_{2k}}$ 아래에서 서로 대응하고,
$n_{i_0} = m$, $n_{i_{2n}} = m'$이다. 이것이 $s = s'$를 함의함을
$n$에 관한 귀납법으로 보이겠다. 기초 단계 $n = 0$은 자명하다.
$n \geq 1$이라고 가정하자. $\mathcal{I}$에 관한 가정을 사용하여
가환 도표
$$
\xymatrix{
& i_1 \ar[ld] \ar[rd] \\
i_0 \ar[rd] & & i_2 \ar[ld] \\
& w
}
$$
를 얻는다. $m$과 $n_{i_2}$는 모두 원소 $n_{i_1}$의 상이므로
$N_w$의 같은 원소로 사상된다. 특히 이 원소는 그 부분집합에 속하는 원소
$m'' \in M_w$이고, $\colim M_i$에서 $s$와 같은 원소를 준다.
이제 사슬
$$
\xymatrix{
&
i_3 \ar[ld] \ar[rd] & &
i_5 \ar[ld] & &
i_{2n-1} \ar[rd] & \\
w & &
i_4 & &
\ldots & &
i_{2n} = i'
}
$$
과 $j \geq 3$에 대한 원소들 $n_{i_j}$을 얻는데, 이 사슬은 처음
사슬보다 길이가 짧다. 이로써 귀납 단계가 증명되고 (1)의 증명이 끝난다.

\medskip\noindent
$G$를 군이라 하고 $\mathcal{I}$를 자기준동형 모노이드가 $G$인
한 대상 범주라 하자. $g_1, g_2 \in G$가 주어지면
$h_1 g_1 = h_2 g_2$를 만족하는 $h_1, h_2 \in G$를 찾을 수 있으므로
$\mathcal{I}$는 보조정리의 조건을 만족한다. $\textit{Ab}$에서
$\mathcal{I}$ 위의 도표 $M$은 $G$-작용을 가진 아벨군 $M$과 같은
데이터이고, $\colim_\mathcal{I} M$은 $M$의 쌍대불변량 $M_G$이다.
$G$를 위수가 $2$인 군으로 취하여
$M = \mathbf{Z}/2\mathbf{Z}$에 자명하게 작용하게 하고, 이를
$N = \mathbf{Z}/2\mathbf{Z} \times \mathbf{Z}/2\mathbf{Z}$의 첫째
직합인자에 사상하자. 여기서 $G$의 자명하지 않은 원소는
$(x, y) \mapsto (x + y, y)$로 작용한다. 그러면 $M_G \to N_G$는 영이다.
\end{proof}

\begin{lemma}
\label{categories-lemma-split-into-directed}
$\mathcal{I}$를 지표 범주, 즉 범주라 하자. 다음을 가정하자.
\begin{enumerate}
\item $\mathcal{I}$의 모든 사상쌍 $a : w \to x$, $b : w \to y$에
대하여 어떤 대상 $z$와 사상들 $c : x \to z$, $d : y \to z$가
존재하여 $c \circ a = d \circ b$이고,
\item 모든 사상쌍 $a, b : x \to y$에 대하여 어떤 사상
$c : y \to z$가 존재하여 $c \circ a = c \circ b$이다.
\end{enumerate}
그러면 $\mathcal{I}$는 서로소 여과 지표 범주들 $\mathcal{I}_j$의
(공집합일 수도 있는) 합집합이다.
\end{lemma}

\begin{proof}
$\mathcal{I}$가 공집합 범주이면 보조정리는 참이다. 그렇지 않으면
어떤 $z$와 사상들 $x \to z$, $y \to z$가 존재할 때 $x \sim y$라고
하여 $\mathcal{I}$의 대상들 위에 관계를 정의한다. 보조정리의 첫째
가정에 의해 이는 동치관계이다. 따라서 $\Ob(\mathcal{I})$는
동치류들의 서로소 합집합이다. 이 동치류들에 대응하는 충만한
부분범주들을 $\mathcal{I}_j$라 하자. 나머지는 정의에서 명백하다.
\end{proof}

\begin{lemma}
\label{categories-lemma-almost-directed-commutes-equalizers}
$\mathcal{I}$를 위의 Lemma \ref{categories-lemma-split-into-directed}의 가정을
만족하는 지표 범주라 하자. 그러면 $\mathcal{I}$ 위의 여극한은 집합의
섬유곱과 등화자(더 일반적으로는 유한 연결 극한)와 가환한다.
\end{lemma}

\begin{proof}
Lemma \ref{categories-lemma-split-into-directed}에 의해 각 $\mathcal{I}_j$가
여과이도록 $\mathcal{I} = \coprod \mathcal{I}_j$로 쓸 수 있다.
Lemma \ref{categories-lemma-directed-commutes}에 의해 $\mathcal{I}_j$의 여극한은
등화자 및 섬유곱과 가환한다. 따라서 집합의 범주에서 등화자와
섬유곱이 쌍대곱(공집합 쌍대곱을 포함한다)과 가환함을 보이면 충분하다.
달리 말해, 집합 $J$, 집합들 $A_j, B_j, C_j$, 그리고 $j \in J$에
대한 집합 사상들 $A_j \to B_j$, $C_j \to B_j$가 주어지면
$$
(\coprod\nolimits_{j \in J} A_j)
\times_{(\coprod\nolimits_{j \in J} B_j)}
(\coprod\nolimits_{j \in J} C_j)
=
\coprod\nolimits_{j \in J} A_j \times_{B_j} C_j
$$
임을 보여야 하고, $a_j, a'_j : A_j \to B_j$가 주어지면
$$
\text{Equalizer}(
\coprod\nolimits_{j \in J} a_j,
\coprod\nolimits_{j \in J} a'_j)
=
\coprod\nolimits_{j \in J}
\text{Equalizer}(a_j, a'_j)
$$
임을 보여야 한다. 이는 $J = \emptyset$인 경우에도 참이다.
세부사항은 생략한다.
\end{proof}



\section{쌍대여과 극한}
\label{categories-section-codirected-limits}

\noindent
쌍대여과 도표 위의 극한은 계산하거나 기술하기가 더 쉽다.
정의는 다음과 같다.

\begin{definition}
\label{categories-definition-codirected}
도표 $M : \mathcal{I} \to \mathcal{C}$가 다음 조건을 만족하면
{\it 쌍대유향} 또는 {\it 쌍대여과}라고 한다.
\begin{enumerate}
\item 범주 $\mathcal{I}$에는 적어도 하나의 대상이 있고,
\item $\mathcal{I}$의 모든 대상쌍 $x, y$에 대하여 어떤 대상 $z$와
사상들 $z \to x$, $z \to y$가 존재하며,
\item $\mathcal{I}$의 모든 대상쌍 $x, y$와 $\mathcal{I}$의 모든
사상쌍 $a, b : x \to y$에 대하여 $\mathcal{C}$의 사상으로서
$M(a \circ c) = M(b \circ c)$를 만족하는 $\mathcal{I}$의 사상
$c : w \to x$가 존재한다.
\end{enumerate}
지표 범주 $\mathcal{I}$가 {\it 쌍대유향} 또는 {\it 쌍대여과}라고 함은
$\text{id} : \mathcal{I} \to \mathcal{I}$가 쌍대여과라는 뜻이다
(달리 말해 위의 (3)에서 $M$을 지운다).
\end{definition}

\noindent
쌍대여과 지표 범주를 가진 모든 도표는 쌍대여과이며, 이 상황은 보통
이렇게 나타난다.

\medskip\noindent
집합의 쌍대여과 극한이 일반적인 극한보다 ``쉽다''는 이유의 예로,
유한 비공집합 집합들의 쌍대여과 도표는 비공집합 극한을 가진다는
사실(Lemma \ref{categories-lemma-nonempty-limit})을 언급한다. 이 결과는 유한
비공집합 집합들의 일반적인 극한에 대해서는 성립하지 않는다.












\section{준순서 집합 위의 극한과 여극한}
\label{categories-section-posets-limits}

\noindent
도표의 특수한 경우로 준순서 집합 위의 계가 있다.

\begin{definition}
\label{categories-definition-directed-set}
$I$를 집합이라 하고 $\leq$를 $I$ 위의 이항관계라 하자.
\begin{enumerate}
\item $\leq$가 추이적($i \leq j$이고 $j \leq k$이면
$i \leq k$)이며 반사적(모든 $i \in I$에 대하여 $i \leq i$)이면
이를 {\it 준순서}라고 한다.
\item {\it 준순서 집합}은 준순서가 주어진 집합이다.
\item {\it 유향 집합}은 $I$가 공집합이 아니고
$\forall i, j \in I$에 대하여 $i \leq k, j \leq k$를 만족하는
$k \in I$가 존재하는 준순서 집합 $(I, \leq)$이다.
\item $\leq$가 준순서이고 반대칭적($i \leq j$이고 $j \leq i$이면
$i = j$)이면 이를 {\it 부분순서}라고 한다.
\item {\it 부분순서 집합}은 부분순서가 주어진 집합이다.
\item {\it 유향 부분순서 집합}은 그 순서가 부분순서인 유향 집합이다.
\end{enumerate}
\end{definition}

\noindent
준순서 집합을 말할 때 표기에서 $\leq$를 생략하는 것이 관례이다.
즉 준순서 집합 $(I, \leq)$ 대신 준순서 집합 $I$라고 한다.
준순서 집합 $I$가 주어지면 모든 $i, j \in I$에 대하여
$i \geq j \Leftrightarrow j \leq i$라는 규칙으로 기호 $\geq$를
정의한다. ``부분순서 집합''이라는 말은 때때로 ``poset''으로 줄여 쓴다.

\medskip\noindent
준순서 집합 $I$가 주어지면 다음과 같이 범주를 구성할 수 있다.
대상은 $I$의 원소들이고, $i \leq i'$이면 사상 $i \to i'$가 정확히
하나 있으며 그렇지 않으면 없다. 거꾸로 임의의 두 대상 사이에
화살표가 많아야 하나인 범주 $\mathcal{C}$가 주어지면, 집합
$\Ob(\mathcal{C})$에는
$x \leq y \Leftrightarrow \Mor_\mathcal{C}(x, y) \not = \emptyset$라는
규칙으로 정의되는 준순서가 주어진다.

\begin{definition}
\label{categories-definition-system-over-poset}
$(I, \leq)$를 준순서 집합이라 하고 $\mathcal{C}$를 범주라 하자.
\begin{enumerate}
\item $\mathcal{C}$에서의 {\it $I$ 위의 계}는, 때때로
$\mathcal{C}$에서의 {\it $I$ 위의 귀납계}라고도 하며,
$\mathcal{C}$의 대상들 $M_i$와 모든 $i \leq i'$에 대한 사상
$f_{ii'} : M_i \to M_{i'}$로 주어진다. 이들은 $f_{ii} = \text{id}$이고
$i \leq i' \leq i''$일 때
$f_{ii''} = f_{i'i''} \circ f_{i i'}$를 만족한다.
\item $\mathcal{C}$에서의 {\it $I$ 위의 역계}는, 때때로
$\mathcal{C}$에서의 {\it $I$ 위의 사영계}라고도 하며,
$\mathcal{C}$의 대상들 $M_i$와 모든 $i' \leq i$에 대한 사상
$f_{ii'} : M_i \to M_{i'}$로 주어진다. 이들은 $f_{ii} = \text{id}$이고
$i'' \leq i' \leq i$일 때
$f_{ii''} = f_{i'i''} \circ f_{i i'}$를 만족한다
(부등호의 방향이 뒤집힌 것에 유의하라).
\end{enumerate}
이를 나타내기 위해 $(M_i, f_{ii'})$를 $I$ 위의 (역)계라고 하겠다.
사상 $f_{ii'}$를 때때로 {\it 전이 사상}이라고 한다.
\end{definition}

\noindent
달리 말해 $I$ 위의 계는 위에서 $I$에 대응시킨 범주를
$\mathcal{I}$라 할 때의 도표 $M : \mathcal{I} \to \mathcal{C}$일 뿐이다.
대상은 $I$의 원소들이고, $i \leq i'$일 필요충분조건은
$\mathcal{I}$ 안에 유일한 화살표 $i \to i'$가 존재하는 것이다.
역계는 도표 $M : \mathcal{I}^{opp} \to \mathcal{C}$이다. 이 관점에서는
$I$ 위의 임의의 (역)계의 (여)극한을 취할 수 있다. 그러나 관례상
{\it $I$ 위의 계에서는 여극한만 취하고},
{\it $I$ 위의 역계에서는 극한만 취한다}. 더 정확히 말하면,
$I$ 위의 계 $(M_i, f_{ii'})$가 주어졌을 때 계
$(M_i, f_{ii'})$의 여극한은
$$
\colim_{i \in I} M_i = \colim_\mathcal{I} M,
$$
즉 대응하는 도표의 여극한으로 정의된다. $I$ 위의 역계
$(M_i, f_{ii'})$가 주어졌을 때 역계 $(M_i, f_{ii'})$의 극한은
$$
\lim_{i \in I} M_i = \lim_{\mathcal{I}^{opp}} M,
$$
즉 대응하는 도표의 극한으로 정의된다.

\begin{remark}
\label{categories-remark-preorder-versus-partial-order}
$I$를 준순서 집합이라 하자. $I$로부터 표준적인 부분순서 집합
$\overline{I}$와 순서 보존 사상 $\pi : I \to \overline{I}$를 구성할
수 있다. 구체적으로 다음 규칙으로 $I$ 위의 동치관계 $\sim$를 정의한다.
$$
i \sim j \Leftrightarrow (i \leq j\text{ and }j \leq i).
$$
$\overline{I} = I/\sim$로 놓고 $\pi : I \to \overline{I}$를 몫사상으로
둔다. 마지막으로 $\overline{I}$에는
$\pi(i) \leq \pi(j) \Leftrightarrow i \leq j$를 만족하는 유일한
부분순서가 주어진다. $I$가 유향 집합이면 $\overline{I}$는 유향
부분순서 집합임에 유의하라. $\overline{I}$ 위의 (역)계 $N$이 주어지면
$M_i = N_{\pi(i)}$로 놓아 $I$ 위의 (역)계 $M$을 얻는다. 이 구성은
$I$와 $\overline{I}$ 위의 역계들의 범주 사이의 함자를 정의한다.
실제로 이는 동치이다. 그 이유는 $i \sim j$이면 $I$ 위의 임의의 계
$M$에 대하여 사상들 $M_i \to M_j$와 $M_j \to M_i$가 서로 역인
동형이기 때문이다. 더 정확히 말해, $\pi$의 절단
$s : \overline{I} \to I$를 택하면 위 함자의 준역은 $M$을
$N_{\overline{i}} = M_{s(\overline{i})}$인 $N$으로 보낸다.
마지막으로 이 대응은 계의 여극한과 양립한다. $M$과 $N$이 위와 같이
관련되어 있고 $\colim_{\overline{I}} N$ 또는 $\colim_I M$ 중 하나가
존재하면 다른 하나도 존재하며
$\colim_{\overline{I}} N = \colim_I M$이다. 역계와 역계의 극한에
대해서도 비슷한 결과가 성립한다.
\end{remark}

\noindent
Remark \ref{categories-remark-preorder-versus-partial-order}의 요지는 계의 여극한이나
역계의 극한을 계산할 때 준순서가 부분순서라고 항상 가정해도 된다는
것이다.

\begin{definition}
\label{categories-definition-directed-system}
$I$를 준순서 집합이라 하자. $I$가 유향 집합이면 계(각각 역계)
$(M_i, f_{ii'})$를 {\it 유향계}(각각 {\it 유향 역계})라고 한다
(Definition \ref{categories-definition-directed-set}). 즉 $I$는 공집합이 아니며,
모든 $i_1, i_2 \in I$에 대하여 $i_1 \leq i$이고 $i_2 \leq i$인
$i\in I$가 존재한다.
\end{definition}

\noindent
이 경우 여극한을 때때로 (유감스럽게도) ``직접극한''이라고 부른다.
여기서는 이 마지막 용어를 사용하지 않겠다. 다음 의미에서 여과 범주
위의 도표는 유향계보다 더 일반적이지 않음이 밝혀진다.

\begin{lemma}
\label{categories-lemma-directed-category-system}
$\mathcal{I}$를 여과 지표 범주라 하자. 다음 성질을 만족하는 유향 집합
$I$와 $\mathcal{I}$에서의 $I$ 위의 계 $(x_i, \varphi_{ii'})$가 존재한다.
\begin{enumerate}
\item 모든 범주 $\mathcal{C}$와 $\mathcal{C}$에 값을 가지는 모든 도표
$M : \mathcal{I} \to \mathcal{C}$에 대하여 대응하는 $I$ 위의 계를
$(M(x_i), M(\varphi_{ii'}))$로 나타내자.
$\colim_{i \in I} M(x_i)$가 존재하면 $\colim_\mathcal{I} M$도 존재하고,
Lemma \ref{categories-lemma-functorial-colimit}의 변환
$$
\theta :
\colim_{i \in I} M(x_i)
\longrightarrow
\colim_\mathcal{I} M
$$
은 동형이다.
\item 모든 범주 $\mathcal{C}$와 $\mathcal{C}$에서의 모든 도표
$M : \mathcal{I}^{opp} \to \mathcal{C}$에 대하여 대응하는 $I$ 위의
역계를 $(M(x_i), M(\varphi_{ii'}))$로 나타내자.
$\lim_{i \in I} M(x_i)$가 존재하면 $\lim_{\mathcal{I}^{opp}} M$도 존재하고,
Lemma \ref{categories-lemma-functorial-limit}의 변환
$$
\theta :
\lim_{\mathcal{I}^{opp}} M
\longrightarrow
\lim_{i \in I} M(x_i)
$$
은 동형이다.
\end{enumerate}
\end{lemma}

\begin{proof}
Definition \ref{categories-definition-system-over-poset} 다음의 본문에서 설명했듯이,
준순서 집합을 범주로, 계를 함자로 볼 수 있다. 증명 전체에서 이 두
관점을 자유롭게 오가겠다. 유향 집합에 대응하는 범주
$\mathcal{I}_0$\footnote{실제로 이 구성은 유향 부분순서 집합을
만든다.}와 공종 함자 $M_0 : \mathcal{I}_0 \to \mathcal{I}$를 구성하여
첫째 명제를 증명한다. 그러면 Lemma \ref{categories-lemma-cofinal}에 의해 도표
$M : \mathcal{I} \to \mathcal{C}$의 여극한은 도표
$M \circ M_0 : \mathcal{I}_0 \to \mathcal{C}$의 여극한과 일치하므로
명제가 따른다. 둘째 명제는 첫째 명제의 쌍대이며, $\mathcal{C}$에서의
극한을 $\mathcal{C}^{opp}$에서의 여극한으로 해석하여 증명할 수 있다.
세부사항은 생략한다.

\medskip\noindent
범주 $\mathcal{F}$ 안에 유한한 화살표 집합 $F$가 존재하여
$\mathcal{F}$의 모든 화살표가 $F$의 화살표들을 합성하여 얻어지면
$\mathcal{F}$를 {\em 유한생성}이라고 한다. 특히 이는 $\mathcal{F}$의
대상이 유한 개임을 함의한다. 먼저 $\mathcal{I}$가 $\mathcal{I}$의 모든 유한생성
부분범주를 유일한 끝 대상을 가진 유한생성 부분범주로 확장할 수 있는
경우로 환원하겠다.

\medskip\noindent
$\omega$를 유한 순서수들의 유향 집합이라 하고, 이를 여과 범주로 보자.
곱범주 $\mathcal{I}\times \omega$도 여과이며 사영
$\Pi : \mathcal{I} \times \omega \to \mathcal{I}$는 공종임을 쉽게
확인할 수 있다.

\medskip\noindent
이제 $\mathcal{F}$를 $\mathcal{I}\times \omega$의 임의의 유한생성
부분범주라 하자. 여과 범주의 공리와 $\mathcal{F}$의 유한한 생성자
집합에 관한 간단한 귀납 논증을 사용하면, 도표
$\mathcal{F} \to \mathcal{I} \times \omega$의 쌍대원뿔
$(\{f_i\}, i_\infty)$를 $\mathcal{I} \times \omega$ 안에서 구성할 수
있다. 즉 $\mathcal{F}$의 모든 대상 $i$에 대하여 사상
$f_i : i \to i_\infty$가 있고, $\mathcal{F}$의 각 화살표
$f : i \to i'$에 대하여 $f_i = f_{i'} \circ f$이다.
또한 $i_\infty$에서 $\mathcal{F}$의 대상으로 가는 화살표가 없도록
$i_\infty$를 택할 수 있다. 실제로 $\mathcal{I}$-성분에서는 항등이고
$\omega$-성분에서는 엄밀히 증가하는 화살표를 사상들 $f_i$에 항상
뒤에서 합성할 수 있기 때문이다. 이제 $\mathcal{F}^+$를
$\mathcal{F}$의 모든 대상과 화살표에 대상 $i_\infty$, 항등화살표
$\text{id}_{i_\infty}$, 화살표들 $f_i$를 더한 범주라 하자.
$\mathcal{F}^+$ 안의 $i_\infty$에서 $\mathcal{F}$의 대상으로 가는
화살표가 없으므로 $\mathcal{F}^+$의 화살표 집합은 합성에 닫혀 있고,
따라서 $\mathcal{F}^+$는 실제로 범주이다. 구성에 의해 이는
$i_\infty$를 유일한 끝 대상으로 가지는 $\mathcal{I}$의 유한생성
부분범주이다. Lemma \ref{categories-lemma-cofinal}에 의해 임의의 도표
$M : \mathcal{I} \to \mathcal{C}$의 여극한은 $M\circ\Pi$의 여극한과
일치하므로 원하는 환원을 얻는다.

\medskip\noindent
유일한 끝 대상을 가진 $\mathcal{I}$의 모든 유한생성 부분범주들의
집합에는 포함관계에 의한 자연스러운 순서가 주어진다. 이 집합에 대응하는
범주를 $\mathcal{I}_0$로 취한다. 또한 함자
$M_0 : \mathcal{I}_0 \to \mathcal{I}$가 있는데, 이는
$\mathcal{I}_0$의 화살표 $\mathcal{F} \subset \mathcal{F'}$를
$\mathcal{F}$의 끝 대상에서 $\mathcal{F}'$의 끝 대상으로 가는
유일한 사상으로 보낸다. $\mathcal{I}$의 임의의 두 유한생성 부분범주가
주어지면 이 두 범주로 생성되는 범주도 유한생성이다. $\mathcal{I}$에
대한 가정에 의해 이는 유일한 끝 대상을 가진 $\mathcal{I}$의 어떤
유한생성 부분범주에 포함된다. 따라서 $\mathcal{I}_0$는 유향이다.

\medskip\noindent
마지막으로 $M_0$가 공종임을 확인하자. $\mathcal{I}$의 임의의 대상은
그 대상과 그 항등화살표만으로 이루어진 부분범주의 끝 대상이므로
함자 $M_0$는 대상 위에서 전사이다. 특히 Definition
\ref{categories-definition-cofinal}의 조건 (1)이 성립한다. $\mathcal{I}$의 대상
$i$, $\mathcal{I}_0$의 대상들 $\mathcal{F}_1, \mathcal{F}_2$, 그리고
$\mathcal{I}$의 사상들 $\varphi_1 : i \to M_0(\mathcal{F}_1)$,
$\varphi_2 : i \to M_0(\mathcal{F}_2)$가 주어지면,
$\mathcal{F}_1$, $\mathcal{F}_2$와 사상들 $\varphi_1, \varphi_2$를
포함하는 유일한 끝 대상을 가진 유한생성 범주를
$\mathcal{F}_{12}$로 취할 수 있다. 그 결과인 도표
$$
\xymatrix{
& M_0(\mathcal{F}_{12}) & \\
M_0(\mathcal{F}_{1}) \ar[ru] & & M_0(\mathcal{F}_{2}) \ar[lu] \\
& i \ar[lu] \ar[ru]
}
$$
는 $\mathcal{F}_{12}$ 안에 있고 $M_0(\mathcal{F}_{12})$가 이 범주의
끝 대상이므로 가환한다. 따라서 조건 (2)도 성립하며 증명이 끝난다.
\end{proof}

\begin{remark}
\label{categories-remark-trick-needed}
유한 유향 집합 $(I, \geq)$는 항상 최대 대상 $i_\infty$를 가짐에
유의하라. 따라서 그러한 집합 위의 계 $(M_i, f_{ii'})$의 여극한은
$M_{i_\infty}$와 같다는 의미에서 자명하다. 이와 대조적으로 유한 여과
범주를 지표로 하는 여극한은 자명하지 않을 수 있다. 예를 들어
$\mathcal{I}$를 대상 $i$ 하나와 $e = e \circ e$를 만족하는 자명하지
않은 사상 $e$ 하나를 가진 범주라 하자. 도표
$M : \mathcal{I} \to Sets$의 여극한은 멱등원 $M(e)$의 상이다.
이는 Lemma \ref{categories-lemma-directed-category-system}의 증명에서
$\mathcal{I}\times \omega$로 옮기는 것과 같은 기법이 본질적임을 보여 준다.
\end{remark}

\begin{lemma}
\label{categories-lemma-nonempty-limit}
$S : \mathcal{I} \to \textit{Sets}$가 집합들의 쌍대여과 도표이고 모든
$S_i$가 유한 비공집합이면 $\lim_i S_i$는 공집합이 아니다. 달리 말해
유한 비공집합 집합들의 유향 역계의 극한은 공집합이 아니다.
\end{lemma}

\begin{proof}
두 명제는 Lemma \ref{categories-lemma-directed-category-system}에 의해 동치이다.
$I$를 유향 집합이라 하고 $(S_i)_{i \in I}$를 $I$ 위의 유한 비공집합
집합들의 역계라 하자. {\it 부분계} $T$를 공집합이 아닌 부분집합들
$T_i \subset S_i$의 족 $T = (T_i)_{i \in I}$로 정의하자. 여기서 모든
$i' \geq i$에 대하여 $T_{i'}$는 전이 사상 $S_{i'} \to S_i$ 아래에서
$T_i$ 안으로 사상된다. 부분계들의 집합을 $\mathcal{T}$로 나타내자.
$\mathcal{T}$에 포함관계로 순서를 주자. $T_\alpha$, $\alpha \in A$가
$\mathcal{T}$의 원소들로 이루어진 전순서 족이라고 하자.
$T_\alpha = (T_{\alpha, i})_{i \in I}$라고 쓰자. 그러면
$T_i = \bigcap_{\alpha \in A} T_{\alpha, i}$로 놓아 하계
$T = (T_i)_{i \in I}$를 찾을 수 있다. 모든 $T_{\alpha, i}$가
공집합이 아니고 $T_\alpha$들이 전순서 족을 이루므로 이는 명백히
$S_i$의 유한 비공집합 부분집합이다. 따라서 초른의 보조정리를 적용하면
$\mathcal{T}$가 극소 원소를 가짐을 알 수 있다.

\medskip\noindent
극소 원소 $T \in \mathcal{T}$의 모습을 분석하자. 먼저 사상들
$T_{i'} \to T_i$가 모두 전사임을 관찰하자. 실제로 $I$가 유향 집합이고
$T_i$가 유한하므로 교집합
$T'_i = \bigcap_{i' \geq i} \Im(T_{i'} \to T_i)$는 공집합이 아니다.
따라서 $T' = (T'_i)$는 $T$에 포함되는 부분계이고 극소성에 의해
$T' = T$이다. 마지막으로 각 $i$에 대하여 $T_i$가 한원소집합임을
주장한다. 실제로 $x \in T_i$이면 $i' \geq i$에 대하여
$T'_{i'} = (T_{i'} \to T_i)^{-1}(\{x\})$로 놓고, $j \not \geq i$이면
$T'_j = T_j$로 놓을 수 있다. 위에서 부분계 $T$의 전이 사상들이
전사임을 보았으므로 이것도 부분계이다. 극소성에 의해 $T = T'$이고,
이는 실제로 $T_i$가 한원소집합임을 함의한다. 이는 모든 $i \in I$에
대하여 성립하므로 어떤 $x_i \in S_i$에 대하여 $T_i = \{x_i\}$이며,
모든 $i' \geq i$에 대하여 $S_{i'} \to S_i$ 아래에서
$x_{i'} \mapsto x_i$이다. 달리 말해 $(x_i) \in \lim S_i$이고
보조정리가 증명되었다.
\end{proof}






\section{본질적으로 상수인 계}
\label{categories-section-essentially-constant}

\noindent
$M : \mathcal{I} \to \mathcal{C}$를 범주 $\mathcal{C}$에서의 도표라 하자.
지표 범주 $\mathcal{I}$가 여과라고 가정하자. 이 경우
$\mathcal{C}$의 대상 $X$를 골라내는, 차례로 더 강해지는 세 개념이 있다.
첫째는 단순히
$$
X = \colim_{i \in \mathcal{I}} M_i.
$$
이다. 그러면 $X$에는 쌍대사영들 $M_i \to X$가 주어진다. 더 강한
조건은 $X$가 여극한이고 어떤 $i \in \mathcal{I}$와 사상
$X \to M_i$가 존재하여 합성 $X \to M_i \to X$가 $\text{id}_X$라고
요구하는 것이다. 이보다 더 강한 조건은 다음과 같다.

\begin{definition}
\label{categories-definition-essentially-constant-diagram}
$M : \mathcal{I} \to \mathcal{C}$를 범주 $\mathcal{C}$에서의 도표라 하자.
\begin{enumerate}
\item 지표 범주 $\mathcal{I}$가 여과라고 가정하고
$(X, \{M_i \to X\}_i)$를 $M$의 쌍대원뿔이라 하자. Remark
\ref{categories-remark-cones-and-cocones}를 참조하라. 어떤 $i \in \mathcal{I}$와
사상 $X \to M_i$가 존재하여 다음을 만족하면 $M$이 {\it 값} $X$를
가지는 {\it 본질적으로 상수인} 도표라고 한다.
\begin{enumerate}
\item $X \to M_i \to X$는 $\text{id}_X$이고,
\item 모든 $j$에 대하여 어떤 $k$와 사상들 $i \to k$, $j \to k$가
존재하여 사상 $M_j \to M_k$가 합성
$M_j \to X \to M_i \to M_k$와 같다.
\end{enumerate}
\item 지표 범주 $\mathcal{I}$가 쌍대여과라고 가정하고
$(X, \{X \to M_i\}_i)$를 $M$의 원뿔이라 하자. Remark
\ref{categories-remark-cones-and-cocones}를 참조하라. 어떤 $i \in \mathcal{I}$와
사상 $M_i \to X$가 존재하여 다음을 만족하면 $M$이 {\it 값} $X$를
가지는 {\it 본질적으로 상수인} 도표라고 한다.
\begin{enumerate}
\item $X \to M_i \to X$는 $\text{id}_X$이고,
\item 모든 $j$에 대하여 어떤 $k$와 사상들 $k \to i$, $k \to j$가
존재하여 사상 $M_k \to M_j$가 합성
$M_k \to M_i \to X \to M_j$와 같다.
\end{enumerate}
\end{enumerate}
이 정의를 사용할 때 Lemma
\ref{categories-lemma-essentially-constant-is-limit-colimit}을 염두에 두라.
\end{definition}

\noindent
두 판본 중 어느 것을 뜻하는지는 문맥에서 분명할 것이다. 혼동의 여지가
있으면 첫째 판본에서는 $M$이 {\it ind-대상}으로서 본질적으로 상수라고
하고, 둘째 판본에서는 {\it pro-대상}으로서 본질적으로 상수라고 하여
구별하겠다. 이 용어는 Remarks \ref{categories-remark-ind-category}와
\ref{categories-remark-pro-category}에서 더 설명한다. 실제로 ``본질적으로 상수인
계''라는 용어를 자주 사용할 텐데, 엄밀히 말하면 이는 유향 집합 위의
계에 대해서만 정의되어 있다.

\begin{definition}
\label{categories-definition-essentially-constant-system}
$\mathcal{C}$를 범주라 하자. 유향계 $(M_i, f_{ii'})$에 대응하는
$M$을 함자 $I \to \mathcal{C}$로 보았을 때 본질적으로 상수인 도표를 정의하면,
이를 {\it 본질적으로 상수인 계}라고 한다. 유향 역계
$(M_i, f_{ii'})$에 대응하는 $M$을 함자 $I^{opp} \to \mathcal{C}$로 보았을 때
본질적으로 상수인 역도표를 정의하면, 이를 {\it 본질적으로 상수인
역계}라고 한다.
\end{definition}

\noindent
$(M_i, f_{ii'})$가 본질적으로 상수인 계이고 사상들 $f_{ii'}$가
단사사상이면, 충분히 큰 모든 $i \leq i'$에 대하여 사상들 $f_{ii'}$는
동형이다. 한편 다음 계를 생각하자.
$$
\mathbf{Z}^2 \to \mathbf{Z}^2 \to \mathbf{Z}^2 \to \ldots
$$
사상들은 $(a, b) \mapsto (a + b, 0)$로 주어진다. 이 계는 값
$\mathbf{Z}$를 가지며 본질적으로 상수이지만, 모든 전이 사상은 핵을 가진다.

\medskip\noindent
본질적으로 상수가 아닌 계의 예는 다음과 같다.
$M = \bigoplus_{n \geq 0} \mathbf{Z}$로 놓고
$S : M \to M$을 이동 작용소
$(a_0, a_1, \ldots) \mapsto (a_1, a_2, \ldots)$로 두자. 이 경우
전이 사상이 $S$인 계 $M \to M \to M \to \ldots$는 여극한이 $0$이고
합성 $0 \to M \to 0$은 항등이지만, 이 계는 본질적으로 상수가 아니다.

\medskip\noindent
다음 보조정리는 정의의 정합성을 확인한다.

\begin{lemma}
\label{categories-lemma-essentially-constant-is-limit-colimit}
$M : \mathcal{I} \to \mathcal{C}$를 도표라 하자. $\mathcal{I}$가
여과이고 $M$이 ind-대상으로서 본질적으로 상수이면
$X = \colim M_i$가 존재하고 $M$은 값 $X$를 가지며 본질적으로 상수이다.
$\mathcal{I}$가 쌍대여과이고 $M$이 pro-대상으로서 본질적으로 상수이면
$X = \lim M_i$가 존재하고 $M$은 값 $X$를 가지며 본질적으로 상수이다.
\end{lemma}

\begin{proof}
생략한다. 정의를 익히기 좋은 연습문제이다.
\end{proof}

\begin{remark}
\label{categories-remark-ind-category}
$\mathcal{C}$를 범주라 하자. $\mathcal{C}$의 {\it ind-대상}들로 이루어진
큰 범주 $\text{Ind-}\mathcal{C}$가 존재한다. 구체적으로
$F : \mathcal{I} \to \mathcal{C}$와
$G : \mathcal{J} \to \mathcal{C}$가 $\mathcal{C}$에서의 여과
도표이면 다음과 같이 정의할 수 있다.
$$
\Mor_{\text{Ind-}\mathcal{C}}(F, G) =
\lim_i \colim_j \Mor_\mathcal{C}(F(i), G(j)).
$$
$X$를 $X$에서의 {\it 상수계}로 보내는 표준 함자
$\mathcal{C} \to \text{Ind-}\mathcal{C}$가 있다. 이는 충만충실한
매장이다. 이 언어로 말하면 도표 $F$가 본질적으로 상수일 필요충분조건은
$F$가 상수계와 동형인 것이다. 이 내용이 필요해지면 여기서 보조정리로
정식화하고 증명하겠다.
\end{remark}

\begin{remark}
\label{categories-remark-pro-category}
$\mathcal{C}$를 범주라 하자. $\mathcal{C}$의 {\it pro-대상}들로 이루어진
큰 범주 $\text{Pro-}\mathcal{C}$가 존재한다. 구체적으로
$F : \mathcal{I} \to \mathcal{C}$와
$G : \mathcal{J} \to \mathcal{C}$가 $\mathcal{C}$에서의 쌍대여과
도표이면 다음과 같이 정의할 수 있다.
$$
\Mor_{\text{Pro-}\mathcal{C}}(F, G) =
\lim_j \colim_i \Mor_\mathcal{C}(F(i), G(j)).
$$
$X$를 $X$에서의 {\it 상수계}로 보내는 표준 함자
$\mathcal{C} \to \text{Pro-}\mathcal{C}$가 있다. 이는 충만충실한
매장이다. 이 언어로 말하면 도표 $F$가 본질적으로 상수일 필요충분조건은
$F$가 상수계와 동형인 것이다. 이 내용이 필요해지면 여기서 보조정리로
정식화하고 증명하겠다.
\end{remark}

\begin{example}
\label{categories-example-pro-morphism-inverse-systems}
$\mathcal{C}$를 범주라 하자. $(X_n)$과 $(Y_n)$을 통상적인 순서가 주어진
$\mathbf{N}$ 위의 $\mathcal{C}$에서의 역계라 하자. 그림으로 쓰면
$$
\ldots \to X_3 \to X_2 \to X_1
\quad\text{and}\quad
\ldots \to Y_3 \to Y_2 \to Y_1
$$
이다. $a : (X_n) \to (Y_n)$를 $\mathcal{C}$의 pro-대상들의 사상이라
하자. $a$는 어떤 데이터인가? 각 $n \in \mathbf{N}$에 대하여 어떤
$m(n)$과 사상 $a_n : X_{m(n)} \to Y_n$이 존재해야 한다. 이 사상들은
다음 의미에서 서로 양립해야 한다. 모든 $n' \geq n$에 대하여
도표
$$
\xymatrix{
X_{m(n, n')} \ar[rr] \ar[d] & & X_{m(n)} \ar[d]^{a_n} \\
X_{m(n')} \ar[r]^{a_{n'}} & Y_{n'} \ar[r] & Y_n
}
$$
가 가환하도록 하는 $m(n', n) \geq m(n'), m(n)$이 존재해야 한다.
$m(n)$을 $\max_{k, l \leq n}\{m(n, k), m(k, l)\}$로 바꾸면
$\ldots \geq m(3) \geq m(2) \geq m(1)$과 가환 도표
$$
\xymatrix{
\ldots \ar[r] &
X_{m(3)} \ar[d]^{a_3} \ar[r] &
X_{m(2)} \ar[d]^{a_2} \ar[r] &
X_{m(1)} \ar[d]^{a_1} \\
\ldots \ar[r] &
Y_3 \ar[r] &
Y_2 \ar[r] &
Y_1
}
$$
를 얻는다. $m' \geq m$인 증가 사상
$m' : \mathbf{N} \to \mathbf{N}$이 주어졌을 때
$a'_i : X_{m'(i)} \to X_{m(i)} \to Y_i$로 놓으면 순서쌍
$(m', a')$는 pro-계의 같은 사상을 정의한다. 거꾸로 위와 같은 두
순서쌍 $(m_1, a_1)$과 $(m_1, a_2)$가 주어졌을 때, 이들이 pro-대상들의
같은 사상을 정의할 필요충분조건은 $a'_1 = a'_2$가 되도록 하는
$m' \geq m_1, m_2$를 찾을 수 있는 것이다.
\end{example}

\begin{remark}
\label{categories-remark-pro-category-copresheaves}
$\mathcal{C}$를 범주라 하자. $F : \mathcal{I} \to \mathcal{C}$와
$G : \mathcal{J} \to \mathcal{C}$를 $\mathcal{C}$에서의 쌍대여과
도표라 하자. 다음과 같이 정의되는 함자
$A, B : \mathcal{C} \to \textit{Sets}$를 생각하자.
$$
A(X) = \colim_i \Mor_\mathcal{C}(F(i), X)
\quad\text{and}\quad
B(X) = \colim_j \Mor_\mathcal{C}(G(j), X)
$$
$F$에서 $G$로 가는 pro-계의 사상은 함자의 변환 $t : B \to A$와
같은 데이터라고 주장한다. 실제로 $t$가 주어지면 $B(G(j))$ 안의
$\text{id}_{G(j)}$의 동치류에 $t$를 적용하여 서로 양립하는 원소들의 계
$\xi_j \in A(G(j)) = \colim_i \Mor_\mathcal{C}(F(i), G(j))$를 얻는다.
이는 Remark \ref{categories-remark-pro-category}에서 정의한
$\text{Pro-}\mathcal{C}$의 사상과 정확히 같다. $F$에서 $G$로 가는
pro-대상의 사상으로부터 변환 $B \to A$를 구성하는 과정은 생략한다.
\end{remark}

\begin{lemma}
\label{categories-lemma-image-essentially-constant}
$\mathcal{C}$를 범주라 하자. $M : \mathcal{I} \to \mathcal{C}$를
여과(각각 쌍대여과) 지표 범주 $\mathcal{I}$를 가진 도표라 하고,
$F : \mathcal{C} \to \mathcal{D}$를 함자라 하자. $M$이 ind-대상
(각각 pro-대상)으로서 본질적으로 상수이면
$F \circ M : \mathcal{I} \to \mathcal{D}$도 그러하다.
\end{lemma}

\begin{proof}
$X$가 $M$의 값이면 정의에서 곧바로 $F(X)$가 $F \circ M$의 값임이
따른다.
\end{proof}

\begin{lemma}
\label{categories-lemma-characterize-essentially-constant-ind}
$\mathcal{C}$를 범주라 하자. $M : \mathcal{I} \to \mathcal{C}$를
여과 지표 범주 $\mathcal{I}$를 가진 도표라 하자. 다음은 서로 동치이다.
\begin{enumerate}
\item $M$은 본질적으로 상수인 ind-대상이다.
\item 모든 $\mathcal{C}$의 $W$에 대하여 사상
$\colim_i \Mor_\mathcal{C}(W, M_i) \to \Mor_\mathcal{C}(W, X)$가
전단사이도록 하는 쌍대원뿔 $(X, \{M_i \to X\}_i)$가 존재한다.
\item $X = \colim_i M_i$가 존재하고, 모든 $\mathcal{C}$의 $W$에
대하여 사상
$\colim_i \Mor_\mathcal{C}(W, M_i) \to \Mor_\mathcal{C}(W, X)$가
전단사이다.
\item 어떤 $i$가 $\mathcal{I}$ 안에 있고 사상 $X \to M_i$가 존재하여,
모든 $\mathcal{C}$의 $W$에 대하여 사상
$\Mor_\mathcal{C}(W, X) \to
\colim_{j \in \mathcal{I}} \Mor_\mathcal{C}(W, M_j)$가 전단사이다.
\end{enumerate}
(2), (3), (4)의 경우 본질적으로 상수인 계의 값은 $X$이다.
\end{lemma}

\begin{proof}
(3)이 (2)를 함의함은 명백하다. (2)를 가정하자. 그러면
$\text{id}_X \in \Mor_\mathcal{C}(X, X)$는 어떤 $i$에 대한 사상
$X \to M_i$에서 온다. 즉 $X \to M_i \to X$는 항등이다. 그러면 모든
$W$에 대하여 두 사상
$$
\Mor_\mathcal{C}(W, X)
\longrightarrow
\colim_i \Mor_\mathcal{C}(W, M_i)
\longrightarrow
\Mor_\mathcal{C}(W, X)
$$
은 전단사이다. 여기서 첫째 사상은 위에서 찾은 사상 $X \to M_i$가
유도하고, 합성은 항등이다. 이는 합성
$$
\colim_i \Mor_\mathcal{C}(W, M_i)
\longrightarrow
\Mor_\mathcal{C}(W, X)
\longrightarrow
\colim_i \Mor_\mathcal{C}(W, M_i)
$$
도 항등이라는 뜻이다. $W = M_j$로 놓고 여극한 안의
$\text{id}_{M_j}$에서 시작하면, 충분히 큰 어떤 $k$에 대하여
$M_j \to X \to M_i \to M_k$가 $M_j \to M_k$와 같음을 알 수 있다.
이는 (1)이 성립함을 증명한다.

\medskip\noindent
(4)를 가정하자. $k$를 $\mathcal{I}$의 대상이라 하자. $W = M_k$로
놓으면, 어떤 $j$와 $\mathcal{I}$의 사상들 $k \to j$, $i \to j$가
존재하여 $M_k \to X \to M_i \to M_j$가 $M_k \to M_j$와 같아지게
하는 유일한 사상 $M_k \to X$가 존재함을 안다. 유일성은 쌍대원뿔
$(X, \{M_k \to X\})$을 얻게 함을 보장한다. 이와 같이 (4)가 (2)를
함의함을 알 수 있다. 일부 세부사항은 생략한다.

\medskip\noindent
(1)이 다른 조건들을 함의하는 증명은 생략한다.
\end{proof}

\begin{lemma}
\label{categories-lemma-characterize-essentially-constant-pro}
$\mathcal{C}$를 범주라 하자. $M : \mathcal{I} \to \mathcal{C}$를
쌍대여과 지표 범주 $\mathcal{I}$를 가진 도표라 하자. 다음은 서로 동치이다.
\begin{enumerate}
\item $M$은 본질적으로 상수인 pro-대상이다.
\item 모든 $\mathcal{C}$의 $W$에 대하여 사상
$\colim_{i \in \mathcal{I}^{opp}} \Mor_\mathcal{C}(M_i, W) \to
\Mor_\mathcal{C}(X, W)$가 전단사이도록 하는 원뿔
$(X, \{X \to M_i\})$가 존재한다.
\item $X = \lim_i M_i$가 존재하고 모든 $\mathcal{C}$의 $W$에 대하여 사상
$\colim_{i \in \mathcal{I}^{opp}} \Mor_\mathcal{C}(M_i, W) \to
\Mor_\mathcal{C}(X, W)$가 전단사이다.
\item 어떤 $i$가 $\mathcal{I}$ 안에 있고 사상 $M_i \to X$가 존재하여,
모든 $\mathcal{C}$의 $W$에 대하여 사상 $\Mor_\mathcal{C}(X, W) \to
\colim_{j \in \mathcal{I}^{opp}} \Mor_\mathcal{C}(M_j, W)$가 전단사이다.
\end{enumerate}
(2), (3), (4)의 경우 본질적으로 상수인 계의 값은 $X$이다.
\end{lemma}

\begin{proof}
이 보조정리는 Lemma \ref{categories-lemma-characterize-essentially-constant-ind}의
쌍대이다.
\end{proof}

\begin{lemma}
\label{categories-lemma-cofinal-essentially-constant}
$\mathcal{C}$를 범주라 하자. $H : \mathcal{I} \to \mathcal{J}$를 여과
지표 범주들의 함자라 하자. $H$가 공종이면, 임의의 도표
$M : \mathcal{J} \to \mathcal{C}$가 본질적으로 상수일 필요충분조건은
$M \circ H$가 본질적으로 상수인 것이다.
\end{lemma}

\begin{proof}
이는 Lemmas \ref{categories-lemma-characterize-essentially-constant-ind}와
\ref{categories-lemma-cofinal}에서 형식적으로 따른다.
\end{proof}

\begin{lemma}
\label{categories-lemma-essentially-constant-over-product}
$\mathcal{I}$와 $\mathcal{J}$를 여과 범주라 하고 사영을
$p : \mathcal{I} \times \mathcal{J} \to \mathcal{J}$로 표기하자.
그러면 $\mathcal{I} \times \mathcal{J}$는 여과 범주이고, 도표
$M : \mathcal{J} \to \mathcal{C}$가 본질적으로 상수일 필요충분조건은
$M \circ p : \mathcal{I} \times \mathcal{J} \to \mathcal{C}$가
본질적으로 상수인 것이다.
\end{lemma}

\begin{proof}
$\mathcal{I} \times \mathcal{J}$가 여과 범주라는 확인은 생략한다.
$p$가 공종이므로(확인은 생략한다), 동치는 Lemma
\ref{categories-lemma-cofinal-essentially-constant}에서 따른다.
\end{proof}

\begin{lemma}
\label{categories-lemma-initial-essentially-constant}
$\mathcal{C}$를 범주라 하자. $H : \mathcal{I} \to \mathcal{J}$를
쌍대여과 지표 범주들의 함자라 하자. $H$가 초기이면, 임의의 도표
$M : \mathcal{J} \to \mathcal{C}$가 본질적으로 상수일 필요충분조건은
$M \circ H$가 본질적으로 상수인 것이다.
\end{lemma}

\begin{proof}
이는 Lemmas \ref{categories-lemma-characterize-essentially-constant-pro},
\ref{categories-lemma-initial}, \ref{categories-lemma-cofinal} 및 $\mathcal{I}$가
$\mathcal{J}$에서 초기이면 $\mathcal{I}^{opp}$가
$\mathcal{J}^{opp}$에서 공종이라는 사실에서 형식적으로 따른다.
\end{proof}





\section{완전 함자}
\label{categories-section-exact-functor}

\noindent
이 절에서는 완전 함자를 정의한다.

\begin{definition}
\label{categories-definition-exact}
$F : \mathcal{A} \to \mathcal{B}$를 함자라 하자.
\begin{enumerate}
\item $\mathcal{A}$에 모든 유한 극한이 존재한다고 가정하자.
$F$가 모든 유한 극한과 가환하면 이를 {\it 좌완전}이라고 한다.
\item $\mathcal{A}$에 모든 유한 여극한이 존재한다고 가정하자.
$F$가 모든 유한 여극한과 가환하면 이를 {\it 우완전}이라고 한다.
\item $F$가 좌완전이고 우완전이면 이를 {\it 완전}이라고 한다.
\end{enumerate}
\end{definition}

\begin{lemma}
\label{categories-lemma-characterize-left-exact}
$F : \mathcal{A} \to \mathcal{B}$를 함자라 하자.
$\mathcal{A}$에 모든 유한 극한이 존재한다고 가정하자.
Lemma \ref{categories-lemma-finite-limits-exist}를 보라. 다음은 서로 동치이다.
\begin{enumerate}
\item $F$는 좌완전이다.
\item $F$는 유한 곱 및 등화자와 가환한다.
\item $F$는 $\mathcal{A}$의 끝 대상을 $\mathcal{B}$의 끝 대상으로
보내고 섬유곱과 가환한다.
\end{enumerate}
\end{lemma}

\begin{proof}
Lemma \ref{categories-lemma-limits-products-equalizers}에 의해 (2)는 (1)을 함의한다.
(3)을 가정하자. 끝 대상 위의 섬유곱은 곱이다.
$a, b : A \to B$가 $\mathcal{A}$의 사상들이면 $a, b$의 등화자는
$$
(A \times_{a, B, b} A)\times_{(\text{pr}_1, \text{pr}_2), A \times A, \Delta} A.
$$
이다. 따라서 (3)은 (2)를 함의한다. 끝으로 공극한은 끝 대상이고
섬유곱은 극한이므로 (1)은 (3)을 함의한다.
\end{proof}

\begin{lemma}
\label{categories-lemma-characterize-right-exact}
$F : \mathcal{A} \to \mathcal{B}$를 함자라 하자.
$\mathcal{A}$에 모든 유한 여극한이 존재한다고 가정하자.
Lemma \ref{categories-lemma-colimits-exist}를 보라. 다음은 서로 동치이다.
\begin{enumerate}
\item $F$는 우완전이다.
\item $F$는 유한 쌍대곱 및 쌍대등화자와 가환한다.
\item $F$는 $\mathcal{A}$의 초기 대상을 $\mathcal{B}$의 초기 대상으로
보내고 밂과 가환한다.
\end{enumerate}
\end{lemma}

\begin{proof}
Lemma \ref{categories-lemma-characterize-left-exact}의 쌍대이다.
\end{proof}




\section{수반 함자}
\label{categories-section-adjoint}

\begin{definition}
\label{categories-definition-adjoint}
$\mathcal{C}$, $\mathcal{D}$를 범주라 하자.
$u : \mathcal{C} \to \mathcal{D}$와
$v : \mathcal{D} \to \mathcal{C}$를 함자라 하자.
$X \in \Ob(\mathcal{C})$와 $Y \in \Ob(\mathcal{D})$에 관해 함자적인 전단사
$$
\Mor_\mathcal{D}(u(X), Y)
\longrightarrow
\Mor_\mathcal{C}(X, v(Y))
$$
가 존재하면 $u$를 $v$의 {\it 왼쪽 수반}이라고 하거나,
$v$를 $u$의 {\it 오른쪽 수반}이라고 한다.
\end{definition}

\noindent
다시 말해 이는 함자 $\mathcal{C}^{opp} \times \mathcal{D} \to \textit{Sets}$
사이에 $\Mor_\mathcal{D}(u(-), -)$에서 $\Mor_\mathcal{C}(-, v(-))$로 가는
{\it 주어진} 동형이 있다는 뜻이다. $\mathcal{C}$의 임의의 대상 $X$에 대해
$\text{id}_{u(X)}$에 대응하는 사상 $X \to v(u(X))$를 얻는다. 마찬가지로
$\mathcal{D}$의 임의의 대상 $Y$에 대해 $\text{id}_{v(Y)}$에 대응하는 사상
$u(v(Y)) \to Y$를 얻는다. 이 사상들을 {\it 수반 사상}이라고 한다.
수반 사상은 $X$와 $Y$에 관해 함자적이므로 함자들의 사상
$$
\eta : \text{id}_\mathcal{C} \to v \circ u\quad (\text{unit})
\quad\text{and}\quad
\epsilon : u \circ v \to \text{id}_\mathcal{D}\quad (\text{counit}).
$$
을 얻는다. 또한 $\alpha : u(X) \to Y$와
$\beta : X \to v(Y)$가 사상들이면 다음은 서로 동치이다.
\begin{enumerate}
\item $\alpha$와 $\beta$는 정의의 전단사에 의해 서로 대응한다.
\item $\beta$는 합성 $X \to v(u(X)) \xrightarrow{v(\alpha)} v(Y)$이다.
\item $\alpha$는 합성 $u(X) \xrightarrow{u(\beta)} u(v(Y)) \to Y$이다.
\end{enumerate}
이와 같이 수반 함자의 개념을 수반 사상으로 다시 정식화할 수 있다.

\begin{lemma}
\label{categories-lemma-adjoint-exists}
$u : \mathcal{C} \to \mathcal{D}$를 범주 사이의 함자라 하자.
각 $y \in \Ob(\mathcal{D})$에 대하여 함자
$x \mapsto \Mor_\mathcal{D}(u(x), y)$가 표현 가능하면 $u$는 오른쪽
수반을 갖는다.
\end{lemma}

\begin{proof}
각 $y$에 대하여 대상 $v(y)$와 함자들의 동형
$\Mor_\mathcal{C}(-, v(y)) \to \Mor_\mathcal{D}(u(-), y)$를 택하자.
Yoneda 보조정리(Lemma \ref{categories-lemma-yoneda})에 의해 임의의 사상
$g : y \to y'$에 대하여 함자들의 변환
$$
\Mor_\mathcal{C}(-, v(y)) \to \Mor_\mathcal{D}(u(-), y) \to
\Mor_\mathcal{D}(u(-), y') \to \Mor_\mathcal{C}(-, v(y'))
$$
은 유일한 사상 $v(g) : v(y) \to v(y')$에 대응한다.
$v$가 함자이고 $u$의 오른쪽 수반이라는 확인은 생략한다.
\end{proof}

\begin{lemma}
\label{categories-lemma-left-adjoint-composed-fully-faithful}
\begin{reference}
Bhargav Bhatt, private communication.
\end{reference}
Definition \ref{categories-definition-adjoint}에서와 같이 $u$를 $v$의 왼쪽 수반이라 하자.
\begin{enumerate}
\item $v \circ u$가 충만충실하면 $u$는 충만충실하다.
\item $u \circ v$가 충만충실하면 $v$는 충만충실하다.
\end{enumerate}
\end{lemma}

\begin{proof}
(2)를 증명하자. $u \circ v$가 충만충실하다고 가정하자.
$X$, $Y$가 $\mathcal{D}$ 안에 있다고 하자. 그러면 자연스러운 합성 사상
$$
\Mor(X,Y) \to \Mor(v(X),v(Y)) \to \Mor(u(v(X)), u(v(Y)))
$$
은 전단사이므로 $v$는 적어도 충실하다. 충만충실함을 보이려면 위의 둘째
사상이 단사임을 보여야 한다. 그러나 $u$와 $v$ 사이의 수반에 따르면
$$
\Mor(v(X), v(Y)) \to \Mor(u(v(X)), u(v(Y))) \to \Mor(u(v(X)), Y)
$$
는 전단사이다. 여기서 첫째 사상은 자연스러운 사상이고 둘째 사상은 수반의
쌍대단위 $u(v(Y)) \to Y$에서 온다. 따라서
$\Mor(v(X), v(Y)) \to \Mor(u(v(X)), u(v(Y)))$도 원하는 대로
단사이다. (1)의 증명은 이것의 쌍대이다.
\end{proof}

\begin{lemma}
\label{categories-lemma-adjoint-fully-faithful}
Definition \ref{categories-definition-adjoint}에서와 같이 $u$를 $v$의 왼쪽 수반이라 하자.
그러면
\begin{enumerate}
\item $u$가 충만충실하다 $\Leftrightarrow$ $\text{id} \cong v \circ u$
$\Leftrightarrow$ $\eta : \text{id} \to v \circ u$가 동형이다.
\item $v$가 충만충실하다 $\Leftrightarrow$
$u \circ v \cong \text{id}$ $\Leftrightarrow$
$\epsilon : u \circ v \to \text{id}$가 동형이다.
\end{enumerate}
\end{lemma}

\begin{proof}
(1)을 증명하자.
$u$가 충만충실하다고 가정하자. $\eta_X : X \to v(u(X))$가 동형임을 보이겠다.
임의의 사상 $X' \to v(u(X))$를 택하자. 수반성에 의해 이는 사상
$u(X') \to u(X)$에 대응한다. $u$의 충만충실성에 의해 이는 유일한 사상
$X' \to X$에 대응한다. 따라서 $\eta_X$를 뒤에서 합성하는 것이 전단사
$\Mor(X', X) \to \Mor(X', v(u(X)))$를 정의한다. 그러므로 $\eta_X$는
동형이다. 함자들의 동형 $\text{id} \cong v \circ u$가 존재하면
$v \circ u$는 충만충실하다. Lemma
\ref{categories-lemma-left-adjoint-composed-fully-faithful}에 의해 $u$가 충만충실함을 안다.
위에서 이는 $\eta$가 동형임을 함의한다. 따라서 모든 $3$ 조건은 동치이다
(그리고 이 조건들은 $v \circ u$가 충만충실하다는 조건과도 동치이다).

\medskip\noindent
(2)는 (1)의 쌍대이다.
\end{proof}

\begin{lemma}
\label{categories-lemma-adjoint-exact}
Definition \ref{categories-definition-adjoint}에서와 같이 $u$를 $v$의 왼쪽 수반이라 하자.
\begin{enumerate}
\item $M : \mathcal{I} \to \mathcal{C}$가 도표이고
$\colim_\mathcal{I} M$이 $\mathcal{C}$에 존재한다고 가정하자. 그러면
$u(\colim_\mathcal{I} M) = \colim_\mathcal{I} u \circ M$이다.
다시 말해 $u$는 (표현 가능한) 여극한과 가환한다.
\item $M : \mathcal{I} \to \mathcal{D}$가 도표이고
$\lim_\mathcal{I} M$이 $\mathcal{D}$에 존재한다고 가정하자. 그러면
$v(\lim_\mathcal{I} M) = \lim_\mathcal{I} v \circ M$이다.
다시 말해 $v$는 표현 가능한 극한과 가환한다.
\end{enumerate}
\end{lemma}

\begin{proof}
여극한에서 대상으로 가는 사상은 그 극한의 구성요소들에서 그 대상으로 가는
서로 양립하는 사상들의 계와 같다. Remark \ref{categories-remark-limit-colim}을 보라.
따라서
$$
\begin{matrix}
\Mor_\mathcal{D}(u(\colim_{i \in \mathcal{I}} M_i), Y) &
= & \Mor_\mathcal{C}(\colim_{i \in \mathcal{I}} M_i, v(Y)) \\
& = &
\lim_{i \in \mathcal{I}^{opp}}
\Mor_\mathcal{C}(M_i, v(Y)) \\
& = &
\lim_{i \in \mathcal{I}^{opp}}
\Mor_\mathcal{D}(u(M_i), Y)
\end{matrix}
$$
는 $u(\colim_{i \in \mathcal{I}} M_i)$가 우리가 구하는 여극한임을 증명한다.
다른 명제에도 같은 논증을 적용한다.
\end{proof}

\begin{lemma}
\label{categories-lemma-exact-adjoint}
Definition \ref{categories-definition-adjoint}에서와 같이 $u$를 $v$의 왼쪽 수반이라 하자.
\begin{enumerate}
\item $\mathcal{C}$가 유한 여극한을 가지면 $u$는 우완전이다.
\item $\mathcal{D}$가 유한 극한을 가지면 $v$는 좌완전이다.
\end{enumerate}
\end{lemma}

\begin{proof}
정의와 Lemma \ref{categories-lemma-adjoint-exact}에서 명백하다.
\end{proof}

\begin{lemma}
\label{categories-lemma-unit-counit-relations}
$u : \mathcal{C} \to \mathcal{D}$를 함자
$v : \mathcal{D} \to \mathcal{C}$의 왼쪽 수반이라 하자.
$\eta_X : X \to v(u(X))$를 단위라 하고
$\epsilon_Y : u(v(Y)) \to Y$를 쌍대단위라 하자. 그러면
$$
u(X) \xrightarrow{u(\eta_X)} u(v(u(X))
\xrightarrow{\epsilon_{u(X)}} u(X)
\quad\text{and}\quad
v(Y) \xrightarrow{\eta_{v(Y)}} v(u(v(Y))) \xrightarrow{v(\epsilon_Y)}
v(Y)
$$
는 항등사상이다.
\end{lemma}

\begin{proof}
생략한다.
\end{proof}

\begin{lemma}
\label{categories-lemma-transformation-between-functors-and-adjoints}
$u_1, u_2 : \mathcal{C} \to \mathcal{D}$를 오른쪽 수반
$v_1, v_2 : \mathcal{D} \to \mathcal{C}$를 갖는 함자들이라 하자.
$\beta : u_2 \to u_1$를 함자들의 변환이라 하자.
$\beta^\vee : v_1 \to v_2$를 이에 대응하는 수반 함자들의 변환이라 하자.
그러면
$$
\xymatrix{
u_2 \circ v_1 \ar[r]_\beta \ar[d]_{\beta^\vee} &
u_1 \circ v_1 \ar[d] \\
u_2 \circ v_2 \ar[r] & \text{id}
}
$$
는 가환한다. 여기서 표지 없는 화살표들은 쌍대단위 변환이다.
\end{lemma}

\begin{proof}
이는 $\beta^\vee_D : v_1D \to v_2D$가 유도 사상
$\Mor(C, v_1D) \to \Mor(C, v_2D)$이
$\beta_C : u_2C \to u_1C$가 유도하는 사상
$\Mor(u_1C, D) \to \Mor(u_2C, D)$이 되게 하는 유일한 사상이기 때문이다.
곧 이는 $\beta_{v_1 D}$가 유도하는 사상
$$
\Mor(u_1 v_1 D, D') \to \Mor(u_2 v_1 D, D')
$$
이 $\beta^\vee_{D'}$가 유도하는 사상
$$
\Mor(v_1 D, v_1 D') \to \Mor(v_1 D, v_2 D')
$$
과 같다는 뜻이다. $D' = D$로 놓으면, $\beta_{v_1D}$를 앞에서 합성한
쌍대단위 $u_1 v_1 D \to D$가 수반 아래에서 $\beta^\vee_D$에 대응함을
얻는다. 이는 도표를 $D$에서 평가했을 때 정확히 가환한다는 뜻이다.
\end{proof}

\begin{lemma}
\label{categories-lemma-compose-counits}
$\mathcal{A}$, $\mathcal{B}$, $\mathcal{C}$를 범주라 하자.
$v : \mathcal{A} \to \mathcal{B}$와
$v' : \mathcal{B} \to \mathcal{C}$를 각각 왼쪽 수반 $u$와 $u'$를 갖는
함자들이라 하자. 그러면
\begin{enumerate}
\item 함자 $v'' = v' \circ v$는 $u'' = u \circ u'$와 같은 왼쪽 수반을 갖는다.
\item $\mathcal{A}$ 안의 $X$가 주어지면
\begin{equation}
\label{categories-equation-compose-counits}
\epsilon_X^v \circ u(\epsilon^{v'}_{v(X)}) = \epsilon^{v''}_X :
u''(v''(X)) \to X
\end{equation}
이다. 여기서 $\epsilon$은 수반들의 쌍대단위이다.
\end{enumerate}
\end{lemma}

\begin{proof}
(2)의 식을 풀어 쓰자. 그러면 (1)도 곧바로 증명된다. 먼저 순서쌍
$(u, v)$와 $(u', v')$에 대한 수반들의 쌍대단위는 사상
$\epsilon_X^v : u(v(X)) \to X$와
$\epsilon_Y^{v'} : u'(v'(Y)) \to Y$이다. Definition
\ref{categories-definition-adjoint} 뒤의 논의를 보라. (1)에서와 같이 $u''$와 $v''$를
놓고 모든 것을 풀어 쓰면
$$
u''(v''(X)) = u(u'(v'(v(X)))) \xrightarrow{u(\epsilon_{v(X)}^{v'})}
u(v(X)) \xrightarrow{\epsilon_X^v} X
$$
를 얻으며, 이는 (\ref{categories-equation-compose-counits})의 좌변에 있는 사상이다.
우선 이를 $\epsilon_X^{v''}$로 표기하자. 이것이 수반쌍
$(u'', v'')$의 쌍대단위임을 보이려면 $\mathcal{C}$ 안의 $Z$가 주어졌을 때
사상 $\beta : Z \to v''(X)$를
$\alpha = \epsilon_X^{v''} \circ u''(\beta) : u''(Z) \to X$로 보내는 규칙이
사상들의 집합 사이의 전단사임을 보여야 한다. 이는 $\beta$를
$\epsilon_{v(X)}^{v'} \circ u'(\beta)$로 보내는 규칙과, 다시 이를
$\epsilon_X^v \circ u(\epsilon_{v(X)}^{v'} \circ u'(\beta))$로 보내는 규칙의
합성이기 때문에 참이다. 첫째는 $(u', v')$에 대한 가정에 의해 전단사이고,
둘째는 $(u, v)$에 대한 가정에 의해 전단사이다.
\end{proof}






\section{표현 가능성의 판정법}
\label{categories-section-representable}

\noindent
다음 보조정리는 큰 범주에서 보편 대상을 구성할 때 흔히 유용하다.
Remark \ref{categories-remark-how-to-use-it}의 논의를 보라.

\begin{lemma}
\label{categories-lemma-a-version-of-brown}
$\mathcal{C}$를 극한을 갖는 큰\footnote{Remark
\ref{categories-remark-big-categories}를 보라.} 범주라 하자.
$F : \mathcal{C} \to \textit{Sets}$를 함자라 하자. 다음을 가정하자.
\begin{enumerate}
\item $F$는 극한과 가환한다.
\item $\mathcal{C}$의 대상들의 족 $\{x_i\}_{i \in I}$와 각
$i \in I$에 대한 원소 $f_i \in F(x_i)$가 존재하여,
$y \in \Ob(\mathcal{C})$ 및 $g \in F(y)$가 주어지면 어떤 $i$와 사상
$\varphi : x_i \to y$가 존재하고 $F(\varphi)(f_i) = g$이다.
\end{enumerate}
그러면 $F$는 표현 가능하다. 즉 $y$에 관해 함자적으로
$$
F(y) = \Mor_\mathcal{C}(x, y)
$$
가 되게 하는 $\mathcal{C}$의 대상 $x$가 존재한다.
\end{lemma}

\begin{proof}
$\mathcal{I}$를 대상이 순서쌍 $(x_i, f_i)$이고 사상
$(x_i, f_i) \to (x_{i'}, f_{i'})$이 $F(\varphi)(f_i) = f_{i'}$를 만족하는
$\mathcal{C}$의 사상 $\varphi : x_i \to x_{i'}$인 범주라 하자.
$$
x = \lim_{(x_i, f_i) \in \mathcal{I}} x_i
$$
로 놓자(이는 우리가 찾는 $x$가 아니다. 아래를 보라). 가정에 의해 이 극한은
존재한다. $F$가 극한과 가환하므로
$$
F(x) = \lim_{(x_i, f_i) \in \mathcal{I}} F(x_i).
$$
이다. 따라서 $F$를 사영사상 $x \to x_i$에 적용한 사상 아래에서
$f_i \in F(x_i)$로 가는 보편 원소 $f \in F(x)$가 존재한다.
$f$를 사용하면 함자들의 변환
$$
\xi : \Mor_\mathcal{C}(x, - ) \longrightarrow F(-)
$$
을 얻는다. Section \ref{categories-section-opposite}를 보라. $y$를 $\mathcal{C}$의
임의의 대상이라 하고 $g \in F(y)$라 하자. 가정에 의해 가능하므로
$f_i$가 $g$로 가게 하는 $x_i \to y$를 택하자. 그러면 $F$를 사상들
$$
x \longrightarrow x_i \longrightarrow y
$$
에 적용하면 $f$가 $g$로 간다(첫째 사상은 $x$를 정의하는 극한의
사영사상이다). 따라서 변환 $\xi$는 전사이다.

\medskip\noindent
$F$를 표현하는 대상을 찾기 위하여, $F(\varphi)(f) = f$를 만족하는 모든
자기사상 $\varphi : x \to x$의 등화자를 $e : x' \to x$라 하자.
$F$가 극한과 가환하므로 등화자와 가환하며, $F(x)$의 $f$로 가는
$f' \in F(x')$가 존재함을 안다. $\xi$가 전사이고 $f'$가 $f$로 가므로
$\xi' : \Mor_\mathcal{C}(x', -) \to F(-)$도 전사이다. 끝으로 사상
$a, b : x' \to y$가 $F(a)(f') = F(b)(f')$를 만족한다고 가정하자.
$a = b$임을 보여야 한다. 그 등화자를 $e' : x'' \to x'$라 하자.
다시 $f'$로 가는 $f'' \in F(x'')$를 얻는다.
$F(\psi)(f) = f''$를 만족하는 사상 $\psi : x \to x''$를 택하자.
그러면 $e \circ e' \circ \psi : x \to x$는
$F(e \circ e' \circ \psi)(f) = f$를 만족하는 사상이다. 따라서
$e \circ e' \circ \psi \circ e = e$이다. $e$가 단사사상이므로
$e'$는 전사사상이고, 원하는 대로 $a = b$이다.
\end{proof}

\begin{remark}
\label{categories-remark-how-to-use-it}
위의 보조정리는 어떤 대상 위의 자유 대상을 구성하는 데 흔히 사용된다.
예를 들면 집합 위의 자유 아벨 군, 집합 위의 자유군 등이 있다.
집합 $E$ 위의 자유군의 경우에는 함자
$$
F : \textit{Groups} \to \textit{Sets},\quad
G \longmapsto \text{Map}(E, G)
$$
를 생각하는 것이 그 발상이다. 이 함자는 극한과 가환한다. 대상들의 족으로는
기수도가 $\max(\aleph_0, |E|)$ 이하인 군들 $G_i$와 집합 사상
$E \to G_i$로 이루어진 족 $E \to G_i$를 택할 수 있으며, 이와 같은 구조의
각 동형류가 적어도 한 번 나타나게 한다. 실제로 $E \to G$가 $E$에서 군
$G$로 가는 사상이면 그 상이 생성하는 부분군 $G'$의 기수도는
$\max(\aleph_0, |E|)$ 이하이다. 보조정리에 의해 함자는 표현 가능하므로
$\Mor_{\textit{Groups}}(F_E, G) = \text{Map}(E, G)$가 되는 군 $F_E$가 존재한다.
특히 $F_E$의 항등사상은 사상 $E \to F_E$에 대응하며, 아무 관계도 부과하지
않고 그 상이 $F_E$를 생성함을 보일 수 있다.

\medskip\noindent
또 하나의 전형적인 응용은 극한의 존재를 알고 있을 때 이 보조정리로 여극한을
구성하는 것이다. Topology, Lemma \ref{topology-lemma-limits}에 의해 극한을
갖는 위상공간의 범주를 사용하여 이를 설명하자. 함자
$\mathcal{I} \to \textit{Top}$, $i \mapsto X_i$가 주어졌다고 하자.
그러면
$$
F : \textit{Top} \longrightarrow \textit{Sets},\quad
Y \longmapsto \lim_\mathcal{I} \Mor_{\textit{Top}}(X_i, Y)
$$
를 생각할 수 있다. 이 함자는 극한과 가환한다. 또한 임의의 위상공간
$Y$와 $F(Y)$의 원소 $(\varphi_i : X_i \to Y)$가 주어지면,
기수도가 $|\coprod X_i|$ 이하인 부분공간 $Y' \subset Y$가 존재하고,
사상들 $\varphi_i$는 $Y'$로 간다. 실제로 $\varphi_i$들의 상의 합집합에 유도된
위상을 취하면 된다. 따라서 보조정리의 가정들이 만족됨은 명백하며, 함자
$F$를 표현하는 위상공간 $X$를 얻는다. 이는 정확히 $X$가 도표
$i \mapsto X_i$의 여극한이라는 뜻이다.
\end{remark}

\begin{theorem}[수반 함자 정리]
\label{categories-theorem-adjoint-functor}
$G : \mathcal{C} \to \mathcal{D}$를 큰 범주들의 함자라 하자.
$\mathcal{C}$가 극한을 갖고 $G$가 그 극한들과 가환한다고 가정하자.
또한 $\mathcal{D}$의 모든 대상 $y$에 대하여 $x_i \in \Ob(\mathcal{C})$,
$f_i \in \Mor_\mathcal{D}(y, G(x_i))$인 순서쌍들의 집합
$(x_i, f_i)_{i \in I}$가 존재하여, $x \in \Ob(\mathcal{C})$,
$f \in \Mor_\mathcal{D}(y, G(x))$인 임의의 순서쌍 $(x, f)$에 대해 어떤
$i$와 사상 $h : x_i \to x$가 존재하고 $f = G(h) \circ f_i$라고 가정하자.
그러면 $G$는 왼쪽 수반 $F$를 갖는다.
\end{theorem}

\begin{proof}
가정에 의해 $\mathcal{D}$의 모든 대상 $y$에 대하여 함자
$x \mapsto \Mor_\mathcal{D}(y, G(x))$는 Lemma
\ref{categories-lemma-a-version-of-brown}의 가정들을 만족한다. 따라서 어떤 대상으로
표현되며 이를 $F(y)$라 하자. Yoneda 보조정리(Lemma \ref{categories-lemma-yoneda})를
적용하면 규칙 $y \mapsto F(y)$가 함자가 되고, 구성에 의해 $G$의 수반이다.
세부사항은 생략한다.
\end{proof}





\section{범주론적 콤팩트 대상}
\label{categories-section-compact}

\noindent
범주의 ``작은'' 대상에 대해 조금 살펴본다.

\begin{definition}
\label{categories-definition-compact-object}
$\mathcal{C}$를 큰\footnote{Remark \ref{categories-remark-big-categories}를 보라.}
범주라 하자. $\mathcal{C}$의 대상 $X$가 $\colim_i M_i$가 존재하는 모든
여과 도표 $M : \mathcal{I} \to \mathcal{C}$에 대하여
$$
\Mor_\mathcal{C}(X, \colim_i M_i) =
\colim_i \Mor_\mathcal{C}(X, M_i)
$$
를 만족하면 이를 {\it 범주론적으로 콤팩트}하다고 한다.
\end{definition}

\noindent
흔히 이 정의는 $\mathcal{C}$가 모든 여과 여극한을 갖는다는 가정 아래에서만
주어진다.

\begin{lemma}
\label{categories-lemma-extend-functor-by-colim}
$\mathcal{C}$와 $\mathcal{D}$를 여과 여극한을 갖는 큰 범주라 하자.
$\mathcal{C}' \subset \mathcal{C}$를 $\mathcal{C}$의 범주론적으로 콤팩트한
대상들로 이루어진 작은 충만한 부분범주라 하고, $\mathcal{C}$의 모든 대상이
$\mathcal{C}'$의 대상들의 여과 여극한이라고 하자. 그러면 모든 함자
$F' : \mathcal{C}' \to \mathcal{D}$는 여과 여극한과 가환하는 유일한 확장
$F : \mathcal{C} \to \mathcal{D}$를 갖는다.
\end{lemma}

\begin{proof}
$\mathcal{C}$의 모든 대상 $X$에 대하여 $X$를
$X_i \in \Ob(\mathcal{C}')$인 여과 여극한
$X = \colim X_i$로 쓸 수 있다. 이때 $\mathcal{D}$에서
$$
F(X) = \colim F'(X_i)
$$
로 놓는다. 아래에서 이 구성이 $X$의 여극한 표시의 선택에 의존하지 않음을
보이겠다.

\medskip\noindent
$\mathcal{C}$의 사상 $\alpha : X \to Y$가 주어지고,
$X = \colim_{i \in I} X_i$와 $Y = \colim_{j \in J} Y_i$가
$\mathcal{C}'$의 대상들의 여과 여극한으로 쓰였다고 하자. 각 $i \in I$에
대하여 $X_i$는 $\mathcal{C}$의 범주론적으로 콤팩트한 대상이므로 어떤
$j \in J$와 가환 도표
$$
\xymatrix{
X_i \ar[r] \ar[d] & X \ar[d]^\alpha \\
Y_j \ar[r] & Y
}
$$
를 찾을 수 있다. 그러면 사상 $F'(X_i) \to F'(Y_j) \to F(Y)$를 얻는다.
여기서 둘째 사상은 $F(Y) = \colim F'(Y_j)$로 들어가는 쌍대사영이다.
화살표 $\beta_i : F'(X_i) \to F(Y)$는 $j$의 선택에 의존하지 않는다.
$i \leq i'$에 대하여 합성
$$
F'(X_i) \to F'(X_{i'}) \xrightarrow{\beta_{i'}} F(Y)
$$
은 $\beta_i$와 같다. 따라서 여극한의 보편 성질에 의해 잘 정의된 화살표
$$
F(\alpha) : F(X) = \colim F(X_i) \to F(Y)
$$
를 얻는다. $\alpha' : Y \to Z$가 $\mathcal{C}$의 또 다른 사상이고
$Z = \colim Z_k$도 $\mathcal{C}'$의 대상들의 여과 여극한으로 쓰였다고 하자.
그러면 유도된 사상들 $F(\alpha) : F(X) \to F(Y)$와
$F(\alpha') : F(Y) \to F(Z)$의 합성이 $F(\alpha' \circ \alpha)$임을
보이는 것은 즐거운 연습문제이다. 세부사항은 생략한다.

\medskip\noindent
특히 계들의 여과 여극한으로 주어진 두 표시 $X = \colim X_i$와
$X = \colim X'_{i'}$가 $\mathcal{C}'$에 있다고 하자. 그러면 서로 역인 화살표
$\colim F'(X_i) \to \colim F'(X'_{i'})$와
$\colim F'(X'_{i'}) \to \colim F'(X_i)$를 얻는다. 다시 말해 값 $F(X)$는
$X$를 $\mathcal{C}'$의 대상들의 여과 여극한으로 나타내는 표시의 선택과
무관하게 잘 정의된다. 앞 단락에서 논의한 $F$의 함자성과 합치면 $F$가
함자임을 안다. 또한 $X \in \Ob(\mathcal{C}')$이면 $F(X) = F'(X)$임은
명백하다.

\medskip\noindent
보조정리의 유일성 명제는 $F$가 여과 여극한과 가환함을 보인다면 명백하다
(그렇지 않으면 이 명제가 의미가 없기 때문이다). 이를 보이기 위해
$X = \colim_{\lambda \in \Lambda} X_\lambda$가 $\mathcal{C}$의 여과
여극한이라고 하자. $F$가 함자이므로 사상
$$
\colim_\lambda F(X_\lambda) \longrightarrow F(X)
$$
을 얻는다. 한편 $X = \colim X_i$를 $\mathcal{C}'$의 대상들의 여과
여극한으로 쓰자. 위와 같이 각 $i \in I$에 대하여 어떤
$\lambda \in \Lambda$와 가환 도표
$$
\xymatrix{
X_i \ar[rr] \ar[rd] & & X_\lambda \ar[ld] \\
& X
}
$$
를 택할 수 있다. 위와 같이 이는 전이 사상들과 양립하는 잘 정의된 사상
$F'(X_i) \to \colim_\lambda F(X_\lambda)$를 결정하며, 따라서 사상
$$
F(X) = \colim_i F'(X_i) \longrightarrow \colim_\lambda F(X_\lambda)
$$
을 결정한다. 이 사상은 위 사상의 역이다(세부사항은 생략한다). 따라서 원하는
대로 $F(X) = \colim_\lambda F(X_\lambda)$이다.
\end{proof}







\section{범주의 국소화}
\label{categories-section-localization}

\noindent
이 절의 기본 발상은 범주 $\mathcal{C}$와 화살표들의 집합 $S$가 주어졌을 때,
$S$의 모든 원소가 $S^{-1}\mathcal{C}$에서 가역이 되고 이 성질을 가진 모든
함자 중에서 $F$가 보편적이도록 함자
$F : \mathcal{C} \to S^{-1}\mathcal{C}$를 구성하는 것이다. 이 절의 참고문헌은
\cite[Chapter I, Section 2]{GZ}와
\cite[Chapter II, Section 2]{Verdier}이다.

\begin{definition}
\label{categories-definition-multiplicative-system}
$\mathcal{C}$를 범주라 하자. $\mathcal{C}$의 화살표들의 집합 $S$가 다음
성질들을 가지면 이를 {\it 왼쪽 곱셈계}라고 한다.
\begin{enumerate}
\item[LMS1] $\mathcal{C}$의 모든 대상의 항등사상은 $S$에 속하고,
$S$의 합성 가능한 두 원소의 합성도 $S$에 속한다.
\item[LMS2] $t \in S$인 모든 실선 도표
$$
\xymatrix{
X \ar[d]_t \ar[r]_g & Y \ar@{..>}[d]^s \\
Z \ar@{..>}[r]^f & W
}
$$
를 $s \in S$인 가환 점선 사각형으로 완성할 수 있다.
\item[LMS3] 사상들의 모든 쌍 $f, g : X \to Y$와, 공역이 $X$이고
$f \circ t = g \circ t$를 만족하는 $t \in S$에 대하여, 정의역이 $Y$이고
$s \circ f = s \circ g$를 만족하는 $s \in S$가 존재한다.
\end{enumerate}
$\mathcal{C}$의 화살표들의 집합 $S$가 다음 성질들을 가지면 이를
{\it 오른쪽 곱셈계}라고 한다.
\begin{enumerate}
\item[RMS1] $\mathcal{C}$의 모든 대상의 항등사상은 $S$에 속하고,
$S$의 합성 가능한 두 원소의 합성도 $S$에 속한다.
\item[RMS2] $s \in S$인 모든 실선 도표
$$
\xymatrix{
X \ar@{..>}[d]_t \ar@{..>}[r]_g & Y \ar[d]^s \\
Z \ar[r]^f & W
}
$$
를 $t \in S$인 가환 점선 사각형으로 완성할 수 있다.
\item[RMS3] 사상들의 모든 쌍 $f, g : X \to Y$와, 정의역이 $Y$이고
$s \circ f = s \circ g$를 만족하는 $s \in S$에 대하여, 공역이 $X$이고
$f \circ t = g \circ t$를 만족하는 $t \in S$가 존재한다.
\end{enumerate}
$\mathcal{C}$의 화살표들의 집합 $S$가 왼쪽 곱셈계이면서 오른쪽 곱셈계이면
이를 {\it 곱셈계}라고 한다. 다시 말해 MS1, MS2, MS3가 성립한다는 뜻이다.
여기서 MS1 $=$ LMS1 $+$ RMS1, MS2 $=$ LMS2 $+$ RMS2이고,
MS3 $=$ LMS3 $+$ RMS3이다. (물론 LMS1 $=$ RMS1 $=$ MS1이다.)
\end{definition}

\noindent
이 조건들은 다음과 같이 범주 $S^{-1}\mathcal{C}$를 구성하는 데 유용하다.

\medskip\noindent
{\bf 왼쪽 분수 계산법.}
$\mathcal{C}$를 범주라 하고 $S$를 왼쪽 곱셈계라 하자. 새 범주
$S^{-1}\mathcal{C}$를 다음과 같이 정의한다(이 구성이 성립함은 Lemma
\ref{categories-lemma-left-localization}의 증명에서 확인한다).
\begin{enumerate}
\item $\Ob(S^{-1}\mathcal{C}) = \Ob(\mathcal{C})$로 놓는다.
\item $S^{-1}\mathcal{C}$의 사상 $X \to Y$는 동치 아래에서
$s \in S$인 순서쌍 $(f : X \to Y', s : Y \to Y')$로 주어진다.
(동치는 아래에서 정의한다. 순서쌍 $(f, s)$의 동치류를
$s^{-1}f : X \to Y$로 생각하라.)
\item 두 순서쌍 $(f_1 : X \to Y_1, s_1 : Y \to Y_1)$과
$(f_2 : X \to Y_2, s_2 : Y \to Y_2)$에 대하여, 셋째 순서쌍
$(f_3 : X \to Y_3, s_3 : Y \to Y_3)$와 $\mathcal{C}$의 사상
$u : Y_1 \to Y_3$ 및 $v : Y_2 \to Y_3$가 존재하여 가환 도표
$$
\xymatrix{
 & Y_1 \ar[d]^u & \\
X \ar[ru]^{f_1} \ar[r]^{f_3} \ar[rd]_{f_2} &
Y_3 &
Y \ar[lu]_{s_1} \ar[l]_{s_3} \ar[ld]^{s_2} \\
& Y_2 \ar[u]_v &
}
$$
에 들어맞으면 이 두 순서쌍이 동치라고 한다.
\item 순서쌍 $(f : X \to Y', s : Y \to Y')$와
$(g : Y \to Z', t : Z \to Z')$의 동치류들의 합성은 순서쌍
$(h \circ f : X \to Z'', u \circ t : Z \to Z'')$의 동치류로 정의한다.
여기서 $h$와 $u \in S$는 가정에 의해 존재하는 가환 도표
$$
\xymatrix{
Y \ar[d]_s \ar[r]_g & Z' \ar[d]^u \\
Y' \ar[r]^h & Z''
}
$$
에 들어맞도록 택한다.
\item $S^{-1} \mathcal{C}$의 항등사상 $X \to X$는 순서쌍
$(\text{id} : X \to X, \text{id} : X \to X)$의 동치류이다.
\end{enumerate}

\begin{lemma}
\label{categories-lemma-left-localization}
$\mathcal{C}$를 범주라 하고 $S$를 왼쪽 곱셈계라 하자.
\begin{enumerate}
\item 위에서 순서쌍들에 정의한 관계는 동치관계이다.
\item 위에서 준 합성 규칙은 동치류들 위에서 잘 정의된다.
\item 합성은 결합적이고(항등사상들은 항등 공리를 만족한다), 따라서
$S^{-1}\mathcal{C}$는 범주이다.
\end{enumerate}
\end{lemma}

\begin{proof}
(1)을 증명하자. 두 순서쌍
$p_1 = (f_1 : X \to Y_1, s_1 : Y \to Y_1)$과
$p_2 = (f_2 : X \to Y_2, s_2 : Y \to Y_2)$에 대하여
$a \circ f_1 = f_2$ 및 $a \circ s_1 = s_2$를 만족하는 $\mathcal{C}$의 사상
$a : Y_1 \to Y_2$가 존재하면 이들이 초등적으로 동치라고 하자. 도표로 쓰면
$$
\xymatrix{
X \ar@{=}[d] \ar[r]_{f_1} & Y_1 \ar[d]^a & Y \ar[l]^{s_1} \ar@{=}[d] \\
X \ar[r]^{f_2} & Y_2 & Y \ar[l]_{s_2}
}
$$
이다. 이 성질을 $p_1Ep_2$라고 표기하자. $pEp$이고
$aEb, bEc \Rightarrow aEc$임에 유의하라. (그 이름과 달리 $E$는
동치관계가 아니다.) (1)의 주장은 관계
$p \sim p' \Leftrightarrow \exists q: pEq \wedge p'Eq$가 동치관계라는 것이다
(여기서 $q$는 $p$ 및 $p'$와 같은 조건들을 만족하는 순서쌍이어야 한다).
위의 $E$의 성질들을 사용하는 간단한 형식적 논증으로, 다음을 증명하면
충분함을 알 수 있다.
$p_3Ep_1, p_3Ep_2 \Rightarrow p_1 \sim p_2$.
따라서 $s_i \in S$인 가환 도표
$$
\xymatrix{
 & Y_1 & \\
X \ar[ru]^{f_1} \ar[r]^{f_3} \ar[rd]_{f_2} &
Y_3 \ar[u]_{a_{31}} \ar[d]^{a_{32}} &
Y \ar[lu]_{s_1} \ar[l]_{s_3} \ar[ld]^{s_2} \\
& Y_2 &
}
$$
가 주어졌다고 가정하자. 먼저 LMS2를 적용하여
$a_{24} \in S$인 가환 도표
$$
\xymatrix{
Y \ar[d]_{s_1} \ar[r]_{s_2} & Y_2 \ar@{..>}[d]^{a_{24}} \\
Y_1 \ar@{..>}[r]^{a_{14}} & Y_4
}
$$
를 얻는다. 그러면
$$
a_{14} \circ a_{31} \circ s_3 =
a_{14} \circ s_1 =
a_{24} \circ s_2 =
a_{24} \circ a_{32} \circ s_3.
$$
이다. 따라서 LMS3에 의해 $s_{44} \in S$이고
$s_{44} \circ a_{14} \circ a_{31}
= s_{44} \circ a_{24} \circ a_{32}$가 되게 하는 사상
$s_{44} : Y_4 \to Y'_4$가 존재한다. $Y_4$, $a_{14}$, $a_{24}$를 각각
$Y'_4$, $s_{44} \circ a_{14}$, $s_{44} \circ a_{24}$로 바꾸면
$a_{14} \circ a_{31} = a_{24} \circ a_{32}$라고 가정할 수 있다
(그리고 여전히 $a_{24} \in S$이고
$a_{14} \circ s_1 = a_{24} \circ s_2$이다). 다음과 같이 놓자.
$$
f_4 =
a_{14} \circ f_1 =
a_{14} \circ a_{31} \circ f_3 =
a_{24} \circ a_{32} \circ f_3 =
a_{24} \circ f_2
$$
그리고 $s_4 = a_{14} \circ s_1 = a_{24} \circ s_2$이다. 그러면 도표
$$
\xymatrix{
X \ar@{=}[d] \ar[r]_{f_1} & Y_1 \ar[d]^{a_{14}} & Y \ar[l]^{s_1} \ar@{=}[d] \\
X \ar[r]^{f_4} & Y_4 & Y \ar[l]_{s_4}
}
$$
는 가환하고, LMS1에 의해 $s_4 \in S$이다. 따라서 $p_1 E p_4$이다.
여기서 $p_4 = (f_4, s_4)$이다. 마찬가지로 $p_2 E p_4$이다. 이들을
합치면 $p_1 \sim p_2$를 얻는다.

\medskip\noindent
(2)를 증명하자. $p = (f : X \to Y', s : Y \to Y')$와
$q = (g : Y \to Z', t : Z \to Z')$를 위의 합성 정의에서와 같은
순서쌍들이라 하자. 합성을 위해 $u_2 \in S$인 도표
$$
\xymatrix{
Y \ar[d]_s \ar[r]_g & Z' \ar[d]^{u_2} \\
Y' \ar[r]^{h_2} & Z_2
}
$$
를 택한다. 먼저 순서쌍
$r_2 = (h_2 \circ f : X \to Z_2, u_2 \circ t : Z \to Z_2)$의 동치류가
$(Z_2, h_2, u_2)$의 선택에 의존하지 않음을 보이자.
$(Z_3, h_3, u_3)$가 이에 대응하는 합성
$r_3 = (h_3 \circ f : X \to Z_3, u_3 \circ t : Z \to Z_3)$을 주는 또 다른
선택이라고 하자. 그러면 LMS2에 의해 $u_{34} \in S$인 도표
$$
\xymatrix{
Z' \ar[d]_{u_2} \ar[r]_{u_3} & Z_3 \ar[d]^{u_{34}} \\
Z_2 \ar[r]^{h_{24}} & Z_4
}
$$
를 택할 수 있다. $h_2 \circ s = u_2 \circ g$이고 마찬가지로
$h_3 \circ s = u_3 \circ g$이다. 이제
$$
u_{34} \circ h_3 \circ s
= u_{34} \circ u_3 \circ g
= h_{24} \circ u_2 \circ g
= h_{24} \circ h_2 \circ s.
$$
이다. 따라서 LMS3에 의해
$s_{44} \circ u_{34} \circ h_3 = s_{44} \circ h_{24} \circ h_2$가 되게 하는
$Z'_4$와 $s_{44} : Z_4 \to Z'_4$가 존재한다. $Z_4$, $h_{24}$, $u_{34}$를
각각 $Z'_4$, $s_{44} \circ h_{24}$, $s_{44} \circ u_{34}$로 바꾸면
$u_{34} \circ h_3 = h_{24} \circ h_2$라고 가정할 수 있다. 한편 관계
$u_{34} \circ u_3 = h_{24} \circ u_2$와 $u_{34} \in S$는 계속 성립한다.
이제 $h_4 = u_{34} \circ h_3 = h_{24} \circ h_2$와
$u_4 = u_{34} \circ u_3 = h_{24} \circ u_2$로 놓자. 그러면 가환 도표
$$
\xymatrix{
X \ar@{=}[d] \ar[r]_{h_2\circ f} &
Z_2 \ar[d]^{h_{24}} &
Z \ar[l]^{u_2 \circ t} \ar@{=}[d] \\
X \ar@{=}[d] \ar[r]^{h_4\circ f} &
Z_4 &
Z \ar@{=}[d] \ar[l]_{u_4 \circ t} \\
X \ar[r]^{h_3 \circ f} &
Z_3 \ar[u]^{u_{34}} &
Z \ar[l]_{u_3 \circ t}
}
$$
를 얻는다. 따라서 순서쌍
$r_4 = (h_4 \circ f : X \to Z_4, u_4 \circ t : Z \to Z_4)$를 얻고,
위 도표에 의해 $r_2Er_4$ 및 $r_3Er_4$이다. 그러므로 원하는 대로
$r_2 \sim r_3$이다. 따라서 이제 위에서 얻은 가능한 모든 순서쌍 $r$의
동치류로 $p \circ q$를 정의하는 것이 의미가 있다.

\medskip\noindent
(2)의 증명을 끝내려면 $p_1Ep_2$인 순서쌍들 $p_1, p_2, q$가 주어졌을 때,
합성이 의미가 있는 경우 $p_1 \circ q = p_2 \circ q$ 및
$q \circ p_1 = q \circ p_2$임을 보여야 한다. 다음과 같이 쓰자.
$p_1 = (f_1 : X \to Y_1, s_1 : Y \to Y_1)$,
$p_2 = (f_2 : X \to Y_2, s_2 : Y \to Y_2)$.
또한 $f_2 = a \circ f_1$ 및 $s_2 = a \circ s_1$을 만족하는
$\mathcal{C}$의 사상을 $a : Y_1 \to Y_2$라 하자. 먼저
$q = (g : Y \to Z', t : Z \to Z')$라고 가정하자. 이 경우 왼쪽과 같은
가환 도표를 택하면
$$
\vcenter{
\xymatrix{
Y \ar[d]_{s_2} \ar[r]^g & Z' \ar[d]^u \\
Y_2 \ar[r]^h & Z''
}
}
\quad
\Rightarrow
\quad
\vcenter{
\xymatrix{
Y \ar[d]_{s_1} \ar[r]^g & Z' \ar[d]^u \\
Y_1 \ar[r]^{h \circ a} & Z''
}
}
$$
오른쪽 도표도 가환한다(여기서 $u \in S$이다). 이 도표들을 사용하면 합성
$q \circ p_1$과 $q \circ p_2$가 모두 순서쌍
$(h \circ a \circ f_1 : X \to Z'', u \circ t : Z \to Z'')$의 동치류임을
알 수 있다. 따라서 $q \circ p_1 = q \circ p_2$이다. 다른 경우, 곧
$p_1 \circ q = p_2 \circ q$를 보여야 하는 경우의 증명은 생략한다.
(우리가 다룬 경우와 비슷하다.)

\medskip\noindent
(3)을 증명하자. 합성의 결합성을 증명해야 한다. 실선 도표
$$
\xymatrix{
& & & Z \ar[d] \\
& & Y \ar[d] \ar[r] & Z' \ar@{..>}[d] \\
& X \ar[d] \ar[r] & Y' \ar@{..>}[d] \ar@{..>}[r] & Z'' \ar@{..>}[d] \\
W \ar[r] & X' \ar@{..>}[r] & Y'' \ar@{..>}[r] & Z'''
}
$$
를 생각하자(그 수직 화살표들은 $S$에 속한다). 이 도표는 합성 가능한 세
순서쌍을 준다. LMS2를 사용하여 사각형들을 가환하게 하고 수직 화살표들이
$S$에 속하도록 하는 점선 화살표들을 택할 수 있다. 그러면 세 순서쌍의 합성은
아랫줄의 수평 화살표들과 오른쪽 열의 수직 화살표들을 합성하여 얻은 순서쌍
$(W \to Z''', Z \to Z''')$의 동치류임이 명백하다.

\medskip\noindent
항등 공리의 확인은 독자에게 맡긴다.
\end{proof}

\begin{remark}
\label{categories-remark-motivation-localization}
$S^{-1} \mathcal{C}$를 구성하는 동기는 $S$ 안의 사상들의 역을 인공적으로
만들어(그 대가로 기존 사상들 가운데 일부가 서로 동일시될 수 있다) 이들을
가역이 되도록 ``강제''하는 것이다. 이는 가환환을 곱셈적 부분집합에서
국소화하는 것과 비슷하고, 더 일반적으로는 비가환환을 오른쪽 분모 집합에서
국소화하는 것과 비슷하다(\cite[Section 10A]{Lam}을 보라). 이는 단순한
유사성에 그치지 않는다. $S^{-1} \mathcal{C}$의 구성(더 정확히는 가법 범주
$\mathcal{C}$에 대한 그 구성)은 실제로 후자의 국소화를 일반화한다.
실제로 비가환환은 대상이 하나뿐인 전가법 범주(사상은 환의 원소이다)로 볼 수
있고, 이 환의 곱셈적 부분집합은 LMS1(즉 RMS1)을 만족하는 사상들의 집합
$S$가 된다. 그러면 이 범주와 부분집합 $S$에 대한 RMS2와 RMS3는 오른쪽
분모 집합의 두 조건(``오른쪽 교환 가능''과 ``오른쪽 가역'')으로 옮겨진다
(LMS와 왼쪽 분모 집합도 마찬가지이다). 그리고 적절히 정의된 가법 구조를
갖는 $S^{-1} \mathcal{C}$는 환의 국소화에 대응하는 대상 하나짜리 범주이다.
\end{remark}

\begin{definition}
\label{categories-definition-left-localization-as-fraction}
$\mathcal{C}$를 범주라 하고 $S$를 $\mathcal{C}$의 사상들의 왼쪽 곱셈계라
하자. $\mathcal{C}$의 임의의 사상 $f : X \to Y'$와 $S$의 임의의 사상
$s : Y \to Y'$가 주어졌을 때, 순서쌍
$(f : X \to Y', s : Y \to Y')$의 동치류를 {\it $s^{-1} f$}로 표기한다.
이는 $S^{-1} \mathcal{C}$에서 $X$로부터 $Y$로 가는 사상이다.
\end{definition}

\noindent
이 표기는 암시적이며, 그 암시하는 내용은 참이다. $\mathcal{C}$의 임의의
사상 $f : X \to Y'$와 $S$의 임의의 두 사상
$s : Y \to Y'$ 및 $t : Y' \to Y''$에 대하여
$\left(t \circ s\right)^{-1} \left(t \circ f\right) = s^{-1} f$이다.
또한 $\mathcal{C}$의 임의의 $f : X \to Y'$와 $g : Y' \to Z'$ 및
$S$의 모든 $s : Z \to Z'$에 대하여
$s^{-1} \left(g \circ f\right) = \left(s^{-1} g\right) \circ
\left(\text{id}_{Y'}^{-1} f\right)$이다. 끝으로 $\mathcal{C}$의 임의의
$f : X \to Y'$, $S$의 모든 $s : Y \to Y'$ 및 $S$의 $t : Z \to Y$에 대하여
$\left(s \circ t\right)^{-1} f
= \left(t^{-1} \text{id}_Y\right)
\circ \left(s^{-1} f\right)$이다. 이는 모두 정의에서 명백하다.
우리는 ``공역이 같은 유한 개의 사상들을 공통 분모를 가진 분수로 쓸 수 있다.''

\begin{lemma}
\label{categories-lemma-morphisms-left-localization}
$\mathcal{C}$를 범주라 하고 $S$를 $\mathcal{C}$의 사상들의 왼쪽 곱셈계라
하자. $S^{-1}\mathcal{C}$의 사상들의 임의의 유한 족
$g_i : X_i \to Y$가 주어졌다고 하자($i$로 지표화한다). 그러면 $S$의 원소
$s : Y \to Y'$와 $\mathcal{C}$의 사상들의 족 $f_i : X_i \to Y'$를 찾아,
각 $g_i$가 순서쌍 $(f_i : X_i \to Y', s : Y \to Y')$의 동치류가 되게
할 수 있다.
\end{lemma}

\begin{proof}
각 $i$에 대하여 $g_i$의 대표 $(X_i \to Y_i, s_i : Y \to Y_i)$를 택하자.
각 $i$에 대하여 $a_i \circ s_i = s$인 사상 $a_i : Y_i \to Y'$가 존재하도록
$S$ 안의 사상 $s : Y \to Y'$를 찾을 수 있다면 보조정리가 따른다.
지표가 두 개 $i = 1, 2$라면 Definition
\ref{categories-definition-multiplicative-system}에 의해 가능한 대로 사각형
$$
\xymatrix{
Y \ar[d]_{s_1} \ar[r]_{s_2} & Y_2 \ar[d]^{t_2} \\
Y_1 \ar[r]^{a_1} & Y'
}
$$
을 $t_2 \in S$로 완성하여 이를 할 수 있다. 그러면
$s = t_2 \circ s_2 \in S$가 원하는 사상이다. 사상이 $n > 2$개라면 위의
방법으로 $n - 1$개의 사상인 경우로 줄이고 귀납법을 적용한다.
\end{proof}

\noindent
분모가 같을 때 사상들의 동등성을 간단히 특징지을 수 있다.

\begin{lemma}
\label{categories-lemma-equality-morphisms-left-localization}
$\mathcal{C}$를 범주라 하고 $S$를 $\mathcal{C}$의 사상들의 왼쪽 곱셈계라
하자. $A, B : X \to Y$를 각각
$(f : X \to Y', s : Y \to Y')$와 $(g : X \to Y', s : Y \to Y')$의
동치류인 $S^{-1}\mathcal{C}$의 사상들이라 하자. 다음은 서로 동치이다.
\begin{enumerate}
\item $A = B$
\item $t \circ f = t \circ g$를 만족하는 $S$ 안의 사상
$t : Y' \to Y''$가 존재한다.
\item $a \circ f = a \circ g$ 및 $a \circ s \in S$를 만족하는 사상
$a : Y' \to Y''$가 존재한다.
\end{enumerate}
\end{lemma}

\begin{proof}
$S^{-1}\mathcal{C}$가 범주라는 사실(Lemma \ref{categories-lemma-left-localization})을
사용하고, Definition \ref{categories-definition-left-localization-as-fraction}의 표기와
그 정의 뒤의 논의를 사용하여 $S^{-1}\mathcal{C}$의 몇몇 사상들을
동일시하겠다. 따라서 $A = s^{-1}f$ 및 $B = s^{-1}g$로 쓴다.

\medskip\noindent
$A = B$이면
$(\text{id}_{Y'}^{-1}s) \circ A = (\text{id}_{Y'}^{-1}s) \circ B$이다.
또한 $(\text{id}_{Y'}^{-1}s) \circ A = \text{id}_{Y'}^{-1}f$이고
$(\text{id}_{Y'}^{-1}s) \circ B = \text{id}_{Y'}^{-1}g$이다.
$\text{id}_{Y'}^{-1}f$와 $\text{id}_{Y'}^{-1}g$가 같다는 것은 정의에 의해
$t \in S$인 가환 도표
$$
\xymatrix{
 & Y' \ar[d]^u & \\
X \ar[ru]^f \ar[r]^h \ar[rd]_g &
Z &
Y' \ar[lu]_{\text{id}_{Y'}} \ar[l]_t \ar[ld]^{\text{id}_{Y'}} \\
& Y' \ar[u]_v &
}
$$
가 존재한다는 뜻이다. 특히 $u = v = t \in S$이고
$t \circ f = t\circ g$이다. 따라서 (1)은 (2)를 함의한다.

\medskip\noindent
(2) $\Rightarrow$ (3)은 즉시 따른다. (3)에서와 같은 $a$를 가정하자.
$s' = a \circ s \in S$로 표기하자. 그러면
$\text{id}_{Y''}^{-1}s'$는 범주 $S^{-1}\mathcal{C}$에서 동형이며, 그 역은
$(s')^{-1}\text{id}_{Y''}$이다. 따라서 $A = B$를 확인하려면
$\text{id}_{Y''}^{-1}s' \circ A = \text{id}_{Y''}^{-1}s' \circ B$를
확인하면 충분하다. Definition
\ref{categories-definition-left-localization-as-fraction} 뒤의 글에서 논의한 규칙을 사용하면
다음과 같이 계산된다.
$\text{id}_{Y''}^{-1}s' \circ A =
\text{id}_{Y''}^{-1}(a \circ s) \circ s^{-1}f =
\text{id}_{Y''}^{-1}(a \circ f) =
\text{id}_{Y''}^{-1}(a \circ g) =
\text{id}_{Y''}^{-1}(a \circ s) \circ s^{-1}g =
\text{id}_{Y''}^{-1}s' \circ B$이고, 따라서 (1)이 참이다.
\end{proof}

\begin{remark}
\label{categories-remark-left-localization-morphisms-colimit}
$\mathcal{C}$를 범주라 하고 $S$를 왼쪽 곱셈계라 하자. $\mathcal{C}$의 대상
$Y$가 주어졌을 때, 대상은 $s \in S$인 $s : Y \to Y'$이고 사상은 가환 도표
$$
\xymatrix{
& Y \ar[ld]_s \ar[rd]^t & \\
Y' \ar[rr]^a & & Y''
}
$$
인 범주를 $Y/S$로 표기한다. 여기서 $a : Y' \to Y''$는 임의이다.
범주 $Y/S$가 여과 범주라고 주장한다(Definition
\ref{categories-definition-directed}를 보라). 실제로 LMS1에 의해
$\text{id}_Y : Y \to Y$는 $Y/S$ 안에 있으므로 $Y/S$는 공집합이 아니다.
LMS2에 의해 $s_1 : Y \to Y_1$과 $s_2 : Y \to Y_2$가 주어지면
$t \in S$인 도표
$$
\xymatrix{
Y \ar[d]_{s_1} \ar[r]_{s_2} & Y_2 \ar[d]^t \\
Y_1 \ar[r]^a & Y_3
}
$$
를 찾을 수 있다. 따라서 $s_1 : Y \to Y_1$과 $s_2 : Y \to Y_2$는 모두
$Y/S$에서 $t \circ s_2 : Y \to Y_3$로 가는 사상을 갖는다. 끝으로
$Y/S$에서 $s_1 : Y \to Y_1$로부터 $s_2 : Y \to Y_2$로 가는 두 사상
$a, b$가 주어지면 $a \circ s_1 = b \circ s_1$임을 안다. 따라서 LMS3에 의해
$t \circ a = t \circ b$인 $S$ 안의 $t : Y_2 \to Y_3$가 존재한다.
이제 Lemmas \ref{categories-lemma-morphisms-left-localization}와
\ref{categories-lemma-equality-morphisms-left-localization}을 합치면
\begin{equation}
\label{categories-equation-left-localization-morphisms-colimit}
\Mor_{S^{-1}\mathcal{C}}(X, Y) =
\colim_{(s : Y \to Y') \in Y/S} \Mor_\mathcal{C}(X, Y')
\end{equation}
를 얻는다. $S^{-1}\mathcal{C}$의 사상 집합을 $\mathcal{C}$의 사상 집합들의
여과 여극한으로 나타내는 이 공식은 때때로 유용하다.
\end{remark}

\begin{lemma}
\label{categories-lemma-properties-left-localization}
$\mathcal{C}$를 범주라 하고 $S$를 $\mathcal{C}$의 사상들로 이루어진
왼쪽 곱셈계라 하자.
\begin{enumerate}
\item 규칙 $X \mapsto X$와
$(f : X \to Y) \mapsto (f : X \to Y, \text{id}_Y : Y \to Y)$는
함자 $Q : \mathcal{C} \to S^{-1}\mathcal{C}$를 정의한다.
\item 임의의 $s \in S$에 대하여 사상 $Q(s)$는
$S^{-1}\mathcal{C}$에서 동형이다.
\item $G : \mathcal{C} \to \mathcal{D}$가 모든 $s \in S$에 대하여
$G(s)$가 가역인 임의의 함자라면, $H \circ Q = G$를 만족하는 유일한 함자
$H : S^{-1}\mathcal{C} \to \mathcal{D}$가 존재한다.
\end{enumerate}
\end{lemma}

\begin{proof}
(1)과 (2)는 명백하다. ((2)에서 $Q(s)$의 역은 쌍
$(\text{id}_Y, s)$의 동치류이다.)
(3)을 보이려면 $H(X) = G(X)$로 놓고
$H((f : X \to Y', s : Y \to Y')) = G(s)^{-1} \circ G(f)$로 놓으면 된다.
세부 사항은 생략한다.
\end{proof}

\begin{lemma}
\label{categories-lemma-left-localization-limits}
$\mathcal{C}$를 범주라 하고 $S$를 $\mathcal{C}$의 사상들로 이루어진
왼쪽 곱셈계라 하자. 국소화 함자
$Q : \mathcal{C} \to S^{-1}\mathcal{C}$는 유한 여극한과 가환한다.
\end{lemma}

\begin{proof}
$\mathcal{I}$를 유한 범주라 하고
$\mathcal{I} \to \mathcal{C}$, $i \mapsto X_i$를 여극한이 존재하는
함자라 하자. 그러면 (\ref{categories-equation-left-localization-morphisms-colimit}),
$Y/S$가 여과 범주라는 사실, 그리고 Lemma
\ref{categories-lemma-directed-commutes}를 사용하여 다음을 얻는다.
\begin{align*}
\Mor_{S^{-1}\mathcal{C}}(Q(\colim X_i), Q(Y))
& =
\colim_{(s : Y \to Y') \in Y/S} \Mor_\mathcal{C}(\colim X_i, Y') \\
& =
\colim_{(s : Y \to Y') \in Y/S} \lim_i \Mor_\mathcal{C}(X_i, Y') \\
& =
\lim_i \colim_{(s : Y \to Y') \in Y/S} \Mor_\mathcal{C}(X_i, Y') \\
& =
\lim_i \Mor_{S^{-1}\mathcal{C}}(Q(X_i), Q(Y))
\end{align*}
이 동형은 각 $j \in \Ob(\mathcal{I})$에 대하여 양변에서 집합
$\Mor_{S^{-1}\mathcal{C}}(Q(X_j), Q(Y))$로 가는 사영과 가환한다. 따라서
$Q(\colim X_i)$는 함자 $i \mapsto Q(X_i)$의 여극한에 대한 보편 성질을
만족하므로, 원하는 대로 바로 그 여극한이다.
\end{proof}

\begin{lemma}
\label{categories-lemma-left-localization-diagram}
$\mathcal{C}$를 범주라 하자. $S$를 왼쪽 곱셈계라 하자.
$f : X \to Y$, $f' : X' \to Y'$가 $\mathcal{C}$의 두 사상이고
$$
\xymatrix{
Q(X) \ar[d]_{Q(f)} \ar[r]_a & Q(X') \ar[d]^{Q(f')} \\
Q(Y) \ar[r]^b & Q(Y')
}
$$
가 $S^{-1}\mathcal{C}$의 가환 도표라면, $\mathcal{C}$의 사상
$f'' : X'' \to Y''$과 $\mathcal{C}$의 가환 도표
$$
\xymatrix{
X \ar[d]_f \ar[r]_g & X'' \ar[d]^{f''} & X' \ar[d]^{f'} \ar[l]^s \\
Y \ar[r]^h & Y'' & Y' \ar[l]_t
}
$$
가 존재하여 $s, t \in S$이고 $a = s^{-1}g$, $b = t^{-1}h$이다.
\end{lemma}

\begin{proof}
다음과 같이 사상과 대상을 선택한다. 먼저 $S$ 안의 어떤
$s : X' \to X''$와 어떤 $g : X \to X''$에 대하여 $a = s^{-1}g$로 쓴다.
LMS2에 의해 $S$ 안의 $t : Y' \to Y''$와
$f'' : X'' \to Y''$를 찾아
$$
\xymatrix{
X \ar[d]_f \ar[r]_g & X'' \ar[d]^{f''} & X' \ar[d]^{f'} \ar[l]^s \\
Y & Y'' & Y' \ar[l]_t
}
$$
가 가환하게 할 수 있다. 이제 이 도표에서 $d \circ t \in S$라는 성질을
갖는 사상 $d : Y'' \to Y'''$로 $t$와 $f''$를 모두 뒤에서 합성함으로써
$$
X'' \xrightarrow{f''} Y'' \xleftarrow{t} Y'
$$
의 선택을 반복하여 바꿀 것이다. Remark
\ref{categories-remark-left-localization-morphisms-colimit}에 따르면 그러한 교체 뒤에
$b = t^{-1}h$가 성립하는 사상 $h : Y \to Y''$가 존재한다고 가정할 수
있다\footnote{이를 더 직접적으로 보는 방법은 다음과 같다.
$S$ 안의 어떤 $q : Y' \to Z$와 어떤 $i : Y \to Z$에 대하여
$b = q^{-1}i$로 쓰자. LMS2에 의해 $S$ 안의 $r : Y'' \to Y'''$와
$j : Z \to Y'''$를 찾아 $j \circ q = r \circ t$가 되게 할 수 있다.
이제 $d = r$ 및 $h = j \circ i$로 놓는다.}. 이 시점에는 도표의 왼쪽
사각형이 가환한다는 사실만 아직 모를 뿐, 보조정리에 적힌 나머지는 모두
갖추어져 있다.
그러나 $S^{-1} \mathcal{C}$에서의 합성 정의에 따르면
$b \circ Q\left(f\right)$는 쌍
$(h \circ f : X \to Y'', t : Y' \to Y'')$의 동치류이다
($b$는 쌍 $(h : Y \to Y'', t : Y' \to Y'')$의 동치류이고,
$Q\left(f\right)$는 쌍
$(f : X \to Y, \text{id} : Y \to Y)$의 동치류이기 때문이다). 한편
$Q\left(f'\right) \circ a$는 쌍
$(f'' \circ g : X \to Y'', t : Y' \to Y'')$의 동치류이다
($a$는 쌍 $(g : X \to X'', s : X' \to X'')$의 동치류이고,
$Q\left(f'\right)$는 쌍
$(f' : X' \to Y', \text{id} : Y' \to Y')$의 동치류이기 때문이다).
$b \circ Q\left(f\right) = Q\left(f'\right) \circ a$임을 알고 있으므로,
쌍 $(h \circ f : X \to Y'', t : Y' \to Y'')$와
$(f'' \circ g : X \to Y'', t : Y' \to Y'')$의 동치류가 같다고 결론짓는다.
따라서 Lemma \ref{categories-lemma-equality-morphisms-left-localization}를 사용하면
$d \circ t \in S$이고 $d \circ h \circ f = d \circ f'' \circ g$인
사상 $d : Y'' \to Y'''$를 찾을 수 있다. 그러므로 위에서 설명한 종류의
교체를 한 번 더 하면 원하는 결론을 얻는다.
\end{proof}

\noindent
{\bf 오른쪽 분수 계산법.}
$\mathcal{C}$를 범주라 하고 $S$를 오른쪽 곱셈계라 하자. 다음과 같이
새 범주 $S^{-1}\mathcal{C}$를 정의한다(이 정의가 성립함은 Lemma
\ref{categories-lemma-right-localization}의 증명에서 확인한다).
\begin{enumerate}
\item $\Ob(S^{-1}\mathcal{C}) = \Ob(\mathcal{C})$로 놓는다.
\item $S^{-1}\mathcal{C}$의 사상 $X \to Y$는 $s \in S$인 쌍
$(f : X' \to Y, s : X' \to X)$의 동치류로 주어진다.
(동치관계는 아래에서 정의한다. 쌍 $(f, s)$의 동치류를
$fs^{-1} : X \to Y$라고 생각하라.)
\item 두 쌍 $(f_1 : X_1 \to Y, s_1 : X_1 \to X)$와
$(f_2 : X_2 \to Y, s_2 : X_2 \to X)$가 동치라는 것은,
세 번째 쌍 $(f_3 : X_3 \to Y, s_3 : X_3 \to X)$와
$\mathcal{C}$의 사상 $u : X_3 \to X_1$ 및 $v : X_3 \to X_2$가 존재하여
가환 도표
$$
\xymatrix{
 & X_1 \ar[ld]_{s_1} \ar[rd]^{f_1} & \\
X &
X_3 \ar[l]_{s_3} \ar[u]_u \ar[d]^v \ar[r]^{f_3} &
Y \\
& X_2 \ar[lu]^{s_2} \ar[ru]_{f_2} &
}
$$
에 들어맞는다는 뜻이다.
\item 쌍 $(f : X' \to Y, s : X' \to X)$와
$(g : Y' \to Z, t : Y' \to Y)$의 동치류들의 합성은
$h$와 $u \in S$를 가환 도표
$$
\xymatrix{
X'' \ar[d]_u \ar[r]^h & Y' \ar[d]^t \\
X' \ar[r]^f & Y
}
$$
에 들어맞도록 선택할 때 쌍
$(g \circ h : X'' \to Z, s \circ u : X'' \to X)$의 동치류로 정의한다.
그러한 도표는 가정에 의해 존재한다.
\item $S^{-1} \mathcal{C}$에서 항등사상 $X \to X$는 쌍
$(\text{id} : X \to X, \text{id} : X \to X)$의 동치류이다.
\end{enumerate}

\begin{lemma}
\label{categories-lemma-right-localization}
$\mathcal{C}$를 범주라 하고 $S$를 오른쪽 곱셈계라 하자.
\begin{enumerate}
\item 위에서 쌍들에 관하여 정의한 관계는 동치관계이다.
\item 위에서 준 합성 규칙은 동치류 위에서 잘 정의된다.
\item 합성은 결합적이고 항등사상들은 항등 공리를 만족한다. 따라서
$S^{-1}\mathcal{C}$는 범주이다.
\end{enumerate}
\end{lemma}

\begin{proof}
이 보조정리는 Lemma \ref{categories-lemma-left-localization}의 쌍대이다.
$S$가 왼쪽 곱셈계가 되는 반대범주로 $\mathcal{C}$를 바꾸면 그 보조정리에서
형식적으로 따른다.
\end{proof}

\begin{definition}
\label{categories-definition-right-localization-as-fraction}
$\mathcal{C}$를 범주라 하고 $S$를 $\mathcal{C}$의 사상들로 이루어진
오른쪽 곱셈계라 하자. $\mathcal{C}$의 임의의 사상 $f : X' \to Y$와
$S$의 임의의 사상 $s : X' \to X$가 주어졌을 때, 쌍
$(f : X' \to Y, s : X' \to X)$의 동치류를 {\it $f s^{-1}$}로 표기한다.
이는 $S^{-1} \mathcal{C}$에서 $X$로부터 $Y$로 가는 사상이다.
\end{definition}

\noindent
Definition \ref{categories-definition-left-localization-as-fraction}에 있는 것과 비슷한
(실제로는 쌍대인) 항등식들이 성립한다. 우리는 ``공통 분모를 갖는 분수로
동일한 시작점을 갖는 유한한 사상 모음을 쓸 수 있다.''

\begin{lemma}
\label{categories-lemma-morphisms-right-localization}
$\mathcal{C}$를 범주라 하고 $S$를 $\mathcal{C}$의 사상들로 이루어진
오른쪽 곱셈계라 하자. $S^{-1}\mathcal{C}$의 사상들로 이루어진 임의의 유한한
모음 $g_i : X \to Y_i$가 주어지면($i$로 첨자화한다),
$S$의 원소 $s : X' \to X$와 $\mathcal{C}$의 사상들의 족
$f_i : X' \to Y_i$를 찾아 $g_i$가 쌍
$(f_i : X' \to Y_i, s : X' \to X)$의 동치류가 되게 할 수 있다.
\end{lemma}

\begin{proof}
이 보조정리는 Lemma \ref{categories-lemma-morphisms-left-localization}의 쌍대이며,
나타나는 모든 범주를 그 반대범주로 바꾸면 그 보조정리에서 형식적으로 따른다.
\end{proof}

\noindent
같은 분모를 갖는 사상들의 동등성에는 간단한 판정법이 있다.

\begin{lemma}
\label{categories-lemma-equality-morphisms-right-localization}
$\mathcal{C}$를 범주라 하고 $S$를 $\mathcal{C}$의 사상들로 이루어진
오른쪽 곱셈계라 하자. $A, B : X \to Y$가 $S^{-1}\mathcal{C}$의 사상이며,
각각 $(f : X' \to Y, s : X' \to X)$와
$(g : X' \to Y, s : X' \to X)$의 동치류라 하자. 다음은 서로 동치이다.
\begin{enumerate}
\item $A = B$,
\item $f \circ t = g \circ t$인 $S$ 안의 사상
$t : X'' \to X'$가 존재한다.
\item $f \circ a = g \circ a$이고 $s \circ a \in S$인 사상
$a : X'' \to X'$가 존재한다.
\end{enumerate}
\end{lemma}

\begin{proof}
이는 Lemma \ref{categories-lemma-equality-morphisms-left-localization}의 쌍대이다.
\end{proof}

\begin{remark}
\label{categories-remark-right-localization-morphisms-colimit}
$\mathcal{C}$를 범주라 하고 $S$를 오른쪽 곱셈계라 하자.
$\mathcal{C}$의 대상 $X$가 주어졌을 때, 대상은 $s \in S$인
$s : X' \to X$이고 사상은 가환 도표
$$
\xymatrix{
X' \ar[rd]_s \ar[rr]_a & & X'' \ar[ld]^t \\
& X
}
$$
인 범주를 $S/X$로 표기한다. 여기서 $a : X' \to X''$는 임의이다.
범주 $S/X$는 쌍대여과 범주이다(Definition
\ref{categories-definition-codirected}를 보라). (이는 Remark
\ref{categories-remark-left-localization-morphisms-colimit}의 대응하는 명제와 쌍대이다.)
이제 Lemmas \ref{categories-lemma-morphisms-right-localization}와
\ref{categories-lemma-equality-morphisms-right-localization}을 합치면
\begin{equation}
\label{categories-equation-right-localization-morphisms-colimit}
\Mor_{S^{-1}\mathcal{C}}(X, Y) =
\colim_{(s : X' \to X) \in (S/X)^{opp}} \Mor_\mathcal{C}(X', Y)
\end{equation}
를 얻는다. $S^{-1}\mathcal{C}$의 사상들을 $\mathcal{C}$의 사상들의
여과 여극한으로 나타내는 이 공식은 때때로 유용하다.
\end{remark}

\begin{lemma}
\label{categories-lemma-properties-right-localization}
$\mathcal{C}$를 범주라 하고 $S$를 $\mathcal{C}$의 사상들로 이루어진
오른쪽 곱셈계라 하자.
\begin{enumerate}
\item 규칙 $X \mapsto X$와
$(f : X \to Y) \mapsto (f : X \to Y, \text{id}_X : X \to X)$는
함자 $Q : \mathcal{C} \to S^{-1}\mathcal{C}$를 정의한다.
\item 임의의 $s \in S$에 대하여 사상 $Q(s)$는
$S^{-1}\mathcal{C}$에서 동형이다.
\item $G : \mathcal{C} \to \mathcal{D}$가 모든 $s \in S$에 대하여
$G(s)$가 가역인 임의의 함자라면, $H \circ Q = G$를 만족하는 유일한 함자
$H : S^{-1}\mathcal{C} \to \mathcal{D}$가 존재한다.
\end{enumerate}
\end{lemma}

\begin{proof}
이 보조정리는 Lemma \ref{categories-lemma-properties-left-localization}의 쌍대이며,
나타나는 모든 범주를 그 반대범주로 바꾸면 그 보조정리에서 형식적으로 따른다.
\end{proof}

\begin{lemma}
\label{categories-lemma-right-localization-limits}
$\mathcal{C}$를 범주라 하고 $S$를 $\mathcal{C}$의 사상들로 이루어진
오른쪽 곱셈계라 하자. 국소화 함자
$Q : \mathcal{C} \to S^{-1}\mathcal{C}$는 유한 극한과 가환한다.
\end{lemma}

\begin{proof}
이는 Lemma \ref{categories-lemma-left-localization-limits}와 쌍대이다.
\end{proof}

\begin{lemma}
\label{categories-lemma-right-localization-diagram}
$\mathcal{C}$를 범주라 하고 $S$를 오른쪽 곱셈계라 하자.
$f : X \to Y$, $f' : X' \to Y'$가 $\mathcal{C}$의 두 사상이고
$$
\xymatrix{
Q(X) \ar[d]_{Q(f)} \ar[r]_a & Q(X') \ar[d]^{Q(f')} \\
Q(Y) \ar[r]^b & Q(Y')
}
$$
가 $S^{-1}\mathcal{C}$의 가환 도표라면, $\mathcal{C}$의 사상
$f'' : X'' \to Y''$과 $\mathcal{C}$의 가환 도표
$$
\xymatrix{
X \ar[d]_f & X'' \ar[l]^s \ar[d]^{f''} \ar[r]_g & X' \ar[d]^{f'} \\
Y & Y'' \ar[l]_t \ar[r]^h & Y'
}
$$
가 존재하여 $s, t \in S$이고 $a = gs^{-1}$, $b = ht^{-1}$이다.
\end{lemma}

\begin{proof}
이 보조정리는 Lemma \ref{categories-lemma-left-localization-diagram}의 쌍대이다.
\end{proof}

\noindent
{\bf 곱셈계와 양쪽 분수 계산법.}
$S$가 곱셈계이면 왼쪽 및 오른쪽 분수 계산법은 표준적으로 동형인 범주들을
준다.

\begin{lemma}
\label{categories-lemma-multiplicative-system}
$\mathcal{C}$를 범주라 하고 $S$를 곱셈계라 하자. 왼쪽 분수의 범주와
오른쪽 분수의 범주 $S^{-1}\mathcal{C}$는 표준적으로 동형이다.
\end{lemma}

\begin{proof}
두 분수 범주를 $\mathcal{C}_{left}$, $\mathcal{C}_{right}$로 표기하자.
Lemmas \ref{categories-lemma-properties-left-localization}와
\ref{categories-lemma-properties-right-localization}의 보편 성질에 의해 함자
$\mathcal{C}_{left} \to \mathcal{C}_{right}$와
$\mathcal{C}_{right} \to \mathcal{C}_{left}$를 얻는다.
보편 성질의 유일성 명제에 의해 이 함자들은 서로의 역이다.
\end{proof}

\begin{definition}
\label{categories-definition-saturated-multiplicative-system}
$\mathcal{C}$를 범주라 하고 $S$를 곱셈계라 하자. MS1, MS2, MS3에 더하여
다음도 성립하면 $S$를 {\it 포화되었다}고 한다.
\begin{enumerate}
\item[MS4] 합성 가능한 세 사상 $f, g, h$가 주어졌을 때
$fg, gh \in S$이면 $g \in S$이다.
\end{enumerate}
\end{definition}

\noindent
포화된 곱셈계는 모든 동형을 포함한다. 또한 범주의 합성 가능한 사상
$f, g, h$에 대하여 $fg, gh$가 동형이면 $g$는 동형이다(그때 $g$는
왼쪽 역과 오른쪽 역을 모두 가지므로 가역이기 때문이다).

\begin{lemma}
\label{categories-lemma-what-gets-inverted}
$\mathcal{C}$를 범주라 하고 $S$를 곱셈계라 하자. 국소화 함자를
$Q : \mathcal{C} \to S^{-1}\mathcal{C}$로 표기하자. 집합
$$
\hat S = \{f \in \text{Arrows}(\mathcal{C}) \mid
Q(f) \text{ is an isomorphism}\}
$$
은
$$
S' = \{f \in \text{Arrows}(\mathcal{C}) \mid
\text{there exist }g, h\text{ such that }gf, fh \in S\}
$$
와 같으며, $S$를 포함하는 가장 작은 포화 곱셈계이다. 특히 $S$가
포화되었으면 $\hat S = S$이다.
\end{lemma}

\begin{proof}
$S'$의 원소들은 $S^{-1}\mathcal{C}$에서 왼쪽 역과 오른쪽 역을 모두 갖는
사상으로 가므로 $S \subset S' \subset \hat S$임은 명백하다.
$S'$는 MS4를 만족하고 $\hat S$는 MS1을 만족함에 유의하라. 다음으로
$S' = \hat S$임을 증명한다.

\medskip\noindent
$f \in \hat S$라 하자. $S^{-1}\mathcal{C}$에서의 역사상을
$s^{-1}g = ht^{-1}$로 쓰자. (Lemma
\ref{categories-lemma-multiplicative-system}을 보면 $S^{-1}\mathcal{C}$의 사상을
왼쪽 분수와 오른쪽 분수 양쪽으로 나타낼 수 있다.)
$S^{-1}\mathcal{C}$에서의 관계 $\text{id}_X = s^{-1}gf$는 어떤 사상
$f', u, v$와 어떤 $s' \in S$에 대하여 가환 도표
$$
\xymatrix{
 & X' \ar[d]^u & \\
X \ar[ru]^{gf} \ar[r]^{f'} \ar[rd]_{\text{id}_X} &
X'' &
X \ar[lu]_s \ar[l]_{s'} \ar[ld]^{\text{id}_X} \\
& X \ar[u]_v &
}
$$
가 존재한다는 뜻이다. 따라서 $ugf = s' \in S$이다. 마찬가지로
$\text{id}_Y = fht^{-1}$를 사용하면 어떤 $w$에 대하여 $fhw \in S$임을
증명할 수 있다. 그러므로 $f \in S'$라고 결론짓는다. 따라서
$S' = \hat S$이다. 이제 $S' = \hat S$가 곱셈계임을 증명하기만 하면,
이로부터 $S' = \hat S$가 $S$를 포함하는 가장 작은 포화된 계임이 명백하다.

\medskip\noindent
위의 관찰로 MS1과 MS4는 처리되었으므로, 보조정리의 증명을 끝내려면
LMS2, RMS2, LMS3, RMS3가 $\hat S$에 대하여 성립함을 보여야 한다.
$\hat S$에 대하여 LMS2가 성립함을 확인하자. 실선 도표
$$
\xymatrix{
X \ar[d]_t \ar[r]_g & Y \ar@{..>}[d]^s \\
Z \ar@{..>}[r]^f & W
}
$$
가 주어지고 $t \in \hat S$라 하자. $at \in S$인 사상
$a : Z \to Z'$를 하나 택한다. 그러면 $S$에 대한 LMS2를 사용하여
가환 도표
$$
\xymatrix{
X \ar[d]_t \ar[r]_g & Y \ar[dd]^s \\
Z \ar[d]_a \\
Z' \ar[r]^{f'} & W
}
$$
를 찾을 수 있고, $f = f' \circ a$로 놓으면 원하는 결론을 얻는다.
RMS2의 증명은 이것과 쌍대이다. 끝으로 사상들의 쌍 $f, g : X \to Y$와
공역이 $X$인 $t \in \hat S$가 주어져 $ft = gt$라 하자.
그러면 $tb \in S$인 사상 $b$를 하나 택한다. 이때 $ftb = gtb$이고,
$S$에 대한 LMS3에 의해 시작점이 $Y$인 어떤 $s \in S$가 존재하여
원하는 대로 $sf = sg$이다. RMS3의 증명은 이것과 쌍대이다.
\end{proof}








\section{형식적 성질}
\label{categories-section-formal-cat-cat}

\noindent
이 절에서는 범주들의 $2$-범주가 갖는 몇 가지 형식적 성질을 논의한다.
이는 뒤에서 (엄밀한) $2$-범주를 정의하는 데 이르게 된다.

\medskip\noindent
모든 범주들의 모임을 $\Ob(\textit{Cat})$로 표기하자. 범주들의 각 쌍
$\mathcal{A}, \mathcal{B} \in \Ob(\textit{Cat})$에 대하여 함자들의
``작은'' 범주 $\text{Fun}(\mathcal{A}, \mathcal{B})$가 있다.
다음과 같은 함자 변환들의 합성
$$
\xymatrix{
\mathcal{A}
\rruppertwocell^{F''}{t'}
\ar[rr]_(.3){F'}
\rrlowertwocell_F{t}
& &
\mathcal{B}
}
\text{ composes to }
\xymatrix{
\mathcal{A}
\rrtwocell^{F''}_F{\ \ t \circ t'}
& &
\mathcal{B}
}
$$
을 {\it 수직} 합성이라 한다. 이에 보통 기호 $\circ$를 사용할 것이다.
다음으로 {\it 수평} 합성을 정의한다. 이를 위해 현재의 구조를 조금 더
설명하자.

\medskip\noindent
즉, 범주들의 각 세쌍 $\mathcal{A}$, $\mathcal{B}$, $\mathcal{C}$에 대하여
함자들의 합성에서 오는 합성법
$$
\circ : \Ob(\text{Fun}(\mathcal{B}, \mathcal{C}))
\times
\Ob(\text{Fun}(\mathcal{A}, \mathcal{B}))
\longrightarrow
\Ob(\text{Fun}(\mathcal{A}, \mathcal{C}))
$$
이 있다. 이 합성법은 결합적이고 항등함자는 단위원으로 작용한다. 달리 말해
함자들의 변환을 잊으면 $\textit{Cat}$이 범주를 이룸을 알 수 있다.
이 구조는 함자들 사이의 사상과 어떻게 상호작용하는가?

\medskip\noindent
함자 $F, F' : \mathcal{A} \to \mathcal{B}$의 변환
$t : F \to F'$와 함자 $G : \mathcal{B} \to \mathcal{C}$가 주어졌을 때,
함자들의 변환 $G\circ F \to G \circ F'$를 정의할 수 있다. 이 변환을
${}_Gt$로 표기한다. 모든 $x \in \mathcal{A}$에 대하여 공식
$({}_Gt)_x = G(t_x) : G(F(x)) \to G(F'(x))$로 주어진다.
이렇게 하여 $G$와의 합성은 함자
$$
\text{Fun}(\mathcal{A}, \mathcal{B})
\longrightarrow
\text{Fun}(\mathcal{A}, \mathcal{C}).
$$
가 된다. 이를 보이려면 ${}_G(\text{id}_F) = \text{id}_{G \circ F}$와
${}_G(t_1 \circ t_2) = {}_Gt_1 \circ {}_Gt_2$를 확인하기만 하면 된다.
물론 ${}_{\text{id}_\mathcal{B}}t = t$도 성립한다.

\medskip\noindent
마찬가지로 함자 $G, G' : \mathcal{B} \to \mathcal{C}$의 변환
$s : G \to G'$와 함자 $F : \mathcal{A} \to \mathcal{B}$가 주어졌을 때,
$s_F$를 공식 $(s_F)_x = s_{F(x)} : G(F(x)) \to G'(F(x))$로 주어지는
함자들의 변환 $G\circ F \to G' \circ F$로 정의할 수 있다. 여기서
$x \in \mathcal{A}$이다. 이렇게 하여 $F$와의 합성은 함자
$$
\text{Fun}(\mathcal{B}, \mathcal{C})
\longrightarrow
\text{Fun}(\mathcal{A}, \mathcal{C}).
$$
가 된다. 이를 보이려면 $(\text{id}_G)_F = \text{id}_{G\circ F}$와
$(s_1 \circ s_2)_F = s_{1, F} \circ s_{2, F}$를 확인하기만 하면 된다.
물론 $s_{\text{id}_\mathcal{B}} = s$도 성립한다.

\medskip\noindent
이 구성들은 의미가 통할 때마다 추가로
$$
{}_{G_1}({}_{G_2}t) = {}_{G_1\circ G_2}t,
\ (s_{F_1})_{F_2} = s_{F_1 \circ F_2},
\text{ and }{}_H(s_F) = ({}_Hs)_F
$$
를 만족한다. 끝으로 함자 $F, F' : \mathcal{A} \to \mathcal{B}$와
$G, G' : \mathcal{B} \to \mathcal{C}$ 및 변환 $t : F \to F'$와
$s : G \to G'$가 주어지면 다음 도표가 가환한다.
$$
\xymatrix{
G \circ F \ar[r]^{{}_Gt} \ar[d]_{s_F}
&
G \circ F' \ar[d]^{s_{F'}} \\
G' \circ F \ar[r]_{{}_{G'}t}
&
G' \circ F'
}
$$
달리 말해 ${}_{G'}t \circ s_F = s_{F'}\circ {}_Gt$이다. 이를 증명하려면
임의의 대상 $x \in \Ob(\mathcal{A})$에서 일어나는 일을 살펴보면 된다.
$$
\xymatrix{
G(F(x)) \ar[r]^{G(t_x)} \ar[d]_{s_{F(x)}}
&
G(F'(x)) \ar[d]^{s_{F'(x)}} \\
G'(F(x)) \ar[r]_{G'(t_x)}
&
G'(F'(x))
}
$$
$s$가 함자들의 변환이므로 이 도표는 가환한다. 이 양립관계로부터 수평 합성을
정의할 수 있다.

\begin{definition}
\label{categories-definition-horizontal-composition}
왼쪽과 같은 도표가 주어지면
$$
\xymatrix{
\mathcal{A}
\rtwocell^F_{F'}{t}
&
\mathcal{B}
\rtwocell^G_{G'}{s}
&
\mathcal{C}
}
\text{ gives }
\xymatrix{
\mathcal{A}
\rrtwocell^{G \circ F} _{G' \circ F'}{\ \ s \star t}
& &
\mathcal{C}
}
$$
수평 합성 $s \star t$를 함자들의 변환
${}_{G'}t \circ s_F = s_{F'}\circ {}_Gt$로 정의한다.
\end{definition}

\noindent
이제 이전에 구성한 변환 ${}_Gt$와 $s_F$를 각각
$ {}_Gt = \text{id}_G \star t $ 및 $ s_F = s \star \text{id}_F $로
되찾을 수 있다. 더 나아가 위에서 얻은 모든 규칙은 다음 보조정리에 진술된
성질들의 귀결이다.

\begin{lemma}
\label{categories-lemma-properties-2-cat-cats}
수평 합성과 수직 합성은 다음 성질들을 갖는다.
\begin{enumerate}
\item $\circ$와 $\star$는 결합적이다.
\item 항등변환 $\text{id}_F$는 $\circ$의 단위원이다.
\item 항등함자의 항등변환 $\text{id}_{\text{id}_\mathcal{A}}$는
$\star$와 $\circ$의 단위원이다.
\item 도표
$$
\xymatrix{
\mathcal{A}
\rruppertwocell^F{t}
\ar[rr]_(.3){F'}
\rrlowertwocell_{F''}{t'}
& &
\mathcal{B}
\rruppertwocell^G{s}
\ar[rr]_(.3){G'}
\rrlowertwocell_{G''}{s'}
& &
\mathcal{C}
}
$$
가 주어지면
$ (s' \circ s) \star (t' \circ t) = (s' \star t') \circ (s \star t)$이다.
\end{enumerate}
\end{lemma}

\begin{proof}
마지막 명제를 앞의 표기법으로 바꾸면 다음 등식이 된다.
$$
s'_{F''}
\circ
{}_{G'}t'
\circ
s_{F'}
\circ
{}_Gt
=
(s' \circ s)_{F''}
\circ
{}_G(t' \circ t).
$$
가운데 합성에 위의 결과를 적용하면 좌변을
$
s'_{F''}
\circ
s_{F''}
\circ
{}_Gt'
\circ
{}_Gt
$
로 다시 쓸 수 있고, 이는 우변과 같음을 쉽게 확인할 수 있다.
\end{proof}

\noindent
보조정리의 조건 (4)를 달리 말하면, 함자들의 합성과 함자 변환들의 수평 합성이
함자
$$
(\circ, \star) :
\text{Fun}(\mathcal{B}, \mathcal{C})
\times
\text{Fun}(\mathcal{A}, \mathcal{B})
\longrightarrow
\text{Fun}(\mathcal{A}, \mathcal{C})
$$
를 준다는 뜻이다. 그 정의역은 곱범주이다. Definition
\ref{categories-definition-product-category}를 보라.

\section{2-범주}
\label{categories-section-2-categories}

\noindent
스택의 맥락에서 나타나는 (엄밀한) $2$-범주의 정의를 제시할 것이다. 이를
읽기 전에 Section \ref{categories-section-formal-cat-cat}과 Example
\ref{categories-example-2-1-category-of-categories}를 살펴보라. 기본적으로 이 예를
택하여 그 예의 대상, $1$-사상, $2$-사상이 만족하는 모든 규칙을 적는 것이다.

\begin{definition}
\label{categories-definition-2-category}
(엄밀한) {\it $2$-범주} $\mathcal{C}$는 다음 데이터로 이루어진다.
\begin{enumerate}
\item 대상들의 집합 $\Ob(\mathcal{C})$.
\item 각 쌍 $x, y \in \Ob(\mathcal{C})$에 대하여 하나의 범주
$\Mor_\mathcal{C}(x, y)$. $\Mor_\mathcal{C}(x, y)$의 대상을
{\it $1$-사상}이라 하고 $F : x \to y$로 표기한다. 이 $1$-사상들 사이의
사상을 {\it $2$-사상}이라 하고 $t : F' \to F$로 표기한다.
$\Mor_\mathcal{C}(x, y)$에서 $2$-사상들의 합성을 {\it 수직} 합성이라 하고,
$t : F' \to F$와 $t' : F'' \to F'$에 대하여 $t \circ t'$로 표기한다.
\item 각 세쌍 $x, y, z\in \Ob(\mathcal{C})$에 대하여 함자
$$
(\circ, \star) :
\Mor_\mathcal{C}(y, z) \times \Mor_\mathcal{C}(x, y)
\longrightarrow
\Mor_\mathcal{C}(x, z).
$$
왼쪽의 $1$-사상 쌍 $(F, G)$의 상을 $F$와 $G$의 {\it 합성}이라 하고
$F\circ G$로 표기한다. $2$-사상 쌍 $(t, s)$의 상을 {\it 수평} 합성이라
하고 $t \star s$로 표기한다.
\end{enumerate}
이 데이터는 다음 규칙들을 만족해야 한다.
\begin{enumerate}
\item 대상들의 집합과 $1$-사상들의 집합은 $1$-사상의 합성을 갖추면 범주를
이룬다.
\item $2$-사상의 수평 합성은 결합적이다.
\item 항등 $1$-사상 $\text{id}_x$의 항등 $2$-사상
$\text{id}_{\text{id}_x}$는 수평 합성의 단위원이다.
\end{enumerate}
\end{definition}

\noindent
이는 분명 다루기 썩 편한 종류의 대상은 아니다. 반면 어떻게 다루어야 할지가
아주 명확한 예는 많이 있다. 지금까지 가진 유일한 예는, 대상들이 주어진
범주들의 모음이고 $1$-사상들이 그 범주들 사이의 함자이며 $2$-사상들이
함자들의 자연변환인 $2$-범주이다. Section \ref{categories-section-formal-cat-cat}을
보라. 이 글에서 나타나는 모든 $2$-범주는 이 예의 부분 $2$-범주일 것이다.
부분 $2$-범주라는 말의 뜻은 다음과 같다.

\begin{definition}
\label{categories-definition-sub-2-category}
$\mathcal{C}$를 $2$-범주라 하자. $\mathcal{C}$의 {\it 부분 $2$-범주}
$\mathcal{C}'$는 $\Ob(\mathcal{C})$의 부분집합 $\Ob(\mathcal{C}')$와,
모든 $x, y \in \Ob(\mathcal{C}')$에 대하여 범주
$\Mor_\mathcal{C}(x, y)$의 부분범주 $\Mor_{\mathcal{C}'}(x, y)$로
주어진다. 이들은 연산 $\circ$ ($1$-사상의 합성), $\circ$ ($2$-사상의
수직 합성), $\star$ (수평 합성)와 함께 $2$-범주를 이루어야 한다.
\end{definition}

\begin{remark}
\label{categories-remark-big-2-categories}
큰 $2$-범주. 많은 문헌에서는 $2$-범주가 대상들의 모임을 갖도록 허용한다
(다만 ``모임들의 모임''까지 허용되지는 않기를 바란다). 우리도 다음의
경우에는 이러한 ``큰'' $2$-범주를 허용할 것이다(목록은 진행하면서 갱신한다).
\begin{enumerate}
\item 범주들의 $2$-범주 $\textit{Cat}$.
\item 범주들의 $(2, 1)$-범주 $\textit{Cat}$.
\item 군체들의 $2$-범주 $\textit{Groupoids}$; 이는 $(2, 1)$-범주이다.
\item 고정된 범주 위의 섬유화된 범주들의 $2$-범주.
\item 고정된 범주 위의 섬유화된 범주들의 $(2, 1)$-범주.
\item 고정된 범주 위에서 군체들로 섬유화된 범주들의 $2$-범주; 이는
$(2, 1)$-범주이다.
\item 고정된 사이트 위의 스택들의 $2$-범주.
\item 고정된 사이트 위의 스택들의 $(2, 1)$-범주.
\item 고정된 사이트 위의 군체 스택들의 $2$-범주; 이는 $(2, 1)$-범주이다.
\item 고정된 사이트 위의 세토이드 스택들의 $2$-범주; 이는
$(2, 1)$-범주이다.
\item 고정된 스킴 위의 대수적 스택들의 $2$-범주; 이는 $(2, 1)$-범주이다.
\end{enumerate}
Definition \ref{categories-definition-2-1-category}를 보라. 각 경우에 $2$-범주
$\mathcal{C}$의 대상들의 모임은 고유 모임이지만, 모든 대상
$x, y \in \Ob(C)$에 대하여 범주 $\Mor_\mathcal{C}(x, y)$는 우리의
약속에 따른 의미에서 ``작다''는 점에 유의하라.
\end{remark}

\noindent
Section \ref{categories-section-definition-categories}에서 정의한 범주들의 동치 개념은
다음과 같이 $2$-범주라는 더 일반적인 맥락으로 확장된다.

\begin{definition}
\label{categories-definition-equivalence}
$2$-범주의 두 대상 $x, y$가 {\it 동치}라는 것은, $F \circ G$가
$\text{id}_y$와 $2$-동형이고 $G \circ F$가 $\text{id}_x$와 $2$-동형인
$1$-사상 $F : x \to y$와 $G : y \to x$가 존재한다는 뜻이다.
\end{definition}

\noindent
때로는 범주에서 $2$-범주로 가는 함자가 무엇을 뜻하는지 말해야 한다.

\begin{definition}
\label{categories-definition-functor-into-2-category}
$\mathcal{A}$를 범주라 하고 $\mathcal{C}$를 $2$-범주라 하자.
\begin{enumerate}
\item 따로 언급하지 않는 한 보통 범주에서 $2$-범주로 가는 {\it 함자}는
$2$-사상을 무시한다. 달리 말해, 모든 2-사상을 잊어 2-범주로부터 만든
범주로 가는 ``보통'' 함자이다.
\item $\mathcal{A}$에서 2-범주 $\mathcal{C}$로 가는 {\it 약한 함자},
또는 {\it 유사함자} $\varphi$는 다음 데이터로 주어진다.
\begin{enumerate}
\item 사상 $\varphi : \Ob(\mathcal{A}) \to \Ob(\mathcal{C})$,
\item 각 쌍 $x, y\in \Ob(\mathcal{A})$와 각 사상 $f : x \to y$에 대하여
$1$-사상 $\varphi(f) : \varphi(x) \to \varphi(y)$,
\item 각 $x\in \Ob(A)$에 대하여 $2$-사상
$\alpha_x : \text{id}_{\varphi(x)} \to \varphi(\text{id}_x)$,
\item $\mathcal{A}$의 합성 가능한 각 사상 쌍 $f : x \to y$,
$g : y \to z$에 대하여 $2$-사상
$\alpha_{g, f} : \varphi(g \circ f) \to \varphi(g) \circ \varphi(f)$.
\end{enumerate}
이 데이터에는 다음 조건들이 부과된다.
\begin{enumerate}
\item $2$-사상 $\alpha_x$와 $\alpha_{g, f}$는 모두 동형이다.
\item $\mathcal{A}$의 임의의 사상 $f : x \to y$에 대하여
$\alpha_{\text{id}_y, f} = \alpha_y \star \text{id}_{\varphi(f)}$이다.
$$
\xymatrix{
\varphi(x)
\rrtwocell^{\varphi(f)}_{\varphi(f)}{\ \ \ \ \text{id}_{\varphi(f)}}
& &
\varphi(y)
\rrtwocell^{\text{id}_{\varphi(y)}}_{\varphi(\text{id}_y)}{\alpha_y}
& &
\varphi(y)
}
=
\xymatrix{
\varphi(x)
\rrtwocell^{\varphi(f)}_{\varphi(\text{id}_y) \circ \varphi(f)}{\ \ \ \ \alpha_{\text{id}_y, f}}
& &
\varphi(y)
}
$$
\item $\mathcal{A}$의 임의의 사상 $f : x \to y$에 대하여
$\alpha_{f, \text{id}_x} = \text{id}_{\varphi(f)} \star \alpha_x$이다.
\item $\mathcal{A}$의 합성 가능한 임의의 사상 세쌍
$f : w \to x$, $g : x \to y$, $h : y \to z$에 대하여
$$
(\text{id}_{\varphi(h)} \star \alpha_{g, f})
\circ
\alpha_{h, g \circ f}
=
(\alpha_{h, g} \star \text{id}_{\varphi(f)})
\circ
\alpha_{h \circ g, f}
$$
이다. 달리 말해, 대상들이 $1$-사상이고 화살표들이 $2$-사상인 다음 도표가
가환한다.
$$
\xymatrix{
\varphi(h \circ g \circ f)
\ar[d]_{\alpha_{h, g \circ f}}
\ar[rr]_{\alpha_{h \circ g, f}}
& &
\varphi(h \circ g) \circ \varphi(f)
\ar[d]^{\alpha_{h, g} \star \text{id}_{\varphi(f)}} \\
\varphi(h) \circ \varphi(g \circ f)
\ar[rr]^{\text{id}_{\varphi(h)} \star \alpha_{g, f}}
& &
\varphi(h) \circ \varphi(g) \circ \varphi(f)
}
$$
\end{enumerate}
\end{enumerate}
\end{definition}

\noindent
이번에도 이는 다루기 썩 편한 개념은 아니지만 때로 나타난다. 모든 유사함자는
함자와 동형이라는 정리가 있다. 끝으로 {\it $2$-범주 사이의 함자}와
{\it $2$-범주 사이의 유사함자}라는 개념들이 있다. 마지막 개념은 우리를
$3$-범주의 영역으로 이끈다. 우리는 거의 어떤 대가를 치르더라도 이를
정의해야 하는 상황은 피하고 싶다!

\section{(2, 1)-범주}
\label{categories-section-2-1-categories}

\noindent
어떤 $2$-범주들은 모든 $2$-사상이 동형이라는 성질을 갖는다. 이들은 다음에서
중요한 역할을 하며 다루기도 더 쉽다.

\begin{definition}
\label{categories-definition-2-1-category}
(엄밀한) {\it $(2, 1)$-범주}란 모든 $2$-사상이 동형인 $2$-범주이다.
\end{definition}

\begin{example}
\label{categories-example-2-1-category-of-categories}
$2$-범주 $\textit{Cat}$은, Remark \ref{categories-remark-big-2-categories}를 보라,
함자들의 동형만을 $2$-사상으로 허용함으로써 $(2, 1)$-범주로 만들 수 있다.
\end{example}

\noindent
실제로 더 일반적으로, 임의의 $2$-범주 $\mathcal{C}$는 대상과 $1$-사상은
같고 $2$-사상은 $\mathcal{C}$의 가역인 $2$-사상들인 부분 $2$-범주
$\mathcal{C}'$를 고려함으로써 $(2, 1)$-범주를 낳는다. 이 상황에서는
``{\it $\mathcal{C}'$를 $\mathcal{C}$에 결부된 $(2, 1)$-범주라 하자}''와 같이
말할 것이다. 예를 들어 군체들의 $(2, 1)$-범주란 대상들이 군체이고
$1$-사상들이 함자이며 $2$-사상들이 함자들의 동형인 $2$-범주를 뜻한다.
다만 이는 좋지 않은 예인데, 군체들 사이의 함자들 사이의 변환은 자동으로
동형이기 때문이다!

\begin{remark}
\label{categories-remark-other-2-categories}
따라서 위의 Example \ref{categories-example-2-1-category-of-categories}의 구성에는
군체들의 $2$-범주, 고정된 범주 위에서 군체들로 섬유화된 범주들,
또는 스택들 등을 살펴보는 변형들이 있다.
\end{remark}






\section{2-섬유곱}
\label{categories-section-2-fibre-products}

\noindent
이 절에서는 $2$-섬유곱을 도입한다. $\mathcal{C}$가 2-범주라고 하자. 도표
$$
\xymatrix{
w \ar[r] \ar[d] & y \ar[d] \\
x \ar[r] & z }
$$
에서 두 1-사상 $w \to y \to z$와 $w \to x \to z$가 2-동형이면 이 도표가
2-가환한다고 한다. 2-범주에서는 실제로 가환하는 도표를 요구하는 것보다
도표의 2-가환성을 요구하는 편이 더 자연스럽다. (실제로 엄밀한 2-범주를
아예 사용하지 말아야 한다고 말하는 이도 있을 수 있으며, ``약한'' 2-범주에서는
1-사상들의 가환 도표라는 개념조차 의미가 없다.) 이에 따라 섬유곱 개념도
조정해야 한다.

\medskip\noindent
$\mathcal{C}$를 $2$-범주라 하자. $x, y, z\in \Ob(\mathcal{C})$ 및
$f\in \Mor_\mathcal{C}(x, z)$, $g\in \Mor_{\mathcal C}(y, z)$라 하자.
$f$와 $g$의 2-섬유곱을 정의하기 위해 2-가환 도표
$$
\xymatrix{
& w \ar[r]_a \ar[d]_b & x \ar[d]^{f} \\
& y \ar[r]^{g} & z. }
$$
들을 살펴볼 것이다. 범주들의 경우 섬유곱은 그러한 도표들의 범주에서 종말
대상이다. 이에 따라 2-섬유곱은 2-범주에서의 끝 대상이다(아래 정의를 보라).
{\it $f$와 $g$ 위의 $2$-가환 도표들의 $2$-범주}는 다음과 같이 정의되는
$2$-범주이다.
\begin{enumerate}
\item 대상은 위와 같은 네쌍 $(w, a, b, \phi)$이며, 여기서 $\phi$는 가역인
2-사상 $\phi : f \circ a \to g \circ b$이다.
\item $(w', a', b', \phi')$에서 $(w, a, b, \phi)$로 가는 1-사상은
$(k : w' \to w, \alpha : a' \to a \circ k, \beta : b' \to b \circ k)$로
주어지며, 도표
$$
\xymatrix{
f \circ a'
\ar[rr]_{\text{id}_f \star \alpha}
\ar[d]_{\phi'}
& &
f \circ a \circ k
\ar[d]^{\phi \star \text{id}_k}
\\
g \circ b'
\ar[rr]^{\text{id}_g \star \beta}
& &
g \circ b \circ k
}
$$
가 가환해야 한다.
\item 두 번째 $1$-사상
$(k', \alpha', \beta') : (w'', a'', b'', \phi'') \to
(w', \alpha', \beta', \phi')$가 주어졌을 때, $1$-사상의 합성은 규칙
$$
(k, \alpha, \beta) \circ (k', \alpha', \beta') =
(k \circ k',
(\alpha \star \text{id}_{k'}) \circ \alpha',
(\beta \star \text{id}_{k'}) \circ \beta'),
$$
으로 주어진다.
\item 같은 시작점과 끝점을 갖는 $1$-사상
$(k_i, \alpha_i, \beta_i)$, $i = 1, 2$ 사이의 2-사상은 도표
$$
\xymatrix{
a'
\ar[rd]_{\alpha_2}
\ar[r]_{\alpha_1} &
a \circ k_1
\ar[d]^{\text{id}_a \star \delta} &
&
b \circ k_1
\ar[d]_{\text{id}_b \star \delta} &
b'
\ar[l]^{\beta_1}
\ar[ld]^{\beta_2}
\\
&
a \circ k_2
&
&
b \circ k_2
&
}
$$
들이 가환하도록 하는 2-사상 $\delta : k_1 \to k_2$로 주어진다.
\item $2$-사상들의 수직 합성은 $\mathcal{C}$에서 사상 $\delta$들의 수직
합성으로 주어진다.
\item 도표
$$
\xymatrix{
(w'', a'', b'', \phi'')
\rrtwocell^{(k'_1, \alpha'_1, \beta'_1)}_{(k'_2, \alpha'_2, \beta'_2)}{\delta'}
& &
(w', a', b', \phi')
\rrtwocell^{(k_1, \alpha_1, \beta_1)}_{(k_2, \alpha_2, \beta_2)}{\delta}
& &
(w, a, b, \phi)
}
$$
의 수평 합성은 도표
$$
\xymatrix@C=12pc{
(w'', a'', b'', \phi'')
\rtwocell^{(k_1 \circ k'_1, (\alpha_1 \star \text{id}_{k'_1}) \circ \alpha'_1, (\beta_1 \star \text{id}_{k'_1}) \circ \beta'_1)}_{(k_2 \circ k'_2, (\alpha_2 \star \text{id}_{k'_2}) \circ \alpha'_2, (\beta_2 \star \text{id}_{k'_2}) \circ \beta'_2)}{\ \ \ \delta \star \delta'}
&
(w, a, b, \phi)
}
$$
로 주어진다.
\end{enumerate}
$\mathcal{C}$가 실제로 $(2, 1)$-범주이면, 위의 (2)에 있는 사상
$\alpha$와 $\beta$도 자동으로 동형이다\footnote{실제로 $2$-범주의 경우에는
화살표 $\alpha$, $\beta$의 방향을 뒤집거나 심지어 그중 하나의 방향만 뒤집은
2-가환 도표들의 또 다른 2-범주를 정의할 수도 있을 듯하다. 이것이 뒤에서
$(2, 1)$-범주로 제한하는 이유이다.}. 또한 $\mathcal{C}$가 $(2, 1)$-범주이면
$2$-가환 도표들의 $2$-범주도 $(2, 1)$-범주이다.

\begin{definition}
\label{categories-definition-final-object-2-category}
$(2, 1)$-범주 $\mathcal{C}$의 {\it 끝 대상}은 다음을 만족하는 대상 $x$이다.
\begin{enumerate}
\item 모든 $y \in \Ob(\mathcal{C})$에 대하여 사상 $y \to x$가 존재한다.
\item 임의의 두 사상 $y \to x$는 유일한 2-사상에 의해 동형이다.
\end{enumerate}
\end{definition}

\noindent
아마도 더 일반적인 $2$-범주의 경우에는 끝 대상에 여러 유형이 있을 것이다.
여기에는 들어가고 싶지 않으므로 $(2, 1)$인 경우에만 $2$-섬유곱을 정의한다.

\begin{definition}
\label{categories-definition-2-fibre-products}
$\mathcal{C}$를 $(2, 1)$-범주라 하자. $x, y, z\in \Ob(\mathcal{C})$ 및
$f\in \Mor_\mathcal{C}(x, z)$, $g\in \Mor_{\mathcal C}(y, z)$라 하자.
{\it $f$와 $g$의 2-섬유곱}은 위에서 설명한 2-가환 도표들의 범주에서 종말
대상이다. 2-섬유곱이 존재하면 이를 $x \times_z y\in \Ob(\mathcal{C})$로
표기한다. 도표
$$
\xymatrix{
& x \times_z y \ar[r]^{p} \ar[d]_q & x \ar[d]^{f} \\
& y \ar[r]^{g} & z }
$$
가 2-가환하게 하는 필요한 사상들을
$p\in \Mor_{\mathcal C}(x \times_z y, x)$와
$q\in \Mor_{\mathcal C}(x \times_z y, y)$로 표기하고, 이를 보이는 주어진
가역 2-사상을 $\psi : f \circ p \to g \circ q$로 표기한다.
\end{definition}

\noindent
따라서 다음 보편 성질이 성립한다. 임의의 $w\in \Ob(\mathcal{C})$와 사상
$a \in \Mor_{\mathcal C}(w, x)$ 및 $b \in \Mor_\mathcal{C}(w, y)$가 주어지고
2-동형 $\phi : f \circ a \to g\circ b$가 주어지면, 도표
$$
\xymatrix{
w\ar[rrrd]^a \ar@{-->}[rrd]_\gamma \ar[rrdd]_b & & \\
& & x \times_z y \ar[r]_p \ar[d]_q & x \ar[d]^{f} \\
& & y \ar[r]^{g} & z }
$$
를 2-가환하게 하는 $\gamma \in \Mor_{\mathcal C}(w, x \times_z y)$가
존재하여, $a \to p \circ \gamma$와 $b \to q \circ \gamma$를 적절히
선택하면 도표
$$
\xymatrix{
f \circ a \ar[r] \ar[d]_\phi &
f \circ p \circ \gamma
\ar[d]^{\psi \star \text{id}_\gamma}
\\
g\circ b
\ar[r]
&
g \circ q \circ \gamma
}
$$
가 가환한다. 더욱이 $\gamma$는 동형을 제외하면 유일하다. 물론 정확한
성질들은 이보다 더 정밀하다. 뒤에서 필요한 2-섬유곱의 모든 경우는 다음
범주들의 2-범주 안에서의 2-섬유곱 예에서 나온다.

\begin{example}
\label{categories-example-2-fibre-product-categories}
$\mathcal{A}$, $\mathcal{B}$, $\mathcal{C}$를 범주라 하자. 함자
$F : \mathcal{A} \to \mathcal{C}$와 $G : \mathcal{B} \to \mathcal{C}$가
주어졌다고 하자. 범주 $\mathcal{A} \times_\mathcal{C} \mathcal{B}$를
다음과 같이 정의한다.
\begin{enumerate}
\item $\mathcal{A} \times_\mathcal{C} \mathcal{B}$의 대상은 세쌍
$(A, B, f)$이다. 여기서 $A\in \Ob(\mathcal{A})$, $B\in \Ob(\mathcal{B})$이고
$f : F(A) \to G(B)$는 $\mathcal{C}$의 동형이다.
\item 사상 $(A, B, f) \to (A', B', f')$는 쌍 $(a, b)$로 주어진다.
여기서 $a : A \to A'$는 $\mathcal{A}$의 사상이고 $b : B \to B'$는
$\mathcal{B}$의 사상이며 도표
$$
\xymatrix{
F(A) \ar[r]^f \ar[d]^{F(a)} & G(B) \ar[d]^{G(b)} \\
F(A') \ar[r]^{f'} & G(B')
}
$$
가 가환한다.
\end{enumerate}
또한
$p : \mathcal{A} \times_\mathcal{C}\mathcal{B} \to \mathcal{A}$와
$q : \mathcal{A} \times_\mathcal{C}\mathcal{B} \to \mathcal{B}$를
$$
p(A, B, f) = A, \quad q(A, B, f) = B,
$$
로 놓아 함자들을 정의한다. 달리 말해 이들은 망각함자이다. 함자들의 변환
$\psi : F \circ p \to G \circ q$를 정의한다. 대상
$\xi = (A, B, f)$ 위에서는
$\psi_\xi = f : F(p(\xi)) = F(A) \to G(B) = G(q(\xi))$로 주어진다.
\end{example}

\begin{lemma}
\label{categories-lemma-2-fibre-product-categories}
범주들의 $(2, 1)$-범주에서는 $2$-섬유곱이 존재하며 Example
\ref{categories-example-2-fibre-product-categories}의 구성으로 주어진다.
\end{lemma}

\begin{proof}
보편 성질을 확인하자. $\mathcal{W}$를 범주라 하고,
$a : \mathcal{W} \to \mathcal{A}$와 $b : \mathcal{W} \to \mathcal{B}$를
함자라 하며 $t : F \circ a \to G \circ b$를 함자들의 동형이라 하자.

\medskip\noindent
$W \mapsto (a(W), b(W), t_W)$로 주어지는 함자
$\gamma : \mathcal{W} \to \mathcal{A} \times_\mathcal{C}\mathcal{B}$를
고려하자. (이것이 함자임을 확인하는 과정은 생략한다.) 또한 항등식
$p \circ \gamma = a$와 $q \circ \gamma = b$에서 얻어지는
$\alpha : a \to p \circ \gamma$와 $\beta : b \to q \circ \gamma$를
고려하자. 그러면 $(\gamma, \alpha, \beta)$가
$(W, a, b, t)$에서
$(\mathcal{A} \times_\mathcal{C} \mathcal{B}, p, q, \psi)$로 가는 사상임은
명백하다.

\medskip\noindent
두 번째 그러한 사상
$
(k, \alpha', \beta') :
(W, a, b, t) \to (\mathcal{A} \times_\mathcal{C} \mathcal{B}, p, q, \psi)
$
가 주어졌다고 하자. $\mathcal{W}$의 대상 $W$에 대하여
$k(W) = (a_k(W), b_k(W), t_{k, W})$로 쓰자. 따라서
$p(k(W)) = a_k(W)$이며 나머지도 마찬가지이다. 사상 $\alpha'$는 함자적인
사상 $\alpha' : a(W) \to a_k(W)$에 대응한다. 범주들의 $(2, 1)$-범주에서
일하고 있으므로 실제로 각 사상 $a(W) \to a_k(W)$는 동형이다. 이 사상들과
그에 대응하는 $b(W) \to b_k(W)$를 사용하여 동형
$$
\delta_W :
\gamma(W) = (a(W), b(W), t_W)
\longrightarrow
(a_k(W), b_k(W), t_{k, W}) = k(W).
$$
을 얻을 수 있다. $\delta$가 원하는 대로 $2$-가환 도표들의 $2$-범주에서
$\gamma$와 $k$ 사이의 $2$-동형을 정의함을 확인하는 것은 어렵지 않다.
\end{proof}

\begin{remark}
\label{categories-remark-other-description-2-fibre-product}
$\mathcal{A}$, $\mathcal{B}$, $\mathcal{C}$를 범주라 하자. 함자
$F : \mathcal{A} \to \mathcal{C}$와 $G : \mathcal{B} \to \mathcal{C}$가
주어졌다고 하자. $2$-섬유곱
$\mathcal{A} \times_\mathcal{C} \mathcal{B}$의 조금 더 대칭적인 또 다른
구성은 다음과 같다. 대상은 다섯쌍 $(A, B, C, a, b)$이며, 여기서
$A, B, C$는 $\mathcal{A}, \mathcal{B}, \mathcal{C}$의 대상이고
$a : F(A) \to C$와 $b : G(B) \to C$는 동형이다. 사상
$(A, B, C, a, b) \to (A', B', C', a', b')$는 사상들의 세쌍
$A \to A', B \to B', C \to C'$로 주어지며 $a, b, a', b'$와 양립한다.
이것이 $2$-섬유곱을 낳음을 직접 증명할 수 있다. 그러나 다섯쌍들의 범주에서
Example \ref{categories-example-2-fibre-product-categories}에서 구성한 범주로 가는
동치를 함자 $(A, B, C, a, b) \mapsto (A, B, b^{-1} \circ a)$가 준다고
관찰하는 편이 더 쉽다.
\end{remark}

\begin{lemma}
\label{categories-lemma-functoriality-2-fibre-product}
$$
\xymatrix{
& \mathcal{Y} \ar[d]_I \ar[rd]^K & \\
\mathcal{X} \ar[r]^H \ar[rd]^L &
\mathcal{Z} \ar[rd]^M & \mathcal{B} \ar[d]^G \\
& \mathcal{A} \ar[r]^F & \mathcal{C}
}
$$
를 범주들의 $2$-가환 도표라 하자. 동형들의 선택
$\alpha : G \circ K \to M \circ I$와
$\beta : M \circ H \to F \circ L$은 이 상황에 결부된 $2$-섬유곱들의 사상
$$
\mathcal{X} \times_\mathcal{Z} \mathcal{Y}
\longrightarrow
\mathcal{A} \times_\mathcal{C} \mathcal{B}
$$
을 결정한다.
\end{lemma}

\begin{proof}
대상에 대해서는 함자
$$
(X, Y, \phi) \longmapsto (L(X), K(Y),
\alpha^{-1}_Y \circ M(\phi) \circ \beta^{-1}_X)
$$
를, 사상에 대해서는
$$
(a, b) \longmapsto (L(a), K(b))
$$
를 사용하면 된다.
\end{proof}

\begin{lemma}
\label{categories-lemma-equivalence-2-fibre-product}
Lemma \ref{categories-lemma-functoriality-2-fibre-product}와 같은 가정을 하자.
\begin{enumerate}
\item $K$와 $L$이 충실하면 사상
$\mathcal{X} \times_\mathcal{Z} \mathcal{Y} \to
\mathcal{A} \times_\mathcal{C} \mathcal{B}$는 충실하다.
\item $K$와 $L$이 충만충실하고 $M$이 충실하면 사상
$\mathcal{X} \times_\mathcal{Z} \mathcal{Y} \to
\mathcal{A} \times_\mathcal{C} \mathcal{B}$는 충만충실하다.
\item $K$와 $L$이 동치이고 $M$이 충만충실하면 사상
$\mathcal{X} \times_\mathcal{Z} \mathcal{Y} \to
\mathcal{A} \times_\mathcal{C} \mathcal{B}$는 동치이다.
\end{enumerate}
\end{lemma}

\begin{proof}
$(X, Y, \phi)$와 $(X', Y', \phi')$를
$\mathcal{X} \times_\mathcal{Z} \mathcal{Y}$의 대상이라 하자.
$Z = H(X)$로 놓고 $\phi$를 통해 이를 $I(Y)$와 동일시한다. 또한
$\alpha_X$를 통해 $M(Z)$를 $F(L(X))$와 동일시하고, $\beta_Y$를 통해
$M(Z)$를 $G(K(Y))$와 동일시한다. $Z' = H(X')$와 $M(Z')$에 대해서도
마찬가지이다. 사상 위의 사상은
$$
\xymatrix{
\Mor_\mathcal{X}(X, X')
\times_{\Mor_\mathcal{Z}(Z, Z')}
\Mor_\mathcal{Y}(Y, Y')
\ar[d] \\
\Mor_\mathcal{A}(L(X), L(X'))
\times_{\Mor_\mathcal{C}(M(Z), M(Z'))}
\Mor_\mathcal{B}(K(Y), K(Y'))
}
$$
이다. 그러므로 (1)과 (2)가 따른다. 더욱이 $K$와 $L$이 동치이고 $M$이
충만충실하면, 임의의 대상 $(A, B, \phi)$는 다음 이유로 본질적 상 안에 있다.
$L(X) \cong A$이고 $K(Y) \cong B$인 $X$, $Y$를 택하라. 그러면 $M$의
충만충실성에 의해 동형 $H(X) \cong I(Y)$를 찾을 수 있다. 일부 세부 사항은
생략한다.
\end{proof}

\begin{lemma}
\label{categories-lemma-associativity-2-fibre-product}
$$
\xymatrix{
\mathcal{A} \ar[rd] & & \mathcal{C} \ar[ld] \ar[rd] & & \mathcal{E} \ar[ld] \\
& \mathcal{B} & & \mathcal{D}
}
$$
를 범주와 함자들의 도표라 하자. 그러면 범주들의 표준 동형
$$
(\mathcal{A} \times_\mathcal{B} \mathcal{C}) \times_\mathcal{D} \mathcal{E}
\cong
\mathcal{A} \times_\mathcal{B} (\mathcal{C} \times_\mathcal{D} \mathcal{E})
$$
이 존재한다.
\end{lemma}

\begin{proof}
의미를 알 것이라 생각하고, 함자
$$
((A, C, \phi), E, \psi)
\longmapsto
(A, (C, E, \psi), \phi)
$$
를 사용하면 된다.
\end{proof}

\noindent
이제부터 범주가 2개보다 많은 섬유곱을 다룰 때에는 괄호를 쓰지 않는다.

\begin{lemma}
\label{categories-lemma-triple-2-fibre-product-pr02}
$$
\xymatrix{
\mathcal{A} \ar[rd] & & \mathcal{C} \ar[ld] \ar[rd] & & \mathcal{E} \ar[ld] \\
& \mathcal{B} \ar[rd]_F & & \mathcal{D} \ar[ld]^G \\
& & \mathcal{F} &
}
$$
를 범주와 함자들의 가환 도표라 하자. 그러면 범주들의 표준 함자
$$
\text{pr}_{02} :
\mathcal{A} \times_\mathcal{B} \mathcal{C} \times_\mathcal{D} \mathcal{E}
\longrightarrow
\mathcal{A} \times_\mathcal{F} \mathcal{E}
$$
가 존재한다.
\end{lemma}

\begin{proof}
$\mathcal{A} \times_\mathcal{B} \mathcal{C}
\times_\mathcal{D} \mathcal{E}$를
$(\mathcal{A} \times_\mathcal{B} \mathcal{C})
\times_\mathcal{D} \mathcal{E}$로 쓰면, 의미를 알 것이라 생각하고 함자
$$
((A, C, \phi), E, \psi)
\longmapsto
(A, E, G(\psi) \circ F(\phi))
$$
를 사용하면 된다.
\end{proof}

\begin{lemma}
\label{categories-lemma-2-fibre-product-erase-factor}
$$
\mathcal{A} \to
\mathcal{B} \leftarrow \mathcal{C} \leftarrow \mathcal{D}
$$
를 범주와 함자들의 도표라 하자. 그러면 범주들의 표준 동형
$$
\mathcal{A} \times_\mathcal{B} \mathcal{C} \times_\mathcal{C} \mathcal{D}
\cong
\mathcal{A} \times_\mathcal{B} \mathcal{D}
$$
이 존재한다.
\end{lemma}

\begin{proof}
생략한다.
\end{proof}

\noindent
이는 이러한 $2$-섬유곱을 보통 섬유곱처럼 다룰 수 있다는 뜻이라고 주장한다.
뒤에서 실제로 나타나는 보조정리들을 몇 가지 더 제시한다.

\begin{lemma}
\label{categories-lemma-diagonal-1}
$$
\xymatrix{
\mathcal{C}_3 \ar[r] \ar[d] & \mathcal{S} \ar[d]^\Delta \\
\mathcal{C}_1 \times \mathcal{C}_2 \ar[r]^{G_1 \times G_2} &
\mathcal{S} \times \mathcal{S}
}
$$
를 범주들의 $2$-섬유곱이라 하자. 그러면 표준 동형
$\mathcal{C}_3 \cong
\mathcal{C}_1 \times_{G_1, \mathcal{S}, G_2} \mathcal{C}_2$가 존재한다.
\end{lemma}

\begin{proof}
$\mathcal{C}_3$가 Example \ref{categories-example-2-fibre-product-categories}에서 구성한
범주 $(\mathcal{C}_1 \times \mathcal{C}_2)\times_{\mathcal{S} \times \mathcal{S}}
\mathcal{S}$라고 가정해도 된다. 따라서 대상은 세쌍
$((X_1, X_2), S, \phi)$이며, 여기서
$\phi = (\phi_1, \phi_2) : (G_1(X_1), G_2(X_2)) \to (S, S)$는 동형이다.
그러므로 이 대상에 세쌍 $(X_1, X_2, \phi_2^{-1} \circ \phi_1)$을 대응시킬
수 있다. 반대로 $(X_1, X_2, \psi)$가
$\mathcal{C}_1 \times_{G_1, \mathcal{S}, G_2} \mathcal{C}_2$의 대상이면,
이에 세쌍 $((X_1, X_2), G_2(X_2), (\psi, \text{id}_{G_2(X_2)}))$을 대응시킬
수 있다. 이 구성들이 서로 역인 함자들을 준다고 주장한다. 사상을 처리하는
방법과 이들이 서로 역임을 보이는 과정은 생략한다.
\end{proof}

\begin{lemma}
\label{categories-lemma-diagonal-2}
$$
\xymatrix{
\mathcal{C}' \ar[r] \ar[d] & \mathcal{S} \ar[d]^\Delta \\
\mathcal{C} \ar[r]^{(G_1, G_2)} &
\mathcal{S} \times \mathcal{S}
}
$$
를 범주들의 $2$-섬유곱이라 하자. 그러면 표준 동형
$$
\mathcal{C}' \cong
(\mathcal{C} \times_{G_1, \mathcal{S}, G_2} \mathcal{C})
\times_{(p, q), \mathcal{C} \times \mathcal{C}, \Delta}
\mathcal{C}.
$$
이 존재한다.
\end{lemma}

\begin{proof}
우변의 대상은 $((C_1, C_2, \phi), C_3, \psi)$로 주어진다. 여기서
$\phi : G_1(C_1) \to G_2(C_2)$는 동형이고
$\psi = (\psi_1, \psi_2) : (C_1, C_2) \to (C_3, C_3)$는 동형이다.
따라서 이에 $\mathcal{C}'$의 대상인 세쌍
$(C_3, G_1(C_1), (G_1(\psi_1^{-1}), \phi^{-1} \circ G_2(\psi_2^{-1})))$을
대응시킬 수 있다. 세부 사항은 생략한다.
\end{proof}

\begin{lemma}
\label{categories-lemma-fibre-product-after-map}
$\mathcal{A} \to \mathcal{C}$, $\mathcal{B} \to \mathcal{C}$ 및
$\mathcal{C} \to \mathcal{D}$를 범주들 사이의 함자라 하자. 그러면 도표
$$
\xymatrix{
\mathcal{A} \times_\mathcal{C} \mathcal{B} \ar[d] \ar[r] &
\mathcal{A} \times_\mathcal{D} \mathcal{B} \ar[d] \\
\mathcal{C} \ar[r]^-{\Delta_{\mathcal{C}/\mathcal{D}}} \ar[r] &
\mathcal{C} \times_\mathcal{D} \mathcal{C}
}
$$
는 $2$-섬유곱 도표이다.
\end{lemma}

\begin{proof}
생략한다.
\end{proof}

\begin{lemma}
\label{categories-lemma-base-change-diagonal}
$$
\xymatrix{
\mathcal{U} \ar[d] \ar[r] & \mathcal{V} \ar[d] \\
\mathcal{X} \ar[r] & \mathcal{Y}
}
$$
를 범주들의 $2$-섬유곱이라 하자. 그러면 도표
$$
\xymatrix{
\mathcal{U} \ar[d] \ar[r] &
\mathcal{U} \times_\mathcal{V} \mathcal{U} \ar[d] \\
\mathcal{X} \ar[r] &
\mathcal{X} \times_\mathcal{Y} \mathcal{X}
}
$$
는 $2$-카르테시안이다.
\end{lemma}

\begin{proof}
이는 $2$-섬유곱을 갖는 임의의 $(2, 1)$-범주에서 성립하는 순수하게
$2$-범주론적인 명제이다. 구체적으로 다음 동치들의 사슬에서 따른다.
\begin{align*}
\mathcal{X} \times_{(\mathcal{X} \times_\mathcal{Y} \mathcal{X})}
(\mathcal{U} \times_\mathcal{V} \mathcal{U})
& =
\mathcal{X} \times_{(\mathcal{X} \times_\mathcal{Y} \mathcal{X})}
((\mathcal{X} \times_\mathcal{Y} \mathcal{V})
\times_\mathcal{V} (\mathcal{X} \times_\mathcal{Y} \mathcal{V})) \\
& =
\mathcal{X} \times_{(\mathcal{X} \times_\mathcal{Y} \mathcal{X})}
(\mathcal{X} \times_\mathcal{Y} \mathcal{X}
\times_\mathcal{Y} \mathcal{V}) \\
& =
\mathcal{X} \times_\mathcal{Y} \mathcal{V} = \mathcal{U}
\end{align*}
Lemmas \ref{categories-lemma-associativity-2-fibre-product}과
\ref{categories-lemma-2-fibre-product-erase-factor}를 보라.
\end{proof}







\section{범주 위의 범주들}
\label{categories-section-categories-over-categories}

\noindent
이 절에서는 함자 $p : \mathcal{S} \to \mathcal{C}$를 다룬다.
우리는 $\mathcal{S}$가 위에 있고 $\mathcal{C}$가 아래에 있다고 생각한다.
무엇을 뜻하는지 명확히 하기 위해 $\mathcal{C}$ 위의 범주들로 이루어진
$2$-범주를 정의한다.

\begin{definition}
\label{categories-definition-categories-over-C}
$\mathcal{C}$를 범주라 하자.
{\it $\mathcal{C}$ 위의 범주들의 $2$-범주}는 다음과 같이 정의되는
$2$-범주이다.
\begin{enumerate}
\item 그 대상은 함자 $p : \mathcal{S} \to \mathcal{C}$이다.
\item 그 $1$-사상 $(\mathcal{S}, p) \to (\mathcal{S}', p')$은
$p' \circ G = p$를 만족하는 함자 $G : \mathcal{S} \to \mathcal{S}'$이다.
\item
$G, H : (\mathcal{S}, p) \to (\mathcal{S}', p')$에 대한 그 $2$-사상
$t : G \to H$는 모든 $x \in \Ob(\mathcal{S})$에 대하여
$p'(t_x) = \text{id}_{p(x)}$를 만족하는 함자의 사상이다.
\end{enumerate}
이 상황에서 $(\mathcal{S}, p)$와 $(\mathcal{S}', p')$ 사이의
$1$-사상들의 범주를
$$
\Mor_{\textit{Cat}/\mathcal{C}}(\mathcal{S}, \mathcal{S}')
$$
로 나타내겠다.
\end{definition}

\noindent
이 $2$-범주에서 수평 합성과 수직 합성을
Section \ref{categories-section-formal-cat-cat}에서 $\textit{Cat}$에 대해 한 것과
똑같이 정의한다. $2$-범주의 공리들은 $\textit{Cat}$에서 성립하는 것과
같은 이유로 성립한다. 또한 $\textit{Cat}$에서 공리들이 성립한다는
사실을 이용하여, 예컨대 ``$\mathcal{C}$ 위의 $2$-사상들을 수직 합성하면
다시 $\mathcal{C}$ 위의 $2$-사상이 된다''와 같은 것을 확인해도 된다.
이는 명백하다.

\medskip\noindent
공간 사이 사상의 섬유와 마찬가지로 여기에도 섬유 범주라는 개념과
이 상황에 관련된 몇 가지 올림 개념이 있다.

\begin{definition}
\label{categories-definition-fibre-category}
$\mathcal{C}$를 범주라 하고,
$p : \mathcal{S} \to \mathcal{C}$를 $\mathcal{C}$ 위의 범주라 하자.
\begin{enumerate}
\item 대상 $U\in \Ob(\mathcal{C})$ 위의 {\it 섬유 범주}는 대상들이
$$
\Ob(\mathcal{S}_U) = \{x\in \Ob(\mathcal{S}) :
p(x) = U\}
$$
이고 사상들이
$$
\Mor_{\mathcal{S}_U}(x, y) = \{ \phi \in \Mor_\mathcal{S}(x, y) :
p(\phi) = \text{id}_U\}.
$$
인 범주 $\mathcal{S}_U$이다.
\item 대상 $U \in \Ob(\mathcal{C})$의 {\it 올림}은
$p(x) = U$를 만족하는 대상 $x\in \Ob(\mathcal{S})$, 즉
$x\in \Ob(\mathcal{S}_U)$이다. 또한 때때로
{\it $x$가 $U$ 위에 놓인다}고 말한다.
\item 마찬가지로 $\mathcal{C}$의 사상 $f : V \to U$의 {\it 올림}은
$p(\phi) = f$를 만족하는 $\mathcal{S}$의 사상 $\phi : y \to x$이다.
때때로 {\it $\phi$가 $f$ 위에 놓인다}고 말한다.
\end{enumerate}
\end{definition}

\noindent
여기서 몇 가지 사실을 관찰할 수 있다. 예를 들어
$F : (\mathcal{S}, p) \to (\mathcal{S}', p')$가
$\mathcal{C}$ 위의 범주들의 $1$-사상이면 $F$는 섬유 범주들의 함자
$F : \mathcal{S}_U \to \mathcal{S}'_U$를 유도한다.
$2$-사상에 대해서도 마찬가지이다.

\medskip\noindent
다음은 $\mathcal{C}$ 위의 범주들의 $(2, 1)$-범주에서 $2$-섬유곱을
기술하는 필수적인 보조정리이다.

\begin{lemma}
\label{categories-lemma-2-product-categories-over-C}
$\mathcal{C}$를 범주라 하자. $\mathcal{C}$ 위의 범주들의
$(2, 1)$-범주는 2-섬유곱을 갖는다.
$F : \mathcal{X} \to \mathcal{S}$와
$G : \mathcal{Y} \to \mathcal{S}$가 $\mathcal{C}$ 위의 범주들의
사상이라고 하자. 명시적인 2-섬유곱
$\mathcal{X} \times_\mathcal{S}\mathcal{Y}$는 다음과 같이 주어진다.
\begin{enumerate}
\item $\mathcal{X} \times_\mathcal{S} \mathcal{Y}$의 대상은 사중항
$(U, x, y, f)$이다. 여기서 $U \in \Ob(\mathcal{C})$,
$x\in \Ob(\mathcal{X}_U)$, $y\in \Ob(\mathcal{Y}_U)$이고,
$f : F(x) \to G(y)$는 $\mathcal{S}_U$에서의 동형사상이다.
\item 사상 $(U, x, y, f) \to (U', x', y', f')$은 쌍
$(a, b)$으로 주어진다. 여기서 $a : x \to x'$는 $\mathcal{X}$의
사상이고 $b : y \to y'$는 $\mathcal{Y}$의 사상이며 다음을 만족한다.
\begin{enumerate}
\item $a$와 $b$가 유도하는 사상 $U \to U'$이 같고,
\item 도표
$$
\xymatrix{
F(x) \ar[r]^f \ar[d]^{F(a)} & G(y) \ar[d]^{G(b)} \\
F(x') \ar[r]^{f'} & G(y')
}
$$
가 가환한다.
\end{enumerate}
\end{enumerate}
이 경우 함자 $p : \mathcal{X} \times_\mathcal{S}\mathcal{Y} \to \mathcal{X}$와
$q : \mathcal{X} \times_\mathcal{S}\mathcal{Y} \to \mathcal{Y}$는 망각 함자이다.
변환 $\psi : F \circ p \to G \circ q$는 대상
$\xi = (U, x, y, f)$에서
$\psi_\xi = f : F(p(\xi)) = F(x) \to G(y) = G(q(\xi))$로 주어진다.
\end{lemma}

\begin{proof}
보편 성질을 확인하자. $p_\mathcal{W} : \mathcal{W}\to \mathcal{C}$를
$\mathcal{C}$ 위의 범주라 하고, $X : \mathcal{W} \to \mathcal{X}$와
$Y : \mathcal{W} \to \mathcal{Y}$를 $\mathcal{C}$ 위의 함자라 하며,
$t : F \circ X \to G \circ Y$를 $\mathcal{C}$ 위의 함자들의
동형사상이라 하자. 원하는 함자
$\gamma : \mathcal{W} \to \mathcal{X} \times_\mathcal{S} \mathcal{Y}$는
$W \mapsto (p_\mathcal{W}(W), X(W), Y(W), t_W)$로 주어진다.
세부 사항은 생략한다. Lemma \ref{categories-lemma-2-fibre-product-categories}와 비교하라.
\end{proof}

\begin{example}
\label{categories-example-different-2-fibre-products}
Lemma \ref{categories-lemma-2-product-categories-over-C}에서 주어진 범주 위의
범주들의 $2$-섬유곱 구성과 Lemma \ref{categories-lemma-2-fibre-product-categories}에서
주어진 범주들의 2-섬유곱 구성은
(Example \ref{categories-example-2-fibre-product-categories}에서처럼) 일반적으로
서로 동치가 아닌 결과를 낸다. 실제로 $\mathcal{S}$를 대상 하나와
화살표 두 개를 갖는 준군 범주라 하고, $\mathcal{X}$를 대상 하나를 갖는
이산 범주라 하자. 범주로서 $2$-섬유곱
$\mathcal{X} \times_\mathcal{S} \mathcal{X}$을 취하면 대상 두 개를 갖는
이산 범주가 된다. 그러나 이들을 모두 $\mathcal{S}$ 위의 범주로 보면,
$\mathcal{S}$ 위의 범주로서 $2$-섬유곱
$\mathcal{X} \times_\mathcal{S} \mathcal{X}$은 대상 하나를 갖는
이산 범주이다. 차이는 (Lemma \ref{categories-lemma-2-product-categories-over-C}의
표기에서) 첫 번째 상황에서는 임의의 비교 동형사상 $f$를 선택할 수
있지만 두 번째 상황에서는 항등 화살표만 선택할 수 있다는 점이다.
\end{example}

\begin{lemma}
\label{categories-lemma-fibre-2-fibre-product-categories-over-C}
$\mathcal{C}$를 범주라 하자.
$f : \mathcal{X} \to \mathcal{S}$와
$g : \mathcal{Y} \to \mathcal{S}$를 $\mathcal{C}$ 위의 범주들의
사상이라 하자. $\mathcal{C}$의 임의의 대상 $U$에 대하여 섬유 범주들의
다음 항등식이 성립한다.
$$
\left(\mathcal{X} \times_\mathcal{S}\mathcal{Y}\right)_U
=
\mathcal{X}_U \times_{\mathcal{S}_U} \mathcal{Y}_U
$$
\end{lemma}

\begin{proof}
생략한다.
\end{proof}







\section{섬유화된 범주들}
\label{categories-section-fibred-categories}

\noindent
섬유화된 범주에 관하여 아주 간단히 논의할 필요가 있다.

\medskip\noindent
$p : \mathcal{S} \to \mathcal{C}$를 $\mathcal{C}$ 위의 범주라 하자.
$p(x) = U$인 대상 $x \in \mathcal{S}$와 사상 $f : V \to U$가 주어졌을 때
일종의 ``섬유곱 $V \times_U x$''(또는 $V \to U$에 의한 $x$의
{\it 기저변환})을 취해 볼 수 있다. 구체적으로 ``$V \times_U x$''로
향하는 대상 $z \in \mathcal{S}$로부터의 사상은
$p(\varphi) = f \circ g$를 만족하는 쌍 $(\varphi, g)$,
즉 $\varphi : z \to x$, $g : p(z) \to V$로 주어져야 한다.
그림으로 나타내면 다음과 같다.
$$
\xymatrix{
z \ar@{~>}[d]^p \ar@{-}[r] &
? \ar[r] \ar@{~>}[d]^p &
x \ar@{~>}[d]^p \\
p(z) \ar[r] & V \ar[r]^f & U
}
$$
그러한 사상 $V \times_U x \to x$가 존재하면 이를 강한 카르테시안
사상이라고 한다.

\begin{definition}
\label{categories-definition-cartesian-over-C}
$\mathcal{C}$를 범주라 하고,
$p : \mathcal{S} \to \mathcal{C}$를 $\mathcal{C}$ 위의 범주라 하자.
{\it 강한 카르테시안 사상}, 더 정확히는
{\it 강한 $\mathcal{C}$-카르테시안 사상}이란 다음 성질을 갖는
$\mathcal{S}$의 사상 $\varphi : y \to x$이다. 모든
$z \in \Ob(\mathcal{S})$에 대하여
$$
\Mor_\mathcal{S}(z, y)
\longrightarrow
\Mor_\mathcal{S}(z, x)
\times_{\Mor_\mathcal{C}(p(z), p(x))}
\Mor_\mathcal{C}(p(z), p(y)),
$$
$\psi \longmapsto (\varphi \circ \psi, p(\psi))$로 주어지는 사상이
전단사이다.
\end{definition}

\noindent
Yoneda Lemma \ref{categories-lemma-yoneda}에 따라 다음에 유의하자.
$U \in \Ob(\mathcal{C})$ 위에 놓인 $x \in \Ob(\mathcal{S})$와
$\mathcal{C}$의 사상 $f : V \to U$가 주어졌을 때,
$p(\varphi) = f$인 강한 카르테시안 사상 $\varphi : y \to x$가
존재한다면 $(y, \varphi)$는 유일한 동형사상을 제외하고 유일하다.
이는 위 정의에서 명백하다. 실제로 함자
$$
z
\longmapsto
\Mor_\mathcal{S}(z, x)
\times_{\Mor_\mathcal{C}(p(z), U)}
\Mor_\mathcal{C}(p(z), V)
$$
는 자료 $(x, U, f : V \to U)$에만 의존한다. 따라서 때때로
$f$의 올림인 강한 카르테시안 사상을 $V \times_U x \to x$ 또는
$f^*x \to x$로 나타내겠다.

\begin{lemma}
\label{categories-lemma-composition-cartesian}
$\mathcal{C}$를 범주라 하고,
$p : \mathcal{S} \to \mathcal{C}$를 $\mathcal{C}$ 위의 범주라 하자.
\begin{enumerate}
\item 두 강한 카르테시안 사상의 합성은 강한 카르테시안 사상이다.
\item $\mathcal{S}$의 모든 동형사상은 강한 카르테시안 사상이다.
\item $p(\varphi)$가 동형사상인 강한 카르테시안 사상 $\varphi$는
동형사상이다.
\end{enumerate}
\end{lemma}

\begin{proof}
(1)의 증명. $\varphi : y \to x$와 $\psi : z \to y$를 강한
카르테시안 사상이라 하자. $t$를 $\mathcal{S}$의 임의의 대상이라 하자.
그러면
\begin{align*}
& \Mor_\mathcal{S}(t, z) \\
& =
\Mor_\mathcal{S}(t, y)
\times_{\Mor_\mathcal{C}(p(t), p(y))}
\Mor_\mathcal{C}(p(t), p(z)) \\
& =
\Mor_\mathcal{S}(t, x)
\times_{\Mor_\mathcal{C}(p(t), p(x))}
\Mor_\mathcal{C}(p(t), p(y))
\times_{\Mor_\mathcal{C}(p(t), p(y))}
\Mor_\mathcal{C}(p(t), p(z)) \\
& =
\Mor_\mathcal{S}(t, x)
\times_{\Mor_\mathcal{C}(p(t), p(x))}
\Mor_\mathcal{C}(p(t), p(z))
\end{align*}
이므로 $z \to x$는 강한 카르테시안 사상이다.

\medskip\noindent
(2)의 증명. $y \to x$를 동형사상이라 하자. 그러면
$p(y) \to p(x)$도 동형사상이다. 따라서 다음 사상은 전단사이다.
\[
\Mor_\mathcal{C}(p(z), p(y)) \to
\Mor_\mathcal{C}(p(z), p(x))
\]
그러므로
\[
\Mor_\mathcal{S}(z, x)
\times_{\Mor_\mathcal{C}(p(z), p(x))}
\Mor_\mathcal{C}(p(z), p(y))
\]
는 $\Mor_\mathcal{S}(z, x)$와 전단사이다.
Definition \ref{categories-definition-cartesian-over-C}에 표시된 사상도
$y \to x$가 동형사상이므로 전단사이고, 따라서 $y \to x$는 강한
카르테시안 사상이다.

\medskip\noindent
(3)의 증명. $\varphi : y \to x$를 강한 카르테시안 사상이라 하고
$p(\varphi) : p(y) \to p(x)$가 동형사상이라고 하자. 정의에서
$z = x$로 놓으면 $(\text{id}_x, p(\varphi)^{-1})$은 유일한 사상
$\chi : x \to y$에서 온다. $\chi$가 $\varphi$의 역이라는 확인은
생략한다.
\end{proof}

\begin{lemma}
\label{categories-lemma-cartesian-over-cartesian}
$F : \mathcal{A} \to \mathcal{B}$와 $G : \mathcal{B} \to \mathcal{C}$를
합성 가능한 범주 사이의 함자들이라 하자. $x \to y$를
$\mathcal{A}$의 사상이라 하자. $x \to y$가 강한
$\mathcal{B}$-카르테시안 사상이고 $F(x) \to F(y)$가 강한
$\mathcal{C}$-카르테시안 사상이면 $x \to y$는 강한
$\mathcal{C}$-카르테시안 사상이다.
\end{lemma}

\begin{proof}
정의에서 바로 따른다.
\end{proof}

\begin{lemma}
\label{categories-lemma-strongly-cartesian-fibre-product}
$\mathcal{C}$를 범주라 하고,
$p : \mathcal{S} \to \mathcal{C}$를 $\mathcal{C}$ 위의 범주라 하자.
$x \to y$와 $z \to y$를 $\mathcal{S}$의 사상들이라 하자.
다음을 가정하자.
\begin{enumerate}
\item $x \to y$는 강한 카르테시안 사상이다.
\item $p(x) \times_{p(y)} p(z)$가 존재한다.
\item $p(w) = p(x) \times_{p(y)} p(z)$이고
$p(a) = \text{pr}_2 : p(x) \times_{p(y)} p(z) \to p(z)$인
$\mathcal{S}$의 강한 카르테시안 사상 $a : w \to z$가 존재한다.
\end{enumerate}
그러면 섬유곱 $x \times_y z$가 존재하고 $w$와 동형이다.
\end{lemma}

\begin{proof}
$x \to y$가 강한 카르테시안 사상이므로 $p(b) = \text{pr}_1$인
유일한 사상 $b : w \to x$가 존재한다. $w$가 섬유곱임을 보이기 위해
다음을 계산한다.
\begin{align*}
& \Mor_\mathcal{S}(t, w) \\
& = \Mor_\mathcal{S}(t, z)
\times_{\Mor_\mathcal{C}(p(t), p(z))}
\Mor_\mathcal{C}(p(t), p(w)) \\
& = \Mor_\mathcal{S}(t, z)
\times_{\Mor_\mathcal{C}(p(t), p(z))}
(\Mor_\mathcal{C}(p(t), p(x))
\times_{\Mor_\mathcal{C}(p(t), p(y))}
\Mor_\mathcal{C}(p(t), p(z))) \\
& = \Mor_\mathcal{S}(t, z)
\times_{\Mor_\mathcal{C}(p(t), p(y))}
\Mor_\mathcal{C}(p(t), p(x)) \\
& = \Mor_\mathcal{S}(t, z)
\times_{\Mor_\mathcal{S}(t, y)}
\Mor_\mathcal{S}(t, y)
\times_{\Mor_\mathcal{C}(p(t), p(y))}
\Mor_\mathcal{C}(p(t), p(x)) \\
& = \Mor_\mathcal{S}(t, z)
\times_{\Mor_\mathcal{S}(t, y)}
\Mor_\mathcal{S}(t, x)
\end{align*}
이것이 원하는 바이다. 첫 번째 등식은 $a : w \to z$가 강한
카르테시안 사상이므로 성립하고, 마지막 등식은 $x \to y$가 강한
카르테시안 사상이므로 성립한다.
\end{proof}

\begin{definition}
\label{categories-definition-fibred-category}
$\mathcal{C}$를 범주라 하고,
$p : \mathcal{S} \to \mathcal{C}$를 $\mathcal{C}$ 위의 범주라 하자.
$U \in \Ob(\mathcal{C})$ 위에 놓인 임의의
$x \in \Ob(\mathcal{S})$와 $\mathcal{C}$의 임의의 사상
$f : V \to U$가 주어졌을 때, $f$ 위에 놓이는 강한 카르테시안 사상
$f^*x \to x$가 존재하면 $\mathcal{S}$를
{\it $\mathcal{C}$ 위의 섬유화된 범주}라고 한다.
\end{definition}

\noindent
$p : \mathcal{S} \to \mathcal{C}$가 섬유화된 범주라고 하자.
정의에서처럼 각 $f : V \to U$와 $x\in \Ob(\mathcal{S}_U)$에 대하여
$f$ 위에 놓인 강한 카르테시안 사상 $f^\ast x \to x$를 선택할 수 있다.
선택 공리에 의해 모든 $f: V \to U = p(x)$에 대한 $f^*x \to x$를
동시에 선택할 수 있다. $\mathcal{S}_U$의 모든 사상
$\phi : x \to x'$와 $f : V \to U$에 대하여 다음 도표가 가환하도록 하는
$\mathcal{S}_V$의 유일한 사상
$f^\ast \phi : f^\ast x \to f^\ast x'$가 존재한다고 주장한다.
$$
\xymatrix{
f^\ast x \ar[r]_{f^\ast \phi} \ar[d] & f^\ast x' \ar[d] \\
x \ar[r]^{\phi} & x' }
$$
실제로 $f^*x' \to x'$가 강한 카르테시안 사상이므로 이 화살표는
존재하고 유일하다. 이 화살표의 유일성에 의해 이제 사상에도 정의된
$f^\ast$는 함자 $ f^\ast : \mathcal{S}_U \to \mathcal{S}_V$가 된다.

\begin{definition}
\label{categories-definition-pullback-functor-fibred-category}
$p : \mathcal{S} \to \mathcal{C}$가 섬유화된 범주라고 하자.
\begin{enumerate}
\item $p : \mathcal{S} \to \mathcal{C}$에 대한
{\it 당김의 선택}\footnote{이는 아마 표준 용어가 아닐 것이다.
일부 문헌에서는 이를 ``절개(cleavage)''라고 부르지만 적절하지 않은
이미지를 떠올리게 한다. ``쪼갬(cleaving)''이 더 나은 말일 수도 있다.
이와 관련된 개념으로 ``분할(splitting)''이 있으나, 많은 문헌에서
``분할''은 합성 가능한 임의의 사상 쌍에 대해
$g^*f^* = (f \circ g)^*$가 되도록 하는 당김의 선택을 뜻한다.
Definition \ref{categories-definition-split-fibred-category}도 비교하라.}이란
$\mathcal{C}$의 임의의 사상 $f: V \to U$와 임의의
$x \in \Ob(\mathcal{S}_U)$에 대하여 $f$ 위에 놓인 강한 카르테시안
사상 $f^\ast x \to x$를 하나씩 선택한 것이다.
\item 당김의 선택이 주어졌을 때, $\mathcal{C}$의 임의의 사상
$f : V \to U$에 대하여 위에서 기술한 함자
$f^* : \mathcal{S}_U \to \mathcal{S}_V$를 (위에서 한 선택
$f^*x \to x$에 대응하는) {\it 당김 함자}라고 한다.
\end{enumerate}
\end{definition}

\noindent
물론 당김의 선택이 $\text{id}_U^*x = x$를 만족한다고 언제든 가정할
수 있다. 다만 실제로는 당김에 추가 조건을 부과하지 않으면 쓸모없는
성질이다.

\begin{lemma}
\label{categories-lemma-fibred}
$p : \mathcal{S} \to \mathcal{C}$가 섬유화된 범주이고,
$p : \mathcal{S} \to \mathcal{C}$에 대한 당김의 선택이 주어졌다고 하자.
\begin{enumerate}
\item 합성 가능한 임의의 사상 쌍 $f : V \to U$,
$g : W \to V$에 대하여 함자 $\mathcal{S}_U \to \mathcal{S}_W$ 사이의
유일한 동형사상
$$
\alpha_{g, f} :
(f \circ g)^\ast
\longrightarrow
g^\ast \circ f^\ast
$$
이 존재하여, 모든 $y\in \Ob(\mathcal{S}_U)$에 대해 다음 도표를
가환하게 한다.
$$
\xymatrix{
g^\ast f^\ast y \ar[r]
&
f^\ast y \ar[d] \\
(f \circ g)^\ast y \ar[r]
\ar[u]^{(\alpha_{g, f})_y}
&
y
}
$$
\item $f = \text{id}_U$이면 함자
$\mathcal{S}_U \to \mathcal{S}_U$ 사이의 표준 동형사상
$\alpha_U : \text{id} \to (\text{id}_U)^*$가 존재한다.
\item 사중항
$(U \mapsto \mathcal{S}_U, f \mapsto f^*, \alpha_{g, f}, \alpha_U)$는
$\mathcal{C}^{opp}$에서 범주들의 $(2, 1)$-범주로 가는 유사함자를
정의한다. Definition \ref{categories-definition-functor-into-2-category}를 보라.
\end{enumerate}
\end{lemma}

\begin{proof}
사실 (1)의 가환 도표가 섬유 범주 $\mathcal{S}_W$의 사상
$(\alpha_{g, f})_y$를 유일하게 결정한다는 것은 명백하다.
사상 $(f \circ g)^*y \to y$와 합성
$g^*f^*y \to f^*y \to y$가 모두 $f \circ g$를 올리는 강한
카르테시안 사상이므로 이는 동형사상이다
(Definition \ref{categories-definition-cartesian-over-C} 뒤의 논의와
Lemma \ref{categories-lemma-composition-cartesian}를 보라). 같은 방식으로
$\text{id}_x : x \to x$는 분명히 $\text{id}_U$ 위에서 강한
카르테시안 사상이고(여기서 $U = p(x)$), 따라서 동형사상
$(\alpha_U)_x : x \to (\text{id}_U)^*x$가 존재한다.
(물론 $f$가 항등사상일 때마다 $f^*x = x$라고 미리 가정할 수도 있지만,
일반성을 위해 이를 가정하지 않는 편이 낫다.) 이렇게 얻은
$\alpha_{g, f}$와 $\alpha_U$가 함자들의 변환이라는 확인은 생략한다.
(3)의 확인도 생략한다.
\end{proof}

\begin{lemma}
\label{categories-lemma-fibred-equivalent}
$\mathcal{C}$를 범주라 하고, $\mathcal{S}_1$, $\mathcal{S}_2$를
$\mathcal{C}$ 위의 범주들이라 하자. $\mathcal{S}_1$과
$\mathcal{S}_2$가 $\mathcal{C}$ 위의 범주로서 동치라고 하자.
그러면 $\mathcal{S}_1$이 $\mathcal{C}$ 위에서 섬유화된 범주일
필요충분조건은 $\mathcal{S}_2$가 $\mathcal{C}$ 위에서 섬유화된
범주인 것이다.
\end{lemma}

\begin{proof}
주어진 함자들을 $p_i : \mathcal{S}_i \to \mathcal{C}$로 나타내자.
$F : \mathcal{S}_1 \to \mathcal{S}_2$,
$G : \mathcal{S}_2 \to \mathcal{S}_1$을 $\mathcal{C}$ 위의 함자들이라
하고, $i : F \circ G \to \text{id}_{\mathcal{S}_2}$,
$j : G \circ F \to \text{id}_{\mathcal{S}_1}$를 $\mathcal{C}$ 위의
함자들의 동형사상이라 하자. 이 경우 $F$가 강한 카르테시안 사상을
강한 카르테시안 사상으로 보낸다고 주장한다. 실제로
$\varphi : y \to x$가 $\mathcal{S}_1$에서 강한 카르테시안 사상이라고
하자. $f : V \to U$를 $p_1(\varphi)$와 같게 놓자.
$z' \in \Ob(\mathcal{S}_2)$이고 $W = p_2(z')$라 하며,
$p_2(\psi') = f \circ g$를 만족하는 $g : W \to V$와
$\psi' : z' \to F(x)$가 주어졌다고 하자. 그러면
$$
\psi = j \circ G(\psi') : G(z') \to G(F(x)) \to x
$$
는 $p_1(\psi) = f \circ g$를 만족하는 $\mathcal{S}_1$의 사상이다.
따라서 가정에 의해 $g$ 위에 놓이고
$\psi = \varphi \circ \xi$를 만족하는 유일한 사상
$\xi : G(z') \to y$가 존재한다. 이는 다시 $g$ 위에 놓이고
$\psi' = F(\varphi) \circ \xi'$를 만족하는 사상
$$
\xi' = F(\xi) \circ i^{-1} : z' \to F(G(z')) \to F(y)
$$
을 준다. $\xi'$의 유일성에 대한 확인은 생략한다.
\end{proof}

\noindent
Lemma \ref{categories-lemma-fibred-equivalent}의 결론은 동치가 강한 카르테시안
사상을 강한 카르테시안 사상으로 보낸다는 것이다. 그러나
$\mathcal{C}$ 위의 섬유화된 범주 사이의 임의의 함자에 대해서는
그렇지 않을 수 있다. 따라서 섬유화된 범주들의 $2$-범주를 다음과
같이 정의한다.

\begin{definition}
\label{categories-definition-fibred-categories-over-C}
$\mathcal{C}$를 범주라 하자.
{\it $\mathcal{C}$ 위의 섬유화된 범주들의 $2$-범주}는
$\mathcal{C}$ 위의 범주들의 $2$-범주의 부분 $2$-범주로
(Definition \ref{categories-definition-categories-over-C}를 보라), 다음과 같이
정의된다.
\begin{enumerate}
\item 그 대상은 섬유화된 범주 $p : \mathcal{S} \to \mathcal{C}$이다.
\item 그 $1$-사상 $(\mathcal{S}, p) \to (\mathcal{S}', p')$은
$p' \circ G = p$를 만족하고 $G$가 강한 카르테시안 사상을 강한
카르테시안 사상으로 보내는 함자 $G : \mathcal{S} \to \mathcal{S}'$이다.
\item
$G, H : (\mathcal{S}, p) \to (\mathcal{S}', p')$에 대한 그 $2$-사상
$t : G \to H$는 모든 $x \in \Ob(\mathcal{S})$에 대하여
$p'(t_x) = \text{id}_{p(x)}$를 만족하는 함자의 사상이다.
\end{enumerate}
이 상황에서 $(\mathcal{S}, p)$와 $(\mathcal{S}', p')$ 사이의
$1$-사상들의 범주를
$$
\Mor_{\textit{Fib}/\mathcal{C}}(\mathcal{S}, \mathcal{S}')
$$
로 나타내겠다.
\end{definition}

\noindent
$1$-사상에 부과된 조건에 유의하라. 또한 이것은 진정한 $2$-범주이지
$(2, 1)$-범주가 아니라는 점에도 유의하라. 따라서 $2$-섬유곱을 취할
때에는 먼저 이에 수반되는 $(2, 1)$-범주로 옮겨 간다.

\begin{lemma}
\label{categories-lemma-2-product-fibred-categories-over-C}
$\mathcal{C}$를 범주라 하자. $\mathcal{C}$ 위의 섬유화된 범주들의
$(2, 1)$-범주는 2-섬유곱을 가지며, 이는
Lemma \ref{categories-lemma-2-product-categories-over-C}에서와 같이 기술된다.
\end{lemma}

\begin{proof}
본질적으로 보여야 할 것은 다음이다.
$F : \mathcal{X} \to \mathcal{S}$와
$G : \mathcal{Y} \to \mathcal{S}$가 $\mathcal{C}$ 위의 섬유화된
범주들의 사상일 때, Lemma \ref{categories-lemma-2-product-categories-over-C}에서
기술한 범주 $\mathcal{X} \times_\mathcal{S} \mathcal{Y}$가
섬유화되어 있음을 보여야 한다.
$\mathcal{X} \times_\mathcal{S} \mathcal{Y}$가 충분히 많은 강한
카르테시안 사상을 가짐을 보이자. 대상
$(U, x, y, \phi)$가 $\mathcal{X} \times_\mathcal{S} \mathcal{Y}$에
속하고, $f : V \to U$가 $\mathcal{C}$의 사상이라고 하자.
$\mathcal{X}$에서 $f$ 위에 놓인 강한 카르테시안 사상
$a : f^*x \to x$와 $\mathcal{Y}$에서 $f$ 위에 놓인 강한
카르테시안 사상 $b : f^*y \to y$를 선택하자. 가정에 의해
$F(a)$와 $G(b)$는 강한 카르테시안 사상이다.
$\phi : F(x) \to G(y)$가 동형사상이므로 강한 카르테시안 사상의
유일성에 의해
$G(b) \circ f^*\phi = \phi \circ F(a)$를 만족하는 유일한 동형사상
$f^*\phi : F(f^*x) \to G(f^*y)$를 얻는다. 다시 말해
$(a, b) : (V, f^*x, f^*y, f^*\phi) \to (U, x, y, \phi)$는
$\mathcal{X} \times_\mathcal{S} \mathcal{Y}$의 사상이다. 이것이 강한
카르테시안 사상이라는 확인(그리고 이것들이 실제로 모든 강한
카르테시안 사상이라는 확인)은 생략한다.
\end{proof}

\begin{lemma}
\label{categories-lemma-cute}
$\mathcal{C}$를 범주라 하고 $U \in \Ob(\mathcal{C})$라 하자.
$p : \mathcal{S} \to \mathcal{C}$가 섬유화된 범주이고 $p$가
$p' : \mathcal{S} \to \mathcal{C}/U$를 통해 분해되면,
$p' : \mathcal{S} \to \mathcal{C}/U$는 섬유화된 범주이다.
\end{lemma}

\begin{proof}
$\varphi : x' \to x$가 $p$에 관하여 강한 카르테시안 사상이라고 하자.
$\varphi$가 $p'$에 관해서도 강한 카르테시안 사상이라고 주장한다.
$g = p'(\varphi)$로 놓자. 그러면 어떤 사상 $f : V \to U$와
$f' : V' \to U$에 대하여 $g : V'/U \to V/U$이다.
$z \in \Ob(\mathcal{S})$라 하고 $p'(z) = (W \to U)$로 놓자.
$\varphi$가 $p'$에 관하여 강한 카르테시안 사상임을 보이려면
$$
\Mor_\mathcal{S}(z, x')
\longrightarrow
\Mor_\mathcal{S}(z, x)
\times_{\Mor_{\mathcal{C}/U}(W/U, V/U)}
\Mor_{\mathcal{C}/U}(W/U, V'/U),
$$
$\psi' \longmapsto (\varphi \circ \psi', p'(\psi'))$로 주어지는 사상이
전단사임을 보여야 한다. 우변의 원소 $(\psi, h)$가 주어졌다고 하자.
그러면 특히 $g \circ h = p(\psi)$이고, $\varphi$가 강한 카르테시안
사상이라는 조건에 의해 $\psi = \varphi \circ \psi'$와
$p(\psi') = h$를 만족하는 유일한 사상 $\psi' : z \to x'$를 얻는다.
좋다. 이제 $p'(\psi') : W/U \to V/U$는 이에 대응하는 사상
$W \to V$가 $h$인 사상이므로, $\mathcal{C}/U$의 사상으로서
$h$와 같다. 따라서 $\psi'$는 주어진 쌍 $(\psi, h)$로 보내지는
유일한 사상 $z \to x'$이다. 이로써 주장이 증명되었다.

\medskip\noindent
끝으로 $g : V'/U \to V/U$와 $p'(x) = V/U$인 $x$가 주어졌다고 하자.
$p : \mathcal{S} \to \mathcal{C}$가 섬유화된 범주이므로
$p(\varphi) = g$인 강한 카르테시안 사상 $\varphi : x' \to x$가
존재한다. 위와 같은 논증으로부터
$p'(\varphi) = g : V'/U \to V/U$가 따른다. 또한 위에서 보았듯이
$\varphi$는 강한 카르테시안 사상이다. 따라서
Definition \ref{categories-definition-fibred-category}의 조건들이 성립하며
증명이 끝난다.
\end{proof}

\begin{lemma}
\label{categories-lemma-fibred-over-fibred}
$\mathcal{A} \to \mathcal{B} \to \mathcal{C}$를 범주 사이의 함자들이라
하자. $\mathcal{A}$가 $\mathcal{B}$ 위에서 섬유화되어 있고
$\mathcal{B}$가 $\mathcal{C}$ 위에서 섬유화되어 있으면,
$\mathcal{A}$는 $\mathcal{C}$ 위에서 섬유화되어 있다.
\end{lemma}

\begin{proof}
정의들과 Lemma \ref{categories-lemma-cartesian-over-cartesian}에서 따른다.
\end{proof}

\begin{lemma}
\label{categories-lemma-fibred-category-representable-goes-up}
$p : \mathcal{S} \to \mathcal{C}$를 섬유화된 범주라 하자.
$x \to y$와 $z \to y$를 $\mathcal{S}$의 사상이라 하고,
$x \to y$가 강한 카르테시안 사상이라고 하자.
$p(x) \times_{p(y)} p(z)$가 존재하면 $x \times_y z$가 존재하고,
$p(x \times_y z) = p(x) \times_{p(y)} p(z)$이며,
$x \times_y z \to z$는 강한 카르테시안 사상이다.
\end{lemma}

\begin{proof}
$\text{pr}_2 : p(x) \times_{p(y)} p(z) \to p(z)$ 위에 놓인 강한
카르테시안 사상 $\text{pr}_2^*z \to z$를 선택하자. 그러면
Lemma \ref{categories-lemma-strongly-cartesian-fibre-product}에 의해
$\text{pr}_2^*z = x \times_y z$이다.
\end{proof}

\begin{lemma}
\label{categories-lemma-ameliorate-morphism-fibred-categories}
$\mathcal{C}$를 범주라 하고, $F : \mathcal{X} \to \mathcal{Y}$를
$\mathcal{C}$ 위의 섬유화된 범주들의 $1$-사상이라 하자.
다음을 만족하는 $\mathcal{C}$ 위의 섬유화된 범주들의 $1$-사상들이
존재한다.
$$
\xymatrix{
\mathcal{X} \ar@<1ex>[r]^u &
\mathcal{X}' \ar[r]^v \ar@<1ex>[l]^w & \mathcal{Y}
}
$$
$F = v \circ u$이고,
\begin{enumerate}
\item $u : \mathcal{X} \to \mathcal{X}'$는 충실충만하며,
\item $w$는 $u$의 왼쪽 수반이고,
\item $v : \mathcal{X}' \to \mathcal{Y}$는 섬유화된 범주이다.
\end{enumerate}
\end{lemma}

\begin{proof}
구조 함자들을 $p : \mathcal{X} \to \mathcal{C}$와
$q : \mathcal{Y} \to \mathcal{C}$로 나타내자. $\mathcal{X}'$를 다음과
같이 명시적으로 구성한다. $\mathcal{X}'$의 대상은 사중항
$(U, x, y, f)$이다. 여기서 $x \in \Ob(\mathcal{X}_U)$,
$y \in \Ob(\mathcal{Y}_U)$이고 $f : y \to F(x)$는
$\mathcal{Y}_U$의 사상이다. 사상
$(a, b) : (U, x, y, f) \to (U', x', y', f')$은
$p(a) = q(b) : U \to U'$ 및 $f' \circ b = F(a) \circ f$를 만족하는
$a : x \to x'$와 $b : y \to y'$로 주어진다.

\medskip\noindent
$p$와 $q$ 각각에 대해 당김을 선택하고 둘 모두에 같은 표기를 쓰자.
$(U, x, y, f)$를 대상, $h : V \to U$를 사상이라 하자.
주어진 강한 카르테시안 사상 $h^*x \to x$와 $h^*y \to y$에서 오는
사상 $c : (V, h^*x, h^*y, h^*f) \to (U, x, y, f)$를 생각하자.
$c$가 $\mathcal{C}$ 위의 $\mathcal{X}'$에서 강한 카르테시안
사상이라고 주장한다. 실제로 $\mathcal{X}'$의 대상
$(W, x', y', f')$, $W \to U$ 위에 놓인 사상
$(a, b) : (W, x', y', f') \to (U, x, y, f)$, 그리고 $h$를 통한
$W \to U$의 분해 $W \to V \to U$가 주어졌다고 하자.
$h^*x \to x$와 $h^*y \to y$가 강한 카르테시안 사상이므로 주어진
사상 $W \to V$ 위에 놓인 사상 $a' : x' \to h^*x$와
$b' : y' \to h^*y$를 얻는다. 다음 도표를 생각하자.
$$
\xymatrix{
y' \ar[d]_{f'} \ar[r] & h^*y \ar[r] \ar[d]_{h^*f} & y \ar[d]_f \\
F(x') \ar[r] & F(h^*x) \ar[r] & F(x)
}
$$
바깥 직사각형과 오른쪽 정사각형은 가환한다. $F$가 섬유화된
범주들의 $1$-사상이므로 사상 $F(h^*x) \to F(x)$는 강한 카르테시안
사상이다. 따라서 강한 카르테시안 사상의 보편 성질에 의해 왼쪽
정사각형도 가환한다. 이로써 $\mathcal{X}'$가 $\mathcal{C}$ 위에서
섬유화되어 있음이 증명된다.

\medskip\noindent
함자 $u : \mathcal{X} \to \mathcal{X}'$는
$x \mapsto (p(x), x, F(x), \text{id})$로 주어지며 충실충만하다.
함자 $\mathcal{X}' \to \mathcal{Y}$는
$(U, x, y, f) \mapsto y$로 주어진다. 함자
$w : \mathcal{X}' \to \mathcal{X}$는 $(U, x, y, f) \mapsto x$로
주어진다. $\mathcal{C}$ 위의 $\mathcal{X}'$의 강한 카르테시안
사상에 관한 위 기술에 의해 이 함자들은 각각 $\mathcal{C}$ 위의
섬유화된 범주들의 $1$-사상이다. $w$와 $u$가 수반이라는 것은
$$
\Mor_\mathcal{X}(x, x') =
\Mor_{\mathcal{X}'}((U, x, y, f), (p(x'), x', F(x'), \text{id})),
$$
라는 뜻이며, 이는 정의에서 바로 따른다.

\medskip\noindent
끝으로 $\mathcal{X}' \to \mathcal{Y}$가 섬유화된 범주임을 보여야 한다.
$c : y' \to y$를 $\mathcal{Y}$의 사상이라 하고
$(U, x, y, f)$를 $y$ 위에 놓인 $\mathcal{X}'$의 대상이라 하자.
$V = q(y')$ 및 $h = q(c) : V \to U$로 놓자. $a : h^*x \to x$와
$b : h^*y \to y$를 $h$를 덮는 강한 카르테시안 사상들이라 하자.
$F$가 섬유화된 범주들의 $1$-사상이므로 강한 카르테시안 사상
$F(a) : F(h^*x) \to F(x)$와 함께 $h^*F(x) = F(h^*x)$로 식별할 수 있다.
$b : h^*y \to y$의 보편 성질에 의해 $c = b \circ c'$를 만족하는
$\mathcal{Y}_V$의 사상 $c' : y' \to h^*y$가 존재한다. 다음 사상이
$$
(a, c) : (V, h^*x, y', h^*f \circ c') \longrightarrow (U, x, y, f)
$$
$\mathcal{Y}$ 위의 $\mathcal{X}'$에서 강한 카르테시안 사상이라고
주장한다. 이를 보기 위해 $(W, x_1, y_1, f_1)$을 $\mathcal{X}'$의
대상이라 하고, $(a_1, b_1) : (W, x_1, y_1, f_1) \to (U, x, y, f)$를
사상이라 하며, 어떤 사상 $b_1' : y_1 \to y'$에 대하여
$b_1 = c \circ b_1'$라고 하자. 그러면
$$
(a_1', b_1') : (W, x_1, y_1, f_1) \longrightarrow (V, h^*x, y', h^*f \circ c')
$$
가 원하는 사상이다. 여기서 $a_1' : x_1 \to h^*x$는 주어진 사상
$q(b_1') : W \to V$ 위에 놓이고 $a_1 = a \circ a_1'$를 만족하는
유일한 사상이다.
\end{proof}





\section{관성}
\label{categories-section-inertia}

\noindent
범주 $\mathcal{C}$ 위의 섬유화된 범주
$p : \mathcal{S} \to \mathcal{C}$와
$p' : \mathcal{S}' \to \mathcal{C}$, 그리고 $1$-사상
$F : \mathcal{S} \to \mathcal{S}'$가 주어지면, $\mathcal{C}$ 위의
섬유화된 범주들의 $(2, 1)$-범주에서 대각 사상
$$
\Delta = \Delta_{\mathcal{S}/\mathcal{S}'} :
\mathcal{S} \longrightarrow \mathcal{S} \times_{\mathcal{S}'} \mathcal{S}
$$
을 얻는다.

\begin{lemma}
\label{categories-lemma-inertia-fibred-category}
$\mathcal{C}$를 범주라 하자.
$p : \mathcal{S} \to \mathcal{C}$와
$p' : \mathcal{S}' \to \mathcal{C}$를 섬유화된 범주들이라 하자.
$F : \mathcal{S} \to \mathcal{S}'$를 $\mathcal{C}$ 위의 섬유화된
범주들의 $1$-사상이라 하자. 다음과 같은 $\mathcal{C}$ 위의 범주
$\mathcal{I}_{\mathcal{S}/\mathcal{S}'}$를 생각하자.
\begin{enumerate}
\item 그 대상은 $x \in \Ob(\mathcal{S})$이고
$F(\alpha) = \text{id}$인 자기동형사상 $\alpha : x \to x$인 쌍
$(x, \alpha)$이다.
\item 사상 $(x, \alpha) \to (y, \beta)$은 다음 도표를 가환하게 하는
사상 $\phi : x \to y$로 주어진다.
$$
\xymatrix{
x\ar[r]_\phi\ar[d]_\alpha &
y\ar[d]^{\beta} \\
x\ar[r]^\phi &
y \\
}
$$
\item 함자 $\mathcal{I}_{\mathcal{S}/\mathcal{S}'} \to \mathcal{C}$는
$(x, \alpha) \mapsto p(x)$로 주어진다.
\end{enumerate}
그러면
\begin{enumerate}
\item $\mathcal{C}$ 위의 범주들의 $(2, 1)$-범주에서 동치
$$
\mathcal{I}_{\mathcal{S}/\mathcal{S}'} \longrightarrow
\mathcal{S}
\times_{\Delta, (\mathcal{S} \times_{\mathcal{S}'} \mathcal{S}), \Delta}
\mathcal{S}
$$
가 존재하고,
\item $\mathcal{I}_{\mathcal{S}/\mathcal{S}'}$는 $\mathcal{C}$ 위의
섬유화된 범주이다.
\end{enumerate}
\end{lemma}

\begin{proof}
(2)는 (1)과 Lemmas \ref{categories-lemma-2-product-fibred-categories-over-C} 및
\ref{categories-lemma-fibred-equivalent}에서 따른다. 따라서 (1)만 증명하면 충분하다.
Lemma \ref{categories-lemma-2-product-fibred-categories-over-C}의 $2$-섬유곱 구성을
더 언급하지 않고 사용하겠다. 특히
$\mathcal{S}
\times_{\Delta, (\mathcal{S} \times_{\mathcal{S}'} \mathcal{S}), \Delta}
\mathcal{S}$의 대상은 사중항 $(x, y, (\iota, \kappa))$이다. 여기서
$x$와 $y$는 $\mathcal{S}$의 대상이고,
$(\iota, \kappa) : (x, x, \text{id}_{F(x)}) \to (y, y, \text{id}_{F(y)})$는
$\mathcal{S} \times_{\mathcal{S}'} \mathcal{S}$의 동형사상이다.
이는 곧 $\iota, \kappa : x \to y$가 동형사상이고
$F(\iota) = F(\kappa)$라는 뜻이다. 다음 함자를 생각하자.
$$
I_{\mathcal{S}/\mathcal{S}'}
\longrightarrow
\mathcal{S}
\times_{\Delta, (\mathcal{S} \times_{\mathcal{S}'} \mathcal{S}), \Delta}
\mathcal{S}
$$
이 함자는 좌변의 대상 $(x, \alpha)$에 우변의 대상
$(x, x, (\alpha, \text{id}_x))$를 대응시키고, 좌변의 사상 $\phi$에
우변의 사상 $(\phi, \phi)$를 대응시킨다. 그 사상의 준역은 다음
함자로 주어진다고 주장한다.
$$
\mathcal{S}
\times_{\Delta, (\mathcal{S} \times_{\mathcal{S}'} \mathcal{S}), \Delta}
\mathcal{S}
\longrightarrow
I_{\mathcal{S}/\mathcal{S}'}
$$
이 함자는 좌변의 대상 $(x, y, (\iota, \kappa))$에 우변의 대상
$(x, \kappa^{-1} \circ \iota)$를 대응시키고, 좌변의 사상
$(\phi, \phi') : (x, y, (\iota, \kappa)) \to (z, w, (\lambda, \mu))$에
사상 $\phi$를 대응시킨다. 실제로 위의 두 함자를 합성하여 얻는
$I_{\mathcal{S}/\mathcal{S}'}$의 자기함자는 엄밀히 항등함자이고,
$\mathcal{S}
\times_{\Delta, (\mathcal{S} \times_{\mathcal{S}'} \mathcal{S}), \Delta}
\mathcal{S}$에서 유도되는 자기함자는 자연 동형사상
$$
(\text{id}_x, \kappa) :
(x, x, (\kappa^{-1} \circ \iota, \text{id}_x))
\longrightarrow
(x, y, (\iota, \kappa)).
$$
을 통해 항등함자와 동형이다. 일부 세부 사항은 생략한다.
\end{proof}

\begin{definition}
\label{categories-definition-inertia-fibred-category}
$\mathcal{C}$를 범주라 하자.
\begin{enumerate}
\item $F : \mathcal{S} \to \mathcal{S}'$를 $\mathcal{C}$ 위의
섬유화된 범주들의 $1$-사상이라 하자. $\mathcal{S}'$ 위의
{\it $\mathcal{S}$의 상대 관성}은 Lemma
\ref{categories-lemma-inertia-fibred-category}의 섬유화된 범주
$\mathcal{I}_{\mathcal{S}/\mathcal{S}'} \to \mathcal{C}$이다.
\item {\it $\mathcal{S}$의 관성 섬유화 범주
$\mathcal{I}_\mathcal{S}$}란
$\mathcal{I}_\mathcal{S} = \mathcal{I}_{\mathcal{S}/\mathcal{C}}$를 뜻한다.
\end{enumerate}
\end{definition}

\noindent
$\mathcal{C}$ 위의 섬유화된 범주들의 표준 $1$-사상
\begin{equation}
\label{categories-equation-inertia-structure-map}
\mathcal{I}_{\mathcal{S}/\mathcal{S}'} \longrightarrow \mathcal{S}
\quad\text{and}\quad
\mathcal{I}_\mathcal{S} \longrightarrow \mathcal{S}
\end{equation}
이 있음에 유의하라. Lemma \ref{categories-lemma-inertia-fibred-category}의 기술로
보면 이들은 단순히 대상 $(x, \alpha)$를 대상 $x$로 보내고 사상
$\phi : (x, \alpha) \to (y, \beta)$를 사상 $\phi : x \to y$로 보낸다.
또한 규칙 $x \mapsto (x, \text{id}_x)$와
$(\phi : x \to y) \mapsto \phi$로 정의되는 {\it 중립 단면}
\begin{equation}
\label{categories-equation-neutral-section}
e : \mathcal{S} \to \mathcal{I}_{\mathcal{S}/\mathcal{S}'}
\quad\text{and}\quad
e : \mathcal{S} \to \mathcal{I}_\mathcal{S}
\end{equation}
이 있다. 이는 (\ref{categories-equation-inertia-structure-map})의 오른쪽 역이다.
$2$-가환 정사각형
$$
\xymatrix{
\mathcal{S}_1 \ar[d]_{F_1} \ar[r]_G & \mathcal{S}_2 \ar[d]^{F_2} \\
\mathcal{S}'_1 \ar[r]^{G'} & \mathcal{S}'_2
}
$$
이 주어지면 규칙 $(x, \alpha) \mapsto (G(x), G(\alpha))$와
$\phi \mapsto G(\phi)$로 정의되는 {\it 함자성 사상}
\begin{equation}
\label{categories-equation-functorial}
\mathcal{I}_{\mathcal{S}_1/\mathcal{S}'_1}
\longrightarrow
\mathcal{I}_{\mathcal{S}_2/\mathcal{S}'_2}
\quad\text{and}\quad
\mathcal{I}_{\mathcal{S}_1}
\longrightarrow
\mathcal{I}_{\mathcal{S}_2}
\end{equation}
이 있다. 특히 항상 비교 사상
\begin{equation}
\label{categories-equation-comparison}
\mathcal{I}_{\mathcal{S}/\mathcal{S}'}
\longrightarrow
\mathcal{I}_\mathcal{S}
\end{equation}
이 있고, 위의 모든 사상은 이 사상과 양립한다.

\begin{lemma}
\label{categories-lemma-relative-inertia-as-fibre-product}
$F : \mathcal{S} \to \mathcal{S}'$를 범주 $\mathcal{C}$ 위에서 섬유화된
범주들의 $1$-사상이라 하자. 그러면 도표
$$
\xymatrix{
\mathcal{I}_{\mathcal{S}/\mathcal{S}'}
\ar[d]_{F \circ (\ref{categories-equation-inertia-structure-map})}
\ar[rr]_{(\ref{categories-equation-comparison})} & &
\mathcal{I}_\mathcal{S} \ar[d]^{(\ref{categories-equation-functorial})} \\
\mathcal{S}' \ar[rr]^e & &
\mathcal{I}_{\mathcal{S}'}
}
$$
는 $2$-섬유곱이다.
\end{lemma}

\begin{proof}
생략한다.
\end{proof}





\section{군체로 섬유화된 범주들}
\label{categories-section-fibred-groupoids}

\noindent
이 절에서는 군체로 섬유화된 범주를 어떻게 이해하는지 설명하고,
그것이 본질적으로 군체들의 $(2, 1)$-범주에 값을 갖는 함자와 같음을
살펴본다.

\begin{definition}
\label{categories-definition-fibred-groupoids}
$p : \mathcal{S} \to \mathcal{C}$를 함자라 하자. 다음 두 조건이
성립하면 $\mathcal{S}$가 $\mathcal{C}$ 위에서
{\it 군체로 섬유화되어 있다}고 한다.
\begin{enumerate}
\item $\mathcal{C}$의 모든 사상 $f : V \to U$와 $U$의 모든 올림
$x$에 대하여, 끝점이 $x$인 $f$의 올림 $\phi : y \to x$가 존재한다.
\item 임의의 사상 쌍 $\phi : y \to x$, $ \psi : z \to x$와
$p(\phi) \circ f = p(\psi)$를 만족하는 임의의 사상
$f : p(z) \to p(y)$에 대하여, $\phi \circ \chi = \psi$를 만족하는
$f$의 유일한 올림 $\chi : z \to y$가 존재한다.
\end{enumerate}
\end{definition}

\noindent
조건 (2)를 달리 말하면, 함자 $p$를 적용할 때 아래의 가환 도표에 있는
점선 화살표들의 집합 사이에 전단사가 주어진다는 뜻이다.
$$
\xymatrix{
y \ar[r] & x & p(y) \ar[r] & p(x) \\
z \ar@{-->}[u] \ar[ru] & & p(z) \ar@{-->}[u]\ar[ru] & \\
}
$$
두 번째 조건을 이해하는 또 다른 방법은 다음과 같다.
$g : W \to V$와 $f : V \to U$가 $\mathcal{C}$의 사상들이라고 하자.
$x \in \Ob(\mathcal{S}_U)$라 하자. 첫 번째 조건에 의해 $f$를
$ \phi : y \to x$로 올리고, 이어서 $g$를 $\psi : z \to y$로 올릴 수
있다. 이 두 단계 과정을 거치는 대신 $g \circ f$를
$\gamma : z' \to x$로 직접 올릴 수 있다. 이로써 도표의 실선
화살표들을 얻는다.
\begin{equation}
\label{categories-equation-fibred-groupoids}
\vcenter{
\xymatrix{
z' \ar@{-->}[d]\ar[rrd]^\gamma & & \\
z \ar@{-->}[u] \ar[r]^\psi \ar@{~>}[d]^p &
y \ar[r]^\phi \ar@{~>}[d]^p &
x \ar@{~>}[d]^p
\\
W \ar[r]^g & V \ar[r]^f & U \\
}
}
\end{equation}
여기서 물결 화살표는 사상이 아니라 함자 $p$를 나타낸다.
두 번째 조건을 화살표 $\phi \circ \psi$, $\gamma$와
$\text{id}_W$에 적용하면 $\gamma \circ \chi = \phi \circ \psi$를
만족하는 $\mathcal{S}_W$의 유일한 사상 $\chi : z \to z'$가 존재한다.
마찬가지로 유일한 사상 $z' \to z$가 존재한다. 유일성에 의해 사상
$z' \to z$와 $z\to z'$는 서로 역이며, 다시 말해 동형사상이다.

\medskip\noindent
이 논의로부터 군체로 섬유화된 범주가 섬유화된 범주와 매우 밀접한
관련이 있음은 명백할 것이다. 정확한 결과는 다음과 같다.

\begin{lemma}
\label{categories-lemma-fibred-groupoids}
$p : \mathcal{S} \to \mathcal{C}$를 함자라 하자. 다음은 동치이다.
\begin{enumerate}
\item $p : \mathcal{S} \to \mathcal{C}$는 군체로 섬유화된 범주이고,
\item 모든 섬유 범주는 군체이며 $\mathcal{S}$는 $\mathcal{C}$ 위의
섬유화된 범주이다.
\end{enumerate}
더 나아가 이 경우 $\mathcal{S}$의 모든 사상은 강한 카르테시안
사상이다. 또한 모든 $f: V \to U = p(x)$에 대하여 $f$ 위에 놓인
$f^\ast x \to x$가 주어졌을 때, Lemma \ref{categories-lemma-fibred}에서 구성한
자료 $(U \mapsto \mathcal{S}_U, f \mapsto f^*, \alpha_{f, g}, \alpha_U)$는
$\mathcal{C}^{opp}$에서 군체들의 $(2, 1)$-범주로 가는 유사함자를
정의한다.
\end{lemma}

\begin{proof}
$p : \mathcal{S} \to \mathcal{C}$가 군체로 섬유화되어 있다고 하자.
모든 $U \in \Ob(\mathcal{C})$에 대한 섬유 범주 $\mathcal{S}_U$가
군체임을 보이려면, $\mathcal{S}_U$의 모든 $f : y \to x$에 대하여
역사상을 제시해야 한다. 왼쪽 도표($\mathcal{S}_U$ 안의 도표)는
$p$에 의해 오른쪽 도표로 보내진다.
$$
\xymatrix{
y \ar[r]^f & x & U \ar[r]^{\text{id}_U} & U \\
x \ar@{-->}[u] \ar[ru]_{\text{id}_x} & &
U \ar@{-->}[u]\ar[ru]_{\text{id}_U} & \\
}
$$
오직 $\text{id}_U$만 오른쪽 도표를 가환하게 하므로, 왼쪽 도표를
가환하게 하는 유일한 $g : x \to y$가 존재하고 따라서
$fg = \text{id}_x$이다. 비슷한 논증으로 $gh = \text{id}_y$인
유일한 $h : y \to x$가 존재한다. 그러면
$fgh = f : y \to x$이다. $fg = \text{id}_x$이므로 $h = f$이다.
Definition \ref{categories-definition-fibred-groupoids}의 조건 (2)는 정확히
$\mathcal{S}$의 모든 사상이 강한 카르테시안 사상이라는 뜻이다.
따라서 Definition \ref{categories-definition-fibred-groupoids}의 조건 (1)에 의해
$\mathcal{S}$는 $\mathcal{C}$ 위의 섬유화된 범주이다.

\medskip\noindent
반대로 모든 섬유 범주가 군체이고 $\mathcal{S}$가 $\mathcal{C}$ 위의
섬유화된 범주라고 하자. Definition \ref{categories-definition-fibred-groupoids}의
조건 (1)과 (2)를 확인해야 한다. 첫 번째 조건은 자명하게 따른다.
Definition \ref{categories-definition-fibred-groupoids}의 조건 (2)에서처럼
$\phi : y \to x$, $\psi : z \to x$ 및
$p(\phi) \circ f = p(\psi)$를 만족하는 $f : p(z) \to p(y)$가
주어졌다고 하자. $U = p(x)$, $V = p(y)$, $W = p(z)$,
$p(\phi) = g : V \to U$, $p(\psi) = h : W \to U$로 쓰자.
$g$ 위에 놓인 강한 카르테시안 사상 $g^*x \to x$를 선택하자.
그러면 $\mathcal{S}_V$의 사상 $i : y \to g^*x$를 얻으며, 따라서 이는
동형사상이다. 또한 $g^*x \to x$가 강한 카르테시안 사상이므로 쌍
$(\psi, f)$에 대응하는 사상 $j : z \to g^*x$를 얻는다.
이제 $\chi = i^{-1} \circ j$가 해임을 확인할 수 있다.

\medskip\noindent
(1) $\Rightarrow$ (2)의 증명에서 $\mathcal{S}$의 모든 사상이 강한
카르테시안 사상임을 보았다. 마지막 명제는 Lemma \ref{categories-lemma-fibred}에서
바로 따른다.
\end{proof}

\begin{lemma}
\label{categories-lemma-fibred-gives-fibred-groupoids}
$\mathcal{C}$를 범주라 하고,
$p : \mathcal{S} \to \mathcal{C}$를 섬유화된 범주라 하자.
$\mathcal{S}'$를 다음과 같이 정의되는 $\mathcal{S}$의 부분범주라 하자.
\begin{enumerate}
\item $\Ob(\mathcal{S}') = \Ob(\mathcal{S})$이고,
\item $x, y \in \Ob(\mathcal{S}')$에 대하여 $\mathcal{S}'$에서
$x$와 $y$ 사이의 사상들의 집합은 $\mathcal{S}$에서 $x$와 $y$ 사이의
강한 카르테시안 사상들의 집합이다.
\end{enumerate}
$p' : \mathcal{S}' \to \mathcal{C}$를 $p$의 $\mathcal{S}'$에 대한
제한이라 하자. 그러면 $p' : \mathcal{S}' \to \mathcal{C}$는 군체로
섬유화되어 있다.
\end{lemma}

\begin{proof}
Lemma \ref{categories-lemma-composition-cartesian}에 의해 $\mathcal{S}$의 모든
대상의 항등사상이 강한 카르테시안 사상이고, 강한 카르테시안 사상들의
합성도 강한 카르테시안 사상이므로 이 구성은 잘 정의된다.
Definition \ref{categories-definition-fibred-groupoids}의 첫 번째 올림 성질은
섬유화된 범주에서 임의의 사상 $f : V \to U$와 $U$ 위에 놓인 $x$가
주어지면 $f$ 위에 놓인 강한 카르테시안 사상
$\varphi : y \to x$가 존재한다는 조건에서 따른다.
$\mathcal{C}$ 위의 범주 $p' : \mathcal{S}' \to \mathcal{C}$에 대해
Definition \ref{categories-definition-fibred-groupoids}의 두 번째 올림 성질을
확인하자. 이를 위해 Definition \ref{categories-definition-fibred-groupoids} 뒤의
논의와 같이 논증한다. 따라서 Diagram \ref{categories-equation-fibred-groupoids}의
사상 $\phi$, $\psi$, $\gamma$는 $\mathcal{S}$의 강한 카르테시안
사상이다. 그러므로 $\gamma$와 $\phi \circ \psi$는 $\mathcal{S}$의
강한 카르테시안 사상들이고 $\mathcal{C}$의
같은 화살표 위에 놓이고 $\mathcal{S}$에서 끝점이 같은 강한
카르테시안 사상들이다. Definition \ref{categories-definition-cartesian-over-C}
뒤의 논의에 의해 이 두 화살표는 원하는 대로 동형이다
(여기서 Lemma \ref{categories-lemma-composition-cartesian}에 의해 $\mathcal{S}$의
모든 동형사상이 강한 카르테시안 사상이라는 것도 사용한다).
\end{proof}

\begin{example}
\label{categories-example-group-homomorphism-fibreedingroupoids}
군 준동형 $p : G \to H$는 Example
\ref{categories-example-group-homomorphism-functor}에서와 같은 함자
$p : \mathcal{S}\to\mathcal{C}$를 낳는다. 이 함자
$p : \mathcal{S}\to\mathcal{C}$가 군체로 섬유화되어 있을
필요충분조건은 $p$가 전사인 것이다. (유일한) 대상
$U\in \Ob(\mathcal{C})$ 위의 섬유 범주 $\mathcal{S}_U$는 Example
\ref{categories-example-group-groupoid}에서처럼 $p$의 핵에 대응하는 범주이다.
\end{example}

\noindent
$p : \mathcal{S} \to \mathcal{C}$가 주어졌다고 하자. 모든
$U \in \Ob(\mathcal{C})$에 대하여 섬유 범주 $\mathcal{S}_U$가
군체이면 $\mathcal{S}$는 반드시 $\mathcal{C}$ 위에서 군체로
섬유화되어 있는가? 다음과 같이 답이 부정적임을 알 수 있다.
군체로 섬유화된 범주 $p : \mathcal{S} \to \mathcal{C}$에서 시작하자.
$\mathcal{S}$에서 어떤 대상의 항등사상으로 보내지지 않는 사상들을
바꾸어도 범주 $\mathcal{S}_U$는 바뀌지 않는다. 따라서 올림의 존재성과
유일성 조건을 깨뜨릴 수 있다. 한 예는 $G \to H$가 전사가 아닐 때의
Example \ref{categories-example-group-homomorphism-fibreedingroupoids}의 함자이다.
다른 예를 하나 들자.

\begin{example}
\label{categories-example-not-fibred-in-groupoids-but-fibre-cats-are}
$\Ob(\mathcal{C}) = \{A, B, T\}$,
$\Mor_\mathcal{C}(A, B) = \{f\}$, $\Mor_\mathcal{C}(B, T) = \{g\}$,
$\Mor_\mathcal{C}(A, T) = \{h\} = \{gf\}, $이고 각 대상의 항등사상도
있다고 하자. 이 범주의 그림은 아래 도표를 보라. 이제
$\Ob(\mathcal{S}) = \{A', B', T'\}$,
$\Mor_\mathcal{S}(A', B') = \emptyset$,
$\Mor_\mathcal{S}(B', T') = \{g'\}$,
$\Mor_\mathcal{S}(A', T') = \{h'\}, $이고 항등사상들도 있다고 하자.
함자 $p : \mathcal{S} \to \mathcal{C}$는 자명하다. 그러면 모든
$U \in \Ob(\mathcal{C})$에 대하여 $\mathcal{S}_U$는 대상 하나와 그
대상의 항등사상으로 이루어진 범주이므로 군체이지만, 사상
$f: A \to B$는 올릴 수 없다. 마찬가지로
$\Mor_\mathcal{S}(A', B') = \{f'_1, f'_2\}$이고
$ \Mor_\mathcal{S}(A', T') = \{h'\} = \{g'f'_1 \} = \{g'f'_2\}$라고
정하면 섬유 범주들은 그대로이고 아래 도표의 $f: A \to B$는 두 개의
올림을 갖는다.
$$
\xymatrix{
B' \ar[r]^{g'} & T' & & B \ar[r]^g & T & \\
A' \ar@{-->}[u]^{??} \ar[ru]_{h'} & & \ar@{}[u]^{above} &
A \ar[u]^f \ar[ru]_{gf = h} & \\
}
$$
\end{example}

\noindent
나중에는 ``$\mathcal{C}$ 위에서 군체로 섬유화된 모든 범주는 분할된
범주와 동치이다'' 또는 ``섬유 범주들이 집합과 같은, 군체로 섬유화된
모든 범주는 집합들로 섬유화된 범주와 동치이다''와 같은 명제를 말하고
싶을 것이다. 동치의 개념은 우리가 다루는 $2$-범주에 의존한다.

\begin{definition}
\label{categories-definition-categories-fibred-in-groupoids-over-C}
$\mathcal{C}$를 범주라 하자.
{\it $\mathcal{C}$ 위에서 군체로 섬유화된 범주들의 $2$-범주}는
$\mathcal{C}$ 위의 섬유화된 범주들의 $2$-범주의 부분 $2$-범주이며
(Definition \ref{categories-definition-fibred-categories-over-C}를 보라), 다음과
같이 정의된다.
\begin{enumerate}
\item 그 대상은 군체로 섬유화된 범주
$p : \mathcal{S} \to \mathcal{C}$이다.
\item 그 $1$-사상 $(\mathcal{S}, p) \to (\mathcal{S}', p')$은
$p' \circ G = p$를 만족하는 함자 $G : \mathcal{S} \to \mathcal{S}'$이다
(모든 사상이 강한 카르테시안 사상이므로 $G$는 자동으로 이를 보존한다).
\item
$G, H : (\mathcal{S}, p) \to (\mathcal{S}', p')$에 대한 그 $2$-사상
$t : G \to H$는 모든 $x \in \Ob(\mathcal{S})$에 대하여
$p'(t_x) = \text{id}_{p(x)}$를 만족하는 함자의 사상이다.
\end{enumerate}
\end{definition}

\noindent
모든 $2$-사상은 자동으로 동형사상임에 유의하라! 따라서 이것은 단순한
$2$-범주가 아니라 실제로 $(2, 1)$-범주이다. 다음은 $2$-섬유곱에 관한
필수적인 보조정리이다.

\begin{lemma}
\label{categories-lemma-2-product-fibred-categories}
$\mathcal{C}$를 범주라 하자. $\mathcal{C}$ 위에서 군체로 섬유화된
범주들의 $2$-범주는 2-섬유곱을 가지며, 이는
Lemma \ref{categories-lemma-2-product-categories-over-C}에서와 같이 기술된다.
\end{lemma}

\begin{proof}
Lemma \ref{categories-lemma-2-product-fibred-categories-over-C}에 의해
Lemma \ref{categories-lemma-2-product-categories-over-C}에서 기술한 섬유곱은
섬유화된 범주이다. 따라서 Lemma \ref{categories-lemma-fibred-groupoids}를
참조하면 그 섬유 범주들이 군체임을 증명하는 것으로 충분하다.
Lemma \ref{categories-lemma-fibre-2-fibre-product-categories-over-C}에 의해
군체들의 $2$-섬유곱이 군체임을 보이면 충분한데, 이는 명백하다
(예를 들어 Lemma \ref{categories-lemma-2-fibre-product-categories}의 구성을 보라).
\end{proof}

\begin{remark}
\label{categories-remark-2-product-fibred-categories-same}
$\mathcal{C}$를 범주라 하자. $f : \mathcal{X} \to \mathcal{S}$와
$g : \mathcal{Y} \to \mathcal{S}$를 $\mathcal{C}$ 위에서 군체로
섬유화된 범주들의 $1$-사상이라 하자.
$p : \mathcal{S} \to \mathcal{C}$를 주어진 함자라 하자.
Lemma \ref{categories-lemma-2-product-fibred-categories}의 $2$-섬유곱은
Example \ref{categories-example-2-fibre-product-categories}의 것과 (범주로서)
표준적으로 동치라고 주장한다. 앞의 범주의 대상은
$p(\alpha) = \text{id}_U$인 사중항 $(U, x, y, \alpha)$이고
(Lemma \ref{categories-lemma-2-product-categories-over-C}를 보라), 뒤의 범주의
대상은 사중항 $(x, y, \alpha)$이다
(Example \ref{categories-example-2-fibre-product-categories}를 보라).
두 범주 사이의 동치는 규칙
$(U, x, y, \alpha) \mapsto (x, y, \alpha)$와
$(x, y, \alpha) \mapsto (p(f(x)), x, y', \alpha')$로 주어진다.
여기서 $\alpha' = g(\gamma)^{-1} \circ \alpha$이고,
$\gamma : y' \to y$는 화살표
$p(\alpha) : p(f(x)) \to p(g(y))$의 올림이다. 세부 사항은 생략한다.
\end{remark}

\begin{lemma}
\label{categories-lemma-equivalence-fibred-categories}
$p : \mathcal{S}\to \mathcal{C}$와
$p' : \mathcal{S'}\to \mathcal{C}$를 군체로 섬유화된 범주들이라 하고,
$G : \mathcal{S}\to \mathcal {S}'$가 $\mathcal{C}$ 위의 함자라고 하자.
\begin{enumerate}
\item $G$가 충실한 것(각각 충실충만한 것, 동치인 것)의 필요충분조건은
각 $U\in\Ob(\mathcal{C})$에 대하여 유도된 함자
$G_U : \mathcal{S}_U\to \mathcal{S}'_U$가 충실한 것
(각각 충실충만한 것, 동치인 것)이다.
\item $G$가 동치이면 $G$는 $\mathcal{C}$ 위에서 군체로 섬유화된
범주들의 $2$-범주에서의 동치이다.
\end{enumerate}
\end{lemma}

\begin{proof}
$x, y$를 같은 대상 $U$ 위에 놓인 $\mathcal{S}$의 대상들이라 하자.
다음 가환 도표를 생각하자.
$$
\xymatrix{
\Mor_\mathcal{S}(x, y) \ar[rd]_p \ar[rr]_G & &
\Mor_{\mathcal{S}'}(G(x), G(y)) \ar[ld]^{p'} \\
& \Mor_\mathcal{C}(U, U) &
}
$$
이 도표로부터 $G$가 충실하면(각각 충실충만하면) 각 $G_U$도
그러하다는 것이 명백하다.

\medskip\noindent
$G$가 동치라고 하자. $\mathcal{S}'$의 모든 대상 $x'$에 대하여
$G(x)$가 $x'$와 동형인 $\mathcal{S}$의 대상 $x$가 존재한다.
$x'$가 $U'$ 위에 놓이고 $x$가 $U$ 위에 놓인다고 하자. 그러면
$\mathcal{C}$에 동형사상 $f : U' \to U$가 존재한다. 구체적으로 이는
동형사상 $x' \to G(x)$에 $p'$를 적용한 것이다. 군체로 섬유화된
범주의 공리에 의해 $f$ 위에 놓인 $\mathcal{S}$의 화살표
$f^*x \to x$가 존재한다. 따라서 $p'(\alpha) = \text{id}_{U'}$인
동형사상 $\alpha : x' \to G(f^*x)$가 존재한다
(이번에는 $\mathcal{S}'$에 대한 공리에 의한다). 결국
$\mathcal{S}'$의 모든 대상 $x'$에 대하여 $\mathcal{S}$의 대상
$o_{x'}$와 $p'(\alpha_{x'}) = \text{id}_{p'(x')}$인 동형사상
$\alpha_{x'} : x' \to G(o_{x'})$로 이루어진 쌍
$(o_{x'}, \alpha_{x'})$를 선택할 수 있다. 이 지점부터는 통상적인
방식으로 진행하여(Lemma \ref{categories-lemma-equivalence-categories}의 증명을
보라) 역함자 $F : \mathcal{S}' \to \mathcal{S}$를 구성한다.
대상에서는 $x' \mapsto o_{x'}$로 보내고, 사상
$\varphi' : x' \to y'$는
$\alpha_{y'}^{-1} \circ G(\varphi_{\varphi'}) \circ \alpha_{x'} = \varphi'$
를 만족하는 유일한 화살표
$\varphi_{\varphi'} : o_{x'} \to o_{y'}$로 보낸다. 이 선택 아래에서
$F$는 $\mathcal{C}$ 위의 함자이다. $G \circ F$와 $F \circ G$가 각각의
항등함자와 $2$-동형이라는 확인은 생략한다
($\mathcal{C}$ 위에서 군체로 섬유화된 범주들의 $2$-범주에서).

\medskip\noindent
모든 $U\in\Ob(\mathcal C)$에 대하여 $G_U$가 충실하다고
(각각 충실충만하다고) 하자. $G$가 충실함(각각 충실충만함)을 보이려면
임의의 대상 $x, y\in\Ob(\mathcal{S})$에 대하여 $G$가
$\Mor_\mathcal{S}(x, y)$와
$\Mor_{\mathcal{S}'}(G(x), G(y))$ 사이에 단사(각각 전단사)를
유도함을 보여야 한다. $U = p(x)$와 $V = p(y)$로 놓자. 임의의 사상
$f : U \to V$에 대하여 $G$가 $f$ 위에 놓인 사상 $x \to y$에서
$f$ 위에 놓인 사상 $G(x) \to G(y)$로 가는 단사(각각 전단사)를
유도함을 보이는 것으로 충분하다. 이제 $f : U \to V$를 고정하자.
$f^*y \to y$로 당김을 나타내자. 그러면
$G(f^*y) \to G(y)$도 당김이다. $f$ 위에 놓인 $x$에서 $y$로 가는
사상들의 집합은 $\text{id}_U$ 위에 놓인 $x$와 $f^*y$ 사이의 사상들의
집합과 전단사이다(군체로 섬유화된 범주의 두 번째 공리에 의한다).
마찬가지로 $f$ 위에 놓인 $G(x)$에서 $G(y)$로 가는 사상들의 집합은
$\text{id}_U$ 위에 놓인 $G(x)$와 $G(f^*y)$ 사이의 사상들의 집합과
전단사이다. 따라서 $G_U$가 충실하다는(각각 충실충만하다는) 사실로부터
원하는 결과가 따른다.

\medskip\noindent
끝으로 모든 $U$에 대하여 $G_U$가 동치라고 하자. 그러면 이는
충실충만하고 본질적으로 전사이다. 앞에서 이로부터 $G$가 충실충만함을
보았으므로, 이 함자가 동치임을 보이려면 본질적으로 전사임을 보이면 된다.
이는 명백하다. 실제로 $z'\in \Ob(\mathcal{S}')$이면
$U = p'(z')$이고 $z'\in \Ob(\mathcal{S}'_U)$이다.
$G_U$가 본질적으로 전사이므로 어떤 $z\in \Ob(\mathcal{S}_U)$에 대한
$G_U(z)$ 꼴의 대상과 $\mathcal{S}'_U$에서 $z'$가 동형임을 안다.
$\mathcal{S}'_U$의 사상은 $\mathcal{S}'$의 사상이므로 $z'$는
$\mathcal{S}'$에서 $G(z)$와 동형이다.
\end{proof}

\begin{lemma}
\label{categories-lemma-fully-faithful-diagonal-equivalence}
$\mathcal{C}$를 범주라 하자. $p : \mathcal{S}\to \mathcal{C}$와
$p' : \mathcal{S'}\to \mathcal{C}$를 군체로 섬유화된 범주들이라 하자.
$G : \mathcal{S}\to \mathcal {S}'$를 $\mathcal{C}$ 위의 함자라 하자.
$G$가 충실충만할 필요충분조건은 대각 사상
$$
\Delta_G :
\mathcal{S}
\longrightarrow
\mathcal{S} \times_{G, \mathcal{S}', G} \mathcal{S}
$$
이 동치인 것이다.
\end{lemma}

\begin{proof}
Lemma \ref{categories-lemma-equivalence-fibred-categories}에 의해 $\mathcal{C}$의
대상 $U$ 위의 섬유 범주들만 살펴보면 충분하다. 우변의 대상은 삼중항
$(x, x', \alpha)$이며, 여기서
$\alpha : G(x) \to G(x')$는 $\mathcal{S}'_U$의 사상이다. 함자
$\Delta_G$는 $\mathcal{S}_U$의 대상 $x$를 삼중항
$(x, x, \text{id}_{G(x)})$로 보낸다. $(x, x', \alpha)$가
$\Delta_G$의 본질적 상에 속할 필요충분조건은 어떤
$\mathcal{S}_U$의 사상 $\beta : x \to x'$에 대하여
$\alpha = G(\beta)$인 것이다(세부 사항은 생략한다). 따라서
$\Delta_G$가 동치가 되려면 모든 $\alpha$가 어떤 사상
$\beta : x \to x'$의 상이어야 하고, 서로 다른 두 사상
$\beta, \beta' : x \to x'$는 서로 다른 사상 $G(\beta), G(\beta')$를
주어야 한다. 이로써 보조정리가 증명된다.
\end{proof}

\begin{lemma}
\label{categories-lemma-morphisms-equivalent-fibred-groupoids}
$\mathcal{C}$를 범주라 하자. $\mathcal{S}_i$, $i = 1, 2, 3, 4$를
$\mathcal{C}$ 위에서 군체로 섬유화된 범주들이라 하자.
$\varphi : \mathcal{S}_1 \to \mathcal{S}_2$와
$\psi : \mathcal{S}_3 \to \mathcal{S}_4$가 $\mathcal{C}$ 위의
동치라고 하자. 그러면
$$
\Mor_{\textit{Cat}/\mathcal{C}}(\mathcal{S}_2, \mathcal{S}_3)
\longrightarrow
\Mor_{\textit{Cat}/\mathcal{C}}(\mathcal{S}_1, \mathcal{S}_4),
\quad \alpha \longmapsto \psi \circ \alpha \circ \varphi
$$
는 범주들의 동치이다.
\end{lemma}

\begin{proof}
이는 일반적인 사실이며 모든 $2$-범주에서 성립한다.
\end{proof}

\begin{lemma}
\label{categories-lemma-inertia-fibred-groupoids}
$\mathcal{C}$를 범주라 하자.
$p : \mathcal{S} \to \mathcal{C}$가 군체로 섬유화되어 있으면 관성
섬유화 범주 $\mathcal{I}_\mathcal{S} \to \mathcal{C}$도 그러하다.
\end{lemma}

\begin{proof}
Lemma \ref{categories-lemma-inertia-fibred-category}의 구성에서 명백하다. 또는 같은
보조정리의 동치
$I_\mathcal{S} \to \mathcal{S}
\times_{\Delta, \mathcal{S} \times_\mathcal{C} \mathcal{S}, \Delta}\mathcal{S}$
를 사용하고 Lemma \ref{categories-lemma-2-product-fibred-categories}를 적용해도 된다.
\end{proof}

\begin{lemma}
\label{categories-lemma-cute-groupoids}
$\mathcal{C}$를 범주라 하고 $U \in \Ob(\mathcal{C})$라 하자.
$p : \mathcal{S} \to \mathcal{C}$가 군체로 섬유화된 범주이고 $p$가
$p' : \mathcal{S} \to \mathcal{C}/U$를 통해 분해되면,
$p' : \mathcal{S} \to \mathcal{C}/U$는 군체로 섬유화되어 있다.
\end{lemma}

\begin{proof}
Lemma \ref{categories-lemma-cute}에서 $p'$가 섬유화된 범주임을 이미 보았다.
따라서 Lemma \ref{categories-lemma-fibred-groupoids}에 의해 섬유 범주들이
군체임을 보이는 것으로 충분하다. $V \in \Ob(\mathcal{C})$에 대하여
$$
\mathcal{S}_V = \coprod\nolimits_{f : V \to U} \mathcal{S}_{(f : V \to U)}
$$
이다. 여기서 좌변은 $p$의 섬유 범주이고 우변은 $p'$의 섬유 범주들의
분리합집합이다. 따라서 결과가 따른다.
\end{proof}

\begin{lemma}
\label{categories-lemma-fibred-in-groupoids-over-fibred-in-groupoids}
$\mathcal{A} \to \mathcal{B} \to \mathcal{C}$를 범주 사이의 함자들이라
하자. $\mathcal{A}$가 $\mathcal{B}$ 위에서 군체로 섬유화되어 있고
$\mathcal{B}$가 $\mathcal{C}$ 위에서 군체로 섬유화되어 있으면,
$\mathcal{A}$는 $\mathcal{C}$ 위에서 군체로 섬유화되어 있다.
\end{lemma}

\begin{proof}
정의에서 직접 증명할 수 있다. 그러나 여기서는
Lemma \ref{categories-lemma-fibred-groupoids}의 판정법을 사용하겠다.
Lemma \ref{categories-lemma-fibred-over-fibred}에 의해 $\mathcal{A}$는
$\mathcal{C}$ 위에서 섬유화되어 있다. 증명을 끝내기 위해
$\mathcal{C}$의 $U$에 대한 섬유 범주 $\mathcal{A}_U$가 군체임을
보이자. 실제로 $x \to y$가 $\mathcal{A}_U$의 사상이면 그
$\mathcal{B}$에서의 상은 $\mathcal{B}_U$가 군체이므로 동형사상이다.
그러면 $x \to y$도 동형사상이다. 예를 들어 이는
Lemma \ref{categories-lemma-composition-cartesian} 및 $\mathcal{A}$의 모든 사상이
강한 $\mathcal{B}$-카르테시안 사상이라는 사실에서 따른다
(Lemma \ref{categories-lemma-fibred-groupoids}를 보라).
\end{proof}

\begin{lemma}
\label{categories-lemma-fibred-groupoids-fibre-product-goes-up}
$p : \mathcal{S} \to \mathcal{C}$를 군체로 섬유화된 범주라 하자.
$x \to y$와 $z \to y$를 $\mathcal{S}$의 사상들이라 하자.
$p(x) \times_{p(y)} p(z)$가 존재하면 $x \times_y z$가 존재하고
$p(x \times_y z) = p(x) \times_{p(y)} p(z)$이다.
\end{lemma}

\begin{proof}
Lemma \ref{categories-lemma-fibred-category-representable-goes-up}에서 따른다.
\end{proof}

\begin{lemma}
\label{categories-lemma-ameliorate-morphism-categories-fibred-groupoids}
$\mathcal{C}$를 범주라 하고, $F : \mathcal{X} \to \mathcal{Y}$를
$\mathcal{C}$ 위에서 군체로 섬유화된 범주들의 $1$-사상이라 하자.
$\mathcal{C}$ 위에서 군체로 섬유화된 범주들의 $1$-사상에 의한 분해
$\mathcal{X} \to \mathcal{X}' \to \mathcal{Y}$가 존재하여,
$\mathcal{X} \to \mathcal{X}'$는 $\mathcal{C}$ 위의 동치이고
$\mathcal{X}'$는 $\mathcal{Y}$ 위에서 군체로 섬유화된 범주가 된다.
\end{lemma}

\begin{proof}
구조 함자들을 $p : \mathcal{X} \to \mathcal{C}$와
$q : \mathcal{Y} \to \mathcal{C}$로 나타내자. $\mathcal{X}'$를 다음과
같이 명시적으로 구성한다. $\mathcal{X}'$의 대상은 사중항
$(U, x, y, f)$이다. 여기서 $x \in \Ob(\mathcal{X}_U)$,
$y \in \Ob(\mathcal{Y}_U)$이고 $f : F(x) \to y$는
$\mathcal{Y}_U$의 동형사상이다. 사상
$(a, b) : (U, x, y, f) \to (U', x', y', f')$은
$p(a) = q(b)$ 및 $f' \circ F(a) = b \circ f$를 만족하는
$a : x \to x'$와 $b : y \to y'$로 주어진다. 다시 말해
$\mathcal{X}' = \mathcal{X} \times_{F, \mathcal{Y}, \text{id}} \mathcal{Y}$
이며, 이는 Lemma \ref{categories-lemma-2-product-categories-over-C}의 $2$-섬유곱
구성이다. Lemma \ref{categories-lemma-2-product-fibred-categories}에 의해
$\mathcal{X}'$는 $\mathcal{C}$ 위에서 군체로 섬유화된 범주이고,
$\mathcal{X}' \to \mathcal{Y}$는 $\mathcal{C}$ 위의 범주들의
사상이다. 함자 $\mathcal{X} \to \mathcal{X}'$는 대상에서
$x \mapsto (p(x), x, F(x), \text{id}_{F(x)})$로, 사상에서
$(a : x \to x') \mapsto (a, F(a))$로 정한다. 합성
$\mathcal{X} \to \mathcal{X}' \to \mathcal{Y}$가 $F$와 같음은 명백하다.
$\mathcal{X} \to \mathcal{X}'$가 $\mathcal{C}$ 위의 섬유화된 범주들의
동치라는 확인은 생략한다.

\medskip\noindent
끝으로 $\mathcal{X}' \to \mathcal{Y}$가 군체로 섬유화된 범주임을
보여야 한다. $b : y' \to y$를 $\mathcal{Y}$의 사상이라 하고,
$(U, x, y, f)$를 $y$ 위에 놓인 $\mathcal{X}'$의 대상이라 하자.
$\mathcal{X}$가 $\mathcal{C}$ 위에서 군체로 섬유화되어 있으므로
$U' = q(y') \to q(y) = U$ 위에 놓인 사상 $a : x' \to x$를 찾을 수
있다. $\mathcal{Y}$가 $\mathcal{C}$ 위에서 군체로 섬유화되어 있고
$F(x') \to F(x)$와 $y' \to y$가 모두 같은 사상 $U' \to U$ 위에
놓이므로, $f \circ F(a) = b \circ f'$를 만족하고
$\text{id}_{U'}$ 위에 놓인 $f' : F(x') \to y'$를 찾을 수 있다.
따라서 $(a, b) : (U', x', y', f') \to (U, x, y, f)$를 얻는다.
이는 Definition \ref{categories-definition-fibred-groupoids}의 첫 번째 조건 (1)을
확인한다. (2)를 확인하기 위해
$(a, b) : (U', x', y', f') \to (U, x, y, f)$와
$(a', b') : (U'', x'', y'', f'') \to (U, x, y, f)$를
$\mathcal{X}'$의 사상들이라 하고, $b' \circ b'' = b$를 만족하는
$\mathcal{Y}$의 사상 $b'' : y' \to y''$가 주어졌다고 하자.
$f'' \circ F(a'') = b'' \circ f'$ 및
$(a', b') \circ (a'', b'') = (a, b)$를 만족하는 유일한 사상
$a'' : x' \to x''$가 존재함을 보여야 한다. $\mathcal{X}$가 군체로
섬유화되어 있으므로 $a' \circ a'' = a$와 $p(a'') = q(b'')$를
만족하는 유일한 사상 $a'' : x' \to x''$가 존재한다.
$\mathcal{Y}$가 군체로 섬유화되어 있으므로 $F(a'')$은
$F(a') \circ F(a'') = F(a)$와 $q(F(a'')) = q(b'')$를 만족하는
유일한 사상 $F(x') \to F(x'')$이다. 이 사실과 주어진 관계
$f \circ F(a) = b \circ f'$ 및 $f \circ F(a') = b' \circ f''$로부터
$f'' \circ F(a'') = b'' \circ f'$가 따른다.
\end{proof}

\begin{lemma}
\label{categories-lemma-amelioration-unique}
$\mathcal{C}$를 범주라 하고, $F : \mathcal{X} \to \mathcal{Y}$를
$\mathcal{C}$ 위에서 군체로 섬유화된 범주들의 $1$-사상이라 하자.
$2$-가환 도표
$$
\xymatrix{
\mathcal{X}' \ar[rd]_f &
\mathcal{X} \ar[l]^a \ar[d]^F \ar[r]_b &
\mathcal{X}'' \ar[ld]^g \\
& \mathcal{Y}
}
$$
가 주어졌다고 하자. 여기서 $a$와 $b$는 $\mathcal{C}$ 위의 범주들의
동치이고 $f$와 $g$는 군체로 섬유화된 범주들이다. 그러면
$\mathcal{Y}$ 위의 범주들의 동치 $h : \mathcal{X}'' \to \mathcal{X}'$가
존재하여, $h \circ b$는 $\mathcal{C}$ 위의 범주들의 $1$-사상으로서
$a$와 $2$-동형이다. 위 도표가 실제로 가환하면, $h \circ b$가
$\mathcal{Y}$ 위의 범주들의 $1$-사상으로서 $a$와 $2$-동형이 되도록
정할 수 있다.
\end{lemma}

\begin{proof}
$\mathcal{Y}$ 위의 $\mathcal{X}'$와 $\mathcal{X}''$가 모두
$\mathcal{Y}$ 위에서 군체로 섬유화된 범주
$\mathcal{X} \times_{F, \mathcal{Y}, \text{id}} \mathcal{Y}$와 동치임을
보이겠다. Lemma
\ref{categories-lemma-ameliorate-morphism-categories-fibred-groupoids}의 증명을 보라.
$\mathcal{C}$ 위의 범주들의 $2$-범주에서 준역
$b^{-1} : \mathcal{X}'' \to \mathcal{X}$를 선택하자. 도표의 오른쪽
삼각형이 $2$-가환이므로
$$
\xymatrix{
\mathcal{X} \ar[d]_F & \mathcal{X}'' \ar[l]^{b^{-1}} \ar[d]^g \\
\mathcal{Y} & \mathcal{Y} \ar[l]
}
$$
가 $2$-가환이다. 따라서 $2$-섬유곱의 보편 성질에 의해 $1$-사상
$c : \mathcal{X}'' \to
\mathcal{X} \times_{F, \mathcal{Y}, \text{id}} \mathcal{Y}$를 얻는다.
더 나아가 $c$는 $\mathcal{Y}$ 위의 범주들의 사상이고(!) 동치이다
($b$가 동치라는 가정과 Lemma \ref{categories-lemma-equivalence-2-fibre-product}를
보라). 따라서 Lemma \ref{categories-lemma-equivalence-fibred-categories}에 의해
$c$는 $\mathcal{Y}$ 위에서 군체로 섬유화된 범주들의 $2$-범주에서의
동치이다.

\medskip\noindent
이제 $c \circ b$와 Lemma
\ref{categories-lemma-ameliorate-morphism-categories-fibred-groupoids}의 증명에서
구성한 함자
$d : \mathcal{X} \to
\mathcal{X} \times_{F, \mathcal{Y}, \text{id}} \mathcal{Y}$,
$x \mapsto (p(x), x, F(x), \text{id}_{F(x)})$ 사이의 $2$-동형을
구성해야 한다. $\alpha : F \to g \circ b$와
$\beta : b^{-1} \circ b \to \text{id}$를 $\mathcal{C}$ 위의 범주들의
$1$-사상 사이의 $2$-동형이라고 하자. $c \circ b$는 대상에서 규칙
$$
x \mapsto (p(x), b^{-1}(b(x)), g(b(x)), \alpha_x \circ F(\beta_x))
$$
으로 주어진다는 점에 유의하라. 그러면
$$
(\beta_x, \alpha_x) :
(p(x), x, F(x), \text{id}_{F(x)})
\longrightarrow
(p(x), b^{-1}(b(x)), g(b(x)), \alpha_x \circ F(\beta_x))
$$
는 우리의 $2$-사상 $d \to b \circ c$를 주는 함자적 동형사상이다.
끝으로 도표가 가환하면 모든 $x$에 대해 $\alpha_x$가 항등이고,
이 $2$-사상은 $\mathcal{Y}$ 위의 범주들의 $2$-범주에서의
$2$-사상임을 알 수 있다.
\end{proof}


































\section{범주들의 준층}
\label{categories-section-presheaves-categories}

\noindent
이 절에서는 섬유화된 범주의 개념을 이와 밀접하게 관련된
``범주들의 준층''이라는 개념과 비교한다. 기본 구성은 다음 예에서
설명한다.

\begin{example}
\label{categories-example-functor-categories}
$\mathcal{C}$를 범주라 하자. $F : \mathcal{C}^{opp} \to \textit{Cat}$가
범주들의 $2$-범주로 가는 함자라고 하자. Definition
\ref{categories-definition-functor-into-2-category}를 보라. $\mathcal{C}$의 사상
$f : V \to U$에 대하여 $F(U)$에서 $F(V)$로 가는 함자 $F(f) = f^\ast$를
암시적으로 이렇게 표기하겠다. 이로부터 $\mathcal{C}$ 위의 섬유화된
범주 $\mathcal{S}_F$를 다음과 같이 구성할 수 있다. 다음과 같이 정의하자.
$$
\Ob(\mathcal{S}_F) =
\{(U, x) \mid U\in \Ob(\mathcal{C}), x\in \Ob(F(U))\}.
$$
$(U, x), (V, y) \in \Ob(\mathcal{S}_F)$에 대하여 다음과 같이 정의한다.
\begin{align*}
\Mor_{\mathcal{S}_F}((V, y), (U, x)) & =
\{ (f, \phi) \mid f \in \Mor_\mathcal{C}(V, U),
\phi \in \Mor_{F(V)}(y, f^\ast x)\} \\
& =
\coprod\nolimits_{f \in \Mor_\mathcal{C}(V, U)}
\Mor_{F(V)}(y, f^\ast x)
\end{align*}
합성을 정의하기 위해 $\mathcal{C}$의 합성 가능한 사상 쌍에 대한
$g^\ast \circ f^\ast = (f \circ g)^\ast$를 사용한다
($2$-범주로 가는 함자의 정의에 의한다). 구체적으로
$\psi : z \to g^\ast y$와 $ \phi : y \to f^\ast x$의 합성을
$ g^\ast(\phi) \circ \psi$로 정의한다. 함자
$p_F : \mathcal{S}_F \to \mathcal{C}$는 규칙 $(U, x) \mapsto U$로
주어진다. 이것이 실제로 섬유화된 범주인지 확인하자.
$\mathcal{C}$의 $f: V \to U$와 $U$의 올림 $(U, x)$가 주어졌을 때,
$(f, \text{id}_{f^\ast x}): (V, {f^\ast x}) \to (U, x)$가 $f$의
강한 카르테시안 올림이라고 주장한다. 왼쪽 도표의 $h$가 오른쪽의
$(h, \nu)$를 결정함을 보여야 한다.
$$
\xymatrix{
V \ar[r]^f &
U &
(V, f^*x) \ar[r]^{(f, \text{id}_{f^*x})} &
(U, x) \\
W \ar@{-->}[u]^h \ar[ru]_g & &
(W, z) \ar@{-->}[u]^{(h, \nu)} \ar[ru]_{(g, \psi)} &
}
$$
$\nu = \psi$로 취하면 된다. 이는 $f \circ h = g$이고 따라서
$g^*x = h^*f^*x$이므로 성립한다. 더 나아가 이것은 오른쪽 도표를
가환하게 하는 유일한 올림이다.
\end{example}

\begin{definition}
\label{categories-definition-split-fibred-category}
$\mathcal{C}$를 범주라 하자. $F : \mathcal{C}^{opp} \to \textit{Cat}$가
범주들의 $2$-범주로 가는 함자라고 하자. Example
\ref{categories-example-functor-categories}에서 구성한 섬유화된 범주를
$p_F : \mathcal{S}_F \to \mathcal{C}$로 쓰겠다.
{\it 분할된 섬유화 범주}란 이 범주들 가운데 하나인
{\it $\mathcal{S}_F$}와 $\mathcal{C}$ 위에서 동형인(!) 섬유화된
범주를 말한다.
\end{definition}

\begin{lemma}
\label{categories-lemma-when-split}
$\mathcal{C}$를 범주라 하고, $\mathcal{S}$를 $\mathcal{C}$ 위의
섬유화된 범주라 하자. 그러면 $\mathcal{S}$가 분할될 필요충분조건은
어떤 당김의 선택에 대하여(Definition
\ref{categories-definition-pullback-functor-fibred-category}를 보라) 당김 함자
$(f \circ g)^*$와 $g^* \circ f^*$가 같은 것이다.
\end{lemma}

\begin{proof}
정의에서 바로 따른다.
\end{proof}

\begin{lemma}
\label{categories-lemma-fibred-strict}
$ p : \mathcal{S} \to \mathcal{C}$를 섬유화된 범주라 하자.
$\mathcal{S}$가 $\mathcal{C}$ 위의 섬유화된 범주들의 $2$-범주에서
$\mathcal{S}_F$와 동치가 되는 반변함자
$F : \mathcal{C} \to \textit{Cat}$가 존재한다. 다시 말해 모든
섬유화된 범주는 분할된 범주와 동치이다.
\end{lemma}

\begin{proof}
당김을 하나씩 선택하자(Definition
\ref{categories-definition-pullback-functor-fibred-category}를 보라).
Lemma \ref{categories-lemma-fibred}에 의해 $\mathcal{C}$의 모든 사상 $f$에 대한
당김 함자 $f^*$를 얻는다.

\medskip\noindent
새 범주 $\mathcal{S}'$를 다음과 같이 구성한다. $\mathcal{S}'$의
대상은 $\mathcal{C}$의 사상 $f : V \to U$와 $U$ 위의
$\mathcal{S}$의 대상 $x$, 즉 $x\in \Ob(\mathcal{S}_U)$로 이루어진 쌍
$(x, f)$이다. 함자 $p' : \mathcal{S}' \to \mathcal{C}$는 쌍
$(x, f)$를 사상 $f$의 시작점으로 보낸다. 다시 말해
$p'(x, f : V\to U) = V$이다. 사상
$\varphi : (x_1, f_1: V_1 \to U_1) \to (x_2, f_2 : V_2 \to U_2)$는
$p(\varphi) = g$를 만족하는 사상 $g : V_1 \to V_2$와
$\varphi : f_1^\ast x_1 \to f_2^\ast x_2$로 이루어진 쌍
$(\varphi, g)$으로 주어진다. 합성 법칙은 아무 문제 없이 정의된다.
합성 가능한 임의의 사상 쌍에 대하여
$(\varphi, g) \circ (\psi, h) =
(\varphi \circ \psi, g\circ h)$로 놓는다. $U$ 위의 $x$를
쌍 $(x, \text{id}_U)$로 보내는 자연스러운 함자
$\mathcal{S} \to \mathcal{S}'$가 있다.

\medskip\noindent
이제 $p'$가 $\mathcal{S}'$를 $\mathcal{C}$ 위의 섬유화된 범주로
만드는지, 그리고 $\mathcal{S} \to \mathcal{S}'$가 강한 카르테시안
사상을 강한 카르테시안 사상으로 보내는 $\mathcal{C}$ 위의 범주들의
동치인지 확인해야 한다. 이 확인들은 생략한다.

\medskip\noindent
끝으로 $g : V' \to V$와 $f : V \to U$일 때 대상에서
$g^\ast(x, f) = (x, f \circ g)$로 놓음으로써 $\mathcal{S}'$ 위에
당김 함자를 정의할 수 있다. $\mathcal{S}'_V$의 사상들 사이의 사상
$(\varphi, \text{id}_V) : (x_1, f_1) \to (x_2, f_2)$에 대해서는
$g^\ast(\varphi, \text{id}_V) =
(g^\ast\varphi, \text{id}_{V'})$로 놓는다. 여기서 Lemma
\ref{categories-lemma-fibred}의 유일한 식별
$g^\ast f_i^\ast x_i = (f_i \circ g)^\ast x_i$를 사용하여
$g^\ast\varphi$를 $(f_1 \circ g)^\ast x_1$에서
$(f_2 \circ g)^\ast x_2$로 가는 사상으로 생각한다. 이 당김 함자들
$g^\ast$가 $g_1^\ast \circ g_2^\ast = (g_2\circ g_1)^\ast$를
만족함은 명백하다. 다시 말해 $\mathcal{S}'$는 원하는 대로 분할된다.
\end{proof}



















\section{군체들의 준층}
\label{categories-section-presheaves-groupoids}

\noindent
이 절에서는 군체로 섬유화된 범주의 개념을 이와 밀접하게 관련된
``군체들의 준층''이라는 개념과 비교한다. 기본 구성은 다음 예에서
설명한다.

\begin{example}
\label{categories-example-functor-groupoids}
이 예는 ``범주들의 준층'' 대신 ``군체들의 준층''을 다루는
Example \ref{categories-example-functor-categories}의 유사물이다. 결과물은
섬유화된 범주 대신 군체로 섬유화된 범주가 된다.
$F : \mathcal{C}^{opp} \to \textit{Groupoids}$가 군체들의 범주로 가는
함자라고 하자. Definition \ref{categories-definition-functor-into-2-category}를 보라.
$\mathcal{C}$의 사상 $f : V \to U$에 대하여 $F(U)$에서 $F(V)$로
가는 함자 $F(f) = f^\ast$를 암시적으로 이렇게 표기하겠다.
$\mathcal{C}$ 위에서 군체로 섬유화된 범주 $\mathcal{S}_F$를 다음과
같이 구성한다. 다음과 같이 정의하자.
$$
\Ob(\mathcal{S}_F) =
\{(U, x) \mid U\in \Ob(\mathcal{C}), x\in \Ob(F(U))\}.
$$
$(U, x), (V, y) \in \Ob(\mathcal{S}_F)$에 대하여 다음과 같이 정의한다.
\begin{align*}
\Mor_{\mathcal{S}_F}((V, y), (U, x))
& =
\{ (f, \phi) \mid f \in \Mor_\mathcal{C}(V, U),
\phi \in \Mor_{F(V)}(y, f^\ast x)\} \\
& =
\coprod\nolimits_{f \in \Mor_\mathcal{C}(V, U)}
\Mor_{F(V)}(y, f^\ast x)
\end{align*}
합성을 정의하기 위해 $\mathcal{C}$의 합성 가능한 사상 쌍에 대한
$g^\ast \circ f^\ast = (f \circ g)^\ast$를 사용한다
($2$-범주로 가는 함자의 정의에 의한다). 구체적으로
$\psi : z \to g^\ast y$와 $ \phi : y \to f^\ast x$의 합성을
$ g^\ast(\phi) \circ \psi$로 정의한다. 함자
$p_F : \mathcal{S}_F \to \mathcal{C}$는 규칙 $(U, x) \mapsto U$로
주어진다. 모든 $U$에 대하여 $F(U)$가 군체라는 조건에 의해
$\mathcal{S}_F$는 $\mathcal{C}$ 위에서 군체로 섬유화되어 있다.
실제로 Example \ref{categories-example-functor-categories}에서 $\mathcal{S}_F$가
섬유화된 범주임을 이미 보았으므로 Lemma \ref{categories-lemma-fibred-groupoids}를
적용할 수 있다. 그러나 Definition \ref{categories-definition-fibred-groupoids}의
조건 (1), (2)를 다음과 같이 직접 증명할 수도 있다. (1) 사상의 올림은
존재한다. 실제로 $\mathcal{C}$의 $f: V \to U$와 $U$ 위에 놓인
$\mathcal{S}_F$의 대상 $(U, x)$가 주어지면
$(f, \text{id}_{f^\ast x}): (V, {f^\ast x}) \to (U, x)$는 $f$의
올림이다. (2) 다음 실선 도표들이 주어졌다고 하자.
$$
\xymatrix{
V \ar[r]^f & U & (V, y) \ar[r]^{(f, \phi)} & (U, x) \\
W \ar@{-->}[u]^h \ar[ru]_g & &
(W, z) \ar@{-->}[u]^{(h, \nu)} \ar[ru]_{(g, \psi)} & \\
}
$$
그러면 점선 화살표에 대해서는
$\nu = (h^\ast \phi)^{-1} \circ \psi$이다. 따라서 $h$가 주어지면
$\nu$가 존재하고, 역의 유일성에 의해 유일하다.
\end{example}

\begin{definition}
\label{categories-definition-split-category-fibred-in-groupoids}
$\mathcal{C}$를 범주라 하자.
$F : \mathcal{C}^{opp} \to \textit{Groupoids}$가 군체들의 $2$-범주로
가는 함자라고 하자. Example \ref{categories-example-functor-groupoids}에서 구성한
군체로 섬유화된 범주를 $p_F : \mathcal{S}_F \to \mathcal{C}$로
쓰겠다. {\it 분할된 군체로 섬유화된 범주}란 이 범주들 가운데 하나인
{\it $\mathcal{S}_F$}와 $\mathcal{C}$ 위에서 동형인(!), 군체로
섬유화된 범주를 말한다.
\end{definition}

\begin{lemma}
\label{categories-lemma-fibred-groupoids-strict}
$ p : \mathcal{S} \to \mathcal{C}$를 군체로 섬유화된 범주라 하자.
$\mathcal{S}$가 $\mathcal{C}$ 위에서 $\mathcal{S}_F$와 동치가 되는
반변함자 $F : \mathcal{C} \to \textit{Groupoids}$가 존재한다.
다시 말해 군체로 섬유화된 모든 범주는 분할된 범주와 동치이다.
\end{lemma}

\begin{proof}
당김을 하나씩 선택하자(Definition
\ref{categories-definition-pullback-functor-fibred-category}를 보라).
Lemmas \ref{categories-lemma-fibred}와 \ref{categories-lemma-fibred-groupoids}에 의해
$\mathcal{C}$의 모든 사상 $f$에 대한 당김 함자 $f^*$를 얻는다.

\medskip\noindent
새 범주 $\mathcal{S}'$를 다음과 같이 구성한다. $\mathcal{S}'$의
대상은 $\mathcal{C}$의 사상 $f : V \to U$와 $U$ 위의
$\mathcal{S}$의 대상 $x$, 즉 $x\in \Ob(\mathcal{S}_U)$로 이루어진 쌍
$(x, f)$이다. 함자 $p' : \mathcal{S}' \to \mathcal{C}$는 쌍
$(x, f)$를 사상 $f$의 시작점으로 보낸다. 다시 말해
$p'(x, f : V\to U) = V$이다. 사상
$\varphi : (x_1, f_1: V_1 \to U_1) \to (x_2, f_2 : V_2 \to U_2)$는
$p(\varphi) = g$를 만족하는 사상 $g : V_1 \to V_2$와
$\varphi : f_1^\ast x_1 \to f_2^\ast x_2$로 이루어진 쌍
$(\varphi, g)$으로 주어진다. 합성 가능한 임의의 사상 쌍에 대한
합성 법칙은 $(\varphi, g) \circ (\psi, h) =
(\varphi \circ \psi, g\circ h)$로 아무 문제 없이 정의된다.
$U$ 위의 $x$를 쌍 $(x, \text{id}_U)$로 보내는 자연스러운 함자
$\mathcal{S} \to \mathcal{S}'$가 있다.

\medskip\noindent
이제 $p'$가 $\mathcal{S}'$를 $\mathcal{C}$ 위에서 군체로 섬유화된
범주로 만드는지, 그리고 $\mathcal{S} \to \mathcal{S}'$가
$\mathcal{C}$ 위의 범주들의 동치인지 확인해야 한다. 이 확인들은
생략한다.

\medskip\noindent
끝으로 $g : V' \to V$와 $f : V \to U$일 때 대상에서
$g^\ast(x, f) = (x, f \circ g)$로 놓음으로써 $\mathcal{S}'$ 위에
당김 함자를 정의할 수 있다. $\mathcal{S}'_V$의 사상들 사이의 사상
$(\varphi, \text{id}_V) : (x_1, f_1) \to (x_2, f_2)$에 대해서는
$g^\ast(\varphi, \text{id}_V) =
(g^\ast\varphi, \text{id}_{V'})$로 놓는다. 여기서 Lemma
\ref{categories-lemma-fibred-groupoids}의 유일한 식별
$g^\ast f_i^\ast x_i = (f_i \circ g)^\ast x_i$를 사용하여
$g^\ast\varphi$를 $(f_1 \circ g)^\ast x_1$에서
$(f_2 \circ g)^\ast x_2$로 가는 사상으로 생각한다. 이 당김 함자들
$g^\ast$가 $g_1^\ast \circ g_2^\ast = (g_2\circ g_1)^\ast$를
만족함은 명백하다. 다시 말해 $\mathcal{S}'$는 원하는 대로 분할된다.
\end{proof}

\noindent
이 보조정리의 다른 증명은 Section
\ref{categories-section-representable-1-morphisms}에서 보게 될 것이다.








\section{집합들로 섬유화된 범주들}
\label{categories-section-fibred-in-sets}

\begin{definition}
\label{categories-definition-discrete}
항등사상만을 사상으로 갖는 범주를 {\it 이산 범주}라고 한다.
\end{definition}

\noindent
이산 범주에는 흥미로운 정보가 단 하나, 즉 대상들의 집합만 있다.
따라서 때때로 이산 범주와 집합을 같은 것으로 생각한다.

\begin{definition}
\label{categories-definition-category-fibred-sets}
$\mathcal{C}$를 범주라 하자.
{\it 집합들로 섬유화된 범주}, 또는
{\it 이산 범주들로 섬유화된 범주}란 모든 섬유 범주가 이산 범주인,
군체로 섬유화된 범주를 말한다.
\end{definition}

\noindent
집합들로 섬유화된 범주와 준층 사이의 관계를 명확히 하고자 한다
(Definition \ref{categories-definition-presheaf}를 보라). 이를 위해 먼저 다음을
정의하는 것이 자연스럽다.

\begin{definition}
\label{categories-definition-categories-fibred-in-sets-over-C}
$\mathcal{C}$를 범주라 하자.
{\it $\mathcal{C}$ 위에서 집합들로 섬유화된 범주들의 $2$-범주}는
$\mathcal{C}$ 위에서 군체로 섬유화된 범주들의 범주의 부분 $2$-범주이며
(Definition \ref{categories-definition-categories-fibred-in-groupoids-over-C}를 보라),
다음과 같이 정의된다.
\begin{enumerate}
\item 그 대상은 집합들로 섬유화된 범주
$p : \mathcal{S} \to \mathcal{C}$이다.
\item 그 $1$-사상 $(\mathcal{S}, p) \to (\mathcal{S}', p')$은
$p' \circ G = p$를 만족하는 함자 $G : \mathcal{S} \to \mathcal{S}'$이다
(모든 사상이 강한 카르테시안 사상이므로 $G$는 자동으로 이를 보존한다).
\item
$G, H : (\mathcal{S}, p) \to (\mathcal{S}', p')$에 대한 그 $2$-사상
$t : G \to H$는 모든 $x \in \Ob(\mathcal{S})$에 대하여
$p'(t_x) = \text{id}_{p(x)}$를 만족하는 함자의 사상이다.
\end{enumerate}
\end{definition}

\noindent
모든 $2$-사상은 자동으로 동형사상임에 유의하라. 따라서 이 $2$-범주는
실제로 $(2, 1)$-범주이다. 다음은 $2$-섬유곱의 존재에 관한 필수적인
보조정리이다.

\begin{lemma}
\label{categories-lemma-2-product-categories-fibred-sets}
$\mathcal{C}$를 범주라 하자. $\mathcal{C}$ 위에서 집합들로 섬유화된
범주들의 2-범주는 2-섬유곱을 갖는다. 더 정확히 말해, 집합들로
섬유화된 범주들에서 시작하면 Lemma
\ref{categories-lemma-2-product-categories-over-C}에서 기술한 2-섬유곱은 다시
집합들로 섬유화된 범주를 준다.
\end{lemma}

\begin{proof}
생략한다.
\end{proof}

\begin{example}
\label{categories-example-presheaf}
이 예는 ``범주들의 준층'' 대신 준층을 다루는
Examples \ref{categories-example-functor-categories}와
\ref{categories-example-functor-groupoids}의 유사물이다. 결과물은 섬유화된 범주
대신 집합들로 섬유화된 범주가 된다.
$F : \mathcal{C}^{opp} \to \textit{Sets}$를 준층이라 하자.
$\mathcal{C}$의 사상 $f : V \to U$에 대하여
$F(f) = f^\ast : F(U) \to F(V)$로 암시적으로 표기하겠다.
$\mathcal{C}$ 위에서 집합들로 섬유화된 범주 $\mathcal{S}_F$를 다음과
같이 구성한다. 다음과 같이 정의하자.
$$
\Ob(\mathcal{S}_F) =
\{(U, x) \mid U \in \Ob(\mathcal{C}), x \in \Ob(F(U))\}.
$$
$(U, x), (V, y) \in \Ob(\mathcal{S}_F)$에 대하여 다음과 같이 정의한다.
\begin{align*}
\Mor_{\mathcal{S}_F}((V, y), (U, x))
& =
\{f \in \Mor_\mathcal{C}(V, U) \mid f^*x = y\}
\end{align*}
합성은 $\mathcal{C}$의 합성에서 물려받는다. 이는 $\mathcal{C}$의 합성
가능한 사상 쌍에 대하여 $g^\ast \circ f^\ast = (f \circ g)^\ast$이므로
잘 정의된다. 함자 $p_F : \mathcal{S}_F \to \mathcal{C}$는 규칙
$(U, x) \mapsto U$로 주어진다. 모든 섬유 범주
$\mathcal{S}_{F, U}$가 바탕집합 $F(U)$를 갖는 이산 범주이고,
Example \ref{categories-example-functor-groupoids}에서 $\mathcal{S}_F$가 군체로
섬유화된 범주임을 이미 보았으므로, $\mathcal{S}_F$는 집합들로
섬유화되어 있다는 결론을 얻는다.
\end{example}

\begin{lemma}
\label{categories-lemma-2-category-fibred-sets}
\begin{slogan}
집합들로 섬유화된 범주는 정확히 준층이다.
\end{slogan}
$\mathcal{C}$를 범주라 하자. 집합들로 섬유화된 범주 사이의 유일한
$2$-사상들은 항등사상이다. 다시 말해 집합들로 섬유화된 범주들의
$2$-범주는 범주이다. 더 나아가 범주들의 동치
$$
\left\{
\begin{matrix}
\text{the category of presheaves}\\
\text{of sets over }\mathcal{C}
\end{matrix}
\right\}
\leftrightarrow
\left\{
\begin{matrix}
\text{the category of categories}\\
\text{fibred in sets over }\mathcal{C}
\end{matrix}
\right\}
$$
가 존재한다. 왼쪽에서 오른쪽으로 가는 함자는
Example \ref{categories-example-presheaf}에서 논의한 구성 $F \to \mathcal{S}_F$이다.
오른쪽에서 왼쪽으로 가는 함자는 $p : \mathcal{S} \to \mathcal{C}$에
대상들의 준층 $U \mapsto \Ob(\mathcal{S}_U)$를 대응시킨다.
\end{lemma}

\begin{proof}
섬유 범주의 유일한 사상들이 항등사상이므로 첫 번째 주장은 명백하다.

\medskip\noindent
$p :
\mathcal{S} \to \mathcal{C}$가 집합들로 섬유화되어 있다고 하자.
$f : V \to U$를 $\mathcal{C}$의 사상이라 하고
$x \in \Ob(\mathcal{S}_U)$라 하자. 그러면 대상 $f^\ast x$에 대한
선택은 정확히 하나이다. 따라서 Lemma \ref{categories-lemma-fibred-groupoids}에서와
같은 $f, g$에 대하여 $(f \circ g)^\ast x = g^\ast(f^\ast x)$이다.
그러므로 대응 $U \mapsto \Ob(\mathcal{S}_U)$와 $f \mapsto f^\ast$를
$\mathcal{C}$ 위의 준층으로 생각할 수 있다.
\end{proof}

\noindent
다음은 집합들로 섬유화된 범주의 중요한 예이다.

\begin{example}
\label{categories-example-fibred-category-from-functor-of-points}
$\mathcal{C}$를 범주라 하고 $X \in \Ob(\mathcal{C})$라 하자.
표현 가능한 준층 $h_X = \Mor_\mathcal{C}(-, X)$를 생각하자
(Example \ref{categories-example-hom-functor}를 보라). 한편 Example
\ref{categories-example-category-over-X}의 범주 $p : \mathcal{C}/X \to \mathcal{C}$를
생각하자. 섬유 범주 $(\mathcal{C}/X)_U$의 대상은 사상
$h : U \to X$이고 사상은 항등사상뿐이다. 따라서 Lemma
\ref{categories-lemma-2-category-fibred-sets}의 대응 아래에서
$$
h_X \longleftrightarrow \mathcal{C}/X.
$$
이다. 다시 말해 범주 $\mathcal{C}/X$는 Example
\ref{categories-example-presheaf}에서 $h_X$에 수반시킨 범주
$\mathcal{S}_{h_X}$와 표준적으로 동치이다.
\end{example}

\noindent
이 때문에 군체로 섬유화된 범주들의 2-범주에서 ``표현 가능한'' 대상을
그에 수반된 준층이 표현 가능한, 집합들로 섬유화된 범주로 정의하고
싶어진다. 그러나 우리는 동치 아래에서 불변인 개념을 원하므로 이는
사용하기에 좋은 정의가 아니다. 이를 정밀하게 만들기 위해 어떤
군체로 섬유화된 범주가 집합들로 섬유화된 범주와 동치인지 정확히
연구한다.









\section{세토이드로 섬유화된 범주들}
\label{categories-section-fibred-in-setoids}

\begin{definition}
\label{categories-definition-setoid}
모든 대상이 정확히 하나의 자기동형사상, 즉 항등사상만을 갖는 군체인
범주를 {\it 세토이드}\footnote{스테로이드를 맞은 집합!?}라고 부르자.
\end{definition}

\noindent
$C$가 동치관계 $\sim$을 갖춘 집합이면 다음과 같이 세토이드
$\mathcal{C}$를 만들 수 있다. $\Ob(\mathcal{C}) = C$로 놓고,
$x \sim y$인 경우가 아니면 $\Mor_\mathcal{C}(x, y) = \emptyset$로,
그 경우에는 $\Mor_\mathcal{C}(x, y) = \{1\}$로 놓는다. $\sim$의
추이성에 의해 사상들을 합성할 수 있다. 역으로 임의의 세토이드 범주는
그 대상들 위에 동치관계(동형)를 정의하며, 방금 기술한 절차로부터 그
범주를 유일한 동형사상까지 복원할 수 있다. 여기서 이는 동치까지라는
뜻이 아니다.

\medskip\noindent
이산 범주는 세토이드이다. 임의의 세토이드 $\mathcal{C}$로부터 그와
동치인 이산 범주를 만드는 표준적 절차가 있다. 즉
$\Ob(\mathcal{C})$를 동형류들의 집합으로 바꾸고 항등사상들을
덧붙인다. 동치관계를 갖춘 집합의 관점에서는 이는 그 동치관계에 의한
몫을 취하는 것에 해당한다.

\begin{definition}
\label{categories-definition-category-fibred-setoids}
$\mathcal{C}$를 범주라 하자. {\it 세토이드로 섬유화된 범주}란 모든
섬유 범주가 세토이드인, 군체로 섬유화된 범주를 말한다.
\end{definition}

\noindent
아래에서 세토이드로 섬유화된 범주와 집합들로 섬유화된 범주의 관계를
명확히 하겠다.

\begin{definition}
\label{categories-definition-categories-fibred-in-setoids-over-C}
$\mathcal{C}$를 범주라 하자.
{\it $\mathcal{C}$ 위에서 세토이드로 섬유화된 범주들의 $2$-범주}는
$\mathcal{C}$ 위에서 군체로 섬유화된 범주들의 범주의 부분 $2$-범주이며
(Definition \ref{categories-definition-categories-fibred-in-groupoids-over-C}를 보라),
다음과 같이 정의된다.
\begin{enumerate}
\item 그 대상은 세토이드로 섬유화된 범주
$p : \mathcal{S} \to \mathcal{C}$이다.
\item 그 $1$-사상 $(\mathcal{S}, p) \to (\mathcal{S}', p')$은
$p' \circ G = p$를 만족하는 함자 $G : \mathcal{S} \to \mathcal{S}'$이다
(모든 사상이 강한 카르테시안 사상이므로 $G$는 자동으로 이를 보존한다).
\item
$G, H : (\mathcal{S}, p) \to (\mathcal{S}', p')$에 대한 그 $2$-사상
$t : G \to H$는 모든 $x \in \Ob(\mathcal{S})$에 대하여
$p'(t_x) = \text{id}_{p(x)}$를 만족하는 함자의 사상이다.
\end{enumerate}
\end{definition}

\noindent
모든 $2$-사상은 자동으로 동형사상임에 유의하라. 따라서 이 $2$-범주는
실제로 $(2, 1)$-범주이다.

\noindent
다음은 $2$-섬유곱의 존재에 관한 필수적인 보조정리이다.

\begin{lemma}
\label{categories-lemma-2-product-categories-fibred-setoids}
$\mathcal{C}$를 범주라 하자. $\mathcal{C}$ 위에서 세토이드로 섬유화된
범주들의 2-범주는 2-섬유곱을 갖는다. 더 정확히 말해, 세토이드로
섬유화된 범주들에서 시작하면 Lemma
\ref{categories-lemma-2-product-categories-over-C}에서 기술한 2-섬유곱은 다시
세토이드로 섬유화된 범주를 준다.
\end{lemma}

\begin{proof}
생략한다.
\end{proof}

\begin{lemma}
\label{categories-lemma-setoid-fibres}
$\mathcal{C}$를 범주라 하고 $\mathcal{S}$를 $\mathcal{C}$ 위의 범주라
하자.
\begin{enumerate}
\item $\mathcal{S}'$가 $\mathcal{C}$ 위에서 집합들로 섬유화되어 있고
$\mathcal{S} \to \mathcal{S}'$가 $\mathcal{C}$ 위의 동치이면 다음이
성립한다.
\begin{enumerate}
\item $\mathcal{S}$는 $\mathcal{C}$ 위에서 세토이드로 섬유화된다.
\item 각 $U \in \Ob(\mathcal{C})$에 대하여 사상
$\Ob(\mathcal{S}_U) \to \Ob(\mathcal{S}'_U)$는 공역을 정의역의
동형류들의 집합과 식별한다.
\end{enumerate}
\item $p : \mathcal{S} \to \mathcal{C}$가 세토이드로 섬유화된
범주이면, 집합들로 섬유화된 범주 $p' : \mathcal{S}' \to \mathcal{C}$와
$\mathcal{C}$ 위의 동치 $\text{can} : \mathcal{S} \to \mathcal{S}'$가
존재한다.
\end{enumerate}
\end{lemma}

\begin{proof}
(2)를 증명하자. 범주 $\mathcal{S}'$의 대상은 쌍 $(U, \xi)$이며, 여기서
$U \in \Ob(\mathcal{C})$이고 $\xi$는 $\mathcal{S}_U$의 대상들의
동형류이다. 사상 $(U, \xi) \to (V , \psi)$는 사상 $x \to y$로
주어지며, 여기서 $x \in \xi$이고 $y \in \psi$이다. 두 사상
$x \to y$와 $x' \to y'$가 같은 사상 $U \to V$를 유도하고,
$\mathcal{S}_U$의 동형사상 $x \to x'$ 및 $\mathcal{S}_V$의 동형사상
$y \to y'$의 어떤 선택에 대하여 합성 $x \to x' \to y'$와
$x \to y \to y'$가 일치하면 이 두 사상을 식별한다. 구성에 의해
$\mathcal{S} \to \mathcal{S}'$에서 대상들과 사상들 위의 전사 사상이
존재한다. $\mathcal{S}'$의 사상들의 합성은
$\mathcal{S} \to \mathcal{S}'$를 함자로 만드는 유일한 법칙으로
정의한다. 몇몇 세부사항은 생략한다.
\end{proof}

\noindent
따라서 세토이드로 섬유화된 범주들은 정확히 집합들로 섬유화된 범주와
동치인, 군체로 섬유화된 범주들이다. 더 나아가 Lemma
\ref{categories-lemma-2-category-fibred-sets}에 의해 집합들로 섬유화된 범주들의
동치는 동형사상이다.

\begin{lemma}
\label{categories-lemma-2-category-fibred-setoids}
$\mathcal{C}$를 범주라 하자. Lemma \ref{categories-lemma-setoid-fibres} (2)의 구성은
함자
$$
F :
\left\{
\begin{matrix}
\text{the 2-category of categories}\\
\text{fibred in setoids over }\mathcal{C}
\end{matrix}
\right\}
\longrightarrow
\left\{
\begin{matrix}
\text{the category of categories}\\
\text{fibred in sets over }\mathcal{C}
\end{matrix}
\right\}
$$
를 준다(Definition \ref{categories-definition-functor-into-2-category}를 보라).
이 함자는 다음 의미에서 동치이다.
\begin{enumerate}
\item $F(f) = F(g)$를 만족하는 임의의 두 1-사상
$f, g : \mathcal{S}_1 \to \mathcal{S}_2$에 대하여 유일한 2-동형사상
$f \to g$가 존재한다.
\item 임의의 사상 $h : F(\mathcal{S}_1) \to F(\mathcal{S}_2)$에 대하여
$F(f) = h$를 만족하는 1-사상 $f : \mathcal{S}_1 \to \mathcal{S}_2$가
존재한다.
\item 집합들로 섬유화된 임의의 범주 $\mathcal{S}$는
$F(\mathcal{S})$와 같다.
\end{enumerate}
특히 규칙 $F_i(U) = \Ob(\mathcal{S}_{i, U})/\cong$로
$F_i \in \textit{PSh}(\mathcal{C})$를 정의하면
$$
\Mor_{\textit{Cat}/\mathcal{C}}(\mathcal{S}_1, \mathcal{S}_2)
\Big/
2\text{-isomorphism}
=
\Mor_{\textit{PSh}(\mathcal{C})}(F_1, F_2)
$$
를 얻는다. 더 정확히 말해 임의의 사상 $\phi : F_1 \to F_2$가 주어지면,
대상들의 동형류 위에서 $\phi$를 유도하며 유일한 $2$-동형사상까지
유일한 $1$-사상 $f : \mathcal{S}_1 \to \mathcal{S}_2$가 존재한다.
\end{lemma}

\begin{proof}
Lemma \ref{categories-lemma-2-category-fibred-sets}에 의해 $F$의 공역은 범주이므로
이 주장은 의미가 있다. Lemma \ref{categories-lemma-setoid-fibres} (2)의 구성은
$\mathcal{S}$에, $U$ 위의 값이 $\mathcal{S}_U$ 안의 동형류들의 집합인
집합들로 섬유화된 범주를 대응시킨다. 따라서 이 구성이 표시한 함자를
정의함은 명백하다. (1)에서와 같이
$F(f) = F(g)$를 만족하는 $f, g : \mathcal{S}_1 \to \mathcal{S}_2$를
잡자. $\mathcal{C}$의 각 대상 $U$와 $\mathcal{S}_{1, U}$의 각 대상
$x$에 대하여 가정에 의해 $f(x) \cong g(x)$이다. $\mathcal{S}_2$는
세토이드로 섬유화되어 있으므로 $\mathcal{S}_{2, U}$ 안에 유일한
동형사상 $t_x : f(x) \to g(x)$가 존재한다. 규칙 $x \mapsto t_x$가
원하는 $2$-동형사상 $f \to g$를 줌은 명백하다. (2)와 (3)의 증명은
생략한다. 마지막 주장을 보려면 Lemma
\ref{categories-lemma-2-category-fibred-sets}를 사용하여 우변이
$\Mor_{\textit{Cat}/\mathcal{C}}(F(\mathcal{S}_1), F(\mathcal{S}_2))$와
같음을 확인하고 위의 (1)과 (2)를 적용한다.
\end{proof}

\noindent
다음은 군체로 섬유화된 모든 범주 가운데 세토이드로 섬유화된 범주를
특징짓는 또 다른 방법이다.

\begin{lemma}
\label{categories-lemma-characterize-fibred-setoids-inertia}
$\mathcal{C}$를 범주라 하고
$p : \mathcal{S} \to \mathcal{C}$를 군체로 섬유화된 범주라 하자.
다음은 서로 동치이다.
\begin{enumerate}
\item $p : \mathcal{S} \to \mathcal{C}$는 세토이드로 섬유화된
범주이다.
\item 표준 $1$-사상 $\mathcal{I}_\mathcal{S} \to \mathcal{S}$는
(\ref{categories-equation-inertia-structure-map})를 보라,
$\mathcal{C}$ 위의 범주들의 동치이다.
\end{enumerate}
\end{lemma}

\begin{proof}
(2)를 가정하자. 범주 $\mathcal{I}_\mathcal{S}$의 대상은
$(x, \alpha)$이며, 여기서 $x \in \mathcal{S}$이고 이를테면
$p(x) = U$이며, $\alpha : x \to x$는 $\mathcal{S}_U$의 사상이다.
따라서 $\mathcal{I}_\mathcal{S} \to \mathcal{S}$가 $\mathcal{C}$ 위의
동치이면 임의의 두 대상 $(x, \alpha)$와 $(x, \alpha')$는
$U$ 위의 $\mathcal{I}_\mathcal{S}$의 섬유 범주에서 동형이다.
$\mathcal{I}_\mathcal{S}$의 사상 정의를 살펴보면 $\alpha$와
$\alpha'$는 $x$의 자기동형사상군 안에서 서로 켤레임을 알 수 있다.
그러므로 $\alpha' = \text{id}_x$로 놓으면 $x$의 모든 자기동형사상이
항등사상과 같다는 결론을 얻는다. $\mathcal{S} \to \mathcal{C}$는
군체로 섬유화되어 있으므로 이는 $\mathcal{S} \to \mathcal{C}$가
세토이드로 섬유화되어 있음을 뜻한다. (1) $\Rightarrow$ (2)의 증명은
생략한다.
\end{proof}

\begin{lemma}
\label{categories-lemma-category-fibred-setoids-presheaves-products}
$\mathcal{C}$를 범주라 하자. 세토이드로 섬유화된 범주에 준층을
대응시키는 Lemma \ref{categories-lemma-2-category-fibred-setoids}의 구성은 곱과
호환된다. 즉 $2$-섬유곱 $\mathcal{X} \times_\mathcal{Y} \mathcal{Z}$에
대응하는 준층은 $\mathcal{X}, \mathcal{Y}, \mathcal{Z}$에 대응하는
준층들의 섬유곱이다.
\end{lemma}

\begin{proof}
$U \in \Ob(\mathcal{C})$라 하자. 보조정리의 내용은 단지
$$
\Ob((\mathcal{X} \times_\mathcal{Y} \mathcal{Z})_U)/\!\cong
\quad \text{equals} \quad
\Ob(\mathcal{X}_U)/\!\cong
\ \times_{\Ob(\mathcal{Y}_U)/\!\cong}
\ \Ob(\mathcal{Z}_U)/\!\cong
$$
라는 것이며, 그 증명은 생략한다. (단, 범주 $\mathcal{Y}_U$가
세토이드가 아니면 일반적으로 이는 참이 아님에 유의하라.)
\end{proof}









\section{표현 가능한, 군체로 섬유화된 범주들}
\label{categories-section-representable-fibred-groupoids}

\noindent
다음은 표현 가능한, 군체로 섬유화된 범주의 정의이다. 약속한 대로 이
정의는 동치 아래에서 불변이다.

\begin{definition}
\label{categories-definition-representable-fibred-category}
$\mathcal{C}$를 범주라 하자. 군체로 섬유화된 범주
$p : \mathcal{S} \to \mathcal{C}$에 대하여, $\mathcal{C}$의 대상 $X$와
동치 $j : \mathcal{S} \to \mathcal{C}/X$가 존재하면 이 범주를
{\it 표현 가능하다}고 한다. 여기서 동치는 $\mathcal{C}$ 위에서 군체로
섬유화된 범주들의 $2$-범주에서의 동치이다.
\end{definition}

\noindent
통상적인 표기 남용은 동치 $j$를 언급하지 않고
{\it $X$가 $\mathcal{S}$를 표현한다}고 말하는 것이다. 이것이 무엇을
함의하는지 자세히 적어 보자.

\begin{lemma}
\label{categories-lemma-characterize-representable-fibred-category}
$\mathcal{C}$를 범주라 하고
$p : \mathcal{S} \to \mathcal{C}$를 군체로 섬유화된 범주라 하자.
\begin{enumerate}
\item $\mathcal{S}$가 표현 가능할 필요충분조건은 다음 조건들이
성립하는 것이다.
\begin{enumerate}
\item $\mathcal{S}$는 세토이드로 섬유화되어 있다.
\item 준층 $U \mapsto \Ob(\mathcal{S}_U)/\cong$는 표현 가능하다.
\end{enumerate}
\item $\mathcal{S}$가 표현 가능하면, $j$가 동치
$j : \mathcal{S} \to \mathcal{C}/X$인 쌍 $(X, j)$는 동형까지 유일하게
결정된다.
\end{enumerate}
\end{lemma}

\begin{proof}
첫 번째 주장은 Lemma \ref{categories-lemma-setoid-fibres}에서 즉시 따른다.
두 번째 주장을 위해 $j' : \mathcal{S} \to \mathcal{C}/X'$가 같은 성질을
갖는 두 번째 쌍이라고 하자.
$j' \circ t' \cong \text{id}_{\mathcal{C}/X'}$ 및
$t' \circ j' \cong \text{id}_\mathcal{S}$를 만족하는 $1$-사상
$t' : \mathcal{C}/X' \to \mathcal{S}$를 택하자. 그러면
$j \circ t' : \mathcal{C}/X' \to \mathcal{C}/X$는 동치이다. 따라서
Lemma \ref{categories-lemma-2-category-fibred-sets}에 의해 이는 동형사상이다.
그러므로 Yoneda Lemma \ref{categories-lemma-yoneda}에 의해(예를 들어 Example
\ref{categories-example-fibred-category-from-functor-of-points}를 거쳐) 이는
동형사상 $X' \to X$로 주어진다.
\end{proof}

\begin{lemma}
\label{categories-lemma-morphisms-representable-fibred-categories}
$\mathcal{C}$를 범주라 하자. $\mathcal{X}$와 $\mathcal{Y}$를
$\mathcal{C}$ 위에서 군체로 섬유화된 범주들이라 하자.
$\mathcal{X}$와 $\mathcal{Y}$가 각각 $\mathcal{C}$의 대상 $X$와 $Y$로
표현된다고 가정하자. 그러면
$$
\Mor_{\textit{Cat}/\mathcal{C}}(\mathcal{X}, \mathcal{Y})
\Big/
2\text{-isomorphism}
=
\Mor_\mathcal{C}(X, Y)
$$
이다. 더 정확히 말해 $\phi : X \to Y$가 주어지면, 대상들의 동형류
위에서 $\phi$를 유도하며 유일한 $2$-동형사상까지 유일한 $1$-사상
$f : \mathcal{X} \to \mathcal{Y}$가 존재한다.
\end{lemma}

\begin{proof}
Example \ref{categories-example-fibred-category-from-functor-of-points}에 의해
$\mathcal{C}/X = \mathcal{S}_{h_X}$이고
$\mathcal{C}/Y = \mathcal{S}_{h_Y}$이다. Lemma
\ref{categories-lemma-2-category-fibred-setoids}에 의해
$$
\Mor_{\textit{Cat}/\mathcal{C}}(\mathcal{X}, \mathcal{Y})
\Big/
2\text{-isomorphism}
=
\Mor_{\textit{PSh}(\mathcal{C})}(h_X, h_Y)
$$
이다. Yoneda Lemma \ref{categories-lemma-yoneda}에 의해
$\Mor_{\textit{PSh}(\mathcal{C})}(h_X, h_Y)
= \Mor_\mathcal{C}(X, Y)$이다.
\end{proof}









\section{2-Yoneda 보조정리}
\label{categories-section-2-yoneda}

\noindent
$\mathcal{C}$를 범주라 하자. $\mathcal{C}$ 위의 섬유 범주들의
$2$-범주는 Definition \ref{categories-definition-fibred-categories-over-C}에서
구성하고 정의하였다. $\mathcal{S}$와 $\mathcal{S}'$가
$\mathcal{C}$ 위의 섬유 범주들이면
$$
\Mor_{\textit{Fib}/\mathcal{C}}(\mathcal{S}, \mathcal{S}')
$$
는 이 $2$-범주에서 $1$-사상들의 범주를 나타낸다. 다음은 섬유 범주의
맥락에서 Yoneda 보조정리의 $2$-범주 유사물이다.

\begin{lemma}[섬유 범주에 대한 2-Yoneda 보조정리]
\label{categories-lemma-yoneda-2category-fibred}
$\mathcal{C}$를 범주라 하자. $\mathcal{S} \to \mathcal{C}$를
$\mathcal{C}$ 위의 섬유 범주라 하고 $U \in \Ob(\mathcal{C})$라 하자.
함자
$$
\Mor_{\textit{Fib}/\mathcal{C}}(\mathcal{C}/U, \mathcal{S})
\longrightarrow
\mathcal{S}_U
$$
는 $G \mapsto G(\text{id}_U)$로 주어지며 동치이다.
\end{lemma}

\begin{proof}
$\mathcal{S}$에 대한 당김들을 하나씩 선택하자(Definition
\ref{categories-definition-pullback-functor-fibred-category}를 보라). 다음과 같이
함자
$$
\mathcal{S}_U
\longrightarrow
\Mor_{\textit{Fib}/\mathcal{C}}(\mathcal{C}/U, \mathcal{S})
$$
를 정의한다. $x \in \Ob(\mathcal{S}_U)$가 주어졌을 때 이에 수반된
함자는 다음과 같다.
\begin{enumerate}
\item 대상 위에서는 $(f : V \to U) \mapsto f^*x$이다.
\item 사상 위에서는 화살표 $(g : V'/U \to V/U)$를 합성
$$
(f \circ g)^*x \xrightarrow{(\alpha_{g, f})_x} g^*f^*x \rightarrow f^*x
$$
으로 보낸다. 여기서 $\alpha_{g, f}$는 Lemma \ref{categories-lemma-fibred}에서와
같다.
\end{enumerate}
이 함자가 보조정리의 함자에 대한 역임의 확인은 생략한다.
\end{proof}

\noindent
$\mathcal{C}$를 범주라 하자. $\mathcal{C}$ 위에서 군체로 섬유화된
범주들의 $2$-범주는 $\mathcal{C}$ 위의 범주들의 $2$-범주의 ``충만''한
부분 $2$-범주이다(Definition
\ref{categories-definition-categories-fibred-in-groupoids-over-C}를 보라).
따라서 $\mathcal{S}$와 $\mathcal{S}'$가 $\mathcal{C}$ 위에서 군체로
섬유화되어 있으면
$$
\Mor_{\textit{Cat}/\mathcal{C}}(\mathcal{S}, \mathcal{S}')
$$
는 이 $2$-범주에서 $1$-사상들의 범주를 나타낸다(Definition
\ref{categories-definition-categories-over-C}를 보라). 이 범주들은 모두 군체이다.
Definition \ref{categories-definition-categories-fibred-in-groupoids-over-C} 뒤의
논의를 보라. 다음은 Yoneda 보조정리의 $2$-범주 유사물이다.

\begin{lemma}[2-Yoneda 보조정리]
\label{categories-lemma-yoneda-2category}
$\mathcal{S}\to \mathcal{C}$를 군체로 섬유화된 범주라 하고
$U \in \Ob(\mathcal{C})$라 하자. 함자
$$
\Mor_{\textit{Cat}/\mathcal{C}}(\mathcal{C}/U, \mathcal{S})
\longrightarrow
\mathcal{S}_U
$$
는 $G \mapsto G(\text{id}_U)$로 주어지며 동치이다.
\end{lemma}

\begin{proof}
$\mathcal{S}$에 대한 당김들을 하나씩 선택하자(Definition
\ref{categories-definition-pullback-functor-fibred-category}를 보라). 다음과 같이
함자
$$
\mathcal{S}_U
\longrightarrow
\Mor_{\textit{Cat}/\mathcal{C}}(\mathcal{C}/U, \mathcal{S})
$$
를 정의한다. $x \in \Ob(\mathcal{S}_U)$가 주어졌을 때 이에 수반된
함자는 다음과 같다.
\begin{enumerate}
\item 대상 위에서는 $(f : V \to U) \mapsto f^*x$이다.
\item 사상 위에서는 화살표 $(g : V'/U \to V/U)$를 합성
$$
(f \circ g)^*x \xrightarrow{(\alpha_{g, f})_x} g^*f^*x \rightarrow f^*x
$$
으로 보낸다. 여기서 $\alpha_{g, f}$는 Lemma
\ref{categories-lemma-fibred-groupoids}에서와 같다.
\end{enumerate}
이 함자가 보조정리의 함자에 대한 역임의 확인은 생략한다.
\end{proof}

\begin{remark}
\label{categories-remark-alternative-fibred-groupoids-strict}
$2$-Yoneda 보조정리를 사용하여 Lemma
\ref{categories-lemma-fibred-groupoids-strict}의 다른 증명을 줄 수 있다.
$p : \mathcal{S} \to \mathcal{C}$를 군체로 섬유화된 범주라 하자.
다음과 같이 $\mathcal{C}$에서 군체들의 범주로 가는 반변함자 $F$를
정의한다. $U\in \Ob(\mathcal{C})$에 대하여
$$
F(U) = \Mor_{\textit{Cat}/\mathcal{C}}(\mathcal{C}/U, \mathcal{S}).
$$
로 놓는다. $f : U \to V$이면 유도된 함자
$\mathcal{C}/U \to \mathcal{C}/V$는 사상 $F(f) : F(V) \to F(U)$를
유도한다. $F$가 함자임은 명백하다. Example
\ref{categories-example-functor-groupoids}로부터 수반되는 군체로 섬유화된 범주를
$\mathcal{S}'$라 하자. $\mathcal{C}$ 위의 자명한 함자
$G : \mathcal{S}' \to \mathcal{S}$가 있다. 이는 $U \in \Ob(\mathcal{C})$이고
$x \in F(U)$인 쌍 $(U, x)$를 $x(\text{id}_U) \in \mathcal{S}$로 보낸다.
이제 Lemma \ref{categories-lemma-yoneda-2category}에 의해 각 $U$에 대하여
$$
G_U : \mathcal{S}'_U = F(U)=
\Mor_{\textit{Cat}/\mathcal{C}}(\mathcal{C}/U, \mathcal{S})
\to
\mathcal{S}_U
$$
는 동치이다. 그러므로 Lemma
\ref{categories-lemma-equivalence-fibred-categories}에 의해 $G$는
$\mathcal{S}$와 $\mathcal{S}'$ 사이의 동치이다.
\end{remark}





\section{표현 가능한 1-사상}
\label{categories-section-representable-1-morphisms}

\noindent
$\mathcal{C}$를 범주라 하자. 이 절에서는 $\mathcal{C}$ 위에서 군체로
섬유화된 범주들 사이의 $1$-사상이 표현 가능하다는 것이 무엇을
뜻하는지 설명한다.

\medskip\noindent
$\mathcal{C}$를 범주라 하고 $\mathcal{X}$와 $\mathcal{Y}$를
$\mathcal{C}$ 위에서 군체로 섬유화된 범주들이라 하자.
$U \in \Ob(\mathcal{C})$라 하자. $F : \mathcal{X} \to \mathcal{Y}$와
$G : \mathcal{C}/U \to \mathcal{Y}$를 $\mathcal{C}$ 위에서 군체로
섬유화된 범주들의 $1$-사상이라 하자. $2$-섬유곱
$$
\xymatrix{
(\mathcal{C}/U) \times_\mathcal{Y} \mathcal{X} \ar[r] \ar[d] &
\mathcal{X} \ar[d]^F \\
\mathcal{C}/U \ar[r]^G &
\mathcal{Y}
}
$$
을 기술하고자 한다. $y = G(\text{id}_U) \in \mathcal{Y}_U$라 하자.
$\mathcal{Y}$에 대한 당김들을 하나씩 선택하자(Definition
\ref{categories-definition-pullback-functor-fibred-category}를 보라). 그러면
Lemma \ref{categories-lemma-yoneda-2category}와 그 증명에 의해 $G$는 함자
$(f : V \to U) \mapsto f^*y$와 동형이다. Lemma
\ref{categories-lemma-2-product-categories-over-C}에 따라
$(\mathcal{C}/U) \times_\mathcal{Y} \mathcal{X}$의 대상을 네쌍
$(V, f : V \to U, x, \phi)$으로 생각할 수 있다. 위의 $G$에 대한
기술을 사용하면 $\phi$를 $\mathcal{Y}_V$ 안의 동형사상
$\phi : f^*y \to F(x)$로 생각할 수 있다.

\begin{lemma}
\label{categories-lemma-identify-fibre-product}
위의 상황에서 $\mathcal{C}/U$의 대상 $f : V \to U$ 위에 놓인
$(\mathcal{C}/U) \times_\mathcal{Y} \mathcal{X}$의 섬유 범주는 다음과
같이 기술된다.
\begin{enumerate}
\item 대상은 쌍 $(x, \phi)$이며, 여기서
$x \in \Ob(\mathcal{X}_V)$이고
$\phi : f^*y \to F(x)$는 $\mathcal{Y}_V$의 사상이다.
\item $(x, \phi)$에서 $(x', \phi')$로 가는 사상들의 집합은
$F(\psi) = \phi' \circ \phi^{-1}$를 만족하는 $\mathcal{X}_V$의 사상
$\psi : x \to x'$들의 집합이다.
\end{enumerate}
\end{lemma}

\begin{proof}
위의 논의를 보라.
\end{proof}

\begin{lemma}
\label{categories-lemma-prepare-representable-map-stack-in-groupoids}
$\mathcal{C}$를 범주라 하고 $\mathcal{X}$와 $\mathcal{Y}$를
$\mathcal{C}$ 위에서 군체로 섬유화된 범주들이라 하자.
$F : \mathcal{X} \to \mathcal{Y}$를 $1$-사상이라 하고
$G : \mathcal{C}/U \to \mathcal{Y}$를 $1$-사상이라 하자. 그러면
$$
(\mathcal{C}/U) \times_\mathcal{Y} \mathcal{X}
\longrightarrow
\mathcal{C}/U
$$
는 군체로 섬유화된 범주이다.
\end{lemma}

\begin{proof}
합성
$$
(\mathcal{C}/U) \times_\mathcal{Y} \mathcal{X}
\longrightarrow
\mathcal{C}/U
\longrightarrow
\mathcal{C}
$$
이 군체로 섬유화된 범주임을 Lemma
\ref{categories-lemma-2-product-fibred-categories}에서 이미 보았다. 이제 보조정리는
Lemma \ref{categories-lemma-cute-groupoids}에서 따른다.
\end{proof}

\begin{definition}
\label{categories-definition-representable-map-categories-fibred-in-groupoids}
$\mathcal{C}$를 범주라 하고 $\mathcal{X}$와 $\mathcal{Y}$를
$\mathcal{C}$ 위에서 군체로 섬유화된 범주들이라 하자.
$F : \mathcal{X} \to \mathcal{Y}$를 $1$-사상이라 하자. 모든
$U \in \Ob(\mathcal{C})$와 임의의 $G : \mathcal{C}/U \to \mathcal{Y}$에
대하여 군체로 섬유화된 범주
$$
(\mathcal{C}/U) \times_\mathcal{Y} \mathcal{X}
\longrightarrow
\mathcal{C}/U
$$
가 표현 가능하면 $F$를 {\it 표현 가능하다}고 하며, 또는
{\it $\mathcal{X}$가 $\mathcal{Y}$ 위에서 상대적으로 표현 가능하다}고
한다.
\end{definition}

\begin{lemma}
\label{categories-lemma-spell-out-representable-map-stack-in-groupoids}
$\mathcal{C}$를 범주라 하고 $\mathcal{X}$와 $\mathcal{Y}$를
$\mathcal{C}$ 위에서 군체로 섬유화된 범주들이라 하자.
$F : \mathcal{X} \to \mathcal{Y}$를 $1$-사상이라 하자. $F$가 표현
가능하면 섬유 범주들 사이의 모든 함자
$$
F_U : \mathcal{X}_U \longrightarrow \mathcal{Y}_U
$$
는 충실하다.
\end{lemma}

\begin{proof}
Lemma \ref{categories-lemma-identify-fibre-product}의 섬유 범주에 대한 기술과 Lemma
\ref{categories-lemma-characterize-representable-fibred-category}의 표현 가능한 섬유
범주에 대한 특징짓기에서 명백하다.
\end{proof}

\begin{lemma}
\label{categories-lemma-criterion-representable-map-stack-in-groupoids}
$\mathcal{C}$를 범주라 하고 $\mathcal{X}$와 $\mathcal{Y}$를
$\mathcal{C}$ 위에서 군체로 섬유화된 범주들이라 하자.
$F : \mathcal{X} \to \mathcal{Y}$를 $1$-사상이라 하자.
$\mathcal{Y}$에 대한 당김들을 하나씩 선택하자. 다음을 가정하자.
\begin{enumerate}
\item 섬유 범주들 사이의 각 함자
$F_U : \mathcal{X}_U \longrightarrow \mathcal{Y}_U$는 충실하다.
\item 각 $U$와 각 $y \in \mathcal{Y}_U$에 대하여 준층
$$
(f : V \to U)
\longmapsto
\{(x, \phi) \mid x \in \mathcal{X}_V, \phi : f^*y \to F(x)\}/\cong
$$
은 $\mathcal{C}/U$ 위의 표현 가능한 준층이다.
\end{enumerate}
그러면 $F$는 표현 가능하다.
\end{lemma}

\begin{proof}
Lemma \ref{categories-lemma-identify-fibre-product}의 섬유 범주에 대한 기술과 Lemma
\ref{categories-lemma-characterize-representable-fibred-category}의 표현 가능한 섬유
범주에 대한 특징짓기에서 명백하다.
\end{proof}

\noindent
다음 보조정리를 진술하기 전에, 군체로 섬유화된 범주들의 $2$-범주는
$(2, 1)$-범주이므로 이 범주가 종대상을 갖는다는 것이 무엇을 뜻하는지
알고 있음에 유의하자(Definition
\ref{categories-definition-final-object-2-category}를 보라). 그리고 이 범주는
종대상 $\text{id} : \mathcal{C} \to \mathcal{C}$를 갖는다. 따라서
$\mathcal{C}$ 위에서 군체로 섬유화된 범주들의 {\it $2$-곱}을
$2$-섬유곱
$$
\mathcal{X} \times \mathcal{Y} :=
\mathcal{X} \times_\mathcal{C} \mathcal{Y}.
$$
으로 정의한다. 이 정의 아래 다음 보조정리는 의미가 있다.

\begin{lemma}
\label{categories-lemma-representable-diagonal-groupoids}
$\mathcal{C}$를 범주라 하고
$\mathcal{S} \to \mathcal{C}$를 군체로 섬유화된 범주라 하자.
$\mathcal{C}$가 대상 쌍들의 곱과 섬유곱을 갖는다고 가정하자. 다음은
서로 동치이다.
\begin{enumerate}
\item 대각사상 $\mathcal{S} \to \mathcal{S} \times \mathcal{S}$는 표현
가능하다.
\item $\mathcal{C}$의 모든 $U$에 대하여 임의의
$G : \mathcal{C}/U \to \mathcal{S}$는 표현 가능하다.
\end{enumerate}
\end{lemma}

\begin{proof}
대각사상이 표현 가능하다고 가정하고 $U, G$가 주어졌다고 하자.
임의의 $V \in \Ob(\mathcal{C})$와 임의의
$G' : \mathcal{C}/V \to \mathcal{S}$를 생각하자.
$\mathcal{C}/U \times \mathcal{C}/V = \mathcal{C}/U \times V$는 표현
가능함에 유의하라. 따라서 섬유곱
$$
\xymatrix{
(\mathcal{C}/U \times V)
\times_{(\mathcal{S} \times \mathcal{S})}
\mathcal{S}
\ar[r] \ar[d] &
\mathcal{S} \ar[d] \\
\mathcal{C}/U \times V \ar[r]^{(G, G')} &
\mathcal{S} \times \mathcal{S}
}
$$
은 가정에 의해 표현 가능하다. 이는 $\mathcal{C}$ 안에
$W \to U \times V$가 존재하여
$$
\xymatrix{
\mathcal{C}/W \ar[d] \ar[r] & \mathcal{S} \ar[d] \\
\mathcal{C}/U \times \mathcal{C}/V \ar[r] & \mathcal{S} \times \mathcal{S}
}
$$
가 카르테시안임을 뜻한다. 이는 원하는 대로
$\mathcal{C}/W \cong \mathcal{C}/U \times_\mathcal{S} \mathcal{C}/V$임을
함의한다(Lemma \ref{categories-lemma-diagonal-1}과 Remark
\ref{categories-remark-2-product-fibred-categories-same}을 보라).

\medskip\noindent
(2)가 성립한다고 가정하자. 임의의 $V \in \Ob(\mathcal{C})$와 임의의
$(G, G') : \mathcal{C}/V \to \mathcal{S} \times \mathcal{S}$를 생각하자.
$\mathcal{C}/V \times_{\mathcal{S} \times \mathcal{S}} \mathcal{S}$가 표현
가능함을 보여야 한다. 우리가 아는 것은
$\mathcal{C}/V \times_{G, \mathcal{S}, G'} \mathcal{C}/V$가 표현
가능하며, 이를테면 $\mathcal{C}/V$ 안의 $a : W \to V$에 의해
표현된다는 것이다. 동치
$$
\mathcal{C}/W \to \mathcal{C}/V \times_{G, \mathcal{S}, G'} \mathcal{C}/V
$$
를 $\mathcal{C}/V$로 가는 두 번째 사영과 합성하면 두 번째 사상
$a' : W \to V$를 얻는다.
$W' = W \times_{(a, a'), V \times V} V$를 생각하자. 동치
$$
\mathcal{C}/W' \cong
\mathcal{C}/V \times_{\mathcal{S} \times \mathcal{S}} \mathcal{S}
$$
가 존재한다. 구체적으로
\begin{eqnarray*}
\mathcal{C}/W' & \cong &
\mathcal{C}/W \times_{(\mathcal{C}/V \times \mathcal{C}/V)} \mathcal{C}/V \\
& \cong &
\left(\mathcal{C}/V \times_{(G, \mathcal{S}, G')} \mathcal{C}/V\right)
\times_{(\mathcal{C}/V \times \mathcal{C}/V)} \mathcal{C}/V \\
& \cong &
\mathcal{C}/V \times_{(\mathcal{S} \times \mathcal{S})} \mathcal{S}
\end{eqnarray*}
이다. 마지막 동형사상에 대해서는 Lemma \ref{categories-lemma-diagonal-2}와
Remark \ref{categories-remark-2-product-fibred-categories-same}을 보라. 이로써
보조정리가 증명된다.
\end{proof}

\noindent
{\bf 문헌 주석:}
이 내용의 일부는 Vistoli의 노트 \cite{Vis2}에서 가져왔다.





\section{모노이달 범주}
\label{categories-section-monoidal}

\noindent
$\mathcal{C}$를 범주라 하자. 함자
$$
\otimes : \mathcal{C} \times \mathcal{C} \longrightarrow \mathcal{C}
$$
가 주어졌다고 하자. 흔히 $\otimes$가 결합법칙을 만족하는지, 그리고
$\otimes$에 대한 단위원이 있는지 알고자 한다.

\medskip\noindent
$(\mathcal{C}, \otimes)$에 대한 {\it 결합성 제약}이란 함자적 동형사상
$$
\phi_{X, Y, Z} : X \otimes (Y \otimes Z) \to (X \otimes Y) \otimes Z
$$
으로서 모든 대상 $X, Y, Z, W$에 대하여 도표
\begin{equation}
\label{categories-equation-pentagon}
\vcenter{
\xymatrix{
X \otimes (Y \otimes ( Z \otimes W)) \ar[r] \ar[d] &
X \otimes ((Y \otimes Z) \otimes W) \ar[d] \\
(X \otimes Y) \otimes (Z \otimes W) \ar[d] &
(X \otimes (Y \otimes Z)) \otimes W \ar[ld] \\
((X \otimes Y) \otimes Z) \otimes W
}
}
\end{equation}
를 가환하게 하는 것이다. 여기서 각 화살표는 $\phi$의 적절한 적용과
$\otimes$의 함자성으로 결정된다.

\medskip\noindent
% P01-E1140: frozen source says "as a above"; canonical target retains only a non-rendering pointer.
\cite[Theorem 3.1]{associativity}에서는 위와 같은 세쌍
$(\mathcal{C}, \otimes, \phi)$가 어떻게 정합적인지 설명한다. 이를 다시
말하면, 이러한 세쌍은 모든 $n \geq 1$에 대하여 잘 정의된 함자들의 체계
$$
\mathcal{C} \times \ldots \times \mathcal{C} \longrightarrow \mathcal{C},
\quad
(X_1, \ldots, X_n) \longmapsto X_1 \otimes \ldots \otimes X_n
$$
와, $n, m \geq 1$에 대하여 함자적 동형사상
$$
(X_1 \otimes \ldots \otimes X_n) \otimes (Y_1 \otimes \ldots \otimes Y_m)
\to
X_1 \otimes \ldots \otimes X_n \otimes Y_1 \otimes \ldots \otimes Y_m
$$
을 준다. 이 동형사상들로 만들 수 있는 모든 도표는 가환한다(위의
오각별 도표가 한 예이다). $n = 5$일 때 이는 식들
$$
\begin{matrix}
(((X_1 \otimes X_2) \otimes X_3) \otimes X_4) \otimes X_5 &
((X_1 \otimes (X_2 \otimes X_3)) \otimes X_4) \otimes X_5 \\
((X_1 \otimes X_2) \otimes (X_3 \otimes X_4)) \otimes X_5 &
(X_1 \otimes ((X_2 \otimes X_3) \otimes X_4)) \otimes X_5 \\
(X_1 \otimes (X_2 \otimes (X_3 \otimes X_4))) \otimes X_5 &
((X_1 \otimes X_2) \otimes X_3) \otimes (X_4 \otimes X_5) \\
(X_1 \otimes (X_2 \otimes X_3)) \otimes (X_4 \otimes X_5) &
(X_1 \otimes X_2) \otimes ((X_3 \otimes X_4) \otimes X_5) \\
(X_1 \otimes X_2) \otimes (X_3 \otimes (X_4 \otimes X_5)) &
X_1 \otimes (((X_2 \otimes X_3) \otimes X_4) \otimes X_5) \\
X_1 \otimes ((X_2 \otimes (X_3 \otimes X_4)) \otimes X_5) &
X_1 \otimes ((X_2 \otimes X_3) \otimes (X_4 \otimes X_5)) \\
X_1 \otimes (X_2 \otimes ((X_3 \otimes X_4) \otimes X_5)) &
X_1 \otimes (X_2 \otimes (X_3 \otimes (X_4 \otimes X_5)))
\end{matrix}
$$
이 모두 $X_1 \otimes \ldots \otimes X_5$와 함자적으로 동형이며, 이
동형사상들이 세 개의 연속한 자리에서 가능한 경우 $\phi$를 사용해 이
함자들 사이에서 얻는 모든 가능한 동형사상과 호환된다는 뜻이다. 아래에서
이를 더 언급하지 않고 사용하겠으며, $\otimes$의 반복을 쓸 때 더는
괄호를 쓰지 않겠다.

\medskip\noindent
위와 같은 세쌍 $(\mathcal{C}, \otimes, \phi)$의 {\it 단위원}이란
$\mathcal{C}$의 대상 $\mathbf{1}$과 함자적 동형사상
$$
l : \mathbf{1} \otimes X \to X
\quad\text{and}\quad
r : X \otimes \mathbf{1} \to X
$$
으로서 모든 대상 $X, Y$에 대하여 도표
\begin{equation}
\label{categories-equation-triangle}
\vcenter{
\xymatrix{
X \otimes (\mathbf{1} \otimes Y) \ar[rr]_\phi \ar[rd]_{\text{id} \otimes l} & &
(X \otimes \mathbf{1}) \otimes Y \ar[ld]^{r \otimes \text{id}} \\
& X \otimes Y
}
}
\end{equation}
를 가환하게 하는 것이다. 다음 보조정리에서처럼 단위원을 흔히 쌍
$(\mathbf{1}, 1)$로 생각하겠다.

\begin{lemma}
\label{categories-lemma-monoidal-unit}
$(\mathcal{C}, \otimes, \phi)$가 위와 같다고 하자. $\mathcal{C}$ 안의
단위원 $(\mathbf{1}, l, r)$들과 다음과 같은 쌍 $(\mathbf{1}, 1)$들
사이에는 일대일 대응이 있다. 여기서 $\mathbf{1}$은 $\mathcal{C}$의
대상이고 $1 : \mathbf{1} \otimes \mathbf{1} \to \mathbf{1}$은
동형사상이며, 함자 $L : X \mapsto \mathbf{1} \otimes X$와
$R : X \mapsto X \otimes \mathbf{1}$은 동치이다.
\end{lemma}

\begin{proof}
% P01-E0014: frozen source misspells "isomorphism"; pointer only.
단위원 $(\mathbf{1}, l, r)$이 주어지면 동형사상
$r : \mathbf{1} \otimes \mathbf{1} \to \mathbf{1}$을 얻으며, $L$과 $R$은
항등함자와 동형이므로 동치이다. 역으로 $L$과 $R$이 동치인 쌍
$(\mathbf{1}, 1)$이 주어졌다고 하자.
$L(l_X) = 1 \otimes \text{id}_X$와
$R(r_X) = \text{id}_X \otimes 1$로 특징지어지는 함자적 동형사상
$l_X : \mathbf{1} \otimes X \to X$와
$r_X : X \otimes \mathbf{1} \to X$를 얻는다. 이제 모든 $X$와 $Y$에
대하여 $X \otimes \mathbf{1} \otimes Y$에서 $X \otimes Y$로 가는
두 화살표 $\text{id}_X \otimes l_Y$와
$r_X \otimes \text{id}_Y$가 같음을 보여야 한다. 이 성질은 $X$와 $Y$의
동형류에만 의존한다. $R$과 $L$은 동치이므로, 어떤 대상 $Z$와 $W$에
대하여 $X = Z \otimes \mathbf{1}$이고
$Y = \mathbf{1} \otimes W$인 경우만 보이면 충분하다. 다시 말해
% P01-E0012: frozen source's final identity factor is under errata review; formula retained exactly.
$$
\text{id}_Z \otimes \text{id}_\mathbf{1} \otimes l_{\mathbf{1} \otimes W}
=
r_{Z \otimes \mathbf{1}} \otimes \text{id}_\mathbf{1} \otimes \text{id}_Y
$$
임을 보여야 한다. 구성에 의해 이 두 사상은 각각
$\text{id}_Z \otimes 1 \otimes \text{id}_\mathbf{1} \otimes \text{id}_W$와
$\text{id}_Z \otimes \text{id}_\mathbf{1} \otimes 1 \otimes \text{id}_W$이다.
따라서 $1 \otimes \text{id}_\mathbf{1} =
\text{id}_\mathbf{1} \otimes 1$임을 보이면 충분하다.

\medskip\noindent
다음이 성립한다.
$\text{id}_\mathbf{1} \otimes 1 = r_\mathbf{1} \otimes \text{id}_\mathbf{1}$
이고
$l_\mathbf{1} \otimes \text{id}_\mathbf{1} = \text{id}_\mathbf{1} \otimes 1$이다.
$\mathbf{1}$의 어떤 자기동형사상 $a, b$에 대하여
$r_\mathbf{1} = a \circ 1$ 및 $l_\mathbf{1} = b \circ 1$로 쓸 수 있다.
따라서
$\text{id}_\mathbf{1} \otimes 1 =
(a \otimes \text{id}_\mathbf{1}) \circ (1 \otimes \text{id}_\mathbf{1})$
이고
$1 \otimes \text{id}_\mathbf{1} =
(\text{id}_\mathbf{1} \otimes b) \circ (\text{id}_\mathbf{1} \otimes 1)$이다.
이제 다음과 같이 쓸 수 있다.
\begin{align*}
(1 \otimes \text{id}_\mathbf{1}) \circ
(\text{id}_\mathbf{1} \otimes 1 \otimes \text{id}_\mathbf{1})
& =
(1 \otimes \text{id}_\mathbf{1}) \circ
(\text{id}_\mathbf{1} \otimes \text{id}_\mathbf{1} \otimes b)
\circ
(\text{id}_\mathbf{1} \otimes \text{id}_\mathbf{1} \otimes 1) \\
& =
(\text{id}_\mathbf{1} \otimes b) \circ
(1 \otimes \text{id}_\mathbf{1}) \circ
(\text{id}_\mathbf{1} \otimes \text{id}_\mathbf{1} \otimes 1) \\
& =
(\text{id}_\mathbf{1} \otimes b) \circ
(1 \otimes 1)
\end{align*}
또한 다음도 성립한다.
\begin{align*}
(1 \otimes \text{id}_\mathbf{1}) \circ
(\text{id}_\mathbf{1} \otimes 1 \otimes \text{id}_\mathbf{1})
& =
(\text{id}_\mathbf{1} \otimes b) \circ
(\text{id}_\mathbf{1} \otimes 1) \circ
(a \otimes \text{id}_\mathbf{1} \otimes \text{id}_\mathbf{1}) \circ
(1 \otimes \text{id}_\mathbf{1} \otimes \text{id}_\mathbf{1}) \\
& =
(\text{id}_\mathbf{1} \otimes b) \circ
(a \otimes \text{id}_\mathbf{1}) \circ
(1 \otimes 1)
\end{align*}
이는 $a \otimes \text{id}_\mathbf{1}$이 항등사상임을 증명하며, 따라서
원하는 대로 $a$가 항등사상이다.

\medskip\noindent
증명을 끝내기 위해, 위 규칙들이 단위원들의 범주(적절히 정의된)와 쌍
$(\mathbf{1}, 1)$들의 범주 사이에서 서로 역인 범주의 동치들을
결정함에 유의하자.
\end{proof}

\begin{lemma}
\label{categories-lemma-monoidal-unit-unique}
$(\mathcal{C}, \otimes, \phi)$가 위와 같다고 하자.
$(\mathbf{1}, 1)$을 단위원이라 하자(Lemma
\ref{categories-lemma-monoidal-unit}를 보라). 그러면
\begin{enumerate}
\item $1 \otimes \text{id}_\mathbf{1} = \text{id}_\mathbf{1} \otimes 1$이다.
\item $\Gamma = \text{Mor}(\mathbf{1}, \mathbf{1})$는 가환 모노이드이다.
\item 모든 $a \in \Gamma$에 대하여
$a = 1 \circ (a \otimes \text{id}_\mathbf{1}) \circ 1^{-1} =
1 \circ (\text{id}_\mathbf{1} \otimes a) \circ 1^{-1}$이다.
\item 다른 임의의 단위원은 유일한 동형사상에 의해
$(\mathbf{1}, 1)$과 동형이다.
\end{enumerate}
\end{lemma}

\begin{proof}
(1)은 Lemma \ref{categories-lemma-monoidal-unit}의 증명에서 보였다.
$a \in \Gamma$에 대하여
\begin{align*}
(1 \circ (a \otimes \text{id}_\mathbf{1})) \otimes \text{id}_\mathbf{1}
& =
(1 \otimes \text{id}_\mathbf{1}) \circ
(a \otimes \text{id}_\mathbf{1} \otimes \text{id}_\mathbf{1}) \\
& =
(\text{id}_\mathbf{1} \otimes 1) \circ
(a \otimes \text{id}_\mathbf{1} \otimes \text{id}_\mathbf{1}) \\
& =
(a \otimes \text{id}_\mathbf{1}) \circ
(\text{id}_\mathbf{1} \otimes 1) \\
& =
(a \otimes \text{id}_\mathbf{1}) \circ
(1 \otimes \text{id}_\mathbf{1}) \\
& =
(a \circ 1) \otimes \text{id}_\mathbf{1}
\end{align*}
이다. 따라서 $1 \circ (a \otimes \text{id}_\mathbf{1}) = a \circ 1$이며,
이는 (3)의 절반을 함의한다. 다른 절반도 같은 방식으로 따른다. (2)를
보기 위해 $a, b \in \Gamma$에 대하여
% P01-E0013/P01-E1141: duplicate frozen composition operator retained exactly.
$$
a \circ b =
1 \circ (a \otimes \text{id}_\mathbf{1}) \circ 1^{-1}
\circ
1 \circ (\text{id}_\mathbf{1} \otimes b) \circ 1^{-1}
=
1 \circ (a \otimes b) \circ \circ 1^{-1}
$$
임을 관찰한다. $b \circ a$도 같은 식으로 계산된다. 따라서 (2)가
성립한다. (4)를 보기 위해 $(\mathbf{1}', 1')$이 두 번째 단위원이라고
하자. $r$과 $l'$를 사용하면 동형사상
$\mathbf{1}' \otimes \mathbf{1} \to \mathbf{1}'$ 및
$\mathbf{1}' \otimes \mathbf{1} \to \mathbf{1}$이 존재한다. 따라서
동형사상 $t : \mathbf{1}' \to \mathbf{1}$이 존재한다. 이제 도표
$$
\xymatrix{
\mathbf{1}' \otimes \mathbf{1}' \ar[r]_{t \otimes t} \ar[d]_{1'} & 
\mathbf{1} \otimes \mathbf{1} \ar[d]^1 \\
\mathbf{1}' \ar[r]^t & \mathbf{1}
}
$$
는 $\mathbf{1}$의 한 자기동형사상 $a$까지 가환한다. $t$를
$a \circ t$로 바꾸면 도표가 가환한다. 힌트: (3)을 사용하여
$1 \circ (a \otimes a) = a^2 \circ 1$임을 보라.
\end{proof}

\begin{definition}
\label{categories-definition-monoidal-category}
$\mathcal{C}$가 범주이고
$\otimes : \mathcal{C} \times \mathcal{C} \to \mathcal{C}$가 함자이며
$\phi$가 결합성 제약인 세쌍 $(\mathcal{C}, \otimes, \phi)$에 단위원
$\mathbf{1}$이 존재하면 이를 {\it 모노이달 범주}라고 한다.
\end{definition}

\noindent
모노이달 범주의 단위원은 항상 $\mathbf{1}$로 쓰고, 선택한 동형사상은
$1 : \mathbf{1} \otimes \mathbf{1} \to \mathbf{1}$로 나타낸다. 쌍
$(\mathbf{1}, 1)$은 유일한 동형사상까지 결정되므로(Lemma
\ref{categories-lemma-monoidal-unit-unique}) 하나를 선택해도 문제가 없다.

\begin{lemma}
\label{categories-lemma-monoidal-coherence}
모노이달 범주 $\mathcal{C}, \otimes, \phi, \mathbf{1}, 1$에서 Lemma
\ref{categories-lemma-monoidal-unit}의 증명과 같은 표기를 쓰면 다음이 성립한다.
\begin{enumerate}
\item 화살표 $1, r_\mathbf{1}, l_\mathbf{1} :
\mathbf{1} \otimes \mathbf{1} \to \mathbf{1}$은 서로 같다.
\item 화살표
$l_X \otimes \text{id}_Y, l_{X \otimes Y} :
\mathbf{1} \otimes X \otimes Y \to X \otimes Y$는 서로 같다.
\item 화살표
$\text{id}_X \otimes r_Y , r_{X \otimes Y} :
X \otimes Y \otimes \mathbf{1} \to X \otimes Y$는 서로 같다.
\end{enumerate}
모노이달 범주는 \cite[Theorem 5.2]{associativity}의 가정들을 만족한다.
\end{lemma}

\begin{proof}
(1)은 Lemma \ref{categories-lemma-monoidal-unit}의 증명에서 보았다. Lemma
\ref{categories-lemma-monoidal-unit}의 증명에서 또한
$l_X$와 $l_{X \otimes Y}$가
$\text{id}_\mathbf{1} \otimes l_{X \otimes Y} =
1 \otimes \text{id}_{X \otimes Y}$ 및
$\text{id}_\mathbf{1} \otimes l_X = 1 \otimes \text{id}_X$를 만족하는
유일한 사상들임을 보았다. (2)는 즉시 따른다. (3)도 비슷하게 증명한다.
(\ref{categories-equation-pentagon})과 (\ref{categories-equation-triangle})의 가환성과 함께 이는
보조정리의 마지막 주장이 성립함을 뜻한다.
\end{proof}

\noindent
\cite[Theorem 5.2]{associativity}에서는 위 보조정리의 다섯쌍
$(\mathcal{C}, \otimes, \phi, \mathbf{1}, 1)$이 어떻게 정합적인지
설명한다. 이를 다시 말하면, 우리 의미의 모노이달 범주에서 모든
$n \geq i \geq 0$에 대하여 함자적 동형사상
$$
X_1 \otimes \ldots \otimes X_i \otimes \mathbf{1}
\otimes X_{i + 1} \otimes \ldots \otimes X_n
\to
X_1 \otimes \ldots \otimes X_n
$$
이 있고, 이 동형사상들로 만들 수 있는 모든 가능한 도표가 가환한다.
예를 들어 $\mathbf{1}$을 두 번 삽입한 것에서 시작하여 이 동형사상들을
서로 다른 순서로 사용해도 같은 사상을 얻는다. $\otimes$의 반복을 쓸 때
괄호를 생략한다는 규약에 더하여, 이제부터는 더 언급하지 않고
$X \otimes \mathbf{1}$과 $\mathbf{1} \otimes X$를 $X$와 식별한다.
또한 ``$\mathcal{C}$를 모노이달 범주라 하자''라고 말할 때
$\otimes, \phi, \mathbf{1}$은 주어진 것으로 이해한다.

\begin{definition}
\label{categories-definition-functor-monoidal-categories}
$\mathcal{C}$와 $\mathcal{C}'$를 모노이달 범주라 하자.
{\it 모노이달 범주의 함자} $F : \mathcal{C} \to \mathcal{C}'$는 표시한
함자 $F$와 $X$ 및 $Y$에 함자적인 동형사상
$$
F(X) \otimes F(Y) \to F(X \otimes Y)
$$
으로 주어지며, 모든 대상 $X$, $Y$, $Z$에 대하여 도표
$$
\xymatrix{
F(X) \otimes (F(Y) \otimes F(Z)) \ar[r] \ar[d] &
F(X) \otimes F(Y \otimes Z) \ar[r] &
F(X \otimes (Y \otimes Z)) \ar[d] \\
(F(X) \otimes F(Y)) \otimes F(Z) \ar[r] &
F(X \otimes Y) \otimes F(Z) \ar[r] &
F((X \otimes Y) \otimes Z)
}
$$
를 가환하게 하고 $F(\mathbf{1})$이 $\mathcal{C}'$의 단위원이 되게 한다.
\end{definition}

\noindent
단위원에 관한 우리의 규약에 의해 $F$가 모노이달 범주의 함자이면 항상
$F(\mathbf{1}) = \mathbf{1}$이라고 가정할 수 있다. 예를 들어
$A \to B$가 환 준동형이면 함자 $M \mapsto M \otimes_A B$는
% P01-E1142: frozen source omits an article before "functor"; pointer only.
$\text{Mod}_A$에서 $\text{Mod}_B$로 가는 모노이달 범주의 함자이다.

\begin{lemma}
\label{categories-lemma-invertible}
$\mathcal{C}$를 모노이달 범주라 하고 $X$를 $\mathcal{C}$의 대상이라
하자. 다음은 서로 동치이다.
\begin{enumerate}
\item 함자 $L : Y \mapsto X \otimes Y$는 동치이다.
\item 함자 $R : Y \mapsto Y \otimes X$는 동치이다.
\item $X \otimes X' \cong X' \otimes X \cong \mathbf{1}$을 만족하는
대상 $X'$이 존재한다.
\end{enumerate}
\end{lemma}

\begin{proof}
(1)을 가정하자. $L(X') = \mathbf{1}$, 즉
$X \otimes X' \cong \mathbf{1}$을 만족하는 $X'$을 택하자. $X'$에
대응하는 함자들을 $L'$과 $R'$로 나타내자. 등식
$X \otimes X' \cong \mathbf{1}$은
$L \circ L' \cong \text{id}$를 함의한다. 따라서 $L'$은 $L$의
준역이어야 한다. 이는 가정에 의해 존재한다. 그러므로
$L' \circ L \cong \text{id}$이고, 따라서
$X' \otimes X \cong \mathbf{1}$이다. 즉 (3)이 성립한다.

\medskip\noindent
(2) $\Rightarrow$ (3)의 증명은 방금 한 논의의 쌍대이다.

\medskip\noindent
(3)을 가정하자. $L'$과 $L$이 서로 준역이고 $R'$과 $R$이 서로
준역임은 명백하다. 따라서 (1)과 (2)가 성립한다.
\end{proof}

\begin{definition}
\label{categories-definition-invertible}
$\mathcal{C}$를 모노이달 범주라 하자. $\mathcal{C}$의 대상 $X$가
Lemma \ref{categories-lemma-invertible}의 서로 동치인 조건 중 하나(따라서 모두)를
만족하면 이를 {\it 가역}이라고 한다.
\end{definition}

\noindent
$F : \mathcal{C} \to \mathcal{C}'$가 모노이달 범주의 함자이면 $F$는
가역 대상을 가역 대상으로 보냄에 유의하라.

\begin{definition}
\label{categories-definition-dual}
모노이달 범주 $(\mathcal{C}, \otimes, \phi)$와 대상 $X$가 주어졌다고
하자. {\it 왼쪽 쌍대}란 대상 $Y$와 사상
$\eta : \mathbf{1} \to X \otimes Y$ 및
$\epsilon : Y \otimes X \to \mathbf{1}$로서 도표들
$$
\vcenter{
\xymatrix{
X \ar[rd]_1 \ar[r]_-{\eta \otimes 1} &
X \otimes Y \otimes X \ar[d]^{1 \otimes \epsilon} \\
& X
}
}
\quad\text{and}\quad
\vcenter{
\xymatrix{
Y \ar[rd]_1 \ar[r]_-{1 \otimes \eta} &
Y \otimes X \otimes Y \ar[d]^{\epsilon \otimes 1} \\
& Y
}
}
$$
을 가환하게 하는 것이다. 이 상황에서 $X$를 $Y$의
{\it 오른쪽 쌍대}라고 한다.
\end{definition}

\noindent
$F : \mathcal{C} \to \mathcal{C}'$가 모노이달 범주의 함자이고 $Y$가
$X$의 왼쪽 쌍대이면 $F(Y)$는 $F(X)$의 왼쪽 쌍대임에 유의하라.

\begin{lemma}
\label{categories-lemma-left-dual}
$\mathcal{C}$를 모노이달 범주라 하자. $Y$가 $X$의 왼쪽 쌍대이면
$$
\Mor(Z' \otimes X, Z) = \Mor(Z', Z \otimes Y)
\quad\text{and}\quad
\Mor(Y \otimes Z', Z) = \Mor(Z', X \otimes Z)
$$
이며, 이는 $Z$와 $Z'$에 함자적이다.
\end{lemma}

\begin{proof}
사상들
$$
\Mor(Z' \otimes X, Z) \to
\Mor(Z' \otimes X \otimes Y, Z \otimes Y) \to
\Mor(Z', Z \otimes Y)
$$
을 생각하자. 여기서 두 번째 화살표에는 $\eta$를 사용한다. 또한 사상열
$$
\Mor(Z', Z \otimes Y) \to
\Mor(Z' \otimes X, Z \otimes Y \otimes X) \to
\Mor(Z' \otimes X, Z)
$$
을 생각하자. 여기서 두 번째 화살표에는 $\epsilon$을 사용한다. 이
화살표들이 서로 역임을 보이기 위해 사상 $a : Z' \to Z \otimes Y$를
생각하자. 합성
$$
Z' \xrightarrow{\text{id}_{Z'} \otimes \eta}
Z' \otimes X \otimes Y \xrightarrow{a \otimes \text{id}_{X \otimes Y}}
Z \otimes Y \otimes X \otimes Y
\xrightarrow{\text{id}_Z \otimes \epsilon \otimes \text{id}_Y}
Z \otimes Y
$$
이 $a$와 같음을 보여야 한다. 첫 두 화살표의 합성은
$(\text{id}_{Z \otimes Y} \otimes \eta) \circ a :
Z' \to Z \otimes Y \to Z \otimes Y \otimes X \otimes Y$와 같다.
이제 $\text{id}_{Z \otimes Y} \otimes \eta$와
$\text{id}_Z \otimes \epsilon \otimes \text{id}_Y$의 합성은 쌍대의
정의에 의해 항등사상과 같다. 다른 합성도 마찬가지이다. 두 번째 등식의
증명은 생략한다.
\end{proof}

\begin{remark}
\label{categories-remark-left-dual-adjoint}
Lemma \ref{categories-lemma-left-dual}는 특히 $Z \mapsto Z \otimes Y$가
$Z' \mapsto Z' \otimes X$의 오른쪽 수반임을 말한다. 특히 수반함자의
유일성에 의해 $X$의 왼쪽 쌍대가 존재한다면 유일한 동형사상까지
유일하다. 역으로 함자 $Z \mapsto Z \otimes Y$가 함자
$Z' \mapsto Z' \otimes X$의 오른쪽 수반이라고 가정하자. 즉 $Z$와
$Z'$ 모두에 함자적인 전단사
$$
\Mor(Z' \otimes X, Z) \longrightarrow \Mor(Z', Z \otimes Y)
$$
가 주어졌다고 하자. 수반의 단위는 $Z$에 함자적인 사상들
$$
\eta_Z : Z \to Z \otimes X \otimes Y
$$
을 만들고, 수반의 쌍대단위는 $Z'$에 함자적인 사상들
$$
\epsilon_{Z'} : Z' \otimes Y \otimes X \to Z'
$$
을 만든다. 특히
$\eta = \eta_\mathbf{1} : \mathbf{1} \to X \otimes Y$와
$\epsilon = \epsilon_\mathbf{1} : Y \otimes X \to \mathbf{1}$을 얻는다.
단위, 쌍대단위 및 수반 동형사상의 관계에 관한 연습으로서 독자는
$$
(\epsilon \otimes \text{id}_Y) \circ \eta_Y = \text{id}_Y
\quad\text{and}\quad
\epsilon_X \circ (\eta \otimes \text{id}_X) = \text{id}_X
$$
임을 보일 수 있다. 그러나 이것만으로는
$(\epsilon \otimes \text{id}_Y) \circ (\text{id}_Y \otimes \eta) =
\text{id}_Y$ 및
$(\text{id}_X \otimes \epsilon) \circ (\eta \otimes \text{id}_X) =
\text{id}_X$임을 보이기에 충분하지 않다. 일반적으로
$\eta_Y = \text{id}_Y \otimes \eta$인지도,
$\epsilon_X = \epsilon \otimes \text{id}_X$인지도 모르기 때문이다.
이를 위해서는 수반 동형사상이 다음 성질을 가지면 충분하다. 모든
$W, Z, Z'$에 대하여 도표
$$
\xymatrix{
\Mor(Z' \otimes X, Z) \ar[r] \ar[d]_{\text{id}_W \otimes -} &
\Mor(Z', Z \otimes Y) \ar[d]^{\text{id}_W \otimes -} \\
\Mor(W \otimes Z' \otimes X, W \otimes Z) \ar[r] &
\Mor(W \otimes Z', W \otimes Z \otimes Y)
}
$$
가 가환한다. 이 성질이 성립하면 {\it 이 수반은 주어진 텐서 구조와
호환된다}고 하겠다. 따라서 $Z \mapsto Z \otimes Y$가 주어진 텐서
구조와 호환되는 $Z' \mapsto Z' \otimes X$의 오른쪽 수반이라는 조건은
왼쪽 쌍대라는 성질의 동치인 정식화이다.
\end{remark}

\begin{lemma}
\label{categories-lemma-tensor-dual}
$\mathcal{C}$를 모노이달 범주라 하자. $Y_i$, $i = 1, 2$가 각각
$X_i$, $i = 1, 2$의 왼쪽 쌍대이면 $Y_2 \otimes Y_1$은
$X_1 \otimes X_2$의 왼쪽 쌍대이다.
\end{lemma}

\begin{proof}
수반의 유일성과 Remark \ref{categories-remark-left-dual-adjoint}에서 따른다.
\end{proof}

\noindent
$(\mathcal{C}, \otimes)$에 대한 {\it 가환성 제약}이란 함자적 동형사상
$$
\psi : X \otimes Y \longrightarrow Y \otimes X
$$
으로서 합성
$$
X \otimes Y \xrightarrow{\psi} Y \otimes X \xrightarrow{\psi} X \otimes Y
$$
이 항등사상이 되게 하는 것이다. 주어진 결합성 제약 $\phi$에 대하여,
모든 대상 $X, Y, Z$에 대해 도표
\begin{equation}
\label{categories-equation-hexagon}
\vcenter{
\xymatrix{
X \otimes (Y \otimes Z) \ar[r]_\phi \ar[d]^\psi &
(X \otimes Y) \otimes Z \ar[r]_\psi &
Z \otimes (X \otimes Y) \ar[d]^\phi \\
X \otimes (Z \otimes Y) \ar[r]^\phi &
(X \otimes Z) \otimes Y \ar[r]^\psi &
(Z \otimes X) \otimes Y
}
}
\end{equation}
가 가환하면 $\psi$가 이 제약과 {\it 호환된다}고 한다.

\begin{definition}
\label{categories-definition-symmetric-monoidal-category}
$\mathcal{C}$가 범주이고
$\otimes : \mathcal{C} \otimes \mathcal{C} \to \mathcal{C}$가 함자이며,
$\phi$가 결합성 제약이고 $\psi$가 $\phi$와 호환되는 가환성 제약인
네쌍 $(\mathcal{C}, \otimes, \phi, \psi)$에 단위원이 존재하면 이를
{\it 대칭 모노이달 범주}라고 한다.
\end{definition}

\noindent
분명히 $(\mathcal{C}, \otimes, \phi, \psi)$가 대칭 모노이달 범주이면
$(\mathcal{C}, \otimes, \phi)$는 모노이달 범주이고, 위에서 논의한
용어와 표기를 사용할 수 있다.

\begin{lemma}
\label{categories-lemma-symmetric-monoidal-coherence}
대칭 모노이달 범주
$\mathcal{C}, \otimes, \phi, \psi, \mathbf{1}, 1$에서 다음이 성립한다.
\begin{enumerate}
\item 화살표 $1 \circ \psi, 1 :
\mathbf{1} \otimes \mathbf{1} \to \mathbf{1}$은 서로 같다.
\item 화살표 $\text{id}_X \otimes l_Y,
(l_X \otimes \text{id}) \circ (\psi \otimes \text{id}_Y):
X \otimes \mathbf{1} \otimes Y \to X \otimes Y$는 서로 같다.
\end{enumerate}
대칭 모노이달 범주는 \cite[Theorem 5.1]{associativity}의 가정들을
만족한다.
\end{lemma}

\begin{proof}
유일한 동형사상 $a : \mathbf{1} \to \mathbf{1}$에 대하여
$\psi = a \otimes \text{id}_\mathbf{1}$로 쓸 수 있다. Lemma
\ref{categories-lemma-monoidal-unit-unique}는
$a \otimes \text{id}_\mathbf{1} = \text{id}_\mathbf{1} \otimes a$임을
함의한다. $\psi$의 함자성은 도표
$$
\xymatrix{
(\mathbf{1} \otimes \mathbf{1}) \otimes \mathbf{1}
\ar[r]_\psi
\ar[d]_{1 \otimes \text{id}_\mathbf{1}} &
\mathbf{1} \otimes (\mathbf{1} \otimes \mathbf{1})
\ar[d]^{\text{id}_\mathbf{1} \otimes 1} \\
\mathbf{1} \otimes \mathbf{1}  \ar[r]^\psi &
\mathbf{1} \otimes \mathbf{1}
}
$$
가 가환함을 말한다. 따라서 위쪽 화살표는
$a \otimes \text{id}_\mathbf{1} \otimes \text{id}_\mathbf{1}$과 같다.
그러므로 $X = Y = Z = \mathbf{1}$일 때의
(\ref{categories-equation-hexagon})은
$a \otimes \text{id}_\mathbf{1} \otimes \text{id}_\mathbf{1}$이 자신의
제곱과 같음을 말한다. 따라서 $a = \text{id}_\mathbf{1}$이다. 이로써
(1)이 증명된다.

\medskip\noindent
(2)는 우리의 모노이달 범주에서 정의역과 공역을 $X$와 식별했을 때
$\psi : X \otimes \mathbf{1} \to \mathbf{1} \otimes X$가 항등사상임을
말한다. 이는 $\mathbf{1}, X, \mathbf{1}$에 대한
(\ref{categories-equation-hexagon}), 즉 도표
$$
\xymatrix{
\mathbf{1} \otimes (X \otimes \mathbf{1}) \ar[r]_\phi \ar[d]^\psi &
(\mathbf{1} \otimes X) \otimes \mathbf{1} \ar[r]_\psi &
\mathbf{1} \otimes (\mathbf{1} \otimes X) \ar[d]^\phi \\
\mathbf{1} \otimes (\mathbf{1} \otimes X) \ar[r]^\phi &
(\mathbf{1} \otimes \mathbf{1}) \otimes X \ar[r]^\psi &
(\mathbf{1} \otimes \mathbf{1}) \otimes X
}
$$
가 가환한다는 사실과, 이 도표에서 하나를 제외한 모든 사상 $\psi$가
항등사상과 같다는 사실에서 따른다.

\medskip\noindent
(1)과 (2)에 더하여 도표들 (\ref{categories-equation-pentagon}),
(\ref{categories-equation-triangle}), (\ref{categories-equation-hexagon})의 가환성과 Lemma
\ref{categories-lemma-monoidal-coherence}의 결과는 보조정리의 마지막 주장을
함의한다.
\end{proof}

\noindent
\cite[Theorem 5.1]{associativity}에서는 위 보조정리의 여섯쌍
$(\mathcal{C}, \otimes, \phi, \psi, \mathbf{1}, 1)$이 어떻게
정합적인지 설명한다. 이를 다시 말하면, 우리 의미의 대칭 모노이달
범주에서 모든 $n \geq 1$과 $\{1, \ldots, n\}$의 순열 $\sigma$에 대하여
함자적 동형사상
$$
X_1 \otimes \ldots \otimes X_n
\to
X_{\sigma(1)} \otimes \ldots \otimes X_{\sigma(n)}
$$
이 있고, 이 동형사상들로 만들 수 있는 모든 가능한 도표가 가환하며,
이 동형사상들은 모노이달 범주의 구조와 호환된다. 예를 들어 한 자리에
$\mathbf{1}$을 삽입할 때의 동형사상들과 호환된다.

\begin{lemma}
\label{categories-lemma-dual-symmetric}
$(\mathcal{C}, \otimes, \phi, \psi)$를 대칭 모노이달 범주라 하자.
$X$를 $\mathcal{C}$의 대상이라 하고 $Y$,
$\eta : \mathbf{1} \to X \otimes Y$, 및
$\epsilon : Y \otimes X \to \mathbf{1}$을 Definition
\ref{categories-definition-dual}에서와 같은 $X$의 왼쪽 쌍대라 하자. 그러면
$\eta' = \psi \circ \eta : \mathbf{1} \to Y \otimes X$와
$\epsilon' = \epsilon \circ \psi : X \otimes Y \to \mathbf{1}$은
$X$를 $Y$의 왼쪽 쌍대로 만든다.
\end{lemma}

\begin{proof}
생략한다. 힌트: 정의를 사용하는 즐거운 연습이다.
\end{proof}

\begin{definition}
\label{categories-definition-functor-symmetric-monoidal-categories}
$\mathcal{C}$와 $\mathcal{C}'$를 대칭 모노이달 범주라 하자.
{\it 대칭 모노이달 범주의 함자} $F : \mathcal{C} \to \mathcal{C}'$는
표시한 함자 $F$와 $X$ 및 $Y$에 함자적인 동형사상
$$
F(X) \otimes F(Y) \to F(X \otimes Y)
$$
으로 주어지며, $F$가 모노이달 범주의 함자이고 모든 대상 $X$와 $Y$에
대하여 도표
$$
\xymatrix{
F(X) \otimes F(Y) \ar[r] \ar[d] &
F(X \otimes Y) \ar[d] \\
F(Y) \otimes F(X) \ar[r] &
F(Y \otimes X)
}
$$
를 가환하게 하는 것이다.
\end{definition}

\begin{remark}
\label{categories-remark-internal-hom-monoidal}
$\mathcal{C}$를 모노이달 범주라 하자. $\mathcal{C}$의 모든 대상 쌍
$X, Y$에 대하여 $\mathcal{C}$의 대상 $hom(X, Y)$가 존재하고
$$
\Mor(X, hom(Y, Z)) = \Mor(X \otimes Y, Z)
$$
가 $X, Y, Z$에 함자적으로 성립하면 $\mathcal{C}$가
{\it 내부 hom}을 갖는다고 한다. Yoneda 보조정리에 의해 쌍함자
$(X, Y) \mapsto hom(X, Y)$가 존재하면 유일한 동형사상까지 결정된다.
내부 hom이 주어지면 다음 표준 사상들을 얻는다.
\begin{enumerate}
\item $hom(X, Y) \otimes X \to Y$,
\item $hom(Y, Z) \otimes hom(X, Y) \to hom(X, Z)$,
\item $Z \otimes hom(X, Y) \to hom(X, Z \otimes Y)$,
\item $Y \to hom(X, Y \otimes X)$,
\item $\mathcal{C}$가 대칭 모노이달인 경우
$hom(Y, Z) \otimes X \to hom(hom(X, Y), Z)$.
\end{enumerate}
(1)의 사상은 $\text{id}_{hom(X, Y)}$를
$\Mor(hom(X, Y), hom(X, Y)) \to \Mor(hom(X, Y) \otimes X, Y)$로 보낸
상이다. $hom(X, Z)$의 정의 성질을 이용해 (2)의 사상을 구성하려면 사상
$$
hom(Y, Z) \otimes hom(X, Y) \otimes X \longrightarrow Z
$$
을 구성해야 한다. (1)에 의해 사상들
$hom(X, Y) \otimes X \to Y$와 $hom(Y, Z) \otimes Y \to Z$가 있으므로
그러한 사상이 존재한다. $hom(X, Z \otimes Y)$의 정의 성질을 이용해
(3)의 사상을 구성하려면 사상
$$
Z \otimes hom(X, Y) \otimes X \to Z \otimes Y
$$
을 구성해야 한다. 여기에는 (1)의 사상을 $a$라 할 때
$\text{id}_Z \otimes a$를 사용한다. (4)의 사상을 구성하기 위해 (3)에
의해 이미 사상 $Y \otimes hom(X, X) \to hom(X, Y \otimes X)$가 있음을
관찰하자. 따라서 사상 $\mathbf{1} \to hom(X, X)$을 구성하면 충분하고,
이를 위해 $\Mor(\mathbf{1}, hom(X, X))$에서 표준 동형사상
$\mathbf{1} \otimes X \to X$에 대응하는 원소를 취한다. 이 표준
동형사상은 $\Mor(\mathbf{1} \otimes X, X)$의 원소이다.
마지막으로 (5)를 보자. $hom(hom(X, Y), Z)$의 보편 성질에 의해 사상
$$
hom(Y, Z) \otimes X \otimes hom(X, Y) \longrightarrow Z
$$
을 구성하면 충분하다. 가환성 제약을 사용하여 마지막 두 텐서 인자를
바꾼 다음 사상들 $hom(X, Y) \otimes X \to Y$와
$hom(Y, Z) \otimes Y \to Z$를 사용하여 이를 구성한다.
\end{remark}





\section{점선 화살표의 범주}
\label{categories-section-dotted-arrows}

\noindent 
$(2,1)$-범주에서 어떤 ``점선 화살표의 범주''들을 논의한다. 이들은
대수적 스택에 관한 여러 올림 판정법을 정식화할 때 나타난다. 예를 들어
Morphisms of Stacks, Section \ref{stacks-morphisms-section-valuative}와
More on Morphisms of Stacks, Section
\ref{stacks-more-morphisms-section-formally-smooth}를 보라.

\begin{definition}
\label{categories-definition-dotted-arrows}
$\mathcal{C}$를 $(2,1)$-범주라 하자. $\mathcal{C}$ 안의 실선으로 된
$2$-가환 도표
\begin{equation}
\label{categories-equation-dotted-arrows}
\vcenter{
\xymatrix{
S \ar[r]_-x \ar[d]_j & X \ar[d]^f \\
T \ar[r]^-y \ar@{..>}[ru] & Y
}
}
\end{equation}
를 생각하자. 이 도표의 $2$-가환성을 나타내는 $2$-동형사상
$$
\gamma : y \circ j \rightarrow f \circ x
$$
을 하나 고정하자. (\ref{categories-equation-dotted-arrows})와 $\gamma$가 주어졌을
때, \emph{점선 화살표}란 사상 $a \colon T \to X$와 $2$-동형사상
$\alpha : a \circ j \to x$, $\beta : y \to f \circ a$로 이루어진 세쌍
$(a, \alpha, \beta)$으로서
$\gamma = (\text{id}_f \star \alpha) \circ (\beta \star \text{id}_j)$를
만족하는 것이다. 다시 말해 도표
$$
\xymatrix{
& f \circ a \circ j \ar[rd]^{\text{id}_f \star \alpha} \\
y \circ j \ar[ru]^{\beta \star \text{id}_j} \ar[rr]^\gamma & &
f \circ x
}
$$
가 가환해야 한다. {\it 점선 화살표의 사상}
$(a, \alpha, \beta) \to (a', \alpha', \beta')$이란
$\alpha = \alpha' \circ (\theta \star \text{id}_j)$ 및
$\beta' = (\text{id}_f \star \theta) \circ \beta$를 만족하는
$2$-화살표 $\theta : a \to a'$이다.
\end{definition}

\noindent
Definition \ref{categories-definition-dotted-arrows}의 상황에는 수반된
\emph{점선 화살표의 범주}가 있다. 이 범주는 군체이다. 일반적으로 이는
$\gamma$에 의존할 수 있다. 다음 두 보조정리는 점선 화살표의 범주가
$f$의 기저변환과 합성에 관하여 잘 행동함을 말한다.

\begin{lemma}
\label{categories-lemma-cat-dotted-arrows-base-change}
$\mathcal{C}$를 $(2,1)$-범주라 하자. $\mathcal{C}$ 안에 $2$-가환 도표
$$
\xymatrix{
S \ar[r]_-{x'} \ar[d]_j &
X' \ar[d]^p \ar[r]_q &
X \ar[d]^f \\
T \ar[r]^-{y'} &
Y' \ar[r]^g &
Y
}
$$
가 주어졌다고 하자. 여기서 오른쪽 사각형은 $2$-동형사상
$\phi \colon g \circ p \to f \circ q$에 관하여 $2$-카르테시안이다.
$2$-화살표 $\gamma' : y' \circ j \to p \circ x'$를 선택하자.
$x = q \circ x'$, $y = g \circ y'$로 놓고
$\gamma : y \circ j \to f \circ x$를 $2$-동형사상
$\gamma = (\phi \star \text{id}_{x'}) \circ (\text{id}_g \star \gamma')$로
놓자. 그러면 왼쪽 사각형과 $\gamma'$에 대한 점선 화살표의 범주
$\mathcal{D}'$는 바깥 직사각형과 $\gamma$에 대한 점선 화살표의 범주
$\mathcal{D}$와 동치이다.
\end{lemma}

\begin{proof}
함자 $\mathcal{D}' \to \mathcal{D}$가 존재한다. 이는 대상 위에서
$(a, \alpha, \beta) \mapsto (q \circ a, \text{id}_q \star \alpha,
(\phi \star \text{id}_a) \circ (\text{id}_g \star \beta))$이고,
화살표 위에서 $\theta \mapsto \text{id}_q \star \theta$이다. 이 함자
$\mathcal{D}' \to \mathcal{D}$가 동치임은 Section
\ref{categories-section-2-fibre-products}에서와 같이 $2$-섬유곱의 보편 성질에서
형식적으로 따른다. 세부사항은 생략한다.
\end{proof}

\begin{lemma}
\label{categories-lemma-cat-dotted-arrows-composition}
$\mathcal{C}$를 $(2,1)$-범주라 하자. $\mathcal{C}$ 안에 실선으로 된
$2$-가환 도표
$$
\xymatrix{
S \ar[r]_-x \ar[dd]_j & X \ar[d]^f \\
& Y \ar[d]^g \\
T \ar[r]^-z \ar@{..>}[ruu] & Z
}
$$
가 주어졌다고 하자. $2$-동형사상
$\gamma \colon z \circ j \to g \circ f \circ x$를 선택하자. 바깥
직사각형과 $\gamma$에 대한 점선 화살표의 범주를 $\mathcal{D}$라 하자.
실선으로 된 사각형
$$
\xymatrix{
S \ar[r]_-{f \circ x} \ar[d]_j & Y \ar[d]^g \\
T \ar[r]^-z \ar@{..>}[ru] & Z
}
$$
과 $\gamma$에 대한 점선 화살표의 범주를 $\mathcal{D}'$라 하자. 그러면
$\mathcal{D}$는 다음 성질을 갖는 한 범주 $\mathcal{D}''$와 동치이다.
함자 $\mathcal{D}'' \to \mathcal{D}'$가 존재하여 $\mathcal{D}''$를
$\mathcal{D}'$ 위에서 군체로 섬유화된 범주로 만들고, 그 섬유 범주들은
꼴
$$
\xymatrix{
S \ar[r]_-x \ar[d]_j & X \ar[d]^f \\
T \ar[r]^-y \ar@{..>}[ru] & Y
}
$$
의 어떤 실선으로 된 사각형과 $2$-동형사상
$y \circ j \to f \circ x$의 어떤 선택에 대한 점선 화살표의 범주와
동형이다.
\end{lemma}

\begin{proof}
범주 $\mathcal{D}''$를 다음과 같이 구성하자. 그 대상은 튜플
$(a,\alpha,\beta,b,\eta)$이며, 여기서 $(a,\alpha,\beta)$는
$\mathcal{D}$의 대상이고 $b \colon T \rightarrow Y$는 $1$-사상이며
$\eta \colon b \rightarrow f \circ a$는 $2$-동형사상이다.
$\mathcal{D}''$ 안의 사상
$(a,\alpha,\beta,b,\eta) \rightarrow (a',\alpha',\beta',b',\eta')$은
쌍 $(\theta_1,\theta_2)$이다. 여기서
$\theta_1 \colon a \rightarrow a'$은 $\mathcal{D}$ 안의 화살표
$(a, \alpha, \beta) \rightarrow (a', \alpha', \beta')$를 정의하고,
$\theta_2 \colon b \rightarrow b'$은 호환 조건
$\eta' \circ \theta_2 = (\text{id}_f \star \theta_1) \circ \eta$를
만족하는 $2$-동형사상이다.

\medskip\noindent
함자 $\mathcal{D}'' \rightarrow \mathcal{D}'$가 존재한다. 이는 대상
위에서
$(a, \alpha, \beta, b, \eta) \mapsto
(b, (\text{id}_f \star \alpha) \circ (\eta \star \text{id}_j),
(\text{id}_g \star \eta^{-1}) \circ \beta)$이고, 화살표 위에서
$(\theta_1,\theta_2) \mapsto \theta_2$이다. 그러면
$\mathcal{D}'' \rightarrow \mathcal{D}'$는 군체로 섬유화되어 있다.

\medskip\noindent
$(y, \delta, \epsilon)$을 $\mathcal{D}'$의 대상이라 하자. 마지막으로
표시한 도표에서 $y \circ j \rightarrow f \circ x$가 $\delta$로 주어졌을
때의 점선 화살표의 범주를 $\mathcal{D}_{y,\delta}$로 쓰자. 함자
$\mathcal{D}_{y,\delta} \rightarrow \mathcal{D}''$가 존재한다. 이는 대상
위에서
$(a, \alpha, \eta) \mapsto
(a, \alpha, (\text{id}_g \star \eta) \circ \epsilon, y, \eta)$이고,
화살표 위에서 $\theta \mapsto (\theta, \text{id}_y)$이다. 이 함자는
$\mathcal{D}_{y,\delta}$에서 $(y,\delta,\epsilon)$ 위에 놓인
$\mathcal{D}'' \rightarrow \mathcal{D}'$의 섬유 범주로 가는 동형사상을
준다.

\medskip\noindent
또한 함자 $\mathcal{D} \rightarrow \mathcal{D}''$가 존재한다. 이는 대상
위에서
$(a,\alpha,\beta) \mapsto (a,\alpha,\beta,f \circ a, \text{id}_{f \circ a})$이고,
화살표 위에서 $\theta \mapsto (\theta, \text{id}_f \star \theta)$이다.
이 함자는 완전충실하고 본질적으로 전사이므로 동치이다. 세부사항은
생략한다.
\end{proof}
